% ============================================================
%  preamble.tex — styling, environments, macros
%
%  Shared by BOTH editions. main-en.tex and main-pl.tex each define
%  \booklang before % ============================================================
%  preamble.tex — styling, environments, macros
%
%  Shared by BOTH editions. main-en.tex and main-pl.tex each define
%  \booklang before % ============================================================
%  preamble.tex — styling, environments, macros
%
%  Shared by BOTH editions. main-en.tex and main-pl.tex each define
%  \booklang before % ============================================================
%  preamble.tex — styling, environments, macros
%
%  Shared by BOTH editions. main-en.tex and main-pl.tex each define
%  \booklang before \input{preamble}; nothing else differs between them at
%  this level. Every user-visible string lives in lang/en.tex or lang/pl.tex,
%  so a label that appears in one edition and not the other is a missing macro
%  rather than a silent divergence.
% ============================================================

% \booklang must be set by the main file. Default to English rather than
% failing, so a stray \input{preamble} still compiles and says what it did.
\providecommand{\booklang}{en}

\usepackage{etoolbox}   % loaded first: the language switch below needs it

% ---------- Page geometry ----------
% Same 17 x 24 cm block as the other two books in this series, so a reader who
% owns one recognises the second. Frames are short, so the measure matters
% more here than in a prose book: a frame that runs past a page turn defeats
% the cover-the-next-frame discipline the whole method rests on.
\usepackage[
  paperwidth=17cm, paperheight=24cm,
  inner=2.2cm, outer=1.8cm, top=2.2cm, bottom=2.4cm,
  head=23pt, headsep=0.7cm, footskip=1.2cm
]{geometry}

% ---------- Encoding & fonts ----------
\usepackage[T1]{fontenc}
\usepackage[utf8]{inputenc}
\IfFileExists{lmodern.sty}{\usepackage{lmodern}}{}
\IfFileExists{newtxtext.sty}{\usepackage{newtxtext}}{}
% amsmath goes here, not down with the other packages, because newtxmath
% patches amsmath's internals and must be loaded AFTER it. The wrong order does
% not fail on a machine without newtxmath -- which is most CI images -- and
% fails on the author's full TeX Live.
%
% amssymb only when newtxmath is absent: newtxmath supplies the AMS symbols
% itself (its amssymbols option is on by default) and loading both clashes.
\usepackage{amsmath}
\IfFileExists{newtxmath.sty}{\usepackage{newtxmath}}{\usepackage{amssymb}}
\IfFileExists{inconsolata.sty}{\usepackage[varqu,scaled=0.95]{inconsolata}}{}
\IfFileExists{lmodern.sty}{\usepackage{microtype}}{\usepackage[expansion=false]{microtype}}

% ---------- Language ----------
% Loading babel with a language whose .ldf is absent is a *fatal* error, not a
% warning, so both the package and each language file are probed before either
% is requested. This is inherited from the sibling books, where a minimal TeX
% installation without texlive-lang-polish aborted the run with no PDF and an
% error message that named neither babel nor the language.
%
% The Polish edition needs Polish as the main language for hyphenation; the
% English edition needs British. Both editions load the other language too,
% because the glossary appendix is bilingual in both.
\IfFileExists{babel.sty}{%
  \IfFileExists{polish.ldf}{%
    \IfFileExists{british.ldf}{%
      \ifdefstring{\booklang}{pl}%
        {\usepackage[british,main=polish]{babel}}%
        {\usepackage[polish,main=british]{babel}}%
    }{\usepackage[polish]{babel}}%
  }{%
    \IfFileExists{british.ldf}{\usepackage[british]{babel}}{}%
  }%
}{}

% babel-polish makes " an active shorthand. listings reads its body verbatim
% and an active character in a verbatim context is a category-code accident
% waiting to happen, so the shorthand is turned off globally. Polish quotation
% marks are produced with \enquote-style macros from lang/pl.tex instead.
\ifdefstring{\booklang}{pl}{\AtBeginDocument{\shorthandoff{"}}}{}

% ---------- Core packages ----------
\usepackage{graphicx}
\usepackage{xcolor}
\usepackage{booktabs}
\usepackage{tabularx}
\usepackage{longtable}
\usepackage{array}
\usepackage{enumitem}
\IfFileExists{mathtools.sty}{\usepackage{mathtools}}{}
\usepackage{listings}
\usepackage[most]{tcolorbox}
\usepackage{fancyhdr}
\usepackage{titlesec}
\usepackage{caption}
\usepackage{imakeidx}
% csquotes with autostyle follows babel's active language, so the identical
% source \enquote{x} sets 'x' in the English edition and the Polish quotation
% marks in the Polish one. This is not a nicety: babel-polish makes " an active
% shorthand, which is turned off above precisely because an active character
% near a listings body is a category-code accident waiting to happen -- so a
% quotation mark has to come from a macro rather than from a keyboard key.
\IfFileExists{csquotes.sty}{\usepackage[autostyle=true]{csquotes}}{%
  \providecommand{\enquote}[1]{``##1''}}
\usepackage[hidelinks]{hyperref}

% ---------- Palette ----------
% Deliberately the same family as the sibling books, shifted green so the three
% spines are distinguishable on a shelf.
\definecolor{mblue}{HTML}{1F4E79}
\definecolor{mgreen}{HTML}{2E6B4F}
\definecolor{mteal}{HTML}{0E7C7B}
\definecolor{mamber}{HTML}{B26A00}
\definecolor{mred}{HTML}{9B2C2C}
\definecolor{mgrey}{HTML}{5A5A5A}
\definecolor{mlight}{HTML}{F4F6F8}
\definecolor{mfaint}{HTML}{FBFCFD}
\definecolor{codebg}{HTML}{FAFAFA}
\definecolor{codeframe}{HTML}{D8DEE4}
\definecolor{kw}{HTML}{0B5394}
\definecolor{str}{HTML}{9B2C2C}
\definecolor{cmt}{HTML}{4F7A4F}
\definecolor{typ}{HTML}{0E7C7B}

\hypersetup{
  colorlinks=true,
  linkcolor=mblue,
  urlcolor=mteal,
  citecolor=mblue,
  pdfauthor={Konrad Cinkusz}
}

% \ifpl{Polish text}{English text} -- for the handful of places where a string
% is structural rather than content and belongs in the shared files:
% structure.tex's part titles, and the main files' front matter wiring.
% Anything longer than a few words belongs in lang/ or in the edition's own
% source, not here.
% \DeclareRobustCommand, not \newcommand. A fragile \ifpl inside a \part title
% is expanded one level as it is written to the .toc, which lands
% "\ifdefstring {en}{pl}{...}{...}" in the file -- and that fails on the next
% run with "Extra }, or forgotten \endgroup", an error that names neither the
% part nor this macro.
\DeclareRobustCommand{\ifpl}[2]{\ifdefstring{\booklang}{pl}{#1}{#2}}

% A robust macro cannot go into a PDF bookmark string either -- hyperref strips
% \ifdefstring and warns once per part title. Tell it which branch to take
% instead, so the bookmarks carry the right language rather than nothing.
\makeatletter
\pdfstringdefDisableCommands{%
  \ifdefstring{\booklang}{pl}{\let\ifpl\@firstoftwo}{\let\ifpl\@secondoftwo}%
}
\makeatother

% ---------- Language strings ----------
% Every user-visible word the macros below typeset comes from here. Nothing in
% programs/ or appendices/ may hard-code an English or a Polish word inside a
% macro argument that the other edition also uses.
\input{lang/\booklang}

% ============================================================
%  PROGRAM NUMBERING
%  ---------------------------------------------------------
%  Part F programs are F1..F13; the main programs restart at 1. That is
%  Stroud's scheme and it is worth reproducing, because the Foundation part is
%  explicitly optional and numbering it separately is what makes skipping it
%  feel permitted rather than delinquent.
%
%  \theHchapter is hyperref's *destination* name, and it must stay unique
%  across the whole document. Resetting the chapter counter without also
%  changing \theHchapter produces duplicate PDF destinations, a pile of
%  warnings, and bookmarks that jump to the wrong program.
% ============================================================
\newcommand{\foundationnumbering}{%
  \setcounter{chapter}{0}%
  \renewcommand{\thechapter}{F\arabic{chapter}}%
  \renewcommand{\theHchapter}{F\arabic{chapter}}%
}
\newcommand{\mainnumbering}{%
  \setcounter{chapter}{0}%
  \renewcommand{\thechapter}{\arabic{chapter}}%
  \renewcommand{\theHchapter}{M\arabic{chapter}}%
}

% \appendix resets the chapter counter and switches \thechapter to letters, but
% it knows nothing about \theHchapter, so appendix A would claim the same PDF
% destination as main program 1. Patch it once, here.
\apptocmd{\appendix}{\renewcommand{\theHchapter}{A\Alph{chapter}}}{}{%
  \PackageWarning{mfabook}{Could not patch \string\appendix\space for hyperref destinations}}

% A program is a chapter. \chaptername is set from the language file so the
% word above the title is "Program" in both editions (it happens to be the
% same word, which is a small piece of luck).
\AtBeginDocument{\renewcommand{\chaptername}{\lblProgram}}

% ---------- Headings ----------
% \raggedright on the title body is load-bearing, not cosmetic: at \Huge on a
% 13 cm text block a long title that cannot break overflows the margin badly.
\titleformat{\chapter}[display]
  {\normalfont\huge\bfseries\color{mblue}\raggedright}
  {\normalfont\normalsize\color{mgrey}\MakeUppercase{\chaptertitlename}\ \thechapter}
  {8pt}{\Huge\raggedright}
\titlespacing*{\chapter}{0pt}{10pt}{28pt}
\titleformat{\section}{\normalfont\Large\bfseries\color{mblue}}{\thesection}{0.8em}{}
\titleformat{\subsection}{\normalfont\large\bfseries\color{mgrey}}{\thesubsection}{0.7em}{}

% head=23pt in the geometry options above, not \setlength{\headheight}:
% fancyhdr computes the height it needs from the running-head font and the
% rule, wants 22.5 pt at this size, and silently overprints the top of the text
% block if it does not get it. Setting it through geometry keeps the text block
% where geometry put it instead of pushing it down.
\pagestyle{fancy}
\fancyhf{}
\fancyhead[LE]{\small\nouppercase{\leftmark}}
% The recto head carries the FRAME number, not the section. In a book of
% numbered frames "where am I" is a frame question: every Summary
% back-reference, every Quiz route and every cross-reference between programs
% is a frame number, and hunting for one by page is the single most annoying
% thing about reading a programmed text. \theframeno at shipout is the last
% frame begun on the page, which is what a reader flicking through wants.
% Guarded: the front matter and the appendices have no frames, and a head
% reading "Frame 0" over the introduction is worse than no head at all.
\fancyhead[RO]{\ifnum\value{frameno}>0\relax\small\lblFrame~\theframeno\fi}
\fancyfoot[LE,RO]{\small\thepage}
\renewcommand{\headrulewidth}{0.4pt}

% This MUST come after \pagestyle{fancy}: fancyhdr installs its own \chaptermark
% at that point and silently overwrites anything defined earlier. The symptom is
% a running-head change that appears to do nothing at all.
%
% The running head carries the *frame* number as well as the program, because
% in a book of numbered frames "where am I" is a frame question, not a page
% question, and the summary's back-references are frame numbers.
% \@chapapp, not \lblProgram. \appendix switches \@chapapp from \chaptername
% to \appendixname, and babel supplies both in the right language; hard-coding
% the word gives an appendix a running head reading "Program A. Answers" over a
% page whose own chapter head reads "APPENDIX A".
\makeatletter
\renewcommand{\chaptermark}[1]{\markboth{\@chapapp\ \thechapter.\ #1}{}}
\makeatother
\renewcommand{\sectionmark}[1]{\markright{\thesection\ #1}}

% ============================================================
%  FRAMES — the unit the whole method is built from
%  ---------------------------------------------------------
%  A program is a stream of numbered frames. Nearly every frame ends by asking
%  the reader for something, and the answer opens the *next* frame, so that the
%  reader can cover what is below and commit to an answer before seeing it.
%
%  The environment is called `fr` rather than `frame` for two reasons. First,
%  \frame is already defined by LaTeX2e (it draws a box in a picture
%  environment) and an environment named `frame` would silently redefine it.
%  Second, this is typed forty to seventy times per program, so it wants to be
%  short in the source.
%
%    \begin{fr}
%    \ans{$x = 3$}            % the answer to the previous frame, optional
%    Teaching text, ending in a question.
%    \end{fr}
% ============================================================

\newcounter{frameno}[chapter]
\renewcommand{\theframeno}{\arabic{frameno}}

% Frames are referenced constantly -- by the Summary, by the Quiz, by the
% "Can you?" checklist and by other programs. \label inside a frame must
% therefore capture the frame number, not the section number.
\makeatletter
\newcommand{\framelabelfix}{%
  \@namedef{theHframeno}{\theHchapter.\arabic{frameno}}%
}
\makeatother

\newcommand{\framerule}{%
  \par\penalty-200\vspace{9pt}%
  \noindent\textcolor{codeframe}{\rule{\linewidth}{0.5pt}}%
  \par\vspace{5pt}%
}

% The frame number as a filled badge, set in the text flow rather than in the
% margin. A margin number would have to alternate sides under `twoside` and
% would be the first thing to overflow on a 17 cm page; a badge cannot.
\newcommand{\framebadge}{%
  \begingroup
  \setlength{\fboxsep}{2.5pt}%
  \colorbox{mblue}{\textcolor{white}{\bfseries\scriptsize\theframeno}}%
  \endgroup\hspace{0.7em}%
}

\newenvironment{fr}{%
  \framerule
  \refstepcounter{frameno}%
  \nopagebreak
  \noindent\framebadge\ignorespaces
}{\par}

% The answer to the previous frame. Set apart and centred so the eye can find
% the edge to cover, and so a reader who has answered can check in one glance
% without reading on.
\newcommand{\ans}[1]{%
  \par\nobreak\vspace{-4pt}%
  \begin{tcolorbox}[
    enhanced, sharp corners, boxrule=0pt, leftrule=2.5pt,
    colback=mfaint, colframe=mgreen,
    left=8pt, right=8pt, top=4pt, bottom=4pt,
    width=0.86\linewidth, center, halign=center,
    before skip=2pt, after skip=7pt]
    #1
  \end{tcolorbox}%
  \nobreak\noindent\ignorespaces
}

% A multi-paragraph or worked answer, where a centred box would be wrong.
\newenvironment{ansblock}{%
  \par\nopagebreak
  \begin{tcolorbox}[
    enhanced, breakable, sharp corners, boxrule=0pt, leftrule=2.5pt,
    colback=mfaint, colframe=mgreen,
    left=8pt, right=8pt, top=5pt, bottom=5pt]
}{%
  \end{tcolorbox}%
  \nopagebreak\par\vspace{2pt}\noindent\ignorespaces
}

% The row of dots. Stroud's signal that the reader is to supply a word, a
% number or an expression. \blank is inline; \dotline occupies its own line,
% which is what a longer derivation step wants.
% \penalty0 offers TeX a break opportunity immediately before the dots. Without
% it the run of dots is glued to the last word of the question, and a long
% question in the Polish edition -- where the words are longer and hyphenate
% less freely -- runs into the margin.
\newcommand{\blank}[1][4em]{\penalty0\makebox[#1]{\dotfill}}
\newcommand{\dotline}{\par\vspace{2pt}\noindent\makebox[\linewidth]{\dotfill}\par\vspace{2pt}}

% "Now one for you to do." The third rung of the scaffolding gradient, and the
% one most often skipped by authors who think they have written an exercise
% when they have written another worked example.
\newcommand{\yourturn}{%
  \par\vspace{4pt}\noindent
  {\bfseries\color{mgreen}\lblYourTurn}\par\nobreak\vspace{2pt}\noindent\ignorespaces
}

% ============================================================
%  THE PROGRAM SKELETON
%  ---------------------------------------------------------
%  Every program carries the same seven parts, in the same order. They are
%  macros rather than a convention so that a missing one is a build-time
%  absence rather than something a reader notices.
% ============================================================

% ---- Learning outcomes, on entry ----------------------------------------
% Stored as well as printed, because "Can you?" at the end of the program must
% be the same list, one for one. Storing it is what makes that structural
% rather than a promise: the checklist is generated from the outcomes, so the
% two cannot drift.
\newcounter{outcomeno}[chapter]

\makeatletter
% The key under which everything belonging to the current program is stored.
% It must expand to the printed program number (F1, F2, ..., 1, 2, ...), which
% is what makes an answer findable from the program that owns it.
\newcommand{\mfa@progkey}{\thechapter}
\newcommand{\outcome}[2]{% #1 = frames that teach it, #2 = the outcome
  \stepcounter{outcomeno}%
  \csgdef{mfa@out@\mfa@progkey @\theoutcomeno}{#2}%
  \csgdef{mfa@outfr@\mfa@progkey @\theoutcomeno}{#1}%
  % The "Can you?" row is built here rather than looped over later. A \loop
  % inside a tabularx body does not survive alignment scanning: TeX needs to
  % see the & and the row break while it reads the cell, and a conditional
  % hides them. Accumulating the rows as tokens at declaration time sidesteps
  % that entirely, and has the better property anyway -- the checklist is
  % literally made of the outcomes, so the two cannot disagree.
  \csgappto{mfa@rows@\mfa@progkey}{%
    #2 & {\footnotesize\color{mgrey}#1}
      & \ratingcell & \ratingcell & \ratingcell & \ratingcell & \ratingcell
    \tabularnewline}%
  \item #2\hfill{\footnotesize\color{mgrey}[\lblFrames~#1]}%
}
\makeatother

\makeatletter
\newenvironment{outcomes}{%
  \setcounter{outcomeno}{0}%
  \csgdef{mfa@rows@\mfa@progkey}{}%
  \begin{tcolorbox}[
    enhanced, breakable, sharp corners,
    colback=mlight, colframe=mblue, boxrule=0pt, leftrule=3pt,
    left=8pt, right=8pt, top=6pt, bottom=6pt,
    fonttitle=\bfseries\small, title={\lblOutcomes}]
  \begin{itemize}[leftmargin=1.3em, itemsep=3pt, topsep=2pt]
}{%
  \end{itemize}
  \end{tcolorbox}
}
\makeatother

% ---- "Can you?" -----------------------------------------------------------
% Generated from the stored outcomes, not written out again by hand. Five
% columns rather than a tick box, because a self-assessment with a midpoint is
% one people actually use and a yes/no one is one they lie to.
\newcommand{\ratingcell}{\textcolor{codeframe}{\rule[-0.15em]{0.8em}{0.8em}}}

\makeatletter
\newcommand{\canyou}{%
  \section*{\lblCanYou}%
  \phantomsection\addcontentsline{toc}{section}{\lblCanYou}%
  \noindent\lblCanYouIntro\par\vspace{8pt}
  \noindent
  \begingroup\small
  \ifcsvoid{mfa@rows@\mfa@progkey}{%
    {\color{mred}\lblNoOutcomes}\par
  }{%
    \begin{tabularx}{\linewidth}{@{}Xcccccc@{}}
    \toprule
    \textbf{\lblOnEntryYouCould} & \textbf{\lblFrames} & 1 & 2 & 3 & 4 & 5 \\
    \midrule
    \csuse{mfa@rows@\mfa@progkey}
    \bottomrule
    \end{tabularx}%
  }%
  \endgroup
  \par\vspace{6pt}\noindent{\footnotesize\color{mgrey}\lblCanYouFooter}\par
}
\makeatother

% ============================================================
%  QUIZ, SUMMARY, EXERCISES
% ============================================================

% ---- Quiz -----------------------------------------------------------------
% Foundation programs only. It is a diagnostic on the way in and the same test
% on the way out, which is what lets one book serve a reader who needs the
% whole program and a reader who needs four frames of it. Each question names
% the frames that teach it, so a wrong answer routes the reader somewhere
% specific rather than back to the beginning.
\newlist{quizlist}{enumerate}{1}
\setlist[quizlist]{label=\arabic*., leftmargin=2.0em, itemsep=5pt, topsep=3pt}

\makeatletter
\newenvironment{quiz}{%
  \section*{\lblQuiz}%
  \phantomsection\addcontentsline{toc}{section}{\lblQuiz}%
  \def\mfa@exkey{Q\arabic{quizlisti}}%
  \noindent\lblQuizIntro\par\vspace{5pt}
  \begin{quizlist}
}{%
  \end{quizlist}
}
\makeatother

% The frames that teach the question just asked. Two forms, because a Quiz
% question that routes to a single frame reading "Frames 22" is the kind of
% small wrongness a reader notices and an author never does.
\newcommand{\teachesat}[1]{\hfill{\footnotesize\color{mgrey}(\lblFrames~#1)}}
\newcommand{\teachesatone}[1]{\hfill{\footnotesize\color{mgrey}(\lblFrame~#1)}}

% ---- Summary --------------------------------------------------------------
% Every item carries the frame number it came from, in square brackets. That
% is the eighth-edition improvement and it is the single cheapest thing in the
% format: it turns a summary into a return index, so a reader who fails one
% line of "Can you?" knows exactly which four frames to reread.
\newenvironment{summarybox}{%
  \section*{\lblSummary}%
  \phantomsection\addcontentsline{toc}{section}{\lblSummary}%
  \begin{tcolorbox}[
    enhanced, breakable, sharp corners,
    colback=mlight, colframe=mblue, boxrule=0pt, leftrule=3pt,
    left=8pt, right=8pt, top=6pt, bottom=6pt]
  \begin{itemize}[leftmargin=1.3em, itemsep=5pt, topsep=2pt]
}{%
  \end{itemize}
  \end{tcolorbox}
}

% \sumitem{12}{text} -- the frame reference is not decoration, it is the point.
\newcommand{\sumitem}[2]{\item #2~{\footnotesize\color{mgrey}[#1]}}

% ---- Answers ---------------------------------------------------------------
%  Answers are authored at the point of use and printed at the back.
%
%  They are carried there by a global macro store rather than by an auxiliary
%  file. The sibling books collect their figure and screenshot manifests with
%  \@starttoc and a custom file extension, which works because those entries
%  are one line of text; it is the wrong mechanism for an answer, because an
%  answer contains displayed maths, and material written to an .aux-style file
%  and read back is expanded at shipout, where fragile commands break in ways
%  whose error messages name neither the answer nor the program.
%
%  \include is sequential within a run, so a macro defined globally while
%  typesetting program F1 is still defined when the back matter is typeset.
%  That is all this needs.
\makeatletter
\newcommand{\mfa@storeanswer}[3]{% #1 program key, #2 item key, #3 body
  \csgdef{mfa@a@#1@#2}{#3}%
  % \listcsxadd, not \listcsgadd: the item must be EXPANDED as it is stored.
  % Storing the token \mfa@exkey means every entry in the list expands to
  % whatever the key happens to be at the time the back matter is typeset,
  % which is the last one.
  \listcsxadd{mfa@akeys@#1}{#2}%
}
% Used inside an exercise item. Prints nothing here; the body reappears at the
% back of the book under this program's heading.
\newcommand{\answerto}[1]{\mfa@storeanswer{\mfa@progkey}{\mfa@exkey}{#1}}
\providecommand{\mfa@exkey}{?}
\makeatother

\newlist{exlist}{enumerate}{1}
\setlist[exlist]{label=\arabic*., leftmargin=2.0em, itemsep=6pt, topsep=3pt}
\newlist{fplist}{enumerate}{1}
\setlist[fplist]{label=\arabic*., leftmargin=2.0em, itemsep=5pt, topsep=3pt}

% These MUST be defined inside \makeatletter. Outside it, \def\mfa@exkey{...}
% parses as \def\mfa @exkey{...} -- a definition of \mfa with a delimited
% parameter text -- which raises no error at all and quietly leaves every
% answer filed under the same key. That failure is invisible until the answers
% section at the back of the book prints one program's answers twice.
\makeatletter
\newenvironment{testexercises}{%
  \section*{\lblTestExercises}%
  \phantomsection\addcontentsline{toc}{section}{\lblTestExercises}%
  \def\mfa@exkey{T\arabic{exlisti}}%
  \noindent\lblTestIntro\par\vspace{5pt}
  \begin{exlist}
}{%
  \end{exlist}
}

\newenvironment{furtherproblems}{%
  \section*{\lblFurther}%
  \phantomsection\addcontentsline{toc}{section}{\lblFurther}%
  \def\mfa@exkey{P\arabic{fplisti}}%
  \noindent\lblFurtherIntro\par\vspace{5pt}
  \begin{fplist}
}{%
  \end{fplist}
}
\makeatother

% ---- The programs register -------------------------------------------------
% \program{Title} opens a program: it is \chapter plus the bookkeeping that
% lets the back matter reproduce this program's answers under its own name.
\makeatletter
\newcommand{\program}[1]{%
  \chapter{#1}%
  \csgdef{mfa@title@\thechapter}{#1}%
  % Expanded, for the same reason as the answer keys above.
  \listxadd{\mfaprograms}{\thechapter}%
}
\newcommand{\mfa@printoneanswer}[2]{% #1 program key, #2 item key
  \item[\textbf{#2}] \csuse{mfa@a@#1@#2}%
}
\newcommand{\mfa@answerprogram}[1]{%
  \subsection*{\lblProgram~#1\quad\textmd{\itshape\csuse{mfa@title@#1}}}%
  \ifcsdef{mfa@akeys@#1}{%
    \begin{list}{}{%
      \setlength{\leftmargin}{3.2em}\setlength{\labelwidth}{2.7em}%
      \setlength{\labelsep}{0.5em}\setlength{\itemsep}{3pt}%
      \setlength{\topsep}{3pt}\setlength{\parsep}{0pt}%
      \renewcommand{\makelabel}[1]{\hfill##1}}
      \forlistcsloop{\mfa@printoneanswer{#1}}{mfa@akeys@#1}
    \end{list}%
  }{%
    \par\noindent{\color{mred}\lblNoAnswers}\par
  }%
}
% The answers body, without a heading of its own, so the appendix that carries
% it can be a numbered appendix like any other and appear in the running heads
% and the ToC in the ordinary way.
\newcommand{\answersbody}{%
  \noindent\lblAnswersIntro\par\vspace{8pt}
  \ifdef{\mfaprograms}{\forlistloop{\mfa@answerprogram}{\mfaprograms}}{%
    \noindent{\color{mred}\lblNoAnswers}\par}%
}
% Unnumbered form, for a draft build that has no appendix wiring yet.
\newcommand{\printanswers}{%
  \chapter*{\lblAnswers}%
  \phantomsection\addcontentsline{toc}{chapter}{\lblAnswers}%
  \markboth{\lblAnswers}{\lblAnswers}%
  \answersbody
}
\makeatother

% ============================================================
%  ADMONITIONS
%  ---------------------------------------------------------
%  A deliberately small, fixed vocabulary. Each one earns its place by doing
%  something no other box does; a seventh would be a smell.
% ============================================================

\newtcolorbox{note}[1][]{
  enhanced, breakable, sharp corners,
  colback=mlight, colframe=mblue, boxrule=0pt, leftrule=3pt,
  left=8pt, right=8pt, top=6pt, bottom=6pt,
  fonttitle=\bfseries\small, title={\lblNote}, #1
}

\newtcolorbox{warning}[1][]{
  enhanced, breakable, sharp corners,
  colback=mlight, colframe=mred, boxrule=0pt, leftrule=3pt,
  left=8pt, right=8pt, top=6pt, bottom=6pt,
  fonttitle=\bfseries\small, title={\lblWarning}, #1
}

% The misconception, stated as a summary after it has been walked into. The
% trap itself belongs in the frames -- elicit the wrong answer, then say
% flatly that it is wrong -- and this box is the thing the reader marks with a
% finger and comes back to.
\newtcolorbox{trapbox}[1][]{
  enhanced, breakable, sharp corners,
  colback=white, colframe=mred, boxrule=0.8pt,
  left=8pt, right=8pt, top=6pt, bottom=6pt,
  fonttitle=\bfseries\small, title={\lblTrap}, #1
}

% Where the mathematics just done shows up in a system the reader has actually
% deployed. This is the box that separates the book from a maths textbook, and
% it is also the one most easily filled with waffle: if it cannot name a
% specific line of a specific system, it should not be written.
\newtcolorbox{aibox}[1][]{
  enhanced, breakable, sharp corners,
  colback=mlight, colframe=mteal, boxrule=0pt, leftrule=3pt,
  left=8pt, right=8pt, top=6pt, bottom=6pt,
  fonttitle=\bfseries\small, title={\lblInAI}, #1
}

% What is being asserted rather than proved, and where to read the proof.
% Stroud's method buys fluency at the cost of rigour, and the honest thing is
% to say so at each point rather than once in an introduction nobody rereads.
\newtcolorbox{rigourbox}[1][]{
  enhanced, breakable, sharp corners,
  colback=white, colframe=mgrey, boxrule=0.6pt,
  left=8pt, right=8pt, top=6pt, bottom=6pt,
  fonttitle=\bfseries\small, title={\lblRigour}, #1
}

% Notation that genuinely differs -- between the Polish and English editions,
% and between disciplines. Matrix calculus alone has two incompatible layout
% conventions, and half the confusion about backpropagation is a transpose
% that came from the other one.
\newtcolorbox{notationbox}[1][]{
  enhanced, breakable, sharp corners,
  colback=mlight, colframe=mamber, boxrule=0pt, leftrule=3pt,
  left=8pt, right=8pt, top=6pt, bottom=6pt,
  fonttitle=\bfseries\small, title={\lblNotation}, #1
}

% The maths analogue of the sibling books' "run this before you trust it".
% A number in this book is either produced by a script under code/ and pulled
% in with \val, or it carries this box. Nothing ships to a reader with one
% still attached.
\newtcolorbox{verifybox}[1][]{
  enhanced, breakable, sharp corners,
  colback=white, colframe=mamber, boxrule=0.8pt,
  left=8pt, right=8pt, top=6pt, bottom=6pt,
  fonttitle=\bfseries\small, title={\lblVerify}, #1
}

\newtcolorbox{exercisebox}[1][]{
  enhanced, breakable, sharp corners,
  colback=white, colframe=mgreen, boxrule=0pt, leftrule=3pt,
  left=8pt, right=8pt, top=6pt, bottom=6pt,
  fonttitle=\bfseries\small, title={\lblExercise}, #1
}

% ============================================================
%  CODE LISTINGS
%  ---------------------------------------------------------
%  Code appears here for one reason only: to check a number. Every listing in
%  this book is runnable and every number it prints is the number in the text.
% ============================================================
\lstdefinelanguage{pythonx}{
  language=Python,
  morekeywords={as, with, assert, match, case},
  morekeywords=[2]{
    float64,float32,float16,bfloat16,int8,uint8,
    np,numpy,array,asarray,frombuffer,nextafter,finfo,iinfo,
    exp,log,log1p,expm1,logaddexp,sum,max,abs,sqrt,mean,var
  },
  sensitive=true
}

\lstdefinestyle{mfacode}{
  basicstyle=\ttfamily\footnotesize,
  keywordstyle=\color{kw}\bfseries,
  keywordstyle=[2]\color{typ},
  stringstyle=\color{str},
  commentstyle=\color{cmt}\itshape,
  numbers=left,
  numberstyle=\ttfamily\tiny\color{mgrey},
  numbersep=8pt,
  backgroundcolor=\color{codebg},
  frame=single,
  rulecolor=\color{codeframe},
  framesep=6pt,
  breaklines=true,
  breakatwhitespace=false,
  postbreak=\mbox{\textcolor{mgrey}{$\hookrightarrow$}\space},
  showstringspaces=false,
  tabsize=4,
  columns=fullflexible,
  keepspaces=true,
  captionpos=b,
  aboveskip=1.2em,
  belowskip=1.2em,
  % Maps the characters most likely to be pasted into a listing by accident.
  %
  % Do NOT add a mapping for U+00A0 (non-breaking space). It looks like the
  % obvious next entry and it makes `listings` abort the run with "Improper
  % alphabetic constant" -- fatal, no PDF, and the message names neither the
  % character nor this line. The ASCII-in-listings convention is the guard
  % against it instead. This trap has now been hit in two of the three books
  % in this series.
  literate={ł}{{\l}}1 {Ł}{{\L}}1 {ą}{{\k a}}1 {ę}{{\k e}}1
           {ś}{{\'s}}1 {ć}{{\'c}}1 {ż}{{\.z}}1 {ź}{{\'z}}1 {ó}{{\'o}}1 {ń}{{\'n}}1
           {Ą}{{\k A}}1 {Ę}{{\k E}}1 {Ś}{{\'S}}1 {Ć}{{\'C}}1
           {Ż}{{\.Z}}1 {Ź}{{\'Z}}1 {Ó}{{\'O}}1 {Ń}{{\'N}}1
           {—}{{-{}-}}2 {–}{{-}}1 {‘}{{'}}1 {’}{{'}}1
           {“}{{"}}1 {”}{{"}}1 {°}{{\textdegree}}1
}

\lstset{style=mfacode}

\lstnewenvironment{python}[1][]
  {\lstset{language=pythonx,style=mfacode,#1}}
  {}

\lstnewenvironment{shellcmd}[1][]
  {\lstset{language=bash,style=mfacode,numbers=none,keywordstyle=\color{mteal}\bfseries,#1}}
  {}

% Terminal transcripts: no line numbers, no keyword colouring. Everything in
% one of these was copied from a real run.
\lstnewenvironment{console}[1][]
  {\lstset{style=mfacode,numbers=none,basicstyle=\ttfamily\scriptsize,
           keywordstyle=\color{black},backgroundcolor=\color{white},#1}}
  {}

\newcommand{\pyfile}[3]{%
  \lstinputlisting[language=pythonx,style=mfacode,caption={#2},label={#3}]{#1}%
}
\newcommand{\pyregion}[4]{%
  \lstinputlisting[
    language=pythonx, style=mfacode,
    linerange={--8<--\ [start:#2]--- 8<--\ [end:#2]},
    includerangemarker=false,
    caption={#3}, label={#4}]{#1}%
}

% Inline literals
\newcommand{\code}[1]{\texttt{\small #1}}
\newcommand{\pkg}[1]{\texttt{\small\color{mblue}#1}}
\newcommand{\api}[1]{\texttt{\small\color{mteal}#1}}

% ============================================================
%  COMPUTED VALUES
%  ---------------------------------------------------------
%  The house rule in the sibling books is "run the listing, or mark it". The
%  analogue here is stronger, because a wrong digit in a maths book is not a
%  broken example, it is a lie the reader will carry.
%
%  Every number that came out of a computation is written by a script under
%  code/ into figures/values/<program>.tex as
%
%      \mfaval{f01.eps64}{2.220446049250313e-16}
%
%  and pulled into the text with \val{f01.eps64}. The script is the source of
%  truth; the book cannot disagree with it, because the book does not contain
%  the digits. A value the scripts no longer produce prints a visible marker
%  and is counted by `make debt`.
%
%  \val also applies the locale's decimal separator, which is how the Polish
%  edition gets its decimal comma without a translator having to remember.
% ============================================================
\IfFileExists{siunitx.sty}{%
  \usepackage{siunitx}%
  \sisetup{
    group-digits = integer,
    group-minimum-digits = 5,
    group-separator = {\,},
    exponent-product = \cdot,
    output-decimal-marker = {\mfadecimalmarker}
  }%
  \newcommand{\mfanum}[1]{\num{##1}}%
}{%
  \newcommand{\mfanum}[1]{##1}%
}

\makeatletter
\newcommand{\mfaval}[2]{\csgdef{mfaval@#1}{#2}\listxadd{\mfavalues}{#1}}
% \num on a non-numeric body is a FATAL siunitx error -- "Invalid number" --
% and under -halt-on-error that is the whole build. A script emitting a word
% where the book expected a number is a plausible mistake.
%
% The check is NOT done here. Deciding in TeX whether a token list is a number
% means parsing it character by character, which is fragile and was got wrong
% once already. The generator knows the answer for free, so it emits \mfaval
% for a number and \mfavaltext for a word, and this end only has to enforce
% which macro reads which.
\newcommand{\val}[1]{%
  \ifcsdef{mfaval@#1}{%
    \ifcsdef{mfavaltext@#1}{%
      \PackageError{mfabook}{Value `#1' is not a number}%
        {\string\val\space passes its body to siunitx, which refuses anything
         that is not a number. Use \string\valtext{#1} instead, or fix the
         script in code/ that emits this key.}%
    }{}%
    \mfanum{\csuse{mfaval@#1}}%
  }{\textcolor{mred}{\textbf{[?\,#1]}}}%
}
% A value that is deliberately a word rather than a number -- a format name, a
% verdict, a unit. Rare, and a visibly different macro so nobody reaches for
% \val and gets an error they cannot place.
\newcommand{\valtext}[1]{%
  \ifcsdef{mfaval@#1}{\csuse{mfaval@#1}}%
    {\textcolor{mred}{\textbf{[?\,#1]}}}%
}
\newcommand{\mfavaltext}[2]{%
  \csgdef{mfaval@#1}{#2}\csgdef{mfavaltext@#1}{}\listxadd{\mfavalues}{#1}%
}

% The raw stored string, for the cases where a number must appear exactly as
% the machine printed it -- inside a console block, or where the digits
% themselves are the point and grouping them would obscure the pattern.
\newcommand{\rawval}[1]{%
  \ifcsdef{mfaval@#1}{\csuse{mfaval@#1}}%
    {\textcolor{mred}{\textbf{[?\,#1]}}}%
}
\makeatother

% figures/values/all.tex is generated by `make numbers` and does nothing but
% \input each program's value file. A build with no values yet still compiles;
% every \val{} in it prints a visible marker instead of a number, which is the
% correct behaviour for a draft and an obvious failure for a release.
\IfFileExists{figures/values/all.tex}{\input{figures/values/all}}{%
  \AtBeginDocument{\PackageWarning{mfabook}{No computed values found. Run `make numbers`.}}}

% ============================================================
%  DIAGRAMS — Mermaid pipeline
%  ---------------------------------------------------------
%  \mermaidfig{key}{caption}{one-line manifest description}
%
%  Source of truth is figures/mermaid/<lang>/<key>.mmd, committed. `make
%  diagrams` renders each to figures/diagrams/<lang>/<key>.pdf, which is build
%  output and is not committed, so a diagram change reviews as a text diff.
%
%  The source is per-language because a diagram with English labels in the
%  Polish edition is exactly the kind of half-translation that makes a
%  bilingual book feel machine-made. CI checks that both languages carry the
%  same set of keys.
%
%  If the rendered PDF is absent the macro draws a placeholder AND typesets the
%  Mermaid source, so a build without mermaid-cli still shows the structure.
% ============================================================
\makeatletter
\newcommand{\listofdiagrams}{%
  \section*{\lblDiagramManifest}%
  \phantomsection\addcontentsline{toc}{section}{\lblDiagramManifest}%
  \@starttoc{dgm}%
}
\makeatother

\newcommand{\mermaidfig}[3]{%
  \addcontentsline{dgm}{subsection}{\protect\numberline{}\texttt{#1.mmd} --- #3}%
  \IfFileExists{figures/diagrams/\booklang/#1.pdf}{%
    \begin{figure}[htbp]
    \centering
    \includegraphics[width=0.95\linewidth,height=0.42\textheight,keepaspectratio]{figures/diagrams/\booklang/#1.pdf}%
    \caption{#2}\label{fig:#1}
    \end{figure}%
  }{%
    % The unrendered case deliberately does NOT float. A Mermaid source
    % typeset verbatim is often taller than a page, and a float cannot break,
    % so wrapping it in `figure` produces an overfull box per diagram and a
    % pile of tcolorbox warnings. In the flow it breaks like any other text.
    \par\medskip
    \begin{tcolorbox}[
      enhanced, breakable, width=\linewidth, sharp corners,
      colback=white, colframe=mgrey, boxrule=0.8pt,
      left=8pt, right=8pt, top=8pt, bottom=8pt]
      {\bfseries\color{mgrey}\small \lblNotRendered}\\[4pt]
      {\ttfamily\scriptsize\color{mgrey} figures/mermaid/\booklang/#1.mmd}\\[6pt]
      \IfFileExists{figures/mermaid/\booklang/#1.mmd}{%
        % ASCII only, like every other listing in this book: the fallback runs
        % the Mermaid source through `listings`, which aborts on a multi-byte
        % character it has no literate mapping for.
        \lstinputlisting[style=mfacode,numbers=none,basicstyle=\ttfamily\scriptsize,
                         backgroundcolor=\color{mlight},frame=none,
                         keywordstyle=\color{black}]{figures/mermaid/\booklang/#1.mmd}%
      }{%
        {\small\color{mred} \lblSourceMissing: \texttt{figures/mermaid/\booklang/#1.mmd}}%
      }%
    \end{tcolorbox}%
    \captionof{figure}{#2}\label{fig:#1}
    \par\medskip
  }%
}

% ============================================================
%  PROGRAM STUBS
%  ---------------------------------------------------------
%  Prints a visible marker so an unwritten program cannot be mistaken for a
%  written one and the page count never flatters the draft. `make debt` counts
%  them, and CI publishes the count on every build.
% ============================================================
\newcommand{\programstub}[1]{%
  \begin{tcolorbox}[
    enhanced, breakable, width=\linewidth, sharp corners,
    colback=mlight, colframe=mred, boxrule=1pt,
    borderline={0.6pt}{3pt}{mred, dashed},
    left=12pt, right=12pt, top=12pt, bottom=12pt]
    {\bfseries\color{mred}\large \lblNotWritten}\\[8pt]
    \small\raggedright #1
  \end{tcolorbox}%
}

% ============================================================
%  MATHEMATICAL OPERATORS
%  ---------------------------------------------------------
%  Language-invariant names live here. The ones that genuinely differ between
%  Polish and English usage -- tan/tg, Var/D^2, gcd/NWD, and the interval
%  brackets -- are defined in lang/en.tex and lang/pl.tex instead, so the
%  notation contract is executed by the build rather than remembered by a
%  translator.
% ============================================================
\DeclareMathOperator*{\argmin}{arg\,min}
\DeclareMathOperator*{\argmax}{arg\,max}
\DeclareMathOperator{\tr}{tr}
\DeclareMathOperator{\rank}{rank}
\DeclareMathOperator{\diag}{diag}
\DeclareMathOperator{\sgn}{sgn}
\DeclareMathOperator{\relu}{ReLU}
\DeclareMathOperator{\softmax}{softmax}
\DeclareMathOperator{\logsumexp}{logsumexp}
\DeclareMathOperator{\fl}{fl}
\DeclareMathOperator{\ulp}{ulp}
\DeclareMathOperator{\round}{round}

% A bare \log is FORBIDDEN in this book, and tools/check_notation.py fails the
% build on one. In Polish textbooks \log means base ten; in machine-learning
% writing it means base e. The same three characters carry opposite meanings to
% the two readerships this book serves, and they collide hardest in exactly the
% programs where it matters most -- entropy in bits and cross-entropy in nats
% are two programs apart. Write \ln for nats and \logb{2} for bits. It costs
% two characters and removes the ambiguity entirely.
\makeatletter
\let\mfa@origlog\log
\renewcommand{\log}{%
  \PackageError{mfabook}{A bare \string\log\space is forbidden in this book}%
    {Write \string\ln\space for natural logarithms or \string\logb{2}\space for
     bits. The base is never left to a convention here; see Appendix B.}%
  \mfa@origlog}
% The one legitimate form: an explicit base.
\newcommand{\logb}[1]{\mathop{\mfa@origlog}\nolimits_{#1}}
% ...and a way to write the forbidden thing when the subject IS the ban, which
% is Appendix B and nowhere else.
\newcommand{\mfalogplain}{\mathop{\mfa@origlog}\nolimits}
\makeatother

\newcommand{\R}{\mathbb{R}}
\newcommand{\N}{\mathbb{N}}
\newcommand{\Z}{\mathbb{Z}}
\newcommand{\Q}{\mathbb{Q}}
\newcommand{\C}{\mathbb{C}}
\newcommand{\Fset}{\mathbb{F}}          % the floating-point numbers, F1's subject
\newcommand{\vect}[1]{\mathbf{#1}}
\newcommand{\mat}[1]{\mathbf{#1}}
\newcommand{\T}{^{\mathsf{T}}}
\newcommand{\norm}[1]{\left\lVert #1 \right\rVert}
\newcommand{\abs}[1]{\left\lvert #1 \right\rvert}
\newcommand{\inner}[2]{\left\langle #1, #2 \right\rangle}
\newcommand{\eps}{\varepsilon}

% ============================================================
%  INDEX
%  ---------------------------------------------------------
%  Two-column, and its entries include \texttt{} names that do not hyphenate.
%  Ragged right keeps multi-word entries off the margin; it cannot help a
%  single token wider than the column, so keep verbatim index entries under
%  about 29 characters and index longer names by concept instead. Both sibling
%  books learned this the same way.
% ============================================================
\apptocmd{\theindex}{\raggedright}{}{%
  \PackageWarning{mfabook}{Could not patch theindex for ragged-right setting}}

\makeindex

% ============================================================
%  PINNED FACTS — single source of truth
%  Change these here; they propagate through both editions.
% ============================================================
\newcommand{\bookedition}{0.1}
% \pinnedon is a user-visible string and lives in lang/<lang>.tex, not here.
% It said "August 2026" on the Polish title page for exactly as long as it
% took somebody to read the Polish title page.
\newcommand{\pyver}{Python 3.13}
\newcommand{\numpyver}{2.3}
; nothing else differs between them at
%  this level. Every user-visible string lives in lang/en.tex or lang/pl.tex,
%  so a label that appears in one edition and not the other is a missing macro
%  rather than a silent divergence.
% ============================================================

% \booklang must be set by the main file. Default to English rather than
% failing, so a stray % ============================================================
%  preamble.tex — styling, environments, macros
%
%  Shared by BOTH editions. main-en.tex and main-pl.tex each define
%  \booklang before \input{preamble}; nothing else differs between them at
%  this level. Every user-visible string lives in lang/en.tex or lang/pl.tex,
%  so a label that appears in one edition and not the other is a missing macro
%  rather than a silent divergence.
% ============================================================

% \booklang must be set by the main file. Default to English rather than
% failing, so a stray \input{preamble} still compiles and says what it did.
\providecommand{\booklang}{en}

\usepackage{etoolbox}   % loaded first: the language switch below needs it

% ---------- Page geometry ----------
% Same 17 x 24 cm block as the other two books in this series, so a reader who
% owns one recognises the second. Frames are short, so the measure matters
% more here than in a prose book: a frame that runs past a page turn defeats
% the cover-the-next-frame discipline the whole method rests on.
\usepackage[
  paperwidth=17cm, paperheight=24cm,
  inner=2.2cm, outer=1.8cm, top=2.2cm, bottom=2.4cm,
  head=23pt, headsep=0.7cm, footskip=1.2cm
]{geometry}

% ---------- Encoding & fonts ----------
\usepackage[T1]{fontenc}
\usepackage[utf8]{inputenc}
\IfFileExists{lmodern.sty}{\usepackage{lmodern}}{}
\IfFileExists{newtxtext.sty}{\usepackage{newtxtext}}{}
% amsmath goes here, not down with the other packages, because newtxmath
% patches amsmath's internals and must be loaded AFTER it. The wrong order does
% not fail on a machine without newtxmath -- which is most CI images -- and
% fails on the author's full TeX Live.
%
% amssymb only when newtxmath is absent: newtxmath supplies the AMS symbols
% itself (its amssymbols option is on by default) and loading both clashes.
\usepackage{amsmath}
\IfFileExists{newtxmath.sty}{\usepackage{newtxmath}}{\usepackage{amssymb}}
\IfFileExists{inconsolata.sty}{\usepackage[varqu,scaled=0.95]{inconsolata}}{}
\IfFileExists{lmodern.sty}{\usepackage{microtype}}{\usepackage[expansion=false]{microtype}}

% ---------- Language ----------
% Loading babel with a language whose .ldf is absent is a *fatal* error, not a
% warning, so both the package and each language file are probed before either
% is requested. This is inherited from the sibling books, where a minimal TeX
% installation without texlive-lang-polish aborted the run with no PDF and an
% error message that named neither babel nor the language.
%
% The Polish edition needs Polish as the main language for hyphenation; the
% English edition needs British. Both editions load the other language too,
% because the glossary appendix is bilingual in both.
\IfFileExists{babel.sty}{%
  \IfFileExists{polish.ldf}{%
    \IfFileExists{british.ldf}{%
      \ifdefstring{\booklang}{pl}%
        {\usepackage[british,main=polish]{babel}}%
        {\usepackage[polish,main=british]{babel}}%
    }{\usepackage[polish]{babel}}%
  }{%
    \IfFileExists{british.ldf}{\usepackage[british]{babel}}{}%
  }%
}{}

% babel-polish makes " an active shorthand. listings reads its body verbatim
% and an active character in a verbatim context is a category-code accident
% waiting to happen, so the shorthand is turned off globally. Polish quotation
% marks are produced with \enquote-style macros from lang/pl.tex instead.
\ifdefstring{\booklang}{pl}{\AtBeginDocument{\shorthandoff{"}}}{}

% ---------- Core packages ----------
\usepackage{graphicx}
\usepackage{xcolor}
\usepackage{booktabs}
\usepackage{tabularx}
\usepackage{longtable}
\usepackage{array}
\usepackage{enumitem}
\IfFileExists{mathtools.sty}{\usepackage{mathtools}}{}
\usepackage{listings}
\usepackage[most]{tcolorbox}
\usepackage{fancyhdr}
\usepackage{titlesec}
\usepackage{caption}
\usepackage{imakeidx}
% csquotes with autostyle follows babel's active language, so the identical
% source \enquote{x} sets 'x' in the English edition and the Polish quotation
% marks in the Polish one. This is not a nicety: babel-polish makes " an active
% shorthand, which is turned off above precisely because an active character
% near a listings body is a category-code accident waiting to happen -- so a
% quotation mark has to come from a macro rather than from a keyboard key.
\IfFileExists{csquotes.sty}{\usepackage[autostyle=true]{csquotes}}{%
  \providecommand{\enquote}[1]{``##1''}}
\usepackage[hidelinks]{hyperref}

% ---------- Palette ----------
% Deliberately the same family as the sibling books, shifted green so the three
% spines are distinguishable on a shelf.
\definecolor{mblue}{HTML}{1F4E79}
\definecolor{mgreen}{HTML}{2E6B4F}
\definecolor{mteal}{HTML}{0E7C7B}
\definecolor{mamber}{HTML}{B26A00}
\definecolor{mred}{HTML}{9B2C2C}
\definecolor{mgrey}{HTML}{5A5A5A}
\definecolor{mlight}{HTML}{F4F6F8}
\definecolor{mfaint}{HTML}{FBFCFD}
\definecolor{codebg}{HTML}{FAFAFA}
\definecolor{codeframe}{HTML}{D8DEE4}
\definecolor{kw}{HTML}{0B5394}
\definecolor{str}{HTML}{9B2C2C}
\definecolor{cmt}{HTML}{4F7A4F}
\definecolor{typ}{HTML}{0E7C7B}

\hypersetup{
  colorlinks=true,
  linkcolor=mblue,
  urlcolor=mteal,
  citecolor=mblue,
  pdfauthor={Konrad Cinkusz}
}

% \ifpl{Polish text}{English text} -- for the handful of places where a string
% is structural rather than content and belongs in the shared files:
% structure.tex's part titles, and the main files' front matter wiring.
% Anything longer than a few words belongs in lang/ or in the edition's own
% source, not here.
% \DeclareRobustCommand, not \newcommand. A fragile \ifpl inside a \part title
% is expanded one level as it is written to the .toc, which lands
% "\ifdefstring {en}{pl}{...}{...}" in the file -- and that fails on the next
% run with "Extra }, or forgotten \endgroup", an error that names neither the
% part nor this macro.
\DeclareRobustCommand{\ifpl}[2]{\ifdefstring{\booklang}{pl}{#1}{#2}}

% A robust macro cannot go into a PDF bookmark string either -- hyperref strips
% \ifdefstring and warns once per part title. Tell it which branch to take
% instead, so the bookmarks carry the right language rather than nothing.
\makeatletter
\pdfstringdefDisableCommands{%
  \ifdefstring{\booklang}{pl}{\let\ifpl\@firstoftwo}{\let\ifpl\@secondoftwo}%
}
\makeatother

% ---------- Language strings ----------
% Every user-visible word the macros below typeset comes from here. Nothing in
% programs/ or appendices/ may hard-code an English or a Polish word inside a
% macro argument that the other edition also uses.
\input{lang/\booklang}

% ============================================================
%  PROGRAM NUMBERING
%  ---------------------------------------------------------
%  Part F programs are F1..F13; the main programs restart at 1. That is
%  Stroud's scheme and it is worth reproducing, because the Foundation part is
%  explicitly optional and numbering it separately is what makes skipping it
%  feel permitted rather than delinquent.
%
%  \theHchapter is hyperref's *destination* name, and it must stay unique
%  across the whole document. Resetting the chapter counter without also
%  changing \theHchapter produces duplicate PDF destinations, a pile of
%  warnings, and bookmarks that jump to the wrong program.
% ============================================================
\newcommand{\foundationnumbering}{%
  \setcounter{chapter}{0}%
  \renewcommand{\thechapter}{F\arabic{chapter}}%
  \renewcommand{\theHchapter}{F\arabic{chapter}}%
}
\newcommand{\mainnumbering}{%
  \setcounter{chapter}{0}%
  \renewcommand{\thechapter}{\arabic{chapter}}%
  \renewcommand{\theHchapter}{M\arabic{chapter}}%
}

% \appendix resets the chapter counter and switches \thechapter to letters, but
% it knows nothing about \theHchapter, so appendix A would claim the same PDF
% destination as main program 1. Patch it once, here.
\apptocmd{\appendix}{\renewcommand{\theHchapter}{A\Alph{chapter}}}{}{%
  \PackageWarning{mfabook}{Could not patch \string\appendix\space for hyperref destinations}}

% A program is a chapter. \chaptername is set from the language file so the
% word above the title is "Program" in both editions (it happens to be the
% same word, which is a small piece of luck).
\AtBeginDocument{\renewcommand{\chaptername}{\lblProgram}}

% ---------- Headings ----------
% \raggedright on the title body is load-bearing, not cosmetic: at \Huge on a
% 13 cm text block a long title that cannot break overflows the margin badly.
\titleformat{\chapter}[display]
  {\normalfont\huge\bfseries\color{mblue}\raggedright}
  {\normalfont\normalsize\color{mgrey}\MakeUppercase{\chaptertitlename}\ \thechapter}
  {8pt}{\Huge\raggedright}
\titlespacing*{\chapter}{0pt}{10pt}{28pt}
\titleformat{\section}{\normalfont\Large\bfseries\color{mblue}}{\thesection}{0.8em}{}
\titleformat{\subsection}{\normalfont\large\bfseries\color{mgrey}}{\thesubsection}{0.7em}{}

% head=23pt in the geometry options above, not \setlength{\headheight}:
% fancyhdr computes the height it needs from the running-head font and the
% rule, wants 22.5 pt at this size, and silently overprints the top of the text
% block if it does not get it. Setting it through geometry keeps the text block
% where geometry put it instead of pushing it down.
\pagestyle{fancy}
\fancyhf{}
\fancyhead[LE]{\small\nouppercase{\leftmark}}
% The recto head carries the FRAME number, not the section. In a book of
% numbered frames "where am I" is a frame question: every Summary
% back-reference, every Quiz route and every cross-reference between programs
% is a frame number, and hunting for one by page is the single most annoying
% thing about reading a programmed text. \theframeno at shipout is the last
% frame begun on the page, which is what a reader flicking through wants.
% Guarded: the front matter and the appendices have no frames, and a head
% reading "Frame 0" over the introduction is worse than no head at all.
\fancyhead[RO]{\ifnum\value{frameno}>0\relax\small\lblFrame~\theframeno\fi}
\fancyfoot[LE,RO]{\small\thepage}
\renewcommand{\headrulewidth}{0.4pt}

% This MUST come after \pagestyle{fancy}: fancyhdr installs its own \chaptermark
% at that point and silently overwrites anything defined earlier. The symptom is
% a running-head change that appears to do nothing at all.
%
% The running head carries the *frame* number as well as the program, because
% in a book of numbered frames "where am I" is a frame question, not a page
% question, and the summary's back-references are frame numbers.
% \@chapapp, not \lblProgram. \appendix switches \@chapapp from \chaptername
% to \appendixname, and babel supplies both in the right language; hard-coding
% the word gives an appendix a running head reading "Program A. Answers" over a
% page whose own chapter head reads "APPENDIX A".
\makeatletter
\renewcommand{\chaptermark}[1]{\markboth{\@chapapp\ \thechapter.\ #1}{}}
\makeatother
\renewcommand{\sectionmark}[1]{\markright{\thesection\ #1}}

% ============================================================
%  FRAMES — the unit the whole method is built from
%  ---------------------------------------------------------
%  A program is a stream of numbered frames. Nearly every frame ends by asking
%  the reader for something, and the answer opens the *next* frame, so that the
%  reader can cover what is below and commit to an answer before seeing it.
%
%  The environment is called `fr` rather than `frame` for two reasons. First,
%  \frame is already defined by LaTeX2e (it draws a box in a picture
%  environment) and an environment named `frame` would silently redefine it.
%  Second, this is typed forty to seventy times per program, so it wants to be
%  short in the source.
%
%    \begin{fr}
%    \ans{$x = 3$}            % the answer to the previous frame, optional
%    Teaching text, ending in a question.
%    \end{fr}
% ============================================================

\newcounter{frameno}[chapter]
\renewcommand{\theframeno}{\arabic{frameno}}

% Frames are referenced constantly -- by the Summary, by the Quiz, by the
% "Can you?" checklist and by other programs. \label inside a frame must
% therefore capture the frame number, not the section number.
\makeatletter
\newcommand{\framelabelfix}{%
  \@namedef{theHframeno}{\theHchapter.\arabic{frameno}}%
}
\makeatother

\newcommand{\framerule}{%
  \par\penalty-200\vspace{9pt}%
  \noindent\textcolor{codeframe}{\rule{\linewidth}{0.5pt}}%
  \par\vspace{5pt}%
}

% The frame number as a filled badge, set in the text flow rather than in the
% margin. A margin number would have to alternate sides under `twoside` and
% would be the first thing to overflow on a 17 cm page; a badge cannot.
\newcommand{\framebadge}{%
  \begingroup
  \setlength{\fboxsep}{2.5pt}%
  \colorbox{mblue}{\textcolor{white}{\bfseries\scriptsize\theframeno}}%
  \endgroup\hspace{0.7em}%
}

\newenvironment{fr}{%
  \framerule
  \refstepcounter{frameno}%
  \nopagebreak
  \noindent\framebadge\ignorespaces
}{\par}

% The answer to the previous frame. Set apart and centred so the eye can find
% the edge to cover, and so a reader who has answered can check in one glance
% without reading on.
\newcommand{\ans}[1]{%
  \par\nobreak\vspace{-4pt}%
  \begin{tcolorbox}[
    enhanced, sharp corners, boxrule=0pt, leftrule=2.5pt,
    colback=mfaint, colframe=mgreen,
    left=8pt, right=8pt, top=4pt, bottom=4pt,
    width=0.86\linewidth, center, halign=center,
    before skip=2pt, after skip=7pt]
    #1
  \end{tcolorbox}%
  \nobreak\noindent\ignorespaces
}

% A multi-paragraph or worked answer, where a centred box would be wrong.
\newenvironment{ansblock}{%
  \par\nopagebreak
  \begin{tcolorbox}[
    enhanced, breakable, sharp corners, boxrule=0pt, leftrule=2.5pt,
    colback=mfaint, colframe=mgreen,
    left=8pt, right=8pt, top=5pt, bottom=5pt]
}{%
  \end{tcolorbox}%
  \nopagebreak\par\vspace{2pt}\noindent\ignorespaces
}

% The row of dots. Stroud's signal that the reader is to supply a word, a
% number or an expression. \blank is inline; \dotline occupies its own line,
% which is what a longer derivation step wants.
% \penalty0 offers TeX a break opportunity immediately before the dots. Without
% it the run of dots is glued to the last word of the question, and a long
% question in the Polish edition -- where the words are longer and hyphenate
% less freely -- runs into the margin.
\newcommand{\blank}[1][4em]{\penalty0\makebox[#1]{\dotfill}}
\newcommand{\dotline}{\par\vspace{2pt}\noindent\makebox[\linewidth]{\dotfill}\par\vspace{2pt}}

% "Now one for you to do." The third rung of the scaffolding gradient, and the
% one most often skipped by authors who think they have written an exercise
% when they have written another worked example.
\newcommand{\yourturn}{%
  \par\vspace{4pt}\noindent
  {\bfseries\color{mgreen}\lblYourTurn}\par\nobreak\vspace{2pt}\noindent\ignorespaces
}

% ============================================================
%  THE PROGRAM SKELETON
%  ---------------------------------------------------------
%  Every program carries the same seven parts, in the same order. They are
%  macros rather than a convention so that a missing one is a build-time
%  absence rather than something a reader notices.
% ============================================================

% ---- Learning outcomes, on entry ----------------------------------------
% Stored as well as printed, because "Can you?" at the end of the program must
% be the same list, one for one. Storing it is what makes that structural
% rather than a promise: the checklist is generated from the outcomes, so the
% two cannot drift.
\newcounter{outcomeno}[chapter]

\makeatletter
% The key under which everything belonging to the current program is stored.
% It must expand to the printed program number (F1, F2, ..., 1, 2, ...), which
% is what makes an answer findable from the program that owns it.
\newcommand{\mfa@progkey}{\thechapter}
\newcommand{\outcome}[2]{% #1 = frames that teach it, #2 = the outcome
  \stepcounter{outcomeno}%
  \csgdef{mfa@out@\mfa@progkey @\theoutcomeno}{#2}%
  \csgdef{mfa@outfr@\mfa@progkey @\theoutcomeno}{#1}%
  % The "Can you?" row is built here rather than looped over later. A \loop
  % inside a tabularx body does not survive alignment scanning: TeX needs to
  % see the & and the row break while it reads the cell, and a conditional
  % hides them. Accumulating the rows as tokens at declaration time sidesteps
  % that entirely, and has the better property anyway -- the checklist is
  % literally made of the outcomes, so the two cannot disagree.
  \csgappto{mfa@rows@\mfa@progkey}{%
    #2 & {\footnotesize\color{mgrey}#1}
      & \ratingcell & \ratingcell & \ratingcell & \ratingcell & \ratingcell
    \tabularnewline}%
  \item #2\hfill{\footnotesize\color{mgrey}[\lblFrames~#1]}%
}
\makeatother

\makeatletter
\newenvironment{outcomes}{%
  \setcounter{outcomeno}{0}%
  \csgdef{mfa@rows@\mfa@progkey}{}%
  \begin{tcolorbox}[
    enhanced, breakable, sharp corners,
    colback=mlight, colframe=mblue, boxrule=0pt, leftrule=3pt,
    left=8pt, right=8pt, top=6pt, bottom=6pt,
    fonttitle=\bfseries\small, title={\lblOutcomes}]
  \begin{itemize}[leftmargin=1.3em, itemsep=3pt, topsep=2pt]
}{%
  \end{itemize}
  \end{tcolorbox}
}
\makeatother

% ---- "Can you?" -----------------------------------------------------------
% Generated from the stored outcomes, not written out again by hand. Five
% columns rather than a tick box, because a self-assessment with a midpoint is
% one people actually use and a yes/no one is one they lie to.
\newcommand{\ratingcell}{\textcolor{codeframe}{\rule[-0.15em]{0.8em}{0.8em}}}

\makeatletter
\newcommand{\canyou}{%
  \section*{\lblCanYou}%
  \phantomsection\addcontentsline{toc}{section}{\lblCanYou}%
  \noindent\lblCanYouIntro\par\vspace{8pt}
  \noindent
  \begingroup\small
  \ifcsvoid{mfa@rows@\mfa@progkey}{%
    {\color{mred}\lblNoOutcomes}\par
  }{%
    \begin{tabularx}{\linewidth}{@{}Xcccccc@{}}
    \toprule
    \textbf{\lblOnEntryYouCould} & \textbf{\lblFrames} & 1 & 2 & 3 & 4 & 5 \\
    \midrule
    \csuse{mfa@rows@\mfa@progkey}
    \bottomrule
    \end{tabularx}%
  }%
  \endgroup
  \par\vspace{6pt}\noindent{\footnotesize\color{mgrey}\lblCanYouFooter}\par
}
\makeatother

% ============================================================
%  QUIZ, SUMMARY, EXERCISES
% ============================================================

% ---- Quiz -----------------------------------------------------------------
% Foundation programs only. It is a diagnostic on the way in and the same test
% on the way out, which is what lets one book serve a reader who needs the
% whole program and a reader who needs four frames of it. Each question names
% the frames that teach it, so a wrong answer routes the reader somewhere
% specific rather than back to the beginning.
\newlist{quizlist}{enumerate}{1}
\setlist[quizlist]{label=\arabic*., leftmargin=2.0em, itemsep=5pt, topsep=3pt}

\makeatletter
\newenvironment{quiz}{%
  \section*{\lblQuiz}%
  \phantomsection\addcontentsline{toc}{section}{\lblQuiz}%
  \def\mfa@exkey{Q\arabic{quizlisti}}%
  \noindent\lblQuizIntro\par\vspace{5pt}
  \begin{quizlist}
}{%
  \end{quizlist}
}
\makeatother

% The frames that teach the question just asked. Two forms, because a Quiz
% question that routes to a single frame reading "Frames 22" is the kind of
% small wrongness a reader notices and an author never does.
\newcommand{\teachesat}[1]{\hfill{\footnotesize\color{mgrey}(\lblFrames~#1)}}
\newcommand{\teachesatone}[1]{\hfill{\footnotesize\color{mgrey}(\lblFrame~#1)}}

% ---- Summary --------------------------------------------------------------
% Every item carries the frame number it came from, in square brackets. That
% is the eighth-edition improvement and it is the single cheapest thing in the
% format: it turns a summary into a return index, so a reader who fails one
% line of "Can you?" knows exactly which four frames to reread.
\newenvironment{summarybox}{%
  \section*{\lblSummary}%
  \phantomsection\addcontentsline{toc}{section}{\lblSummary}%
  \begin{tcolorbox}[
    enhanced, breakable, sharp corners,
    colback=mlight, colframe=mblue, boxrule=0pt, leftrule=3pt,
    left=8pt, right=8pt, top=6pt, bottom=6pt]
  \begin{itemize}[leftmargin=1.3em, itemsep=5pt, topsep=2pt]
}{%
  \end{itemize}
  \end{tcolorbox}
}

% \sumitem{12}{text} -- the frame reference is not decoration, it is the point.
\newcommand{\sumitem}[2]{\item #2~{\footnotesize\color{mgrey}[#1]}}

% ---- Answers ---------------------------------------------------------------
%  Answers are authored at the point of use and printed at the back.
%
%  They are carried there by a global macro store rather than by an auxiliary
%  file. The sibling books collect their figure and screenshot manifests with
%  \@starttoc and a custom file extension, which works because those entries
%  are one line of text; it is the wrong mechanism for an answer, because an
%  answer contains displayed maths, and material written to an .aux-style file
%  and read back is expanded at shipout, where fragile commands break in ways
%  whose error messages name neither the answer nor the program.
%
%  \include is sequential within a run, so a macro defined globally while
%  typesetting program F1 is still defined when the back matter is typeset.
%  That is all this needs.
\makeatletter
\newcommand{\mfa@storeanswer}[3]{% #1 program key, #2 item key, #3 body
  \csgdef{mfa@a@#1@#2}{#3}%
  % \listcsxadd, not \listcsgadd: the item must be EXPANDED as it is stored.
  % Storing the token \mfa@exkey means every entry in the list expands to
  % whatever the key happens to be at the time the back matter is typeset,
  % which is the last one.
  \listcsxadd{mfa@akeys@#1}{#2}%
}
% Used inside an exercise item. Prints nothing here; the body reappears at the
% back of the book under this program's heading.
\newcommand{\answerto}[1]{\mfa@storeanswer{\mfa@progkey}{\mfa@exkey}{#1}}
\providecommand{\mfa@exkey}{?}
\makeatother

\newlist{exlist}{enumerate}{1}
\setlist[exlist]{label=\arabic*., leftmargin=2.0em, itemsep=6pt, topsep=3pt}
\newlist{fplist}{enumerate}{1}
\setlist[fplist]{label=\arabic*., leftmargin=2.0em, itemsep=5pt, topsep=3pt}

% These MUST be defined inside \makeatletter. Outside it, \def\mfa@exkey{...}
% parses as \def\mfa @exkey{...} -- a definition of \mfa with a delimited
% parameter text -- which raises no error at all and quietly leaves every
% answer filed under the same key. That failure is invisible until the answers
% section at the back of the book prints one program's answers twice.
\makeatletter
\newenvironment{testexercises}{%
  \section*{\lblTestExercises}%
  \phantomsection\addcontentsline{toc}{section}{\lblTestExercises}%
  \def\mfa@exkey{T\arabic{exlisti}}%
  \noindent\lblTestIntro\par\vspace{5pt}
  \begin{exlist}
}{%
  \end{exlist}
}

\newenvironment{furtherproblems}{%
  \section*{\lblFurther}%
  \phantomsection\addcontentsline{toc}{section}{\lblFurther}%
  \def\mfa@exkey{P\arabic{fplisti}}%
  \noindent\lblFurtherIntro\par\vspace{5pt}
  \begin{fplist}
}{%
  \end{fplist}
}
\makeatother

% ---- The programs register -------------------------------------------------
% \program{Title} opens a program: it is \chapter plus the bookkeeping that
% lets the back matter reproduce this program's answers under its own name.
\makeatletter
\newcommand{\program}[1]{%
  \chapter{#1}%
  \csgdef{mfa@title@\thechapter}{#1}%
  % Expanded, for the same reason as the answer keys above.
  \listxadd{\mfaprograms}{\thechapter}%
}
\newcommand{\mfa@printoneanswer}[2]{% #1 program key, #2 item key
  \item[\textbf{#2}] \csuse{mfa@a@#1@#2}%
}
\newcommand{\mfa@answerprogram}[1]{%
  \subsection*{\lblProgram~#1\quad\textmd{\itshape\csuse{mfa@title@#1}}}%
  \ifcsdef{mfa@akeys@#1}{%
    \begin{list}{}{%
      \setlength{\leftmargin}{3.2em}\setlength{\labelwidth}{2.7em}%
      \setlength{\labelsep}{0.5em}\setlength{\itemsep}{3pt}%
      \setlength{\topsep}{3pt}\setlength{\parsep}{0pt}%
      \renewcommand{\makelabel}[1]{\hfill##1}}
      \forlistcsloop{\mfa@printoneanswer{#1}}{mfa@akeys@#1}
    \end{list}%
  }{%
    \par\noindent{\color{mred}\lblNoAnswers}\par
  }%
}
% The answers body, without a heading of its own, so the appendix that carries
% it can be a numbered appendix like any other and appear in the running heads
% and the ToC in the ordinary way.
\newcommand{\answersbody}{%
  \noindent\lblAnswersIntro\par\vspace{8pt}
  \ifdef{\mfaprograms}{\forlistloop{\mfa@answerprogram}{\mfaprograms}}{%
    \noindent{\color{mred}\lblNoAnswers}\par}%
}
% Unnumbered form, for a draft build that has no appendix wiring yet.
\newcommand{\printanswers}{%
  \chapter*{\lblAnswers}%
  \phantomsection\addcontentsline{toc}{chapter}{\lblAnswers}%
  \markboth{\lblAnswers}{\lblAnswers}%
  \answersbody
}
\makeatother

% ============================================================
%  ADMONITIONS
%  ---------------------------------------------------------
%  A deliberately small, fixed vocabulary. Each one earns its place by doing
%  something no other box does; a seventh would be a smell.
% ============================================================

\newtcolorbox{note}[1][]{
  enhanced, breakable, sharp corners,
  colback=mlight, colframe=mblue, boxrule=0pt, leftrule=3pt,
  left=8pt, right=8pt, top=6pt, bottom=6pt,
  fonttitle=\bfseries\small, title={\lblNote}, #1
}

\newtcolorbox{warning}[1][]{
  enhanced, breakable, sharp corners,
  colback=mlight, colframe=mred, boxrule=0pt, leftrule=3pt,
  left=8pt, right=8pt, top=6pt, bottom=6pt,
  fonttitle=\bfseries\small, title={\lblWarning}, #1
}

% The misconception, stated as a summary after it has been walked into. The
% trap itself belongs in the frames -- elicit the wrong answer, then say
% flatly that it is wrong -- and this box is the thing the reader marks with a
% finger and comes back to.
\newtcolorbox{trapbox}[1][]{
  enhanced, breakable, sharp corners,
  colback=white, colframe=mred, boxrule=0.8pt,
  left=8pt, right=8pt, top=6pt, bottom=6pt,
  fonttitle=\bfseries\small, title={\lblTrap}, #1
}

% Where the mathematics just done shows up in a system the reader has actually
% deployed. This is the box that separates the book from a maths textbook, and
% it is also the one most easily filled with waffle: if it cannot name a
% specific line of a specific system, it should not be written.
\newtcolorbox{aibox}[1][]{
  enhanced, breakable, sharp corners,
  colback=mlight, colframe=mteal, boxrule=0pt, leftrule=3pt,
  left=8pt, right=8pt, top=6pt, bottom=6pt,
  fonttitle=\bfseries\small, title={\lblInAI}, #1
}

% What is being asserted rather than proved, and where to read the proof.
% Stroud's method buys fluency at the cost of rigour, and the honest thing is
% to say so at each point rather than once in an introduction nobody rereads.
\newtcolorbox{rigourbox}[1][]{
  enhanced, breakable, sharp corners,
  colback=white, colframe=mgrey, boxrule=0.6pt,
  left=8pt, right=8pt, top=6pt, bottom=6pt,
  fonttitle=\bfseries\small, title={\lblRigour}, #1
}

% Notation that genuinely differs -- between the Polish and English editions,
% and between disciplines. Matrix calculus alone has two incompatible layout
% conventions, and half the confusion about backpropagation is a transpose
% that came from the other one.
\newtcolorbox{notationbox}[1][]{
  enhanced, breakable, sharp corners,
  colback=mlight, colframe=mamber, boxrule=0pt, leftrule=3pt,
  left=8pt, right=8pt, top=6pt, bottom=6pt,
  fonttitle=\bfseries\small, title={\lblNotation}, #1
}

% The maths analogue of the sibling books' "run this before you trust it".
% A number in this book is either produced by a script under code/ and pulled
% in with \val, or it carries this box. Nothing ships to a reader with one
% still attached.
\newtcolorbox{verifybox}[1][]{
  enhanced, breakable, sharp corners,
  colback=white, colframe=mamber, boxrule=0.8pt,
  left=8pt, right=8pt, top=6pt, bottom=6pt,
  fonttitle=\bfseries\small, title={\lblVerify}, #1
}

\newtcolorbox{exercisebox}[1][]{
  enhanced, breakable, sharp corners,
  colback=white, colframe=mgreen, boxrule=0pt, leftrule=3pt,
  left=8pt, right=8pt, top=6pt, bottom=6pt,
  fonttitle=\bfseries\small, title={\lblExercise}, #1
}

% ============================================================
%  CODE LISTINGS
%  ---------------------------------------------------------
%  Code appears here for one reason only: to check a number. Every listing in
%  this book is runnable and every number it prints is the number in the text.
% ============================================================
\lstdefinelanguage{pythonx}{
  language=Python,
  morekeywords={as, with, assert, match, case},
  morekeywords=[2]{
    float64,float32,float16,bfloat16,int8,uint8,
    np,numpy,array,asarray,frombuffer,nextafter,finfo,iinfo,
    exp,log,log1p,expm1,logaddexp,sum,max,abs,sqrt,mean,var
  },
  sensitive=true
}

\lstdefinestyle{mfacode}{
  basicstyle=\ttfamily\footnotesize,
  keywordstyle=\color{kw}\bfseries,
  keywordstyle=[2]\color{typ},
  stringstyle=\color{str},
  commentstyle=\color{cmt}\itshape,
  numbers=left,
  numberstyle=\ttfamily\tiny\color{mgrey},
  numbersep=8pt,
  backgroundcolor=\color{codebg},
  frame=single,
  rulecolor=\color{codeframe},
  framesep=6pt,
  breaklines=true,
  breakatwhitespace=false,
  postbreak=\mbox{\textcolor{mgrey}{$\hookrightarrow$}\space},
  showstringspaces=false,
  tabsize=4,
  columns=fullflexible,
  keepspaces=true,
  captionpos=b,
  aboveskip=1.2em,
  belowskip=1.2em,
  % Maps the characters most likely to be pasted into a listing by accident.
  %
  % Do NOT add a mapping for U+00A0 (non-breaking space). It looks like the
  % obvious next entry and it makes `listings` abort the run with "Improper
  % alphabetic constant" -- fatal, no PDF, and the message names neither the
  % character nor this line. The ASCII-in-listings convention is the guard
  % against it instead. This trap has now been hit in two of the three books
  % in this series.
  literate={ł}{{\l}}1 {Ł}{{\L}}1 {ą}{{\k a}}1 {ę}{{\k e}}1
           {ś}{{\'s}}1 {ć}{{\'c}}1 {ż}{{\.z}}1 {ź}{{\'z}}1 {ó}{{\'o}}1 {ń}{{\'n}}1
           {Ą}{{\k A}}1 {Ę}{{\k E}}1 {Ś}{{\'S}}1 {Ć}{{\'C}}1
           {Ż}{{\.Z}}1 {Ź}{{\'Z}}1 {Ó}{{\'O}}1 {Ń}{{\'N}}1
           {—}{{-{}-}}2 {–}{{-}}1 {‘}{{'}}1 {’}{{'}}1
           {“}{{"}}1 {”}{{"}}1 {°}{{\textdegree}}1
}

\lstset{style=mfacode}

\lstnewenvironment{python}[1][]
  {\lstset{language=pythonx,style=mfacode,#1}}
  {}

\lstnewenvironment{shellcmd}[1][]
  {\lstset{language=bash,style=mfacode,numbers=none,keywordstyle=\color{mteal}\bfseries,#1}}
  {}

% Terminal transcripts: no line numbers, no keyword colouring. Everything in
% one of these was copied from a real run.
\lstnewenvironment{console}[1][]
  {\lstset{style=mfacode,numbers=none,basicstyle=\ttfamily\scriptsize,
           keywordstyle=\color{black},backgroundcolor=\color{white},#1}}
  {}

\newcommand{\pyfile}[3]{%
  \lstinputlisting[language=pythonx,style=mfacode,caption={#2},label={#3}]{#1}%
}
\newcommand{\pyregion}[4]{%
  \lstinputlisting[
    language=pythonx, style=mfacode,
    linerange={--8<--\ [start:#2]--- 8<--\ [end:#2]},
    includerangemarker=false,
    caption={#3}, label={#4}]{#1}%
}

% Inline literals
\newcommand{\code}[1]{\texttt{\small #1}}
\newcommand{\pkg}[1]{\texttt{\small\color{mblue}#1}}
\newcommand{\api}[1]{\texttt{\small\color{mteal}#1}}

% ============================================================
%  COMPUTED VALUES
%  ---------------------------------------------------------
%  The house rule in the sibling books is "run the listing, or mark it". The
%  analogue here is stronger, because a wrong digit in a maths book is not a
%  broken example, it is a lie the reader will carry.
%
%  Every number that came out of a computation is written by a script under
%  code/ into figures/values/<program>.tex as
%
%      \mfaval{f01.eps64}{2.220446049250313e-16}
%
%  and pulled into the text with \val{f01.eps64}. The script is the source of
%  truth; the book cannot disagree with it, because the book does not contain
%  the digits. A value the scripts no longer produce prints a visible marker
%  and is counted by `make debt`.
%
%  \val also applies the locale's decimal separator, which is how the Polish
%  edition gets its decimal comma without a translator having to remember.
% ============================================================
\IfFileExists{siunitx.sty}{%
  \usepackage{siunitx}%
  \sisetup{
    group-digits = integer,
    group-minimum-digits = 5,
    group-separator = {\,},
    exponent-product = \cdot,
    output-decimal-marker = {\mfadecimalmarker}
  }%
  \newcommand{\mfanum}[1]{\num{##1}}%
}{%
  \newcommand{\mfanum}[1]{##1}%
}

\makeatletter
\newcommand{\mfaval}[2]{\csgdef{mfaval@#1}{#2}\listxadd{\mfavalues}{#1}}
% \num on a non-numeric body is a FATAL siunitx error -- "Invalid number" --
% and under -halt-on-error that is the whole build. A script emitting a word
% where the book expected a number is a plausible mistake.
%
% The check is NOT done here. Deciding in TeX whether a token list is a number
% means parsing it character by character, which is fragile and was got wrong
% once already. The generator knows the answer for free, so it emits \mfaval
% for a number and \mfavaltext for a word, and this end only has to enforce
% which macro reads which.
\newcommand{\val}[1]{%
  \ifcsdef{mfaval@#1}{%
    \ifcsdef{mfavaltext@#1}{%
      \PackageError{mfabook}{Value `#1' is not a number}%
        {\string\val\space passes its body to siunitx, which refuses anything
         that is not a number. Use \string\valtext{#1} instead, or fix the
         script in code/ that emits this key.}%
    }{}%
    \mfanum{\csuse{mfaval@#1}}%
  }{\textcolor{mred}{\textbf{[?\,#1]}}}%
}
% A value that is deliberately a word rather than a number -- a format name, a
% verdict, a unit. Rare, and a visibly different macro so nobody reaches for
% \val and gets an error they cannot place.
\newcommand{\valtext}[1]{%
  \ifcsdef{mfaval@#1}{\csuse{mfaval@#1}}%
    {\textcolor{mred}{\textbf{[?\,#1]}}}%
}
\newcommand{\mfavaltext}[2]{%
  \csgdef{mfaval@#1}{#2}\csgdef{mfavaltext@#1}{}\listxadd{\mfavalues}{#1}%
}

% The raw stored string, for the cases where a number must appear exactly as
% the machine printed it -- inside a console block, or where the digits
% themselves are the point and grouping them would obscure the pattern.
\newcommand{\rawval}[1]{%
  \ifcsdef{mfaval@#1}{\csuse{mfaval@#1}}%
    {\textcolor{mred}{\textbf{[?\,#1]}}}%
}
\makeatother

% figures/values/all.tex is generated by `make numbers` and does nothing but
% \input each program's value file. A build with no values yet still compiles;
% every \val{} in it prints a visible marker instead of a number, which is the
% correct behaviour for a draft and an obvious failure for a release.
\IfFileExists{figures/values/all.tex}{\input{figures/values/all}}{%
  \AtBeginDocument{\PackageWarning{mfabook}{No computed values found. Run `make numbers`.}}}

% ============================================================
%  DIAGRAMS — Mermaid pipeline
%  ---------------------------------------------------------
%  \mermaidfig{key}{caption}{one-line manifest description}
%
%  Source of truth is figures/mermaid/<lang>/<key>.mmd, committed. `make
%  diagrams` renders each to figures/diagrams/<lang>/<key>.pdf, which is build
%  output and is not committed, so a diagram change reviews as a text diff.
%
%  The source is per-language because a diagram with English labels in the
%  Polish edition is exactly the kind of half-translation that makes a
%  bilingual book feel machine-made. CI checks that both languages carry the
%  same set of keys.
%
%  If the rendered PDF is absent the macro draws a placeholder AND typesets the
%  Mermaid source, so a build without mermaid-cli still shows the structure.
% ============================================================
\makeatletter
\newcommand{\listofdiagrams}{%
  \section*{\lblDiagramManifest}%
  \phantomsection\addcontentsline{toc}{section}{\lblDiagramManifest}%
  \@starttoc{dgm}%
}
\makeatother

\newcommand{\mermaidfig}[3]{%
  \addcontentsline{dgm}{subsection}{\protect\numberline{}\texttt{#1.mmd} --- #3}%
  \IfFileExists{figures/diagrams/\booklang/#1.pdf}{%
    \begin{figure}[htbp]
    \centering
    \includegraphics[width=0.95\linewidth,height=0.42\textheight,keepaspectratio]{figures/diagrams/\booklang/#1.pdf}%
    \caption{#2}\label{fig:#1}
    \end{figure}%
  }{%
    % The unrendered case deliberately does NOT float. A Mermaid source
    % typeset verbatim is often taller than a page, and a float cannot break,
    % so wrapping it in `figure` produces an overfull box per diagram and a
    % pile of tcolorbox warnings. In the flow it breaks like any other text.
    \par\medskip
    \begin{tcolorbox}[
      enhanced, breakable, width=\linewidth, sharp corners,
      colback=white, colframe=mgrey, boxrule=0.8pt,
      left=8pt, right=8pt, top=8pt, bottom=8pt]
      {\bfseries\color{mgrey}\small \lblNotRendered}\\[4pt]
      {\ttfamily\scriptsize\color{mgrey} figures/mermaid/\booklang/#1.mmd}\\[6pt]
      \IfFileExists{figures/mermaid/\booklang/#1.mmd}{%
        % ASCII only, like every other listing in this book: the fallback runs
        % the Mermaid source through `listings`, which aborts on a multi-byte
        % character it has no literate mapping for.
        \lstinputlisting[style=mfacode,numbers=none,basicstyle=\ttfamily\scriptsize,
                         backgroundcolor=\color{mlight},frame=none,
                         keywordstyle=\color{black}]{figures/mermaid/\booklang/#1.mmd}%
      }{%
        {\small\color{mred} \lblSourceMissing: \texttt{figures/mermaid/\booklang/#1.mmd}}%
      }%
    \end{tcolorbox}%
    \captionof{figure}{#2}\label{fig:#1}
    \par\medskip
  }%
}

% ============================================================
%  PROGRAM STUBS
%  ---------------------------------------------------------
%  Prints a visible marker so an unwritten program cannot be mistaken for a
%  written one and the page count never flatters the draft. `make debt` counts
%  them, and CI publishes the count on every build.
% ============================================================
\newcommand{\programstub}[1]{%
  \begin{tcolorbox}[
    enhanced, breakable, width=\linewidth, sharp corners,
    colback=mlight, colframe=mred, boxrule=1pt,
    borderline={0.6pt}{3pt}{mred, dashed},
    left=12pt, right=12pt, top=12pt, bottom=12pt]
    {\bfseries\color{mred}\large \lblNotWritten}\\[8pt]
    \small\raggedright #1
  \end{tcolorbox}%
}

% ============================================================
%  MATHEMATICAL OPERATORS
%  ---------------------------------------------------------
%  Language-invariant names live here. The ones that genuinely differ between
%  Polish and English usage -- tan/tg, Var/D^2, gcd/NWD, and the interval
%  brackets -- are defined in lang/en.tex and lang/pl.tex instead, so the
%  notation contract is executed by the build rather than remembered by a
%  translator.
% ============================================================
\DeclareMathOperator*{\argmin}{arg\,min}
\DeclareMathOperator*{\argmax}{arg\,max}
\DeclareMathOperator{\tr}{tr}
\DeclareMathOperator{\rank}{rank}
\DeclareMathOperator{\diag}{diag}
\DeclareMathOperator{\sgn}{sgn}
\DeclareMathOperator{\relu}{ReLU}
\DeclareMathOperator{\softmax}{softmax}
\DeclareMathOperator{\logsumexp}{logsumexp}
\DeclareMathOperator{\fl}{fl}
\DeclareMathOperator{\ulp}{ulp}
\DeclareMathOperator{\round}{round}

% A bare \log is FORBIDDEN in this book, and tools/check_notation.py fails the
% build on one. In Polish textbooks \log means base ten; in machine-learning
% writing it means base e. The same three characters carry opposite meanings to
% the two readerships this book serves, and they collide hardest in exactly the
% programs where it matters most -- entropy in bits and cross-entropy in nats
% are two programs apart. Write \ln for nats and \logb{2} for bits. It costs
% two characters and removes the ambiguity entirely.
\makeatletter
\let\mfa@origlog\log
\renewcommand{\log}{%
  \PackageError{mfabook}{A bare \string\log\space is forbidden in this book}%
    {Write \string\ln\space for natural logarithms or \string\logb{2}\space for
     bits. The base is never left to a convention here; see Appendix B.}%
  \mfa@origlog}
% The one legitimate form: an explicit base.
\newcommand{\logb}[1]{\mathop{\mfa@origlog}\nolimits_{#1}}
% ...and a way to write the forbidden thing when the subject IS the ban, which
% is Appendix B and nowhere else.
\newcommand{\mfalogplain}{\mathop{\mfa@origlog}\nolimits}
\makeatother

\newcommand{\R}{\mathbb{R}}
\newcommand{\N}{\mathbb{N}}
\newcommand{\Z}{\mathbb{Z}}
\newcommand{\Q}{\mathbb{Q}}
\newcommand{\C}{\mathbb{C}}
\newcommand{\Fset}{\mathbb{F}}          % the floating-point numbers, F1's subject
\newcommand{\vect}[1]{\mathbf{#1}}
\newcommand{\mat}[1]{\mathbf{#1}}
\newcommand{\T}{^{\mathsf{T}}}
\newcommand{\norm}[1]{\left\lVert #1 \right\rVert}
\newcommand{\abs}[1]{\left\lvert #1 \right\rvert}
\newcommand{\inner}[2]{\left\langle #1, #2 \right\rangle}
\newcommand{\eps}{\varepsilon}

% ============================================================
%  INDEX
%  ---------------------------------------------------------
%  Two-column, and its entries include \texttt{} names that do not hyphenate.
%  Ragged right keeps multi-word entries off the margin; it cannot help a
%  single token wider than the column, so keep verbatim index entries under
%  about 29 characters and index longer names by concept instead. Both sibling
%  books learned this the same way.
% ============================================================
\apptocmd{\theindex}{\raggedright}{}{%
  \PackageWarning{mfabook}{Could not patch theindex for ragged-right setting}}

\makeindex

% ============================================================
%  PINNED FACTS — single source of truth
%  Change these here; they propagate through both editions.
% ============================================================
\newcommand{\bookedition}{0.1}
% \pinnedon is a user-visible string and lives in lang/<lang>.tex, not here.
% It said "August 2026" on the Polish title page for exactly as long as it
% took somebody to read the Polish title page.
\newcommand{\pyver}{Python 3.13}
\newcommand{\numpyver}{2.3}
 still compiles and says what it did.
\providecommand{\booklang}{en}

\usepackage{etoolbox}   % loaded first: the language switch below needs it

% ---------- Page geometry ----------
% Same 17 x 24 cm block as the other two books in this series, so a reader who
% owns one recognises the second. Frames are short, so the measure matters
% more here than in a prose book: a frame that runs past a page turn defeats
% the cover-the-next-frame discipline the whole method rests on.
\usepackage[
  paperwidth=17cm, paperheight=24cm,
  inner=2.2cm, outer=1.8cm, top=2.2cm, bottom=2.4cm,
  head=23pt, headsep=0.7cm, footskip=1.2cm
]{geometry}

% ---------- Encoding & fonts ----------
\usepackage[T1]{fontenc}
\usepackage[utf8]{inputenc}
\IfFileExists{lmodern.sty}{\usepackage{lmodern}}{}
\IfFileExists{newtxtext.sty}{\usepackage{newtxtext}}{}
% amsmath goes here, not down with the other packages, because newtxmath
% patches amsmath's internals and must be loaded AFTER it. The wrong order does
% not fail on a machine without newtxmath -- which is most CI images -- and
% fails on the author's full TeX Live.
%
% amssymb only when newtxmath is absent: newtxmath supplies the AMS symbols
% itself (its amssymbols option is on by default) and loading both clashes.
\usepackage{amsmath}
\IfFileExists{newtxmath.sty}{\usepackage{newtxmath}}{\usepackage{amssymb}}
\IfFileExists{inconsolata.sty}{\usepackage[varqu,scaled=0.95]{inconsolata}}{}
\IfFileExists{lmodern.sty}{\usepackage{microtype}}{\usepackage[expansion=false]{microtype}}

% ---------- Language ----------
% Loading babel with a language whose .ldf is absent is a *fatal* error, not a
% warning, so both the package and each language file are probed before either
% is requested. This is inherited from the sibling books, where a minimal TeX
% installation without texlive-lang-polish aborted the run with no PDF and an
% error message that named neither babel nor the language.
%
% The Polish edition needs Polish as the main language for hyphenation; the
% English edition needs British. Both editions load the other language too,
% because the glossary appendix is bilingual in both.
\IfFileExists{babel.sty}{%
  \IfFileExists{polish.ldf}{%
    \IfFileExists{british.ldf}{%
      \ifdefstring{\booklang}{pl}%
        {\usepackage[british,main=polish]{babel}}%
        {\usepackage[polish,main=british]{babel}}%
    }{\usepackage[polish]{babel}}%
  }{%
    \IfFileExists{british.ldf}{\usepackage[british]{babel}}{}%
  }%
}{}

% babel-polish makes " an active shorthand. listings reads its body verbatim
% and an active character in a verbatim context is a category-code accident
% waiting to happen, so the shorthand is turned off globally. Polish quotation
% marks are produced with \enquote-style macros from lang/pl.tex instead.
\ifdefstring{\booklang}{pl}{\AtBeginDocument{\shorthandoff{"}}}{}

% ---------- Core packages ----------
\usepackage{graphicx}
\usepackage{xcolor}
\usepackage{booktabs}
\usepackage{tabularx}
\usepackage{longtable}
\usepackage{array}
\usepackage{enumitem}
\IfFileExists{mathtools.sty}{\usepackage{mathtools}}{}
\usepackage{listings}
\usepackage[most]{tcolorbox}
\usepackage{fancyhdr}
\usepackage{titlesec}
\usepackage{caption}
\usepackage{imakeidx}
% csquotes with autostyle follows babel's active language, so the identical
% source \enquote{x} sets 'x' in the English edition and the Polish quotation
% marks in the Polish one. This is not a nicety: babel-polish makes " an active
% shorthand, which is turned off above precisely because an active character
% near a listings body is a category-code accident waiting to happen -- so a
% quotation mark has to come from a macro rather than from a keyboard key.
\IfFileExists{csquotes.sty}{\usepackage[autostyle=true]{csquotes}}{%
  \providecommand{\enquote}[1]{``##1''}}
\usepackage[hidelinks]{hyperref}

% ---------- Palette ----------
% Deliberately the same family as the sibling books, shifted green so the three
% spines are distinguishable on a shelf.
\definecolor{mblue}{HTML}{1F4E79}
\definecolor{mgreen}{HTML}{2E6B4F}
\definecolor{mteal}{HTML}{0E7C7B}
\definecolor{mamber}{HTML}{B26A00}
\definecolor{mred}{HTML}{9B2C2C}
\definecolor{mgrey}{HTML}{5A5A5A}
\definecolor{mlight}{HTML}{F4F6F8}
\definecolor{mfaint}{HTML}{FBFCFD}
\definecolor{codebg}{HTML}{FAFAFA}
\definecolor{codeframe}{HTML}{D8DEE4}
\definecolor{kw}{HTML}{0B5394}
\definecolor{str}{HTML}{9B2C2C}
\definecolor{cmt}{HTML}{4F7A4F}
\definecolor{typ}{HTML}{0E7C7B}

\hypersetup{
  colorlinks=true,
  linkcolor=mblue,
  urlcolor=mteal,
  citecolor=mblue,
  pdfauthor={Konrad Cinkusz}
}

% \ifpl{Polish text}{English text} -- for the handful of places where a string
% is structural rather than content and belongs in the shared files:
% structure.tex's part titles, and the main files' front matter wiring.
% Anything longer than a few words belongs in lang/ or in the edition's own
% source, not here.
% \DeclareRobustCommand, not \newcommand. A fragile \ifpl inside a \part title
% is expanded one level as it is written to the .toc, which lands
% "\ifdefstring {en}{pl}{...}{...}" in the file -- and that fails on the next
% run with "Extra }, or forgotten \endgroup", an error that names neither the
% part nor this macro.
\DeclareRobustCommand{\ifpl}[2]{\ifdefstring{\booklang}{pl}{#1}{#2}}

% A robust macro cannot go into a PDF bookmark string either -- hyperref strips
% \ifdefstring and warns once per part title. Tell it which branch to take
% instead, so the bookmarks carry the right language rather than nothing.
\makeatletter
\pdfstringdefDisableCommands{%
  \ifdefstring{\booklang}{pl}{\let\ifpl\@firstoftwo}{\let\ifpl\@secondoftwo}%
}
\makeatother

% ---------- Language strings ----------
% Every user-visible word the macros below typeset comes from here. Nothing in
% programs/ or appendices/ may hard-code an English or a Polish word inside a
% macro argument that the other edition also uses.
\input{lang/\booklang}

% ============================================================
%  PROGRAM NUMBERING
%  ---------------------------------------------------------
%  Part F programs are F1..F13; the main programs restart at 1. That is
%  Stroud's scheme and it is worth reproducing, because the Foundation part is
%  explicitly optional and numbering it separately is what makes skipping it
%  feel permitted rather than delinquent.
%
%  \theHchapter is hyperref's *destination* name, and it must stay unique
%  across the whole document. Resetting the chapter counter without also
%  changing \theHchapter produces duplicate PDF destinations, a pile of
%  warnings, and bookmarks that jump to the wrong program.
% ============================================================
\newcommand{\foundationnumbering}{%
  \setcounter{chapter}{0}%
  \renewcommand{\thechapter}{F\arabic{chapter}}%
  \renewcommand{\theHchapter}{F\arabic{chapter}}%
}
\newcommand{\mainnumbering}{%
  \setcounter{chapter}{0}%
  \renewcommand{\thechapter}{\arabic{chapter}}%
  \renewcommand{\theHchapter}{M\arabic{chapter}}%
}

% \appendix resets the chapter counter and switches \thechapter to letters, but
% it knows nothing about \theHchapter, so appendix A would claim the same PDF
% destination as main program 1. Patch it once, here.
\apptocmd{\appendix}{\renewcommand{\theHchapter}{A\Alph{chapter}}}{}{%
  \PackageWarning{mfabook}{Could not patch \string\appendix\space for hyperref destinations}}

% A program is a chapter. \chaptername is set from the language file so the
% word above the title is "Program" in both editions (it happens to be the
% same word, which is a small piece of luck).
\AtBeginDocument{\renewcommand{\chaptername}{\lblProgram}}

% ---------- Headings ----------
% \raggedright on the title body is load-bearing, not cosmetic: at \Huge on a
% 13 cm text block a long title that cannot break overflows the margin badly.
\titleformat{\chapter}[display]
  {\normalfont\huge\bfseries\color{mblue}\raggedright}
  {\normalfont\normalsize\color{mgrey}\MakeUppercase{\chaptertitlename}\ \thechapter}
  {8pt}{\Huge\raggedright}
\titlespacing*{\chapter}{0pt}{10pt}{28pt}
\titleformat{\section}{\normalfont\Large\bfseries\color{mblue}}{\thesection}{0.8em}{}
\titleformat{\subsection}{\normalfont\large\bfseries\color{mgrey}}{\thesubsection}{0.7em}{}

% head=23pt in the geometry options above, not \setlength{\headheight}:
% fancyhdr computes the height it needs from the running-head font and the
% rule, wants 22.5 pt at this size, and silently overprints the top of the text
% block if it does not get it. Setting it through geometry keeps the text block
% where geometry put it instead of pushing it down.
\pagestyle{fancy}
\fancyhf{}
\fancyhead[LE]{\small\nouppercase{\leftmark}}
% The recto head carries the FRAME number, not the section. In a book of
% numbered frames "where am I" is a frame question: every Summary
% back-reference, every Quiz route and every cross-reference between programs
% is a frame number, and hunting for one by page is the single most annoying
% thing about reading a programmed text. \theframeno at shipout is the last
% frame begun on the page, which is what a reader flicking through wants.
% Guarded: the front matter and the appendices have no frames, and a head
% reading "Frame 0" over the introduction is worse than no head at all.
\fancyhead[RO]{\ifnum\value{frameno}>0\relax\small\lblFrame~\theframeno\fi}
\fancyfoot[LE,RO]{\small\thepage}
\renewcommand{\headrulewidth}{0.4pt}

% This MUST come after \pagestyle{fancy}: fancyhdr installs its own \chaptermark
% at that point and silently overwrites anything defined earlier. The symptom is
% a running-head change that appears to do nothing at all.
%
% The running head carries the *frame* number as well as the program, because
% in a book of numbered frames "where am I" is a frame question, not a page
% question, and the summary's back-references are frame numbers.
% \@chapapp, not \lblProgram. \appendix switches \@chapapp from \chaptername
% to \appendixname, and babel supplies both in the right language; hard-coding
% the word gives an appendix a running head reading "Program A. Answers" over a
% page whose own chapter head reads "APPENDIX A".
\makeatletter
\renewcommand{\chaptermark}[1]{\markboth{\@chapapp\ \thechapter.\ #1}{}}
\makeatother
\renewcommand{\sectionmark}[1]{\markright{\thesection\ #1}}

% ============================================================
%  FRAMES — the unit the whole method is built from
%  ---------------------------------------------------------
%  A program is a stream of numbered frames. Nearly every frame ends by asking
%  the reader for something, and the answer opens the *next* frame, so that the
%  reader can cover what is below and commit to an answer before seeing it.
%
%  The environment is called `fr` rather than `frame` for two reasons. First,
%  \frame is already defined by LaTeX2e (it draws a box in a picture
%  environment) and an environment named `frame` would silently redefine it.
%  Second, this is typed forty to seventy times per program, so it wants to be
%  short in the source.
%
%    \begin{fr}
%    \ans{$x = 3$}            % the answer to the previous frame, optional
%    Teaching text, ending in a question.
%    \end{fr}
% ============================================================

\newcounter{frameno}[chapter]
\renewcommand{\theframeno}{\arabic{frameno}}

% Frames are referenced constantly -- by the Summary, by the Quiz, by the
% "Can you?" checklist and by other programs. \label inside a frame must
% therefore capture the frame number, not the section number.
\makeatletter
\newcommand{\framelabelfix}{%
  \@namedef{theHframeno}{\theHchapter.\arabic{frameno}}%
}
\makeatother

\newcommand{\framerule}{%
  \par\penalty-200\vspace{9pt}%
  \noindent\textcolor{codeframe}{\rule{\linewidth}{0.5pt}}%
  \par\vspace{5pt}%
}

% The frame number as a filled badge, set in the text flow rather than in the
% margin. A margin number would have to alternate sides under `twoside` and
% would be the first thing to overflow on a 17 cm page; a badge cannot.
\newcommand{\framebadge}{%
  \begingroup
  \setlength{\fboxsep}{2.5pt}%
  \colorbox{mblue}{\textcolor{white}{\bfseries\scriptsize\theframeno}}%
  \endgroup\hspace{0.7em}%
}

\newenvironment{fr}{%
  \framerule
  \refstepcounter{frameno}%
  \nopagebreak
  \noindent\framebadge\ignorespaces
}{\par}

% The answer to the previous frame. Set apart and centred so the eye can find
% the edge to cover, and so a reader who has answered can check in one glance
% without reading on.
\newcommand{\ans}[1]{%
  \par\nobreak\vspace{-4pt}%
  \begin{tcolorbox}[
    enhanced, sharp corners, boxrule=0pt, leftrule=2.5pt,
    colback=mfaint, colframe=mgreen,
    left=8pt, right=8pt, top=4pt, bottom=4pt,
    width=0.86\linewidth, center, halign=center,
    before skip=2pt, after skip=7pt]
    #1
  \end{tcolorbox}%
  \nobreak\noindent\ignorespaces
}

% A multi-paragraph or worked answer, where a centred box would be wrong.
\newenvironment{ansblock}{%
  \par\nopagebreak
  \begin{tcolorbox}[
    enhanced, breakable, sharp corners, boxrule=0pt, leftrule=2.5pt,
    colback=mfaint, colframe=mgreen,
    left=8pt, right=8pt, top=5pt, bottom=5pt]
}{%
  \end{tcolorbox}%
  \nopagebreak\par\vspace{2pt}\noindent\ignorespaces
}

% The row of dots. Stroud's signal that the reader is to supply a word, a
% number or an expression. \blank is inline; \dotline occupies its own line,
% which is what a longer derivation step wants.
% \penalty0 offers TeX a break opportunity immediately before the dots. Without
% it the run of dots is glued to the last word of the question, and a long
% question in the Polish edition -- where the words are longer and hyphenate
% less freely -- runs into the margin.
\newcommand{\blank}[1][4em]{\penalty0\makebox[#1]{\dotfill}}
\newcommand{\dotline}{\par\vspace{2pt}\noindent\makebox[\linewidth]{\dotfill}\par\vspace{2pt}}

% "Now one for you to do." The third rung of the scaffolding gradient, and the
% one most often skipped by authors who think they have written an exercise
% when they have written another worked example.
\newcommand{\yourturn}{%
  \par\vspace{4pt}\noindent
  {\bfseries\color{mgreen}\lblYourTurn}\par\nobreak\vspace{2pt}\noindent\ignorespaces
}

% ============================================================
%  THE PROGRAM SKELETON
%  ---------------------------------------------------------
%  Every program carries the same seven parts, in the same order. They are
%  macros rather than a convention so that a missing one is a build-time
%  absence rather than something a reader notices.
% ============================================================

% ---- Learning outcomes, on entry ----------------------------------------
% Stored as well as printed, because "Can you?" at the end of the program must
% be the same list, one for one. Storing it is what makes that structural
% rather than a promise: the checklist is generated from the outcomes, so the
% two cannot drift.
\newcounter{outcomeno}[chapter]

\makeatletter
% The key under which everything belonging to the current program is stored.
% It must expand to the printed program number (F1, F2, ..., 1, 2, ...), which
% is what makes an answer findable from the program that owns it.
\newcommand{\mfa@progkey}{\thechapter}
\newcommand{\outcome}[2]{% #1 = frames that teach it, #2 = the outcome
  \stepcounter{outcomeno}%
  \csgdef{mfa@out@\mfa@progkey @\theoutcomeno}{#2}%
  \csgdef{mfa@outfr@\mfa@progkey @\theoutcomeno}{#1}%
  % The "Can you?" row is built here rather than looped over later. A \loop
  % inside a tabularx body does not survive alignment scanning: TeX needs to
  % see the & and the row break while it reads the cell, and a conditional
  % hides them. Accumulating the rows as tokens at declaration time sidesteps
  % that entirely, and has the better property anyway -- the checklist is
  % literally made of the outcomes, so the two cannot disagree.
  \csgappto{mfa@rows@\mfa@progkey}{%
    #2 & {\footnotesize\color{mgrey}#1}
      & \ratingcell & \ratingcell & \ratingcell & \ratingcell & \ratingcell
    \tabularnewline}%
  \item #2\hfill{\footnotesize\color{mgrey}[\lblFrames~#1]}%
}
\makeatother

\makeatletter
\newenvironment{outcomes}{%
  \setcounter{outcomeno}{0}%
  \csgdef{mfa@rows@\mfa@progkey}{}%
  \begin{tcolorbox}[
    enhanced, breakable, sharp corners,
    colback=mlight, colframe=mblue, boxrule=0pt, leftrule=3pt,
    left=8pt, right=8pt, top=6pt, bottom=6pt,
    fonttitle=\bfseries\small, title={\lblOutcomes}]
  \begin{itemize}[leftmargin=1.3em, itemsep=3pt, topsep=2pt]
}{%
  \end{itemize}
  \end{tcolorbox}
}
\makeatother

% ---- "Can you?" -----------------------------------------------------------
% Generated from the stored outcomes, not written out again by hand. Five
% columns rather than a tick box, because a self-assessment with a midpoint is
% one people actually use and a yes/no one is one they lie to.
\newcommand{\ratingcell}{\textcolor{codeframe}{\rule[-0.15em]{0.8em}{0.8em}}}

\makeatletter
\newcommand{\canyou}{%
  \section*{\lblCanYou}%
  \phantomsection\addcontentsline{toc}{section}{\lblCanYou}%
  \noindent\lblCanYouIntro\par\vspace{8pt}
  \noindent
  \begingroup\small
  \ifcsvoid{mfa@rows@\mfa@progkey}{%
    {\color{mred}\lblNoOutcomes}\par
  }{%
    \begin{tabularx}{\linewidth}{@{}Xcccccc@{}}
    \toprule
    \textbf{\lblOnEntryYouCould} & \textbf{\lblFrames} & 1 & 2 & 3 & 4 & 5 \\
    \midrule
    \csuse{mfa@rows@\mfa@progkey}
    \bottomrule
    \end{tabularx}%
  }%
  \endgroup
  \par\vspace{6pt}\noindent{\footnotesize\color{mgrey}\lblCanYouFooter}\par
}
\makeatother

% ============================================================
%  QUIZ, SUMMARY, EXERCISES
% ============================================================

% ---- Quiz -----------------------------------------------------------------
% Foundation programs only. It is a diagnostic on the way in and the same test
% on the way out, which is what lets one book serve a reader who needs the
% whole program and a reader who needs four frames of it. Each question names
% the frames that teach it, so a wrong answer routes the reader somewhere
% specific rather than back to the beginning.
\newlist{quizlist}{enumerate}{1}
\setlist[quizlist]{label=\arabic*., leftmargin=2.0em, itemsep=5pt, topsep=3pt}

\makeatletter
\newenvironment{quiz}{%
  \section*{\lblQuiz}%
  \phantomsection\addcontentsline{toc}{section}{\lblQuiz}%
  \def\mfa@exkey{Q\arabic{quizlisti}}%
  \noindent\lblQuizIntro\par\vspace{5pt}
  \begin{quizlist}
}{%
  \end{quizlist}
}
\makeatother

% The frames that teach the question just asked. Two forms, because a Quiz
% question that routes to a single frame reading "Frames 22" is the kind of
% small wrongness a reader notices and an author never does.
\newcommand{\teachesat}[1]{\hfill{\footnotesize\color{mgrey}(\lblFrames~#1)}}
\newcommand{\teachesatone}[1]{\hfill{\footnotesize\color{mgrey}(\lblFrame~#1)}}

% ---- Summary --------------------------------------------------------------
% Every item carries the frame number it came from, in square brackets. That
% is the eighth-edition improvement and it is the single cheapest thing in the
% format: it turns a summary into a return index, so a reader who fails one
% line of "Can you?" knows exactly which four frames to reread.
\newenvironment{summarybox}{%
  \section*{\lblSummary}%
  \phantomsection\addcontentsline{toc}{section}{\lblSummary}%
  \begin{tcolorbox}[
    enhanced, breakable, sharp corners,
    colback=mlight, colframe=mblue, boxrule=0pt, leftrule=3pt,
    left=8pt, right=8pt, top=6pt, bottom=6pt]
  \begin{itemize}[leftmargin=1.3em, itemsep=5pt, topsep=2pt]
}{%
  \end{itemize}
  \end{tcolorbox}
}

% \sumitem{12}{text} -- the frame reference is not decoration, it is the point.
\newcommand{\sumitem}[2]{\item #2~{\footnotesize\color{mgrey}[#1]}}

% ---- Answers ---------------------------------------------------------------
%  Answers are authored at the point of use and printed at the back.
%
%  They are carried there by a global macro store rather than by an auxiliary
%  file. The sibling books collect their figure and screenshot manifests with
%  \@starttoc and a custom file extension, which works because those entries
%  are one line of text; it is the wrong mechanism for an answer, because an
%  answer contains displayed maths, and material written to an .aux-style file
%  and read back is expanded at shipout, where fragile commands break in ways
%  whose error messages name neither the answer nor the program.
%
%  \include is sequential within a run, so a macro defined globally while
%  typesetting program F1 is still defined when the back matter is typeset.
%  That is all this needs.
\makeatletter
\newcommand{\mfa@storeanswer}[3]{% #1 program key, #2 item key, #3 body
  \csgdef{mfa@a@#1@#2}{#3}%
  % \listcsxadd, not \listcsgadd: the item must be EXPANDED as it is stored.
  % Storing the token \mfa@exkey means every entry in the list expands to
  % whatever the key happens to be at the time the back matter is typeset,
  % which is the last one.
  \listcsxadd{mfa@akeys@#1}{#2}%
}
% Used inside an exercise item. Prints nothing here; the body reappears at the
% back of the book under this program's heading.
\newcommand{\answerto}[1]{\mfa@storeanswer{\mfa@progkey}{\mfa@exkey}{#1}}
\providecommand{\mfa@exkey}{?}
\makeatother

\newlist{exlist}{enumerate}{1}
\setlist[exlist]{label=\arabic*., leftmargin=2.0em, itemsep=6pt, topsep=3pt}
\newlist{fplist}{enumerate}{1}
\setlist[fplist]{label=\arabic*., leftmargin=2.0em, itemsep=5pt, topsep=3pt}

% These MUST be defined inside \makeatletter. Outside it, \def\mfa@exkey{...}
% parses as \def\mfa @exkey{...} -- a definition of \mfa with a delimited
% parameter text -- which raises no error at all and quietly leaves every
% answer filed under the same key. That failure is invisible until the answers
% section at the back of the book prints one program's answers twice.
\makeatletter
\newenvironment{testexercises}{%
  \section*{\lblTestExercises}%
  \phantomsection\addcontentsline{toc}{section}{\lblTestExercises}%
  \def\mfa@exkey{T\arabic{exlisti}}%
  \noindent\lblTestIntro\par\vspace{5pt}
  \begin{exlist}
}{%
  \end{exlist}
}

\newenvironment{furtherproblems}{%
  \section*{\lblFurther}%
  \phantomsection\addcontentsline{toc}{section}{\lblFurther}%
  \def\mfa@exkey{P\arabic{fplisti}}%
  \noindent\lblFurtherIntro\par\vspace{5pt}
  \begin{fplist}
}{%
  \end{fplist}
}
\makeatother

% ---- The programs register -------------------------------------------------
% \program{Title} opens a program: it is \chapter plus the bookkeeping that
% lets the back matter reproduce this program's answers under its own name.
\makeatletter
\newcommand{\program}[1]{%
  \chapter{#1}%
  \csgdef{mfa@title@\thechapter}{#1}%
  % Expanded, for the same reason as the answer keys above.
  \listxadd{\mfaprograms}{\thechapter}%
}
\newcommand{\mfa@printoneanswer}[2]{% #1 program key, #2 item key
  \item[\textbf{#2}] \csuse{mfa@a@#1@#2}%
}
\newcommand{\mfa@answerprogram}[1]{%
  \subsection*{\lblProgram~#1\quad\textmd{\itshape\csuse{mfa@title@#1}}}%
  \ifcsdef{mfa@akeys@#1}{%
    \begin{list}{}{%
      \setlength{\leftmargin}{3.2em}\setlength{\labelwidth}{2.7em}%
      \setlength{\labelsep}{0.5em}\setlength{\itemsep}{3pt}%
      \setlength{\topsep}{3pt}\setlength{\parsep}{0pt}%
      \renewcommand{\makelabel}[1]{\hfill##1}}
      \forlistcsloop{\mfa@printoneanswer{#1}}{mfa@akeys@#1}
    \end{list}%
  }{%
    \par\noindent{\color{mred}\lblNoAnswers}\par
  }%
}
% The answers body, without a heading of its own, so the appendix that carries
% it can be a numbered appendix like any other and appear in the running heads
% and the ToC in the ordinary way.
\newcommand{\answersbody}{%
  \noindent\lblAnswersIntro\par\vspace{8pt}
  \ifdef{\mfaprograms}{\forlistloop{\mfa@answerprogram}{\mfaprograms}}{%
    \noindent{\color{mred}\lblNoAnswers}\par}%
}
% Unnumbered form, for a draft build that has no appendix wiring yet.
\newcommand{\printanswers}{%
  \chapter*{\lblAnswers}%
  \phantomsection\addcontentsline{toc}{chapter}{\lblAnswers}%
  \markboth{\lblAnswers}{\lblAnswers}%
  \answersbody
}
\makeatother

% ============================================================
%  ADMONITIONS
%  ---------------------------------------------------------
%  A deliberately small, fixed vocabulary. Each one earns its place by doing
%  something no other box does; a seventh would be a smell.
% ============================================================

\newtcolorbox{note}[1][]{
  enhanced, breakable, sharp corners,
  colback=mlight, colframe=mblue, boxrule=0pt, leftrule=3pt,
  left=8pt, right=8pt, top=6pt, bottom=6pt,
  fonttitle=\bfseries\small, title={\lblNote}, #1
}

\newtcolorbox{warning}[1][]{
  enhanced, breakable, sharp corners,
  colback=mlight, colframe=mred, boxrule=0pt, leftrule=3pt,
  left=8pt, right=8pt, top=6pt, bottom=6pt,
  fonttitle=\bfseries\small, title={\lblWarning}, #1
}

% The misconception, stated as a summary after it has been walked into. The
% trap itself belongs in the frames -- elicit the wrong answer, then say
% flatly that it is wrong -- and this box is the thing the reader marks with a
% finger and comes back to.
\newtcolorbox{trapbox}[1][]{
  enhanced, breakable, sharp corners,
  colback=white, colframe=mred, boxrule=0.8pt,
  left=8pt, right=8pt, top=6pt, bottom=6pt,
  fonttitle=\bfseries\small, title={\lblTrap}, #1
}

% Where the mathematics just done shows up in a system the reader has actually
% deployed. This is the box that separates the book from a maths textbook, and
% it is also the one most easily filled with waffle: if it cannot name a
% specific line of a specific system, it should not be written.
\newtcolorbox{aibox}[1][]{
  enhanced, breakable, sharp corners,
  colback=mlight, colframe=mteal, boxrule=0pt, leftrule=3pt,
  left=8pt, right=8pt, top=6pt, bottom=6pt,
  fonttitle=\bfseries\small, title={\lblInAI}, #1
}

% What is being asserted rather than proved, and where to read the proof.
% Stroud's method buys fluency at the cost of rigour, and the honest thing is
% to say so at each point rather than once in an introduction nobody rereads.
\newtcolorbox{rigourbox}[1][]{
  enhanced, breakable, sharp corners,
  colback=white, colframe=mgrey, boxrule=0.6pt,
  left=8pt, right=8pt, top=6pt, bottom=6pt,
  fonttitle=\bfseries\small, title={\lblRigour}, #1
}

% Notation that genuinely differs -- between the Polish and English editions,
% and between disciplines. Matrix calculus alone has two incompatible layout
% conventions, and half the confusion about backpropagation is a transpose
% that came from the other one.
\newtcolorbox{notationbox}[1][]{
  enhanced, breakable, sharp corners,
  colback=mlight, colframe=mamber, boxrule=0pt, leftrule=3pt,
  left=8pt, right=8pt, top=6pt, bottom=6pt,
  fonttitle=\bfseries\small, title={\lblNotation}, #1
}

% The maths analogue of the sibling books' "run this before you trust it".
% A number in this book is either produced by a script under code/ and pulled
% in with \val, or it carries this box. Nothing ships to a reader with one
% still attached.
\newtcolorbox{verifybox}[1][]{
  enhanced, breakable, sharp corners,
  colback=white, colframe=mamber, boxrule=0.8pt,
  left=8pt, right=8pt, top=6pt, bottom=6pt,
  fonttitle=\bfseries\small, title={\lblVerify}, #1
}

\newtcolorbox{exercisebox}[1][]{
  enhanced, breakable, sharp corners,
  colback=white, colframe=mgreen, boxrule=0pt, leftrule=3pt,
  left=8pt, right=8pt, top=6pt, bottom=6pt,
  fonttitle=\bfseries\small, title={\lblExercise}, #1
}

% ============================================================
%  CODE LISTINGS
%  ---------------------------------------------------------
%  Code appears here for one reason only: to check a number. Every listing in
%  this book is runnable and every number it prints is the number in the text.
% ============================================================
\lstdefinelanguage{pythonx}{
  language=Python,
  morekeywords={as, with, assert, match, case},
  morekeywords=[2]{
    float64,float32,float16,bfloat16,int8,uint8,
    np,numpy,array,asarray,frombuffer,nextafter,finfo,iinfo,
    exp,log,log1p,expm1,logaddexp,sum,max,abs,sqrt,mean,var
  },
  sensitive=true
}

\lstdefinestyle{mfacode}{
  basicstyle=\ttfamily\footnotesize,
  keywordstyle=\color{kw}\bfseries,
  keywordstyle=[2]\color{typ},
  stringstyle=\color{str},
  commentstyle=\color{cmt}\itshape,
  numbers=left,
  numberstyle=\ttfamily\tiny\color{mgrey},
  numbersep=8pt,
  backgroundcolor=\color{codebg},
  frame=single,
  rulecolor=\color{codeframe},
  framesep=6pt,
  breaklines=true,
  breakatwhitespace=false,
  postbreak=\mbox{\textcolor{mgrey}{$\hookrightarrow$}\space},
  showstringspaces=false,
  tabsize=4,
  columns=fullflexible,
  keepspaces=true,
  captionpos=b,
  aboveskip=1.2em,
  belowskip=1.2em,
  % Maps the characters most likely to be pasted into a listing by accident.
  %
  % Do NOT add a mapping for U+00A0 (non-breaking space). It looks like the
  % obvious next entry and it makes `listings` abort the run with "Improper
  % alphabetic constant" -- fatal, no PDF, and the message names neither the
  % character nor this line. The ASCII-in-listings convention is the guard
  % against it instead. This trap has now been hit in two of the three books
  % in this series.
  literate={ł}{{\l}}1 {Ł}{{\L}}1 {ą}{{\k a}}1 {ę}{{\k e}}1
           {ś}{{\'s}}1 {ć}{{\'c}}1 {ż}{{\.z}}1 {ź}{{\'z}}1 {ó}{{\'o}}1 {ń}{{\'n}}1
           {Ą}{{\k A}}1 {Ę}{{\k E}}1 {Ś}{{\'S}}1 {Ć}{{\'C}}1
           {Ż}{{\.Z}}1 {Ź}{{\'Z}}1 {Ó}{{\'O}}1 {Ń}{{\'N}}1
           {—}{{-{}-}}2 {–}{{-}}1 {‘}{{'}}1 {’}{{'}}1
           {“}{{"}}1 {”}{{"}}1 {°}{{\textdegree}}1
}

\lstset{style=mfacode}

\lstnewenvironment{python}[1][]
  {\lstset{language=pythonx,style=mfacode,#1}}
  {}

\lstnewenvironment{shellcmd}[1][]
  {\lstset{language=bash,style=mfacode,numbers=none,keywordstyle=\color{mteal}\bfseries,#1}}
  {}

% Terminal transcripts: no line numbers, no keyword colouring. Everything in
% one of these was copied from a real run.
\lstnewenvironment{console}[1][]
  {\lstset{style=mfacode,numbers=none,basicstyle=\ttfamily\scriptsize,
           keywordstyle=\color{black},backgroundcolor=\color{white},#1}}
  {}

\newcommand{\pyfile}[3]{%
  \lstinputlisting[language=pythonx,style=mfacode,caption={#2},label={#3}]{#1}%
}
\newcommand{\pyregion}[4]{%
  \lstinputlisting[
    language=pythonx, style=mfacode,
    linerange={--8<--\ [start:#2]--- 8<--\ [end:#2]},
    includerangemarker=false,
    caption={#3}, label={#4}]{#1}%
}

% Inline literals
\newcommand{\code}[1]{\texttt{\small #1}}
\newcommand{\pkg}[1]{\texttt{\small\color{mblue}#1}}
\newcommand{\api}[1]{\texttt{\small\color{mteal}#1}}

% ============================================================
%  COMPUTED VALUES
%  ---------------------------------------------------------
%  The house rule in the sibling books is "run the listing, or mark it". The
%  analogue here is stronger, because a wrong digit in a maths book is not a
%  broken example, it is a lie the reader will carry.
%
%  Every number that came out of a computation is written by a script under
%  code/ into figures/values/<program>.tex as
%
%      \mfaval{f01.eps64}{2.220446049250313e-16}
%
%  and pulled into the text with \val{f01.eps64}. The script is the source of
%  truth; the book cannot disagree with it, because the book does not contain
%  the digits. A value the scripts no longer produce prints a visible marker
%  and is counted by `make debt`.
%
%  \val also applies the locale's decimal separator, which is how the Polish
%  edition gets its decimal comma without a translator having to remember.
% ============================================================
\IfFileExists{siunitx.sty}{%
  \usepackage{siunitx}%
  \sisetup{
    group-digits = integer,
    group-minimum-digits = 5,
    group-separator = {\,},
    exponent-product = \cdot,
    output-decimal-marker = {\mfadecimalmarker}
  }%
  \newcommand{\mfanum}[1]{\num{##1}}%
}{%
  \newcommand{\mfanum}[1]{##1}%
}

\makeatletter
\newcommand{\mfaval}[2]{\csgdef{mfaval@#1}{#2}\listxadd{\mfavalues}{#1}}
% \num on a non-numeric body is a FATAL siunitx error -- "Invalid number" --
% and under -halt-on-error that is the whole build. A script emitting a word
% where the book expected a number is a plausible mistake.
%
% The check is NOT done here. Deciding in TeX whether a token list is a number
% means parsing it character by character, which is fragile and was got wrong
% once already. The generator knows the answer for free, so it emits \mfaval
% for a number and \mfavaltext for a word, and this end only has to enforce
% which macro reads which.
\newcommand{\val}[1]{%
  \ifcsdef{mfaval@#1}{%
    \ifcsdef{mfavaltext@#1}{%
      \PackageError{mfabook}{Value `#1' is not a number}%
        {\string\val\space passes its body to siunitx, which refuses anything
         that is not a number. Use \string\valtext{#1} instead, or fix the
         script in code/ that emits this key.}%
    }{}%
    \mfanum{\csuse{mfaval@#1}}%
  }{\textcolor{mred}{\textbf{[?\,#1]}}}%
}
% A value that is deliberately a word rather than a number -- a format name, a
% verdict, a unit. Rare, and a visibly different macro so nobody reaches for
% \val and gets an error they cannot place.
\newcommand{\valtext}[1]{%
  \ifcsdef{mfaval@#1}{\csuse{mfaval@#1}}%
    {\textcolor{mred}{\textbf{[?\,#1]}}}%
}
\newcommand{\mfavaltext}[2]{%
  \csgdef{mfaval@#1}{#2}\csgdef{mfavaltext@#1}{}\listxadd{\mfavalues}{#1}%
}

% The raw stored string, for the cases where a number must appear exactly as
% the machine printed it -- inside a console block, or where the digits
% themselves are the point and grouping them would obscure the pattern.
\newcommand{\rawval}[1]{%
  \ifcsdef{mfaval@#1}{\csuse{mfaval@#1}}%
    {\textcolor{mred}{\textbf{[?\,#1]}}}%
}
\makeatother

% figures/values/all.tex is generated by `make numbers` and does nothing but
% \input each program's value file. A build with no values yet still compiles;
% every \val{} in it prints a visible marker instead of a number, which is the
% correct behaviour for a draft and an obvious failure for a release.
\IfFileExists{figures/values/all.tex}{% Generated by `make numbers` --- do not edit.
\input{figures/values/f01}
}{%
  \AtBeginDocument{\PackageWarning{mfabook}{No computed values found. Run `make numbers`.}}}

% ============================================================
%  DIAGRAMS — Mermaid pipeline
%  ---------------------------------------------------------
%  \mermaidfig{key}{caption}{one-line manifest description}
%
%  Source of truth is figures/mermaid/<lang>/<key>.mmd, committed. `make
%  diagrams` renders each to figures/diagrams/<lang>/<key>.pdf, which is build
%  output and is not committed, so a diagram change reviews as a text diff.
%
%  The source is per-language because a diagram with English labels in the
%  Polish edition is exactly the kind of half-translation that makes a
%  bilingual book feel machine-made. CI checks that both languages carry the
%  same set of keys.
%
%  If the rendered PDF is absent the macro draws a placeholder AND typesets the
%  Mermaid source, so a build without mermaid-cli still shows the structure.
% ============================================================
\makeatletter
\newcommand{\listofdiagrams}{%
  \section*{\lblDiagramManifest}%
  \phantomsection\addcontentsline{toc}{section}{\lblDiagramManifest}%
  \@starttoc{dgm}%
}
\makeatother

\newcommand{\mermaidfig}[3]{%
  \addcontentsline{dgm}{subsection}{\protect\numberline{}\texttt{#1.mmd} --- #3}%
  \IfFileExists{figures/diagrams/\booklang/#1.pdf}{%
    \begin{figure}[htbp]
    \centering
    \includegraphics[width=0.95\linewidth,height=0.42\textheight,keepaspectratio]{figures/diagrams/\booklang/#1.pdf}%
    \caption{#2}\label{fig:#1}
    \end{figure}%
  }{%
    % The unrendered case deliberately does NOT float. A Mermaid source
    % typeset verbatim is often taller than a page, and a float cannot break,
    % so wrapping it in `figure` produces an overfull box per diagram and a
    % pile of tcolorbox warnings. In the flow it breaks like any other text.
    \par\medskip
    \begin{tcolorbox}[
      enhanced, breakable, width=\linewidth, sharp corners,
      colback=white, colframe=mgrey, boxrule=0.8pt,
      left=8pt, right=8pt, top=8pt, bottom=8pt]
      {\bfseries\color{mgrey}\small \lblNotRendered}\\[4pt]
      {\ttfamily\scriptsize\color{mgrey} figures/mermaid/\booklang/#1.mmd}\\[6pt]
      \IfFileExists{figures/mermaid/\booklang/#1.mmd}{%
        % ASCII only, like every other listing in this book: the fallback runs
        % the Mermaid source through `listings`, which aborts on a multi-byte
        % character it has no literate mapping for.
        \lstinputlisting[style=mfacode,numbers=none,basicstyle=\ttfamily\scriptsize,
                         backgroundcolor=\color{mlight},frame=none,
                         keywordstyle=\color{black}]{figures/mermaid/\booklang/#1.mmd}%
      }{%
        {\small\color{mred} \lblSourceMissing: \texttt{figures/mermaid/\booklang/#1.mmd}}%
      }%
    \end{tcolorbox}%
    \captionof{figure}{#2}\label{fig:#1}
    \par\medskip
  }%
}

% ============================================================
%  PROGRAM STUBS
%  ---------------------------------------------------------
%  Prints a visible marker so an unwritten program cannot be mistaken for a
%  written one and the page count never flatters the draft. `make debt` counts
%  them, and CI publishes the count on every build.
% ============================================================
\newcommand{\programstub}[1]{%
  \begin{tcolorbox}[
    enhanced, breakable, width=\linewidth, sharp corners,
    colback=mlight, colframe=mred, boxrule=1pt,
    borderline={0.6pt}{3pt}{mred, dashed},
    left=12pt, right=12pt, top=12pt, bottom=12pt]
    {\bfseries\color{mred}\large \lblNotWritten}\\[8pt]
    \small\raggedright #1
  \end{tcolorbox}%
}

% ============================================================
%  MATHEMATICAL OPERATORS
%  ---------------------------------------------------------
%  Language-invariant names live here. The ones that genuinely differ between
%  Polish and English usage -- tan/tg, Var/D^2, gcd/NWD, and the interval
%  brackets -- are defined in lang/en.tex and lang/pl.tex instead, so the
%  notation contract is executed by the build rather than remembered by a
%  translator.
% ============================================================
\DeclareMathOperator*{\argmin}{arg\,min}
\DeclareMathOperator*{\argmax}{arg\,max}
\DeclareMathOperator{\tr}{tr}
\DeclareMathOperator{\rank}{rank}
\DeclareMathOperator{\diag}{diag}
\DeclareMathOperator{\sgn}{sgn}
\DeclareMathOperator{\relu}{ReLU}
\DeclareMathOperator{\softmax}{softmax}
\DeclareMathOperator{\logsumexp}{logsumexp}
\DeclareMathOperator{\fl}{fl}
\DeclareMathOperator{\ulp}{ulp}
\DeclareMathOperator{\round}{round}

% A bare \log is FORBIDDEN in this book, and tools/check_notation.py fails the
% build on one. In Polish textbooks \log means base ten; in machine-learning
% writing it means base e. The same three characters carry opposite meanings to
% the two readerships this book serves, and they collide hardest in exactly the
% programs where it matters most -- entropy in bits and cross-entropy in nats
% are two programs apart. Write \ln for nats and \logb{2} for bits. It costs
% two characters and removes the ambiguity entirely.
\makeatletter
\let\mfa@origlog\log
\renewcommand{\log}{%
  \PackageError{mfabook}{A bare \string\log\space is forbidden in this book}%
    {Write \string\ln\space for natural logarithms or \string\logb{2}\space for
     bits. The base is never left to a convention here; see Appendix B.}%
  \mfa@origlog}
% The one legitimate form: an explicit base.
\newcommand{\logb}[1]{\mathop{\mfa@origlog}\nolimits_{#1}}
% ...and a way to write the forbidden thing when the subject IS the ban, which
% is Appendix B and nowhere else.
\newcommand{\mfalogplain}{\mathop{\mfa@origlog}\nolimits}
\makeatother

\newcommand{\R}{\mathbb{R}}
\newcommand{\N}{\mathbb{N}}
\newcommand{\Z}{\mathbb{Z}}
\newcommand{\Q}{\mathbb{Q}}
\newcommand{\C}{\mathbb{C}}
\newcommand{\Fset}{\mathbb{F}}          % the floating-point numbers, F1's subject
\newcommand{\vect}[1]{\mathbf{#1}}
\newcommand{\mat}[1]{\mathbf{#1}}
\newcommand{\T}{^{\mathsf{T}}}
\newcommand{\norm}[1]{\left\lVert #1 \right\rVert}
\newcommand{\abs}[1]{\left\lvert #1 \right\rvert}
\newcommand{\inner}[2]{\left\langle #1, #2 \right\rangle}
\newcommand{\eps}{\varepsilon}

% ============================================================
%  INDEX
%  ---------------------------------------------------------
%  Two-column, and its entries include \texttt{} names that do not hyphenate.
%  Ragged right keeps multi-word entries off the margin; it cannot help a
%  single token wider than the column, so keep verbatim index entries under
%  about 29 characters and index longer names by concept instead. Both sibling
%  books learned this the same way.
% ============================================================
\apptocmd{\theindex}{\raggedright}{}{%
  \PackageWarning{mfabook}{Could not patch theindex for ragged-right setting}}

\makeindex

% ============================================================
%  PINNED FACTS — single source of truth
%  Change these here; they propagate through both editions.
% ============================================================
\newcommand{\bookedition}{0.1}
% \pinnedon is a user-visible string and lives in lang/<lang>.tex, not here.
% It said "August 2026" on the Polish title page for exactly as long as it
% took somebody to read the Polish title page.
\newcommand{\pyver}{Python 3.13}
\newcommand{\numpyver}{2.3}
; nothing else differs between them at
%  this level. Every user-visible string lives in lang/en.tex or lang/pl.tex,
%  so a label that appears in one edition and not the other is a missing macro
%  rather than a silent divergence.
% ============================================================

% \booklang must be set by the main file. Default to English rather than
% failing, so a stray % ============================================================
%  preamble.tex — styling, environments, macros
%
%  Shared by BOTH editions. main-en.tex and main-pl.tex each define
%  \booklang before % ============================================================
%  preamble.tex — styling, environments, macros
%
%  Shared by BOTH editions. main-en.tex and main-pl.tex each define
%  \booklang before \input{preamble}; nothing else differs between them at
%  this level. Every user-visible string lives in lang/en.tex or lang/pl.tex,
%  so a label that appears in one edition and not the other is a missing macro
%  rather than a silent divergence.
% ============================================================

% \booklang must be set by the main file. Default to English rather than
% failing, so a stray \input{preamble} still compiles and says what it did.
\providecommand{\booklang}{en}

\usepackage{etoolbox}   % loaded first: the language switch below needs it

% ---------- Page geometry ----------
% Same 17 x 24 cm block as the other two books in this series, so a reader who
% owns one recognises the second. Frames are short, so the measure matters
% more here than in a prose book: a frame that runs past a page turn defeats
% the cover-the-next-frame discipline the whole method rests on.
\usepackage[
  paperwidth=17cm, paperheight=24cm,
  inner=2.2cm, outer=1.8cm, top=2.2cm, bottom=2.4cm,
  head=23pt, headsep=0.7cm, footskip=1.2cm
]{geometry}

% ---------- Encoding & fonts ----------
\usepackage[T1]{fontenc}
\usepackage[utf8]{inputenc}
\IfFileExists{lmodern.sty}{\usepackage{lmodern}}{}
\IfFileExists{newtxtext.sty}{\usepackage{newtxtext}}{}
% amsmath goes here, not down with the other packages, because newtxmath
% patches amsmath's internals and must be loaded AFTER it. The wrong order does
% not fail on a machine without newtxmath -- which is most CI images -- and
% fails on the author's full TeX Live.
%
% amssymb only when newtxmath is absent: newtxmath supplies the AMS symbols
% itself (its amssymbols option is on by default) and loading both clashes.
\usepackage{amsmath}
\IfFileExists{newtxmath.sty}{\usepackage{newtxmath}}{\usepackage{amssymb}}
\IfFileExists{inconsolata.sty}{\usepackage[varqu,scaled=0.95]{inconsolata}}{}
\IfFileExists{lmodern.sty}{\usepackage{microtype}}{\usepackage[expansion=false]{microtype}}

% ---------- Language ----------
% Loading babel with a language whose .ldf is absent is a *fatal* error, not a
% warning, so both the package and each language file are probed before either
% is requested. This is inherited from the sibling books, where a minimal TeX
% installation without texlive-lang-polish aborted the run with no PDF and an
% error message that named neither babel nor the language.
%
% The Polish edition needs Polish as the main language for hyphenation; the
% English edition needs British. Both editions load the other language too,
% because the glossary appendix is bilingual in both.
\IfFileExists{babel.sty}{%
  \IfFileExists{polish.ldf}{%
    \IfFileExists{british.ldf}{%
      \ifdefstring{\booklang}{pl}%
        {\usepackage[british,main=polish]{babel}}%
        {\usepackage[polish,main=british]{babel}}%
    }{\usepackage[polish]{babel}}%
  }{%
    \IfFileExists{british.ldf}{\usepackage[british]{babel}}{}%
  }%
}{}

% babel-polish makes " an active shorthand. listings reads its body verbatim
% and an active character in a verbatim context is a category-code accident
% waiting to happen, so the shorthand is turned off globally. Polish quotation
% marks are produced with \enquote-style macros from lang/pl.tex instead.
\ifdefstring{\booklang}{pl}{\AtBeginDocument{\shorthandoff{"}}}{}

% ---------- Core packages ----------
\usepackage{graphicx}
\usepackage{xcolor}
\usepackage{booktabs}
\usepackage{tabularx}
\usepackage{longtable}
\usepackage{array}
\usepackage{enumitem}
\IfFileExists{mathtools.sty}{\usepackage{mathtools}}{}
\usepackage{listings}
\usepackage[most]{tcolorbox}
\usepackage{fancyhdr}
\usepackage{titlesec}
\usepackage{caption}
\usepackage{imakeidx}
% csquotes with autostyle follows babel's active language, so the identical
% source \enquote{x} sets 'x' in the English edition and the Polish quotation
% marks in the Polish one. This is not a nicety: babel-polish makes " an active
% shorthand, which is turned off above precisely because an active character
% near a listings body is a category-code accident waiting to happen -- so a
% quotation mark has to come from a macro rather than from a keyboard key.
\IfFileExists{csquotes.sty}{\usepackage[autostyle=true]{csquotes}}{%
  \providecommand{\enquote}[1]{``##1''}}
\usepackage[hidelinks]{hyperref}

% ---------- Palette ----------
% Deliberately the same family as the sibling books, shifted green so the three
% spines are distinguishable on a shelf.
\definecolor{mblue}{HTML}{1F4E79}
\definecolor{mgreen}{HTML}{2E6B4F}
\definecolor{mteal}{HTML}{0E7C7B}
\definecolor{mamber}{HTML}{B26A00}
\definecolor{mred}{HTML}{9B2C2C}
\definecolor{mgrey}{HTML}{5A5A5A}
\definecolor{mlight}{HTML}{F4F6F8}
\definecolor{mfaint}{HTML}{FBFCFD}
\definecolor{codebg}{HTML}{FAFAFA}
\definecolor{codeframe}{HTML}{D8DEE4}
\definecolor{kw}{HTML}{0B5394}
\definecolor{str}{HTML}{9B2C2C}
\definecolor{cmt}{HTML}{4F7A4F}
\definecolor{typ}{HTML}{0E7C7B}

\hypersetup{
  colorlinks=true,
  linkcolor=mblue,
  urlcolor=mteal,
  citecolor=mblue,
  pdfauthor={Konrad Cinkusz}
}

% \ifpl{Polish text}{English text} -- for the handful of places where a string
% is structural rather than content and belongs in the shared files:
% structure.tex's part titles, and the main files' front matter wiring.
% Anything longer than a few words belongs in lang/ or in the edition's own
% source, not here.
% \DeclareRobustCommand, not \newcommand. A fragile \ifpl inside a \part title
% is expanded one level as it is written to the .toc, which lands
% "\ifdefstring {en}{pl}{...}{...}" in the file -- and that fails on the next
% run with "Extra }, or forgotten \endgroup", an error that names neither the
% part nor this macro.
\DeclareRobustCommand{\ifpl}[2]{\ifdefstring{\booklang}{pl}{#1}{#2}}

% A robust macro cannot go into a PDF bookmark string either -- hyperref strips
% \ifdefstring and warns once per part title. Tell it which branch to take
% instead, so the bookmarks carry the right language rather than nothing.
\makeatletter
\pdfstringdefDisableCommands{%
  \ifdefstring{\booklang}{pl}{\let\ifpl\@firstoftwo}{\let\ifpl\@secondoftwo}%
}
\makeatother

% ---------- Language strings ----------
% Every user-visible word the macros below typeset comes from here. Nothing in
% programs/ or appendices/ may hard-code an English or a Polish word inside a
% macro argument that the other edition also uses.
\input{lang/\booklang}

% ============================================================
%  PROGRAM NUMBERING
%  ---------------------------------------------------------
%  Part F programs are F1..F13; the main programs restart at 1. That is
%  Stroud's scheme and it is worth reproducing, because the Foundation part is
%  explicitly optional and numbering it separately is what makes skipping it
%  feel permitted rather than delinquent.
%
%  \theHchapter is hyperref's *destination* name, and it must stay unique
%  across the whole document. Resetting the chapter counter without also
%  changing \theHchapter produces duplicate PDF destinations, a pile of
%  warnings, and bookmarks that jump to the wrong program.
% ============================================================
\newcommand{\foundationnumbering}{%
  \setcounter{chapter}{0}%
  \renewcommand{\thechapter}{F\arabic{chapter}}%
  \renewcommand{\theHchapter}{F\arabic{chapter}}%
}
\newcommand{\mainnumbering}{%
  \setcounter{chapter}{0}%
  \renewcommand{\thechapter}{\arabic{chapter}}%
  \renewcommand{\theHchapter}{M\arabic{chapter}}%
}

% \appendix resets the chapter counter and switches \thechapter to letters, but
% it knows nothing about \theHchapter, so appendix A would claim the same PDF
% destination as main program 1. Patch it once, here.
\apptocmd{\appendix}{\renewcommand{\theHchapter}{A\Alph{chapter}}}{}{%
  \PackageWarning{mfabook}{Could not patch \string\appendix\space for hyperref destinations}}

% A program is a chapter. \chaptername is set from the language file so the
% word above the title is "Program" in both editions (it happens to be the
% same word, which is a small piece of luck).
\AtBeginDocument{\renewcommand{\chaptername}{\lblProgram}}

% ---------- Headings ----------
% \raggedright on the title body is load-bearing, not cosmetic: at \Huge on a
% 13 cm text block a long title that cannot break overflows the margin badly.
\titleformat{\chapter}[display]
  {\normalfont\huge\bfseries\color{mblue}\raggedright}
  {\normalfont\normalsize\color{mgrey}\MakeUppercase{\chaptertitlename}\ \thechapter}
  {8pt}{\Huge\raggedright}
\titlespacing*{\chapter}{0pt}{10pt}{28pt}
\titleformat{\section}{\normalfont\Large\bfseries\color{mblue}}{\thesection}{0.8em}{}
\titleformat{\subsection}{\normalfont\large\bfseries\color{mgrey}}{\thesubsection}{0.7em}{}

% head=23pt in the geometry options above, not \setlength{\headheight}:
% fancyhdr computes the height it needs from the running-head font and the
% rule, wants 22.5 pt at this size, and silently overprints the top of the text
% block if it does not get it. Setting it through geometry keeps the text block
% where geometry put it instead of pushing it down.
\pagestyle{fancy}
\fancyhf{}
\fancyhead[LE]{\small\nouppercase{\leftmark}}
% The recto head carries the FRAME number, not the section. In a book of
% numbered frames "where am I" is a frame question: every Summary
% back-reference, every Quiz route and every cross-reference between programs
% is a frame number, and hunting for one by page is the single most annoying
% thing about reading a programmed text. \theframeno at shipout is the last
% frame begun on the page, which is what a reader flicking through wants.
% Guarded: the front matter and the appendices have no frames, and a head
% reading "Frame 0" over the introduction is worse than no head at all.
\fancyhead[RO]{\ifnum\value{frameno}>0\relax\small\lblFrame~\theframeno\fi}
\fancyfoot[LE,RO]{\small\thepage}
\renewcommand{\headrulewidth}{0.4pt}

% This MUST come after \pagestyle{fancy}: fancyhdr installs its own \chaptermark
% at that point and silently overwrites anything defined earlier. The symptom is
% a running-head change that appears to do nothing at all.
%
% The running head carries the *frame* number as well as the program, because
% in a book of numbered frames "where am I" is a frame question, not a page
% question, and the summary's back-references are frame numbers.
% \@chapapp, not \lblProgram. \appendix switches \@chapapp from \chaptername
% to \appendixname, and babel supplies both in the right language; hard-coding
% the word gives an appendix a running head reading "Program A. Answers" over a
% page whose own chapter head reads "APPENDIX A".
\makeatletter
\renewcommand{\chaptermark}[1]{\markboth{\@chapapp\ \thechapter.\ #1}{}}
\makeatother
\renewcommand{\sectionmark}[1]{\markright{\thesection\ #1}}

% ============================================================
%  FRAMES — the unit the whole method is built from
%  ---------------------------------------------------------
%  A program is a stream of numbered frames. Nearly every frame ends by asking
%  the reader for something, and the answer opens the *next* frame, so that the
%  reader can cover what is below and commit to an answer before seeing it.
%
%  The environment is called `fr` rather than `frame` for two reasons. First,
%  \frame is already defined by LaTeX2e (it draws a box in a picture
%  environment) and an environment named `frame` would silently redefine it.
%  Second, this is typed forty to seventy times per program, so it wants to be
%  short in the source.
%
%    \begin{fr}
%    \ans{$x = 3$}            % the answer to the previous frame, optional
%    Teaching text, ending in a question.
%    \end{fr}
% ============================================================

\newcounter{frameno}[chapter]
\renewcommand{\theframeno}{\arabic{frameno}}

% Frames are referenced constantly -- by the Summary, by the Quiz, by the
% "Can you?" checklist and by other programs. \label inside a frame must
% therefore capture the frame number, not the section number.
\makeatletter
\newcommand{\framelabelfix}{%
  \@namedef{theHframeno}{\theHchapter.\arabic{frameno}}%
}
\makeatother

\newcommand{\framerule}{%
  \par\penalty-200\vspace{9pt}%
  \noindent\textcolor{codeframe}{\rule{\linewidth}{0.5pt}}%
  \par\vspace{5pt}%
}

% The frame number as a filled badge, set in the text flow rather than in the
% margin. A margin number would have to alternate sides under `twoside` and
% would be the first thing to overflow on a 17 cm page; a badge cannot.
\newcommand{\framebadge}{%
  \begingroup
  \setlength{\fboxsep}{2.5pt}%
  \colorbox{mblue}{\textcolor{white}{\bfseries\scriptsize\theframeno}}%
  \endgroup\hspace{0.7em}%
}

\newenvironment{fr}{%
  \framerule
  \refstepcounter{frameno}%
  \nopagebreak
  \noindent\framebadge\ignorespaces
}{\par}

% The answer to the previous frame. Set apart and centred so the eye can find
% the edge to cover, and so a reader who has answered can check in one glance
% without reading on.
\newcommand{\ans}[1]{%
  \par\nobreak\vspace{-4pt}%
  \begin{tcolorbox}[
    enhanced, sharp corners, boxrule=0pt, leftrule=2.5pt,
    colback=mfaint, colframe=mgreen,
    left=8pt, right=8pt, top=4pt, bottom=4pt,
    width=0.86\linewidth, center, halign=center,
    before skip=2pt, after skip=7pt]
    #1
  \end{tcolorbox}%
  \nobreak\noindent\ignorespaces
}

% A multi-paragraph or worked answer, where a centred box would be wrong.
\newenvironment{ansblock}{%
  \par\nopagebreak
  \begin{tcolorbox}[
    enhanced, breakable, sharp corners, boxrule=0pt, leftrule=2.5pt,
    colback=mfaint, colframe=mgreen,
    left=8pt, right=8pt, top=5pt, bottom=5pt]
}{%
  \end{tcolorbox}%
  \nopagebreak\par\vspace{2pt}\noindent\ignorespaces
}

% The row of dots. Stroud's signal that the reader is to supply a word, a
% number or an expression. \blank is inline; \dotline occupies its own line,
% which is what a longer derivation step wants.
% \penalty0 offers TeX a break opportunity immediately before the dots. Without
% it the run of dots is glued to the last word of the question, and a long
% question in the Polish edition -- where the words are longer and hyphenate
% less freely -- runs into the margin.
\newcommand{\blank}[1][4em]{\penalty0\makebox[#1]{\dotfill}}
\newcommand{\dotline}{\par\vspace{2pt}\noindent\makebox[\linewidth]{\dotfill}\par\vspace{2pt}}

% "Now one for you to do." The third rung of the scaffolding gradient, and the
% one most often skipped by authors who think they have written an exercise
% when they have written another worked example.
\newcommand{\yourturn}{%
  \par\vspace{4pt}\noindent
  {\bfseries\color{mgreen}\lblYourTurn}\par\nobreak\vspace{2pt}\noindent\ignorespaces
}

% ============================================================
%  THE PROGRAM SKELETON
%  ---------------------------------------------------------
%  Every program carries the same seven parts, in the same order. They are
%  macros rather than a convention so that a missing one is a build-time
%  absence rather than something a reader notices.
% ============================================================

% ---- Learning outcomes, on entry ----------------------------------------
% Stored as well as printed, because "Can you?" at the end of the program must
% be the same list, one for one. Storing it is what makes that structural
% rather than a promise: the checklist is generated from the outcomes, so the
% two cannot drift.
\newcounter{outcomeno}[chapter]

\makeatletter
% The key under which everything belonging to the current program is stored.
% It must expand to the printed program number (F1, F2, ..., 1, 2, ...), which
% is what makes an answer findable from the program that owns it.
\newcommand{\mfa@progkey}{\thechapter}
\newcommand{\outcome}[2]{% #1 = frames that teach it, #2 = the outcome
  \stepcounter{outcomeno}%
  \csgdef{mfa@out@\mfa@progkey @\theoutcomeno}{#2}%
  \csgdef{mfa@outfr@\mfa@progkey @\theoutcomeno}{#1}%
  % The "Can you?" row is built here rather than looped over later. A \loop
  % inside a tabularx body does not survive alignment scanning: TeX needs to
  % see the & and the row break while it reads the cell, and a conditional
  % hides them. Accumulating the rows as tokens at declaration time sidesteps
  % that entirely, and has the better property anyway -- the checklist is
  % literally made of the outcomes, so the two cannot disagree.
  \csgappto{mfa@rows@\mfa@progkey}{%
    #2 & {\footnotesize\color{mgrey}#1}
      & \ratingcell & \ratingcell & \ratingcell & \ratingcell & \ratingcell
    \tabularnewline}%
  \item #2\hfill{\footnotesize\color{mgrey}[\lblFrames~#1]}%
}
\makeatother

\makeatletter
\newenvironment{outcomes}{%
  \setcounter{outcomeno}{0}%
  \csgdef{mfa@rows@\mfa@progkey}{}%
  \begin{tcolorbox}[
    enhanced, breakable, sharp corners,
    colback=mlight, colframe=mblue, boxrule=0pt, leftrule=3pt,
    left=8pt, right=8pt, top=6pt, bottom=6pt,
    fonttitle=\bfseries\small, title={\lblOutcomes}]
  \begin{itemize}[leftmargin=1.3em, itemsep=3pt, topsep=2pt]
}{%
  \end{itemize}
  \end{tcolorbox}
}
\makeatother

% ---- "Can you?" -----------------------------------------------------------
% Generated from the stored outcomes, not written out again by hand. Five
% columns rather than a tick box, because a self-assessment with a midpoint is
% one people actually use and a yes/no one is one they lie to.
\newcommand{\ratingcell}{\textcolor{codeframe}{\rule[-0.15em]{0.8em}{0.8em}}}

\makeatletter
\newcommand{\canyou}{%
  \section*{\lblCanYou}%
  \phantomsection\addcontentsline{toc}{section}{\lblCanYou}%
  \noindent\lblCanYouIntro\par\vspace{8pt}
  \noindent
  \begingroup\small
  \ifcsvoid{mfa@rows@\mfa@progkey}{%
    {\color{mred}\lblNoOutcomes}\par
  }{%
    \begin{tabularx}{\linewidth}{@{}Xcccccc@{}}
    \toprule
    \textbf{\lblOnEntryYouCould} & \textbf{\lblFrames} & 1 & 2 & 3 & 4 & 5 \\
    \midrule
    \csuse{mfa@rows@\mfa@progkey}
    \bottomrule
    \end{tabularx}%
  }%
  \endgroup
  \par\vspace{6pt}\noindent{\footnotesize\color{mgrey}\lblCanYouFooter}\par
}
\makeatother

% ============================================================
%  QUIZ, SUMMARY, EXERCISES
% ============================================================

% ---- Quiz -----------------------------------------------------------------
% Foundation programs only. It is a diagnostic on the way in and the same test
% on the way out, which is what lets one book serve a reader who needs the
% whole program and a reader who needs four frames of it. Each question names
% the frames that teach it, so a wrong answer routes the reader somewhere
% specific rather than back to the beginning.
\newlist{quizlist}{enumerate}{1}
\setlist[quizlist]{label=\arabic*., leftmargin=2.0em, itemsep=5pt, topsep=3pt}

\makeatletter
\newenvironment{quiz}{%
  \section*{\lblQuiz}%
  \phantomsection\addcontentsline{toc}{section}{\lblQuiz}%
  \def\mfa@exkey{Q\arabic{quizlisti}}%
  \noindent\lblQuizIntro\par\vspace{5pt}
  \begin{quizlist}
}{%
  \end{quizlist}
}
\makeatother

% The frames that teach the question just asked. Two forms, because a Quiz
% question that routes to a single frame reading "Frames 22" is the kind of
% small wrongness a reader notices and an author never does.
\newcommand{\teachesat}[1]{\hfill{\footnotesize\color{mgrey}(\lblFrames~#1)}}
\newcommand{\teachesatone}[1]{\hfill{\footnotesize\color{mgrey}(\lblFrame~#1)}}

% ---- Summary --------------------------------------------------------------
% Every item carries the frame number it came from, in square brackets. That
% is the eighth-edition improvement and it is the single cheapest thing in the
% format: it turns a summary into a return index, so a reader who fails one
% line of "Can you?" knows exactly which four frames to reread.
\newenvironment{summarybox}{%
  \section*{\lblSummary}%
  \phantomsection\addcontentsline{toc}{section}{\lblSummary}%
  \begin{tcolorbox}[
    enhanced, breakable, sharp corners,
    colback=mlight, colframe=mblue, boxrule=0pt, leftrule=3pt,
    left=8pt, right=8pt, top=6pt, bottom=6pt]
  \begin{itemize}[leftmargin=1.3em, itemsep=5pt, topsep=2pt]
}{%
  \end{itemize}
  \end{tcolorbox}
}

% \sumitem{12}{text} -- the frame reference is not decoration, it is the point.
\newcommand{\sumitem}[2]{\item #2~{\footnotesize\color{mgrey}[#1]}}

% ---- Answers ---------------------------------------------------------------
%  Answers are authored at the point of use and printed at the back.
%
%  They are carried there by a global macro store rather than by an auxiliary
%  file. The sibling books collect their figure and screenshot manifests with
%  \@starttoc and a custom file extension, which works because those entries
%  are one line of text; it is the wrong mechanism for an answer, because an
%  answer contains displayed maths, and material written to an .aux-style file
%  and read back is expanded at shipout, where fragile commands break in ways
%  whose error messages name neither the answer nor the program.
%
%  \include is sequential within a run, so a macro defined globally while
%  typesetting program F1 is still defined when the back matter is typeset.
%  That is all this needs.
\makeatletter
\newcommand{\mfa@storeanswer}[3]{% #1 program key, #2 item key, #3 body
  \csgdef{mfa@a@#1@#2}{#3}%
  % \listcsxadd, not \listcsgadd: the item must be EXPANDED as it is stored.
  % Storing the token \mfa@exkey means every entry in the list expands to
  % whatever the key happens to be at the time the back matter is typeset,
  % which is the last one.
  \listcsxadd{mfa@akeys@#1}{#2}%
}
% Used inside an exercise item. Prints nothing here; the body reappears at the
% back of the book under this program's heading.
\newcommand{\answerto}[1]{\mfa@storeanswer{\mfa@progkey}{\mfa@exkey}{#1}}
\providecommand{\mfa@exkey}{?}
\makeatother

\newlist{exlist}{enumerate}{1}
\setlist[exlist]{label=\arabic*., leftmargin=2.0em, itemsep=6pt, topsep=3pt}
\newlist{fplist}{enumerate}{1}
\setlist[fplist]{label=\arabic*., leftmargin=2.0em, itemsep=5pt, topsep=3pt}

% These MUST be defined inside \makeatletter. Outside it, \def\mfa@exkey{...}
% parses as \def\mfa @exkey{...} -- a definition of \mfa with a delimited
% parameter text -- which raises no error at all and quietly leaves every
% answer filed under the same key. That failure is invisible until the answers
% section at the back of the book prints one program's answers twice.
\makeatletter
\newenvironment{testexercises}{%
  \section*{\lblTestExercises}%
  \phantomsection\addcontentsline{toc}{section}{\lblTestExercises}%
  \def\mfa@exkey{T\arabic{exlisti}}%
  \noindent\lblTestIntro\par\vspace{5pt}
  \begin{exlist}
}{%
  \end{exlist}
}

\newenvironment{furtherproblems}{%
  \section*{\lblFurther}%
  \phantomsection\addcontentsline{toc}{section}{\lblFurther}%
  \def\mfa@exkey{P\arabic{fplisti}}%
  \noindent\lblFurtherIntro\par\vspace{5pt}
  \begin{fplist}
}{%
  \end{fplist}
}
\makeatother

% ---- The programs register -------------------------------------------------
% \program{Title} opens a program: it is \chapter plus the bookkeeping that
% lets the back matter reproduce this program's answers under its own name.
\makeatletter
\newcommand{\program}[1]{%
  \chapter{#1}%
  \csgdef{mfa@title@\thechapter}{#1}%
  % Expanded, for the same reason as the answer keys above.
  \listxadd{\mfaprograms}{\thechapter}%
}
\newcommand{\mfa@printoneanswer}[2]{% #1 program key, #2 item key
  \item[\textbf{#2}] \csuse{mfa@a@#1@#2}%
}
\newcommand{\mfa@answerprogram}[1]{%
  \subsection*{\lblProgram~#1\quad\textmd{\itshape\csuse{mfa@title@#1}}}%
  \ifcsdef{mfa@akeys@#1}{%
    \begin{list}{}{%
      \setlength{\leftmargin}{3.2em}\setlength{\labelwidth}{2.7em}%
      \setlength{\labelsep}{0.5em}\setlength{\itemsep}{3pt}%
      \setlength{\topsep}{3pt}\setlength{\parsep}{0pt}%
      \renewcommand{\makelabel}[1]{\hfill##1}}
      \forlistcsloop{\mfa@printoneanswer{#1}}{mfa@akeys@#1}
    \end{list}%
  }{%
    \par\noindent{\color{mred}\lblNoAnswers}\par
  }%
}
% The answers body, without a heading of its own, so the appendix that carries
% it can be a numbered appendix like any other and appear in the running heads
% and the ToC in the ordinary way.
\newcommand{\answersbody}{%
  \noindent\lblAnswersIntro\par\vspace{8pt}
  \ifdef{\mfaprograms}{\forlistloop{\mfa@answerprogram}{\mfaprograms}}{%
    \noindent{\color{mred}\lblNoAnswers}\par}%
}
% Unnumbered form, for a draft build that has no appendix wiring yet.
\newcommand{\printanswers}{%
  \chapter*{\lblAnswers}%
  \phantomsection\addcontentsline{toc}{chapter}{\lblAnswers}%
  \markboth{\lblAnswers}{\lblAnswers}%
  \answersbody
}
\makeatother

% ============================================================
%  ADMONITIONS
%  ---------------------------------------------------------
%  A deliberately small, fixed vocabulary. Each one earns its place by doing
%  something no other box does; a seventh would be a smell.
% ============================================================

\newtcolorbox{note}[1][]{
  enhanced, breakable, sharp corners,
  colback=mlight, colframe=mblue, boxrule=0pt, leftrule=3pt,
  left=8pt, right=8pt, top=6pt, bottom=6pt,
  fonttitle=\bfseries\small, title={\lblNote}, #1
}

\newtcolorbox{warning}[1][]{
  enhanced, breakable, sharp corners,
  colback=mlight, colframe=mred, boxrule=0pt, leftrule=3pt,
  left=8pt, right=8pt, top=6pt, bottom=6pt,
  fonttitle=\bfseries\small, title={\lblWarning}, #1
}

% The misconception, stated as a summary after it has been walked into. The
% trap itself belongs in the frames -- elicit the wrong answer, then say
% flatly that it is wrong -- and this box is the thing the reader marks with a
% finger and comes back to.
\newtcolorbox{trapbox}[1][]{
  enhanced, breakable, sharp corners,
  colback=white, colframe=mred, boxrule=0.8pt,
  left=8pt, right=8pt, top=6pt, bottom=6pt,
  fonttitle=\bfseries\small, title={\lblTrap}, #1
}

% Where the mathematics just done shows up in a system the reader has actually
% deployed. This is the box that separates the book from a maths textbook, and
% it is also the one most easily filled with waffle: if it cannot name a
% specific line of a specific system, it should not be written.
\newtcolorbox{aibox}[1][]{
  enhanced, breakable, sharp corners,
  colback=mlight, colframe=mteal, boxrule=0pt, leftrule=3pt,
  left=8pt, right=8pt, top=6pt, bottom=6pt,
  fonttitle=\bfseries\small, title={\lblInAI}, #1
}

% What is being asserted rather than proved, and where to read the proof.
% Stroud's method buys fluency at the cost of rigour, and the honest thing is
% to say so at each point rather than once in an introduction nobody rereads.
\newtcolorbox{rigourbox}[1][]{
  enhanced, breakable, sharp corners,
  colback=white, colframe=mgrey, boxrule=0.6pt,
  left=8pt, right=8pt, top=6pt, bottom=6pt,
  fonttitle=\bfseries\small, title={\lblRigour}, #1
}

% Notation that genuinely differs -- between the Polish and English editions,
% and between disciplines. Matrix calculus alone has two incompatible layout
% conventions, and half the confusion about backpropagation is a transpose
% that came from the other one.
\newtcolorbox{notationbox}[1][]{
  enhanced, breakable, sharp corners,
  colback=mlight, colframe=mamber, boxrule=0pt, leftrule=3pt,
  left=8pt, right=8pt, top=6pt, bottom=6pt,
  fonttitle=\bfseries\small, title={\lblNotation}, #1
}

% The maths analogue of the sibling books' "run this before you trust it".
% A number in this book is either produced by a script under code/ and pulled
% in with \val, or it carries this box. Nothing ships to a reader with one
% still attached.
\newtcolorbox{verifybox}[1][]{
  enhanced, breakable, sharp corners,
  colback=white, colframe=mamber, boxrule=0.8pt,
  left=8pt, right=8pt, top=6pt, bottom=6pt,
  fonttitle=\bfseries\small, title={\lblVerify}, #1
}

\newtcolorbox{exercisebox}[1][]{
  enhanced, breakable, sharp corners,
  colback=white, colframe=mgreen, boxrule=0pt, leftrule=3pt,
  left=8pt, right=8pt, top=6pt, bottom=6pt,
  fonttitle=\bfseries\small, title={\lblExercise}, #1
}

% ============================================================
%  CODE LISTINGS
%  ---------------------------------------------------------
%  Code appears here for one reason only: to check a number. Every listing in
%  this book is runnable and every number it prints is the number in the text.
% ============================================================
\lstdefinelanguage{pythonx}{
  language=Python,
  morekeywords={as, with, assert, match, case},
  morekeywords=[2]{
    float64,float32,float16,bfloat16,int8,uint8,
    np,numpy,array,asarray,frombuffer,nextafter,finfo,iinfo,
    exp,log,log1p,expm1,logaddexp,sum,max,abs,sqrt,mean,var
  },
  sensitive=true
}

\lstdefinestyle{mfacode}{
  basicstyle=\ttfamily\footnotesize,
  keywordstyle=\color{kw}\bfseries,
  keywordstyle=[2]\color{typ},
  stringstyle=\color{str},
  commentstyle=\color{cmt}\itshape,
  numbers=left,
  numberstyle=\ttfamily\tiny\color{mgrey},
  numbersep=8pt,
  backgroundcolor=\color{codebg},
  frame=single,
  rulecolor=\color{codeframe},
  framesep=6pt,
  breaklines=true,
  breakatwhitespace=false,
  postbreak=\mbox{\textcolor{mgrey}{$\hookrightarrow$}\space},
  showstringspaces=false,
  tabsize=4,
  columns=fullflexible,
  keepspaces=true,
  captionpos=b,
  aboveskip=1.2em,
  belowskip=1.2em,
  % Maps the characters most likely to be pasted into a listing by accident.
  %
  % Do NOT add a mapping for U+00A0 (non-breaking space). It looks like the
  % obvious next entry and it makes `listings` abort the run with "Improper
  % alphabetic constant" -- fatal, no PDF, and the message names neither the
  % character nor this line. The ASCII-in-listings convention is the guard
  % against it instead. This trap has now been hit in two of the three books
  % in this series.
  literate={ł}{{\l}}1 {Ł}{{\L}}1 {ą}{{\k a}}1 {ę}{{\k e}}1
           {ś}{{\'s}}1 {ć}{{\'c}}1 {ż}{{\.z}}1 {ź}{{\'z}}1 {ó}{{\'o}}1 {ń}{{\'n}}1
           {Ą}{{\k A}}1 {Ę}{{\k E}}1 {Ś}{{\'S}}1 {Ć}{{\'C}}1
           {Ż}{{\.Z}}1 {Ź}{{\'Z}}1 {Ó}{{\'O}}1 {Ń}{{\'N}}1
           {—}{{-{}-}}2 {–}{{-}}1 {‘}{{'}}1 {’}{{'}}1
           {“}{{"}}1 {”}{{"}}1 {°}{{\textdegree}}1
}

\lstset{style=mfacode}

\lstnewenvironment{python}[1][]
  {\lstset{language=pythonx,style=mfacode,#1}}
  {}

\lstnewenvironment{shellcmd}[1][]
  {\lstset{language=bash,style=mfacode,numbers=none,keywordstyle=\color{mteal}\bfseries,#1}}
  {}

% Terminal transcripts: no line numbers, no keyword colouring. Everything in
% one of these was copied from a real run.
\lstnewenvironment{console}[1][]
  {\lstset{style=mfacode,numbers=none,basicstyle=\ttfamily\scriptsize,
           keywordstyle=\color{black},backgroundcolor=\color{white},#1}}
  {}

\newcommand{\pyfile}[3]{%
  \lstinputlisting[language=pythonx,style=mfacode,caption={#2},label={#3}]{#1}%
}
\newcommand{\pyregion}[4]{%
  \lstinputlisting[
    language=pythonx, style=mfacode,
    linerange={--8<--\ [start:#2]--- 8<--\ [end:#2]},
    includerangemarker=false,
    caption={#3}, label={#4}]{#1}%
}

% Inline literals
\newcommand{\code}[1]{\texttt{\small #1}}
\newcommand{\pkg}[1]{\texttt{\small\color{mblue}#1}}
\newcommand{\api}[1]{\texttt{\small\color{mteal}#1}}

% ============================================================
%  COMPUTED VALUES
%  ---------------------------------------------------------
%  The house rule in the sibling books is "run the listing, or mark it". The
%  analogue here is stronger, because a wrong digit in a maths book is not a
%  broken example, it is a lie the reader will carry.
%
%  Every number that came out of a computation is written by a script under
%  code/ into figures/values/<program>.tex as
%
%      \mfaval{f01.eps64}{2.220446049250313e-16}
%
%  and pulled into the text with \val{f01.eps64}. The script is the source of
%  truth; the book cannot disagree with it, because the book does not contain
%  the digits. A value the scripts no longer produce prints a visible marker
%  and is counted by `make debt`.
%
%  \val also applies the locale's decimal separator, which is how the Polish
%  edition gets its decimal comma without a translator having to remember.
% ============================================================
\IfFileExists{siunitx.sty}{%
  \usepackage{siunitx}%
  \sisetup{
    group-digits = integer,
    group-minimum-digits = 5,
    group-separator = {\,},
    exponent-product = \cdot,
    output-decimal-marker = {\mfadecimalmarker}
  }%
  \newcommand{\mfanum}[1]{\num{##1}}%
}{%
  \newcommand{\mfanum}[1]{##1}%
}

\makeatletter
\newcommand{\mfaval}[2]{\csgdef{mfaval@#1}{#2}\listxadd{\mfavalues}{#1}}
% \num on a non-numeric body is a FATAL siunitx error -- "Invalid number" --
% and under -halt-on-error that is the whole build. A script emitting a word
% where the book expected a number is a plausible mistake.
%
% The check is NOT done here. Deciding in TeX whether a token list is a number
% means parsing it character by character, which is fragile and was got wrong
% once already. The generator knows the answer for free, so it emits \mfaval
% for a number and \mfavaltext for a word, and this end only has to enforce
% which macro reads which.
\newcommand{\val}[1]{%
  \ifcsdef{mfaval@#1}{%
    \ifcsdef{mfavaltext@#1}{%
      \PackageError{mfabook}{Value `#1' is not a number}%
        {\string\val\space passes its body to siunitx, which refuses anything
         that is not a number. Use \string\valtext{#1} instead, or fix the
         script in code/ that emits this key.}%
    }{}%
    \mfanum{\csuse{mfaval@#1}}%
  }{\textcolor{mred}{\textbf{[?\,#1]}}}%
}
% A value that is deliberately a word rather than a number -- a format name, a
% verdict, a unit. Rare, and a visibly different macro so nobody reaches for
% \val and gets an error they cannot place.
\newcommand{\valtext}[1]{%
  \ifcsdef{mfaval@#1}{\csuse{mfaval@#1}}%
    {\textcolor{mred}{\textbf{[?\,#1]}}}%
}
\newcommand{\mfavaltext}[2]{%
  \csgdef{mfaval@#1}{#2}\csgdef{mfavaltext@#1}{}\listxadd{\mfavalues}{#1}%
}

% The raw stored string, for the cases where a number must appear exactly as
% the machine printed it -- inside a console block, or where the digits
% themselves are the point and grouping them would obscure the pattern.
\newcommand{\rawval}[1]{%
  \ifcsdef{mfaval@#1}{\csuse{mfaval@#1}}%
    {\textcolor{mred}{\textbf{[?\,#1]}}}%
}
\makeatother

% figures/values/all.tex is generated by `make numbers` and does nothing but
% \input each program's value file. A build with no values yet still compiles;
% every \val{} in it prints a visible marker instead of a number, which is the
% correct behaviour for a draft and an obvious failure for a release.
\IfFileExists{figures/values/all.tex}{\input{figures/values/all}}{%
  \AtBeginDocument{\PackageWarning{mfabook}{No computed values found. Run `make numbers`.}}}

% ============================================================
%  DIAGRAMS — Mermaid pipeline
%  ---------------------------------------------------------
%  \mermaidfig{key}{caption}{one-line manifest description}
%
%  Source of truth is figures/mermaid/<lang>/<key>.mmd, committed. `make
%  diagrams` renders each to figures/diagrams/<lang>/<key>.pdf, which is build
%  output and is not committed, so a diagram change reviews as a text diff.
%
%  The source is per-language because a diagram with English labels in the
%  Polish edition is exactly the kind of half-translation that makes a
%  bilingual book feel machine-made. CI checks that both languages carry the
%  same set of keys.
%
%  If the rendered PDF is absent the macro draws a placeholder AND typesets the
%  Mermaid source, so a build without mermaid-cli still shows the structure.
% ============================================================
\makeatletter
\newcommand{\listofdiagrams}{%
  \section*{\lblDiagramManifest}%
  \phantomsection\addcontentsline{toc}{section}{\lblDiagramManifest}%
  \@starttoc{dgm}%
}
\makeatother

\newcommand{\mermaidfig}[3]{%
  \addcontentsline{dgm}{subsection}{\protect\numberline{}\texttt{#1.mmd} --- #3}%
  \IfFileExists{figures/diagrams/\booklang/#1.pdf}{%
    \begin{figure}[htbp]
    \centering
    \includegraphics[width=0.95\linewidth,height=0.42\textheight,keepaspectratio]{figures/diagrams/\booklang/#1.pdf}%
    \caption{#2}\label{fig:#1}
    \end{figure}%
  }{%
    % The unrendered case deliberately does NOT float. A Mermaid source
    % typeset verbatim is often taller than a page, and a float cannot break,
    % so wrapping it in `figure` produces an overfull box per diagram and a
    % pile of tcolorbox warnings. In the flow it breaks like any other text.
    \par\medskip
    \begin{tcolorbox}[
      enhanced, breakable, width=\linewidth, sharp corners,
      colback=white, colframe=mgrey, boxrule=0.8pt,
      left=8pt, right=8pt, top=8pt, bottom=8pt]
      {\bfseries\color{mgrey}\small \lblNotRendered}\\[4pt]
      {\ttfamily\scriptsize\color{mgrey} figures/mermaid/\booklang/#1.mmd}\\[6pt]
      \IfFileExists{figures/mermaid/\booklang/#1.mmd}{%
        % ASCII only, like every other listing in this book: the fallback runs
        % the Mermaid source through `listings`, which aborts on a multi-byte
        % character it has no literate mapping for.
        \lstinputlisting[style=mfacode,numbers=none,basicstyle=\ttfamily\scriptsize,
                         backgroundcolor=\color{mlight},frame=none,
                         keywordstyle=\color{black}]{figures/mermaid/\booklang/#1.mmd}%
      }{%
        {\small\color{mred} \lblSourceMissing: \texttt{figures/mermaid/\booklang/#1.mmd}}%
      }%
    \end{tcolorbox}%
    \captionof{figure}{#2}\label{fig:#1}
    \par\medskip
  }%
}

% ============================================================
%  PROGRAM STUBS
%  ---------------------------------------------------------
%  Prints a visible marker so an unwritten program cannot be mistaken for a
%  written one and the page count never flatters the draft. `make debt` counts
%  them, and CI publishes the count on every build.
% ============================================================
\newcommand{\programstub}[1]{%
  \begin{tcolorbox}[
    enhanced, breakable, width=\linewidth, sharp corners,
    colback=mlight, colframe=mred, boxrule=1pt,
    borderline={0.6pt}{3pt}{mred, dashed},
    left=12pt, right=12pt, top=12pt, bottom=12pt]
    {\bfseries\color{mred}\large \lblNotWritten}\\[8pt]
    \small\raggedright #1
  \end{tcolorbox}%
}

% ============================================================
%  MATHEMATICAL OPERATORS
%  ---------------------------------------------------------
%  Language-invariant names live here. The ones that genuinely differ between
%  Polish and English usage -- tan/tg, Var/D^2, gcd/NWD, and the interval
%  brackets -- are defined in lang/en.tex and lang/pl.tex instead, so the
%  notation contract is executed by the build rather than remembered by a
%  translator.
% ============================================================
\DeclareMathOperator*{\argmin}{arg\,min}
\DeclareMathOperator*{\argmax}{arg\,max}
\DeclareMathOperator{\tr}{tr}
\DeclareMathOperator{\rank}{rank}
\DeclareMathOperator{\diag}{diag}
\DeclareMathOperator{\sgn}{sgn}
\DeclareMathOperator{\relu}{ReLU}
\DeclareMathOperator{\softmax}{softmax}
\DeclareMathOperator{\logsumexp}{logsumexp}
\DeclareMathOperator{\fl}{fl}
\DeclareMathOperator{\ulp}{ulp}
\DeclareMathOperator{\round}{round}

% A bare \log is FORBIDDEN in this book, and tools/check_notation.py fails the
% build on one. In Polish textbooks \log means base ten; in machine-learning
% writing it means base e. The same three characters carry opposite meanings to
% the two readerships this book serves, and they collide hardest in exactly the
% programs where it matters most -- entropy in bits and cross-entropy in nats
% are two programs apart. Write \ln for nats and \logb{2} for bits. It costs
% two characters and removes the ambiguity entirely.
\makeatletter
\let\mfa@origlog\log
\renewcommand{\log}{%
  \PackageError{mfabook}{A bare \string\log\space is forbidden in this book}%
    {Write \string\ln\space for natural logarithms or \string\logb{2}\space for
     bits. The base is never left to a convention here; see Appendix B.}%
  \mfa@origlog}
% The one legitimate form: an explicit base.
\newcommand{\logb}[1]{\mathop{\mfa@origlog}\nolimits_{#1}}
% ...and a way to write the forbidden thing when the subject IS the ban, which
% is Appendix B and nowhere else.
\newcommand{\mfalogplain}{\mathop{\mfa@origlog}\nolimits}
\makeatother

\newcommand{\R}{\mathbb{R}}
\newcommand{\N}{\mathbb{N}}
\newcommand{\Z}{\mathbb{Z}}
\newcommand{\Q}{\mathbb{Q}}
\newcommand{\C}{\mathbb{C}}
\newcommand{\Fset}{\mathbb{F}}          % the floating-point numbers, F1's subject
\newcommand{\vect}[1]{\mathbf{#1}}
\newcommand{\mat}[1]{\mathbf{#1}}
\newcommand{\T}{^{\mathsf{T}}}
\newcommand{\norm}[1]{\left\lVert #1 \right\rVert}
\newcommand{\abs}[1]{\left\lvert #1 \right\rvert}
\newcommand{\inner}[2]{\left\langle #1, #2 \right\rangle}
\newcommand{\eps}{\varepsilon}

% ============================================================
%  INDEX
%  ---------------------------------------------------------
%  Two-column, and its entries include \texttt{} names that do not hyphenate.
%  Ragged right keeps multi-word entries off the margin; it cannot help a
%  single token wider than the column, so keep verbatim index entries under
%  about 29 characters and index longer names by concept instead. Both sibling
%  books learned this the same way.
% ============================================================
\apptocmd{\theindex}{\raggedright}{}{%
  \PackageWarning{mfabook}{Could not patch theindex for ragged-right setting}}

\makeindex

% ============================================================
%  PINNED FACTS — single source of truth
%  Change these here; they propagate through both editions.
% ============================================================
\newcommand{\bookedition}{0.1}
% \pinnedon is a user-visible string and lives in lang/<lang>.tex, not here.
% It said "August 2026" on the Polish title page for exactly as long as it
% took somebody to read the Polish title page.
\newcommand{\pyver}{Python 3.13}
\newcommand{\numpyver}{2.3}
; nothing else differs between them at
%  this level. Every user-visible string lives in lang/en.tex or lang/pl.tex,
%  so a label that appears in one edition and not the other is a missing macro
%  rather than a silent divergence.
% ============================================================

% \booklang must be set by the main file. Default to English rather than
% failing, so a stray % ============================================================
%  preamble.tex — styling, environments, macros
%
%  Shared by BOTH editions. main-en.tex and main-pl.tex each define
%  \booklang before \input{preamble}; nothing else differs between them at
%  this level. Every user-visible string lives in lang/en.tex or lang/pl.tex,
%  so a label that appears in one edition and not the other is a missing macro
%  rather than a silent divergence.
% ============================================================

% \booklang must be set by the main file. Default to English rather than
% failing, so a stray \input{preamble} still compiles and says what it did.
\providecommand{\booklang}{en}

\usepackage{etoolbox}   % loaded first: the language switch below needs it

% ---------- Page geometry ----------
% Same 17 x 24 cm block as the other two books in this series, so a reader who
% owns one recognises the second. Frames are short, so the measure matters
% more here than in a prose book: a frame that runs past a page turn defeats
% the cover-the-next-frame discipline the whole method rests on.
\usepackage[
  paperwidth=17cm, paperheight=24cm,
  inner=2.2cm, outer=1.8cm, top=2.2cm, bottom=2.4cm,
  head=23pt, headsep=0.7cm, footskip=1.2cm
]{geometry}

% ---------- Encoding & fonts ----------
\usepackage[T1]{fontenc}
\usepackage[utf8]{inputenc}
\IfFileExists{lmodern.sty}{\usepackage{lmodern}}{}
\IfFileExists{newtxtext.sty}{\usepackage{newtxtext}}{}
% amsmath goes here, not down with the other packages, because newtxmath
% patches amsmath's internals and must be loaded AFTER it. The wrong order does
% not fail on a machine without newtxmath -- which is most CI images -- and
% fails on the author's full TeX Live.
%
% amssymb only when newtxmath is absent: newtxmath supplies the AMS symbols
% itself (its amssymbols option is on by default) and loading both clashes.
\usepackage{amsmath}
\IfFileExists{newtxmath.sty}{\usepackage{newtxmath}}{\usepackage{amssymb}}
\IfFileExists{inconsolata.sty}{\usepackage[varqu,scaled=0.95]{inconsolata}}{}
\IfFileExists{lmodern.sty}{\usepackage{microtype}}{\usepackage[expansion=false]{microtype}}

% ---------- Language ----------
% Loading babel with a language whose .ldf is absent is a *fatal* error, not a
% warning, so both the package and each language file are probed before either
% is requested. This is inherited from the sibling books, where a minimal TeX
% installation without texlive-lang-polish aborted the run with no PDF and an
% error message that named neither babel nor the language.
%
% The Polish edition needs Polish as the main language for hyphenation; the
% English edition needs British. Both editions load the other language too,
% because the glossary appendix is bilingual in both.
\IfFileExists{babel.sty}{%
  \IfFileExists{polish.ldf}{%
    \IfFileExists{british.ldf}{%
      \ifdefstring{\booklang}{pl}%
        {\usepackage[british,main=polish]{babel}}%
        {\usepackage[polish,main=british]{babel}}%
    }{\usepackage[polish]{babel}}%
  }{%
    \IfFileExists{british.ldf}{\usepackage[british]{babel}}{}%
  }%
}{}

% babel-polish makes " an active shorthand. listings reads its body verbatim
% and an active character in a verbatim context is a category-code accident
% waiting to happen, so the shorthand is turned off globally. Polish quotation
% marks are produced with \enquote-style macros from lang/pl.tex instead.
\ifdefstring{\booklang}{pl}{\AtBeginDocument{\shorthandoff{"}}}{}

% ---------- Core packages ----------
\usepackage{graphicx}
\usepackage{xcolor}
\usepackage{booktabs}
\usepackage{tabularx}
\usepackage{longtable}
\usepackage{array}
\usepackage{enumitem}
\IfFileExists{mathtools.sty}{\usepackage{mathtools}}{}
\usepackage{listings}
\usepackage[most]{tcolorbox}
\usepackage{fancyhdr}
\usepackage{titlesec}
\usepackage{caption}
\usepackage{imakeidx}
% csquotes with autostyle follows babel's active language, so the identical
% source \enquote{x} sets 'x' in the English edition and the Polish quotation
% marks in the Polish one. This is not a nicety: babel-polish makes " an active
% shorthand, which is turned off above precisely because an active character
% near a listings body is a category-code accident waiting to happen -- so a
% quotation mark has to come from a macro rather than from a keyboard key.
\IfFileExists{csquotes.sty}{\usepackage[autostyle=true]{csquotes}}{%
  \providecommand{\enquote}[1]{``##1''}}
\usepackage[hidelinks]{hyperref}

% ---------- Palette ----------
% Deliberately the same family as the sibling books, shifted green so the three
% spines are distinguishable on a shelf.
\definecolor{mblue}{HTML}{1F4E79}
\definecolor{mgreen}{HTML}{2E6B4F}
\definecolor{mteal}{HTML}{0E7C7B}
\definecolor{mamber}{HTML}{B26A00}
\definecolor{mred}{HTML}{9B2C2C}
\definecolor{mgrey}{HTML}{5A5A5A}
\definecolor{mlight}{HTML}{F4F6F8}
\definecolor{mfaint}{HTML}{FBFCFD}
\definecolor{codebg}{HTML}{FAFAFA}
\definecolor{codeframe}{HTML}{D8DEE4}
\definecolor{kw}{HTML}{0B5394}
\definecolor{str}{HTML}{9B2C2C}
\definecolor{cmt}{HTML}{4F7A4F}
\definecolor{typ}{HTML}{0E7C7B}

\hypersetup{
  colorlinks=true,
  linkcolor=mblue,
  urlcolor=mteal,
  citecolor=mblue,
  pdfauthor={Konrad Cinkusz}
}

% \ifpl{Polish text}{English text} -- for the handful of places where a string
% is structural rather than content and belongs in the shared files:
% structure.tex's part titles, and the main files' front matter wiring.
% Anything longer than a few words belongs in lang/ or in the edition's own
% source, not here.
% \DeclareRobustCommand, not \newcommand. A fragile \ifpl inside a \part title
% is expanded one level as it is written to the .toc, which lands
% "\ifdefstring {en}{pl}{...}{...}" in the file -- and that fails on the next
% run with "Extra }, or forgotten \endgroup", an error that names neither the
% part nor this macro.
\DeclareRobustCommand{\ifpl}[2]{\ifdefstring{\booklang}{pl}{#1}{#2}}

% A robust macro cannot go into a PDF bookmark string either -- hyperref strips
% \ifdefstring and warns once per part title. Tell it which branch to take
% instead, so the bookmarks carry the right language rather than nothing.
\makeatletter
\pdfstringdefDisableCommands{%
  \ifdefstring{\booklang}{pl}{\let\ifpl\@firstoftwo}{\let\ifpl\@secondoftwo}%
}
\makeatother

% ---------- Language strings ----------
% Every user-visible word the macros below typeset comes from here. Nothing in
% programs/ or appendices/ may hard-code an English or a Polish word inside a
% macro argument that the other edition also uses.
\input{lang/\booklang}

% ============================================================
%  PROGRAM NUMBERING
%  ---------------------------------------------------------
%  Part F programs are F1..F13; the main programs restart at 1. That is
%  Stroud's scheme and it is worth reproducing, because the Foundation part is
%  explicitly optional and numbering it separately is what makes skipping it
%  feel permitted rather than delinquent.
%
%  \theHchapter is hyperref's *destination* name, and it must stay unique
%  across the whole document. Resetting the chapter counter without also
%  changing \theHchapter produces duplicate PDF destinations, a pile of
%  warnings, and bookmarks that jump to the wrong program.
% ============================================================
\newcommand{\foundationnumbering}{%
  \setcounter{chapter}{0}%
  \renewcommand{\thechapter}{F\arabic{chapter}}%
  \renewcommand{\theHchapter}{F\arabic{chapter}}%
}
\newcommand{\mainnumbering}{%
  \setcounter{chapter}{0}%
  \renewcommand{\thechapter}{\arabic{chapter}}%
  \renewcommand{\theHchapter}{M\arabic{chapter}}%
}

% \appendix resets the chapter counter and switches \thechapter to letters, but
% it knows nothing about \theHchapter, so appendix A would claim the same PDF
% destination as main program 1. Patch it once, here.
\apptocmd{\appendix}{\renewcommand{\theHchapter}{A\Alph{chapter}}}{}{%
  \PackageWarning{mfabook}{Could not patch \string\appendix\space for hyperref destinations}}

% A program is a chapter. \chaptername is set from the language file so the
% word above the title is "Program" in both editions (it happens to be the
% same word, which is a small piece of luck).
\AtBeginDocument{\renewcommand{\chaptername}{\lblProgram}}

% ---------- Headings ----------
% \raggedright on the title body is load-bearing, not cosmetic: at \Huge on a
% 13 cm text block a long title that cannot break overflows the margin badly.
\titleformat{\chapter}[display]
  {\normalfont\huge\bfseries\color{mblue}\raggedright}
  {\normalfont\normalsize\color{mgrey}\MakeUppercase{\chaptertitlename}\ \thechapter}
  {8pt}{\Huge\raggedright}
\titlespacing*{\chapter}{0pt}{10pt}{28pt}
\titleformat{\section}{\normalfont\Large\bfseries\color{mblue}}{\thesection}{0.8em}{}
\titleformat{\subsection}{\normalfont\large\bfseries\color{mgrey}}{\thesubsection}{0.7em}{}

% head=23pt in the geometry options above, not \setlength{\headheight}:
% fancyhdr computes the height it needs from the running-head font and the
% rule, wants 22.5 pt at this size, and silently overprints the top of the text
% block if it does not get it. Setting it through geometry keeps the text block
% where geometry put it instead of pushing it down.
\pagestyle{fancy}
\fancyhf{}
\fancyhead[LE]{\small\nouppercase{\leftmark}}
% The recto head carries the FRAME number, not the section. In a book of
% numbered frames "where am I" is a frame question: every Summary
% back-reference, every Quiz route and every cross-reference between programs
% is a frame number, and hunting for one by page is the single most annoying
% thing about reading a programmed text. \theframeno at shipout is the last
% frame begun on the page, which is what a reader flicking through wants.
% Guarded: the front matter and the appendices have no frames, and a head
% reading "Frame 0" over the introduction is worse than no head at all.
\fancyhead[RO]{\ifnum\value{frameno}>0\relax\small\lblFrame~\theframeno\fi}
\fancyfoot[LE,RO]{\small\thepage}
\renewcommand{\headrulewidth}{0.4pt}

% This MUST come after \pagestyle{fancy}: fancyhdr installs its own \chaptermark
% at that point and silently overwrites anything defined earlier. The symptom is
% a running-head change that appears to do nothing at all.
%
% The running head carries the *frame* number as well as the program, because
% in a book of numbered frames "where am I" is a frame question, not a page
% question, and the summary's back-references are frame numbers.
% \@chapapp, not \lblProgram. \appendix switches \@chapapp from \chaptername
% to \appendixname, and babel supplies both in the right language; hard-coding
% the word gives an appendix a running head reading "Program A. Answers" over a
% page whose own chapter head reads "APPENDIX A".
\makeatletter
\renewcommand{\chaptermark}[1]{\markboth{\@chapapp\ \thechapter.\ #1}{}}
\makeatother
\renewcommand{\sectionmark}[1]{\markright{\thesection\ #1}}

% ============================================================
%  FRAMES — the unit the whole method is built from
%  ---------------------------------------------------------
%  A program is a stream of numbered frames. Nearly every frame ends by asking
%  the reader for something, and the answer opens the *next* frame, so that the
%  reader can cover what is below and commit to an answer before seeing it.
%
%  The environment is called `fr` rather than `frame` for two reasons. First,
%  \frame is already defined by LaTeX2e (it draws a box in a picture
%  environment) and an environment named `frame` would silently redefine it.
%  Second, this is typed forty to seventy times per program, so it wants to be
%  short in the source.
%
%    \begin{fr}
%    \ans{$x = 3$}            % the answer to the previous frame, optional
%    Teaching text, ending in a question.
%    \end{fr}
% ============================================================

\newcounter{frameno}[chapter]
\renewcommand{\theframeno}{\arabic{frameno}}

% Frames are referenced constantly -- by the Summary, by the Quiz, by the
% "Can you?" checklist and by other programs. \label inside a frame must
% therefore capture the frame number, not the section number.
\makeatletter
\newcommand{\framelabelfix}{%
  \@namedef{theHframeno}{\theHchapter.\arabic{frameno}}%
}
\makeatother

\newcommand{\framerule}{%
  \par\penalty-200\vspace{9pt}%
  \noindent\textcolor{codeframe}{\rule{\linewidth}{0.5pt}}%
  \par\vspace{5pt}%
}

% The frame number as a filled badge, set in the text flow rather than in the
% margin. A margin number would have to alternate sides under `twoside` and
% would be the first thing to overflow on a 17 cm page; a badge cannot.
\newcommand{\framebadge}{%
  \begingroup
  \setlength{\fboxsep}{2.5pt}%
  \colorbox{mblue}{\textcolor{white}{\bfseries\scriptsize\theframeno}}%
  \endgroup\hspace{0.7em}%
}

\newenvironment{fr}{%
  \framerule
  \refstepcounter{frameno}%
  \nopagebreak
  \noindent\framebadge\ignorespaces
}{\par}

% The answer to the previous frame. Set apart and centred so the eye can find
% the edge to cover, and so a reader who has answered can check in one glance
% without reading on.
\newcommand{\ans}[1]{%
  \par\nobreak\vspace{-4pt}%
  \begin{tcolorbox}[
    enhanced, sharp corners, boxrule=0pt, leftrule=2.5pt,
    colback=mfaint, colframe=mgreen,
    left=8pt, right=8pt, top=4pt, bottom=4pt,
    width=0.86\linewidth, center, halign=center,
    before skip=2pt, after skip=7pt]
    #1
  \end{tcolorbox}%
  \nobreak\noindent\ignorespaces
}

% A multi-paragraph or worked answer, where a centred box would be wrong.
\newenvironment{ansblock}{%
  \par\nopagebreak
  \begin{tcolorbox}[
    enhanced, breakable, sharp corners, boxrule=0pt, leftrule=2.5pt,
    colback=mfaint, colframe=mgreen,
    left=8pt, right=8pt, top=5pt, bottom=5pt]
}{%
  \end{tcolorbox}%
  \nopagebreak\par\vspace{2pt}\noindent\ignorespaces
}

% The row of dots. Stroud's signal that the reader is to supply a word, a
% number or an expression. \blank is inline; \dotline occupies its own line,
% which is what a longer derivation step wants.
% \penalty0 offers TeX a break opportunity immediately before the dots. Without
% it the run of dots is glued to the last word of the question, and a long
% question in the Polish edition -- where the words are longer and hyphenate
% less freely -- runs into the margin.
\newcommand{\blank}[1][4em]{\penalty0\makebox[#1]{\dotfill}}
\newcommand{\dotline}{\par\vspace{2pt}\noindent\makebox[\linewidth]{\dotfill}\par\vspace{2pt}}

% "Now one for you to do." The third rung of the scaffolding gradient, and the
% one most often skipped by authors who think they have written an exercise
% when they have written another worked example.
\newcommand{\yourturn}{%
  \par\vspace{4pt}\noindent
  {\bfseries\color{mgreen}\lblYourTurn}\par\nobreak\vspace{2pt}\noindent\ignorespaces
}

% ============================================================
%  THE PROGRAM SKELETON
%  ---------------------------------------------------------
%  Every program carries the same seven parts, in the same order. They are
%  macros rather than a convention so that a missing one is a build-time
%  absence rather than something a reader notices.
% ============================================================

% ---- Learning outcomes, on entry ----------------------------------------
% Stored as well as printed, because "Can you?" at the end of the program must
% be the same list, one for one. Storing it is what makes that structural
% rather than a promise: the checklist is generated from the outcomes, so the
% two cannot drift.
\newcounter{outcomeno}[chapter]

\makeatletter
% The key under which everything belonging to the current program is stored.
% It must expand to the printed program number (F1, F2, ..., 1, 2, ...), which
% is what makes an answer findable from the program that owns it.
\newcommand{\mfa@progkey}{\thechapter}
\newcommand{\outcome}[2]{% #1 = frames that teach it, #2 = the outcome
  \stepcounter{outcomeno}%
  \csgdef{mfa@out@\mfa@progkey @\theoutcomeno}{#2}%
  \csgdef{mfa@outfr@\mfa@progkey @\theoutcomeno}{#1}%
  % The "Can you?" row is built here rather than looped over later. A \loop
  % inside a tabularx body does not survive alignment scanning: TeX needs to
  % see the & and the row break while it reads the cell, and a conditional
  % hides them. Accumulating the rows as tokens at declaration time sidesteps
  % that entirely, and has the better property anyway -- the checklist is
  % literally made of the outcomes, so the two cannot disagree.
  \csgappto{mfa@rows@\mfa@progkey}{%
    #2 & {\footnotesize\color{mgrey}#1}
      & \ratingcell & \ratingcell & \ratingcell & \ratingcell & \ratingcell
    \tabularnewline}%
  \item #2\hfill{\footnotesize\color{mgrey}[\lblFrames~#1]}%
}
\makeatother

\makeatletter
\newenvironment{outcomes}{%
  \setcounter{outcomeno}{0}%
  \csgdef{mfa@rows@\mfa@progkey}{}%
  \begin{tcolorbox}[
    enhanced, breakable, sharp corners,
    colback=mlight, colframe=mblue, boxrule=0pt, leftrule=3pt,
    left=8pt, right=8pt, top=6pt, bottom=6pt,
    fonttitle=\bfseries\small, title={\lblOutcomes}]
  \begin{itemize}[leftmargin=1.3em, itemsep=3pt, topsep=2pt]
}{%
  \end{itemize}
  \end{tcolorbox}
}
\makeatother

% ---- "Can you?" -----------------------------------------------------------
% Generated from the stored outcomes, not written out again by hand. Five
% columns rather than a tick box, because a self-assessment with a midpoint is
% one people actually use and a yes/no one is one they lie to.
\newcommand{\ratingcell}{\textcolor{codeframe}{\rule[-0.15em]{0.8em}{0.8em}}}

\makeatletter
\newcommand{\canyou}{%
  \section*{\lblCanYou}%
  \phantomsection\addcontentsline{toc}{section}{\lblCanYou}%
  \noindent\lblCanYouIntro\par\vspace{8pt}
  \noindent
  \begingroup\small
  \ifcsvoid{mfa@rows@\mfa@progkey}{%
    {\color{mred}\lblNoOutcomes}\par
  }{%
    \begin{tabularx}{\linewidth}{@{}Xcccccc@{}}
    \toprule
    \textbf{\lblOnEntryYouCould} & \textbf{\lblFrames} & 1 & 2 & 3 & 4 & 5 \\
    \midrule
    \csuse{mfa@rows@\mfa@progkey}
    \bottomrule
    \end{tabularx}%
  }%
  \endgroup
  \par\vspace{6pt}\noindent{\footnotesize\color{mgrey}\lblCanYouFooter}\par
}
\makeatother

% ============================================================
%  QUIZ, SUMMARY, EXERCISES
% ============================================================

% ---- Quiz -----------------------------------------------------------------
% Foundation programs only. It is a diagnostic on the way in and the same test
% on the way out, which is what lets one book serve a reader who needs the
% whole program and a reader who needs four frames of it. Each question names
% the frames that teach it, so a wrong answer routes the reader somewhere
% specific rather than back to the beginning.
\newlist{quizlist}{enumerate}{1}
\setlist[quizlist]{label=\arabic*., leftmargin=2.0em, itemsep=5pt, topsep=3pt}

\makeatletter
\newenvironment{quiz}{%
  \section*{\lblQuiz}%
  \phantomsection\addcontentsline{toc}{section}{\lblQuiz}%
  \def\mfa@exkey{Q\arabic{quizlisti}}%
  \noindent\lblQuizIntro\par\vspace{5pt}
  \begin{quizlist}
}{%
  \end{quizlist}
}
\makeatother

% The frames that teach the question just asked. Two forms, because a Quiz
% question that routes to a single frame reading "Frames 22" is the kind of
% small wrongness a reader notices and an author never does.
\newcommand{\teachesat}[1]{\hfill{\footnotesize\color{mgrey}(\lblFrames~#1)}}
\newcommand{\teachesatone}[1]{\hfill{\footnotesize\color{mgrey}(\lblFrame~#1)}}

% ---- Summary --------------------------------------------------------------
% Every item carries the frame number it came from, in square brackets. That
% is the eighth-edition improvement and it is the single cheapest thing in the
% format: it turns a summary into a return index, so a reader who fails one
% line of "Can you?" knows exactly which four frames to reread.
\newenvironment{summarybox}{%
  \section*{\lblSummary}%
  \phantomsection\addcontentsline{toc}{section}{\lblSummary}%
  \begin{tcolorbox}[
    enhanced, breakable, sharp corners,
    colback=mlight, colframe=mblue, boxrule=0pt, leftrule=3pt,
    left=8pt, right=8pt, top=6pt, bottom=6pt]
  \begin{itemize}[leftmargin=1.3em, itemsep=5pt, topsep=2pt]
}{%
  \end{itemize}
  \end{tcolorbox}
}

% \sumitem{12}{text} -- the frame reference is not decoration, it is the point.
\newcommand{\sumitem}[2]{\item #2~{\footnotesize\color{mgrey}[#1]}}

% ---- Answers ---------------------------------------------------------------
%  Answers are authored at the point of use and printed at the back.
%
%  They are carried there by a global macro store rather than by an auxiliary
%  file. The sibling books collect their figure and screenshot manifests with
%  \@starttoc and a custom file extension, which works because those entries
%  are one line of text; it is the wrong mechanism for an answer, because an
%  answer contains displayed maths, and material written to an .aux-style file
%  and read back is expanded at shipout, where fragile commands break in ways
%  whose error messages name neither the answer nor the program.
%
%  \include is sequential within a run, so a macro defined globally while
%  typesetting program F1 is still defined when the back matter is typeset.
%  That is all this needs.
\makeatletter
\newcommand{\mfa@storeanswer}[3]{% #1 program key, #2 item key, #3 body
  \csgdef{mfa@a@#1@#2}{#3}%
  % \listcsxadd, not \listcsgadd: the item must be EXPANDED as it is stored.
  % Storing the token \mfa@exkey means every entry in the list expands to
  % whatever the key happens to be at the time the back matter is typeset,
  % which is the last one.
  \listcsxadd{mfa@akeys@#1}{#2}%
}
% Used inside an exercise item. Prints nothing here; the body reappears at the
% back of the book under this program's heading.
\newcommand{\answerto}[1]{\mfa@storeanswer{\mfa@progkey}{\mfa@exkey}{#1}}
\providecommand{\mfa@exkey}{?}
\makeatother

\newlist{exlist}{enumerate}{1}
\setlist[exlist]{label=\arabic*., leftmargin=2.0em, itemsep=6pt, topsep=3pt}
\newlist{fplist}{enumerate}{1}
\setlist[fplist]{label=\arabic*., leftmargin=2.0em, itemsep=5pt, topsep=3pt}

% These MUST be defined inside \makeatletter. Outside it, \def\mfa@exkey{...}
% parses as \def\mfa @exkey{...} -- a definition of \mfa with a delimited
% parameter text -- which raises no error at all and quietly leaves every
% answer filed under the same key. That failure is invisible until the answers
% section at the back of the book prints one program's answers twice.
\makeatletter
\newenvironment{testexercises}{%
  \section*{\lblTestExercises}%
  \phantomsection\addcontentsline{toc}{section}{\lblTestExercises}%
  \def\mfa@exkey{T\arabic{exlisti}}%
  \noindent\lblTestIntro\par\vspace{5pt}
  \begin{exlist}
}{%
  \end{exlist}
}

\newenvironment{furtherproblems}{%
  \section*{\lblFurther}%
  \phantomsection\addcontentsline{toc}{section}{\lblFurther}%
  \def\mfa@exkey{P\arabic{fplisti}}%
  \noindent\lblFurtherIntro\par\vspace{5pt}
  \begin{fplist}
}{%
  \end{fplist}
}
\makeatother

% ---- The programs register -------------------------------------------------
% \program{Title} opens a program: it is \chapter plus the bookkeeping that
% lets the back matter reproduce this program's answers under its own name.
\makeatletter
\newcommand{\program}[1]{%
  \chapter{#1}%
  \csgdef{mfa@title@\thechapter}{#1}%
  % Expanded, for the same reason as the answer keys above.
  \listxadd{\mfaprograms}{\thechapter}%
}
\newcommand{\mfa@printoneanswer}[2]{% #1 program key, #2 item key
  \item[\textbf{#2}] \csuse{mfa@a@#1@#2}%
}
\newcommand{\mfa@answerprogram}[1]{%
  \subsection*{\lblProgram~#1\quad\textmd{\itshape\csuse{mfa@title@#1}}}%
  \ifcsdef{mfa@akeys@#1}{%
    \begin{list}{}{%
      \setlength{\leftmargin}{3.2em}\setlength{\labelwidth}{2.7em}%
      \setlength{\labelsep}{0.5em}\setlength{\itemsep}{3pt}%
      \setlength{\topsep}{3pt}\setlength{\parsep}{0pt}%
      \renewcommand{\makelabel}[1]{\hfill##1}}
      \forlistcsloop{\mfa@printoneanswer{#1}}{mfa@akeys@#1}
    \end{list}%
  }{%
    \par\noindent{\color{mred}\lblNoAnswers}\par
  }%
}
% The answers body, without a heading of its own, so the appendix that carries
% it can be a numbered appendix like any other and appear in the running heads
% and the ToC in the ordinary way.
\newcommand{\answersbody}{%
  \noindent\lblAnswersIntro\par\vspace{8pt}
  \ifdef{\mfaprograms}{\forlistloop{\mfa@answerprogram}{\mfaprograms}}{%
    \noindent{\color{mred}\lblNoAnswers}\par}%
}
% Unnumbered form, for a draft build that has no appendix wiring yet.
\newcommand{\printanswers}{%
  \chapter*{\lblAnswers}%
  \phantomsection\addcontentsline{toc}{chapter}{\lblAnswers}%
  \markboth{\lblAnswers}{\lblAnswers}%
  \answersbody
}
\makeatother

% ============================================================
%  ADMONITIONS
%  ---------------------------------------------------------
%  A deliberately small, fixed vocabulary. Each one earns its place by doing
%  something no other box does; a seventh would be a smell.
% ============================================================

\newtcolorbox{note}[1][]{
  enhanced, breakable, sharp corners,
  colback=mlight, colframe=mblue, boxrule=0pt, leftrule=3pt,
  left=8pt, right=8pt, top=6pt, bottom=6pt,
  fonttitle=\bfseries\small, title={\lblNote}, #1
}

\newtcolorbox{warning}[1][]{
  enhanced, breakable, sharp corners,
  colback=mlight, colframe=mred, boxrule=0pt, leftrule=3pt,
  left=8pt, right=8pt, top=6pt, bottom=6pt,
  fonttitle=\bfseries\small, title={\lblWarning}, #1
}

% The misconception, stated as a summary after it has been walked into. The
% trap itself belongs in the frames -- elicit the wrong answer, then say
% flatly that it is wrong -- and this box is the thing the reader marks with a
% finger and comes back to.
\newtcolorbox{trapbox}[1][]{
  enhanced, breakable, sharp corners,
  colback=white, colframe=mred, boxrule=0.8pt,
  left=8pt, right=8pt, top=6pt, bottom=6pt,
  fonttitle=\bfseries\small, title={\lblTrap}, #1
}

% Where the mathematics just done shows up in a system the reader has actually
% deployed. This is the box that separates the book from a maths textbook, and
% it is also the one most easily filled with waffle: if it cannot name a
% specific line of a specific system, it should not be written.
\newtcolorbox{aibox}[1][]{
  enhanced, breakable, sharp corners,
  colback=mlight, colframe=mteal, boxrule=0pt, leftrule=3pt,
  left=8pt, right=8pt, top=6pt, bottom=6pt,
  fonttitle=\bfseries\small, title={\lblInAI}, #1
}

% What is being asserted rather than proved, and where to read the proof.
% Stroud's method buys fluency at the cost of rigour, and the honest thing is
% to say so at each point rather than once in an introduction nobody rereads.
\newtcolorbox{rigourbox}[1][]{
  enhanced, breakable, sharp corners,
  colback=white, colframe=mgrey, boxrule=0.6pt,
  left=8pt, right=8pt, top=6pt, bottom=6pt,
  fonttitle=\bfseries\small, title={\lblRigour}, #1
}

% Notation that genuinely differs -- between the Polish and English editions,
% and between disciplines. Matrix calculus alone has two incompatible layout
% conventions, and half the confusion about backpropagation is a transpose
% that came from the other one.
\newtcolorbox{notationbox}[1][]{
  enhanced, breakable, sharp corners,
  colback=mlight, colframe=mamber, boxrule=0pt, leftrule=3pt,
  left=8pt, right=8pt, top=6pt, bottom=6pt,
  fonttitle=\bfseries\small, title={\lblNotation}, #1
}

% The maths analogue of the sibling books' "run this before you trust it".
% A number in this book is either produced by a script under code/ and pulled
% in with \val, or it carries this box. Nothing ships to a reader with one
% still attached.
\newtcolorbox{verifybox}[1][]{
  enhanced, breakable, sharp corners,
  colback=white, colframe=mamber, boxrule=0.8pt,
  left=8pt, right=8pt, top=6pt, bottom=6pt,
  fonttitle=\bfseries\small, title={\lblVerify}, #1
}

\newtcolorbox{exercisebox}[1][]{
  enhanced, breakable, sharp corners,
  colback=white, colframe=mgreen, boxrule=0pt, leftrule=3pt,
  left=8pt, right=8pt, top=6pt, bottom=6pt,
  fonttitle=\bfseries\small, title={\lblExercise}, #1
}

% ============================================================
%  CODE LISTINGS
%  ---------------------------------------------------------
%  Code appears here for one reason only: to check a number. Every listing in
%  this book is runnable and every number it prints is the number in the text.
% ============================================================
\lstdefinelanguage{pythonx}{
  language=Python,
  morekeywords={as, with, assert, match, case},
  morekeywords=[2]{
    float64,float32,float16,bfloat16,int8,uint8,
    np,numpy,array,asarray,frombuffer,nextafter,finfo,iinfo,
    exp,log,log1p,expm1,logaddexp,sum,max,abs,sqrt,mean,var
  },
  sensitive=true
}

\lstdefinestyle{mfacode}{
  basicstyle=\ttfamily\footnotesize,
  keywordstyle=\color{kw}\bfseries,
  keywordstyle=[2]\color{typ},
  stringstyle=\color{str},
  commentstyle=\color{cmt}\itshape,
  numbers=left,
  numberstyle=\ttfamily\tiny\color{mgrey},
  numbersep=8pt,
  backgroundcolor=\color{codebg},
  frame=single,
  rulecolor=\color{codeframe},
  framesep=6pt,
  breaklines=true,
  breakatwhitespace=false,
  postbreak=\mbox{\textcolor{mgrey}{$\hookrightarrow$}\space},
  showstringspaces=false,
  tabsize=4,
  columns=fullflexible,
  keepspaces=true,
  captionpos=b,
  aboveskip=1.2em,
  belowskip=1.2em,
  % Maps the characters most likely to be pasted into a listing by accident.
  %
  % Do NOT add a mapping for U+00A0 (non-breaking space). It looks like the
  % obvious next entry and it makes `listings` abort the run with "Improper
  % alphabetic constant" -- fatal, no PDF, and the message names neither the
  % character nor this line. The ASCII-in-listings convention is the guard
  % against it instead. This trap has now been hit in two of the three books
  % in this series.
  literate={ł}{{\l}}1 {Ł}{{\L}}1 {ą}{{\k a}}1 {ę}{{\k e}}1
           {ś}{{\'s}}1 {ć}{{\'c}}1 {ż}{{\.z}}1 {ź}{{\'z}}1 {ó}{{\'o}}1 {ń}{{\'n}}1
           {Ą}{{\k A}}1 {Ę}{{\k E}}1 {Ś}{{\'S}}1 {Ć}{{\'C}}1
           {Ż}{{\.Z}}1 {Ź}{{\'Z}}1 {Ó}{{\'O}}1 {Ń}{{\'N}}1
           {—}{{-{}-}}2 {–}{{-}}1 {‘}{{'}}1 {’}{{'}}1
           {“}{{"}}1 {”}{{"}}1 {°}{{\textdegree}}1
}

\lstset{style=mfacode}

\lstnewenvironment{python}[1][]
  {\lstset{language=pythonx,style=mfacode,#1}}
  {}

\lstnewenvironment{shellcmd}[1][]
  {\lstset{language=bash,style=mfacode,numbers=none,keywordstyle=\color{mteal}\bfseries,#1}}
  {}

% Terminal transcripts: no line numbers, no keyword colouring. Everything in
% one of these was copied from a real run.
\lstnewenvironment{console}[1][]
  {\lstset{style=mfacode,numbers=none,basicstyle=\ttfamily\scriptsize,
           keywordstyle=\color{black},backgroundcolor=\color{white},#1}}
  {}

\newcommand{\pyfile}[3]{%
  \lstinputlisting[language=pythonx,style=mfacode,caption={#2},label={#3}]{#1}%
}
\newcommand{\pyregion}[4]{%
  \lstinputlisting[
    language=pythonx, style=mfacode,
    linerange={--8<--\ [start:#2]--- 8<--\ [end:#2]},
    includerangemarker=false,
    caption={#3}, label={#4}]{#1}%
}

% Inline literals
\newcommand{\code}[1]{\texttt{\small #1}}
\newcommand{\pkg}[1]{\texttt{\small\color{mblue}#1}}
\newcommand{\api}[1]{\texttt{\small\color{mteal}#1}}

% ============================================================
%  COMPUTED VALUES
%  ---------------------------------------------------------
%  The house rule in the sibling books is "run the listing, or mark it". The
%  analogue here is stronger, because a wrong digit in a maths book is not a
%  broken example, it is a lie the reader will carry.
%
%  Every number that came out of a computation is written by a script under
%  code/ into figures/values/<program>.tex as
%
%      \mfaval{f01.eps64}{2.220446049250313e-16}
%
%  and pulled into the text with \val{f01.eps64}. The script is the source of
%  truth; the book cannot disagree with it, because the book does not contain
%  the digits. A value the scripts no longer produce prints a visible marker
%  and is counted by `make debt`.
%
%  \val also applies the locale's decimal separator, which is how the Polish
%  edition gets its decimal comma without a translator having to remember.
% ============================================================
\IfFileExists{siunitx.sty}{%
  \usepackage{siunitx}%
  \sisetup{
    group-digits = integer,
    group-minimum-digits = 5,
    group-separator = {\,},
    exponent-product = \cdot,
    output-decimal-marker = {\mfadecimalmarker}
  }%
  \newcommand{\mfanum}[1]{\num{##1}}%
}{%
  \newcommand{\mfanum}[1]{##1}%
}

\makeatletter
\newcommand{\mfaval}[2]{\csgdef{mfaval@#1}{#2}\listxadd{\mfavalues}{#1}}
% \num on a non-numeric body is a FATAL siunitx error -- "Invalid number" --
% and under -halt-on-error that is the whole build. A script emitting a word
% where the book expected a number is a plausible mistake.
%
% The check is NOT done here. Deciding in TeX whether a token list is a number
% means parsing it character by character, which is fragile and was got wrong
% once already. The generator knows the answer for free, so it emits \mfaval
% for a number and \mfavaltext for a word, and this end only has to enforce
% which macro reads which.
\newcommand{\val}[1]{%
  \ifcsdef{mfaval@#1}{%
    \ifcsdef{mfavaltext@#1}{%
      \PackageError{mfabook}{Value `#1' is not a number}%
        {\string\val\space passes its body to siunitx, which refuses anything
         that is not a number. Use \string\valtext{#1} instead, or fix the
         script in code/ that emits this key.}%
    }{}%
    \mfanum{\csuse{mfaval@#1}}%
  }{\textcolor{mred}{\textbf{[?\,#1]}}}%
}
% A value that is deliberately a word rather than a number -- a format name, a
% verdict, a unit. Rare, and a visibly different macro so nobody reaches for
% \val and gets an error they cannot place.
\newcommand{\valtext}[1]{%
  \ifcsdef{mfaval@#1}{\csuse{mfaval@#1}}%
    {\textcolor{mred}{\textbf{[?\,#1]}}}%
}
\newcommand{\mfavaltext}[2]{%
  \csgdef{mfaval@#1}{#2}\csgdef{mfavaltext@#1}{}\listxadd{\mfavalues}{#1}%
}

% The raw stored string, for the cases where a number must appear exactly as
% the machine printed it -- inside a console block, or where the digits
% themselves are the point and grouping them would obscure the pattern.
\newcommand{\rawval}[1]{%
  \ifcsdef{mfaval@#1}{\csuse{mfaval@#1}}%
    {\textcolor{mred}{\textbf{[?\,#1]}}}%
}
\makeatother

% figures/values/all.tex is generated by `make numbers` and does nothing but
% \input each program's value file. A build with no values yet still compiles;
% every \val{} in it prints a visible marker instead of a number, which is the
% correct behaviour for a draft and an obvious failure for a release.
\IfFileExists{figures/values/all.tex}{\input{figures/values/all}}{%
  \AtBeginDocument{\PackageWarning{mfabook}{No computed values found. Run `make numbers`.}}}

% ============================================================
%  DIAGRAMS — Mermaid pipeline
%  ---------------------------------------------------------
%  \mermaidfig{key}{caption}{one-line manifest description}
%
%  Source of truth is figures/mermaid/<lang>/<key>.mmd, committed. `make
%  diagrams` renders each to figures/diagrams/<lang>/<key>.pdf, which is build
%  output and is not committed, so a diagram change reviews as a text diff.
%
%  The source is per-language because a diagram with English labels in the
%  Polish edition is exactly the kind of half-translation that makes a
%  bilingual book feel machine-made. CI checks that both languages carry the
%  same set of keys.
%
%  If the rendered PDF is absent the macro draws a placeholder AND typesets the
%  Mermaid source, so a build without mermaid-cli still shows the structure.
% ============================================================
\makeatletter
\newcommand{\listofdiagrams}{%
  \section*{\lblDiagramManifest}%
  \phantomsection\addcontentsline{toc}{section}{\lblDiagramManifest}%
  \@starttoc{dgm}%
}
\makeatother

\newcommand{\mermaidfig}[3]{%
  \addcontentsline{dgm}{subsection}{\protect\numberline{}\texttt{#1.mmd} --- #3}%
  \IfFileExists{figures/diagrams/\booklang/#1.pdf}{%
    \begin{figure}[htbp]
    \centering
    \includegraphics[width=0.95\linewidth,height=0.42\textheight,keepaspectratio]{figures/diagrams/\booklang/#1.pdf}%
    \caption{#2}\label{fig:#1}
    \end{figure}%
  }{%
    % The unrendered case deliberately does NOT float. A Mermaid source
    % typeset verbatim is often taller than a page, and a float cannot break,
    % so wrapping it in `figure` produces an overfull box per diagram and a
    % pile of tcolorbox warnings. In the flow it breaks like any other text.
    \par\medskip
    \begin{tcolorbox}[
      enhanced, breakable, width=\linewidth, sharp corners,
      colback=white, colframe=mgrey, boxrule=0.8pt,
      left=8pt, right=8pt, top=8pt, bottom=8pt]
      {\bfseries\color{mgrey}\small \lblNotRendered}\\[4pt]
      {\ttfamily\scriptsize\color{mgrey} figures/mermaid/\booklang/#1.mmd}\\[6pt]
      \IfFileExists{figures/mermaid/\booklang/#1.mmd}{%
        % ASCII only, like every other listing in this book: the fallback runs
        % the Mermaid source through `listings`, which aborts on a multi-byte
        % character it has no literate mapping for.
        \lstinputlisting[style=mfacode,numbers=none,basicstyle=\ttfamily\scriptsize,
                         backgroundcolor=\color{mlight},frame=none,
                         keywordstyle=\color{black}]{figures/mermaid/\booklang/#1.mmd}%
      }{%
        {\small\color{mred} \lblSourceMissing: \texttt{figures/mermaid/\booklang/#1.mmd}}%
      }%
    \end{tcolorbox}%
    \captionof{figure}{#2}\label{fig:#1}
    \par\medskip
  }%
}

% ============================================================
%  PROGRAM STUBS
%  ---------------------------------------------------------
%  Prints a visible marker so an unwritten program cannot be mistaken for a
%  written one and the page count never flatters the draft. `make debt` counts
%  them, and CI publishes the count on every build.
% ============================================================
\newcommand{\programstub}[1]{%
  \begin{tcolorbox}[
    enhanced, breakable, width=\linewidth, sharp corners,
    colback=mlight, colframe=mred, boxrule=1pt,
    borderline={0.6pt}{3pt}{mred, dashed},
    left=12pt, right=12pt, top=12pt, bottom=12pt]
    {\bfseries\color{mred}\large \lblNotWritten}\\[8pt]
    \small\raggedright #1
  \end{tcolorbox}%
}

% ============================================================
%  MATHEMATICAL OPERATORS
%  ---------------------------------------------------------
%  Language-invariant names live here. The ones that genuinely differ between
%  Polish and English usage -- tan/tg, Var/D^2, gcd/NWD, and the interval
%  brackets -- are defined in lang/en.tex and lang/pl.tex instead, so the
%  notation contract is executed by the build rather than remembered by a
%  translator.
% ============================================================
\DeclareMathOperator*{\argmin}{arg\,min}
\DeclareMathOperator*{\argmax}{arg\,max}
\DeclareMathOperator{\tr}{tr}
\DeclareMathOperator{\rank}{rank}
\DeclareMathOperator{\diag}{diag}
\DeclareMathOperator{\sgn}{sgn}
\DeclareMathOperator{\relu}{ReLU}
\DeclareMathOperator{\softmax}{softmax}
\DeclareMathOperator{\logsumexp}{logsumexp}
\DeclareMathOperator{\fl}{fl}
\DeclareMathOperator{\ulp}{ulp}
\DeclareMathOperator{\round}{round}

% A bare \log is FORBIDDEN in this book, and tools/check_notation.py fails the
% build on one. In Polish textbooks \log means base ten; in machine-learning
% writing it means base e. The same three characters carry opposite meanings to
% the two readerships this book serves, and they collide hardest in exactly the
% programs where it matters most -- entropy in bits and cross-entropy in nats
% are two programs apart. Write \ln for nats and \logb{2} for bits. It costs
% two characters and removes the ambiguity entirely.
\makeatletter
\let\mfa@origlog\log
\renewcommand{\log}{%
  \PackageError{mfabook}{A bare \string\log\space is forbidden in this book}%
    {Write \string\ln\space for natural logarithms or \string\logb{2}\space for
     bits. The base is never left to a convention here; see Appendix B.}%
  \mfa@origlog}
% The one legitimate form: an explicit base.
\newcommand{\logb}[1]{\mathop{\mfa@origlog}\nolimits_{#1}}
% ...and a way to write the forbidden thing when the subject IS the ban, which
% is Appendix B and nowhere else.
\newcommand{\mfalogplain}{\mathop{\mfa@origlog}\nolimits}
\makeatother

\newcommand{\R}{\mathbb{R}}
\newcommand{\N}{\mathbb{N}}
\newcommand{\Z}{\mathbb{Z}}
\newcommand{\Q}{\mathbb{Q}}
\newcommand{\C}{\mathbb{C}}
\newcommand{\Fset}{\mathbb{F}}          % the floating-point numbers, F1's subject
\newcommand{\vect}[1]{\mathbf{#1}}
\newcommand{\mat}[1]{\mathbf{#1}}
\newcommand{\T}{^{\mathsf{T}}}
\newcommand{\norm}[1]{\left\lVert #1 \right\rVert}
\newcommand{\abs}[1]{\left\lvert #1 \right\rvert}
\newcommand{\inner}[2]{\left\langle #1, #2 \right\rangle}
\newcommand{\eps}{\varepsilon}

% ============================================================
%  INDEX
%  ---------------------------------------------------------
%  Two-column, and its entries include \texttt{} names that do not hyphenate.
%  Ragged right keeps multi-word entries off the margin; it cannot help a
%  single token wider than the column, so keep verbatim index entries under
%  about 29 characters and index longer names by concept instead. Both sibling
%  books learned this the same way.
% ============================================================
\apptocmd{\theindex}{\raggedright}{}{%
  \PackageWarning{mfabook}{Could not patch theindex for ragged-right setting}}

\makeindex

% ============================================================
%  PINNED FACTS — single source of truth
%  Change these here; they propagate through both editions.
% ============================================================
\newcommand{\bookedition}{0.1}
% \pinnedon is a user-visible string and lives in lang/<lang>.tex, not here.
% It said "August 2026" on the Polish title page for exactly as long as it
% took somebody to read the Polish title page.
\newcommand{\pyver}{Python 3.13}
\newcommand{\numpyver}{2.3}
 still compiles and says what it did.
\providecommand{\booklang}{en}

\usepackage{etoolbox}   % loaded first: the language switch below needs it

% ---------- Page geometry ----------
% Same 17 x 24 cm block as the other two books in this series, so a reader who
% owns one recognises the second. Frames are short, so the measure matters
% more here than in a prose book: a frame that runs past a page turn defeats
% the cover-the-next-frame discipline the whole method rests on.
\usepackage[
  paperwidth=17cm, paperheight=24cm,
  inner=2.2cm, outer=1.8cm, top=2.2cm, bottom=2.4cm,
  head=23pt, headsep=0.7cm, footskip=1.2cm
]{geometry}

% ---------- Encoding & fonts ----------
\usepackage[T1]{fontenc}
\usepackage[utf8]{inputenc}
\IfFileExists{lmodern.sty}{\usepackage{lmodern}}{}
\IfFileExists{newtxtext.sty}{\usepackage{newtxtext}}{}
% amsmath goes here, not down with the other packages, because newtxmath
% patches amsmath's internals and must be loaded AFTER it. The wrong order does
% not fail on a machine without newtxmath -- which is most CI images -- and
% fails on the author's full TeX Live.
%
% amssymb only when newtxmath is absent: newtxmath supplies the AMS symbols
% itself (its amssymbols option is on by default) and loading both clashes.
\usepackage{amsmath}
\IfFileExists{newtxmath.sty}{\usepackage{newtxmath}}{\usepackage{amssymb}}
\IfFileExists{inconsolata.sty}{\usepackage[varqu,scaled=0.95]{inconsolata}}{}
\IfFileExists{lmodern.sty}{\usepackage{microtype}}{\usepackage[expansion=false]{microtype}}

% ---------- Language ----------
% Loading babel with a language whose .ldf is absent is a *fatal* error, not a
% warning, so both the package and each language file are probed before either
% is requested. This is inherited from the sibling books, where a minimal TeX
% installation without texlive-lang-polish aborted the run with no PDF and an
% error message that named neither babel nor the language.
%
% The Polish edition needs Polish as the main language for hyphenation; the
% English edition needs British. Both editions load the other language too,
% because the glossary appendix is bilingual in both.
\IfFileExists{babel.sty}{%
  \IfFileExists{polish.ldf}{%
    \IfFileExists{british.ldf}{%
      \ifdefstring{\booklang}{pl}%
        {\usepackage[british,main=polish]{babel}}%
        {\usepackage[polish,main=british]{babel}}%
    }{\usepackage[polish]{babel}}%
  }{%
    \IfFileExists{british.ldf}{\usepackage[british]{babel}}{}%
  }%
}{}

% babel-polish makes " an active shorthand. listings reads its body verbatim
% and an active character in a verbatim context is a category-code accident
% waiting to happen, so the shorthand is turned off globally. Polish quotation
% marks are produced with \enquote-style macros from lang/pl.tex instead.
\ifdefstring{\booklang}{pl}{\AtBeginDocument{\shorthandoff{"}}}{}

% ---------- Core packages ----------
\usepackage{graphicx}
\usepackage{xcolor}
\usepackage{booktabs}
\usepackage{tabularx}
\usepackage{longtable}
\usepackage{array}
\usepackage{enumitem}
\IfFileExists{mathtools.sty}{\usepackage{mathtools}}{}
\usepackage{listings}
\usepackage[most]{tcolorbox}
\usepackage{fancyhdr}
\usepackage{titlesec}
\usepackage{caption}
\usepackage{imakeidx}
% csquotes with autostyle follows babel's active language, so the identical
% source \enquote{x} sets 'x' in the English edition and the Polish quotation
% marks in the Polish one. This is not a nicety: babel-polish makes " an active
% shorthand, which is turned off above precisely because an active character
% near a listings body is a category-code accident waiting to happen -- so a
% quotation mark has to come from a macro rather than from a keyboard key.
\IfFileExists{csquotes.sty}{\usepackage[autostyle=true]{csquotes}}{%
  \providecommand{\enquote}[1]{``##1''}}
\usepackage[hidelinks]{hyperref}

% ---------- Palette ----------
% Deliberately the same family as the sibling books, shifted green so the three
% spines are distinguishable on a shelf.
\definecolor{mblue}{HTML}{1F4E79}
\definecolor{mgreen}{HTML}{2E6B4F}
\definecolor{mteal}{HTML}{0E7C7B}
\definecolor{mamber}{HTML}{B26A00}
\definecolor{mred}{HTML}{9B2C2C}
\definecolor{mgrey}{HTML}{5A5A5A}
\definecolor{mlight}{HTML}{F4F6F8}
\definecolor{mfaint}{HTML}{FBFCFD}
\definecolor{codebg}{HTML}{FAFAFA}
\definecolor{codeframe}{HTML}{D8DEE4}
\definecolor{kw}{HTML}{0B5394}
\definecolor{str}{HTML}{9B2C2C}
\definecolor{cmt}{HTML}{4F7A4F}
\definecolor{typ}{HTML}{0E7C7B}

\hypersetup{
  colorlinks=true,
  linkcolor=mblue,
  urlcolor=mteal,
  citecolor=mblue,
  pdfauthor={Konrad Cinkusz}
}

% \ifpl{Polish text}{English text} -- for the handful of places where a string
% is structural rather than content and belongs in the shared files:
% structure.tex's part titles, and the main files' front matter wiring.
% Anything longer than a few words belongs in lang/ or in the edition's own
% source, not here.
% \DeclareRobustCommand, not \newcommand. A fragile \ifpl inside a \part title
% is expanded one level as it is written to the .toc, which lands
% "\ifdefstring {en}{pl}{...}{...}" in the file -- and that fails on the next
% run with "Extra }, or forgotten \endgroup", an error that names neither the
% part nor this macro.
\DeclareRobustCommand{\ifpl}[2]{\ifdefstring{\booklang}{pl}{#1}{#2}}

% A robust macro cannot go into a PDF bookmark string either -- hyperref strips
% \ifdefstring and warns once per part title. Tell it which branch to take
% instead, so the bookmarks carry the right language rather than nothing.
\makeatletter
\pdfstringdefDisableCommands{%
  \ifdefstring{\booklang}{pl}{\let\ifpl\@firstoftwo}{\let\ifpl\@secondoftwo}%
}
\makeatother

% ---------- Language strings ----------
% Every user-visible word the macros below typeset comes from here. Nothing in
% programs/ or appendices/ may hard-code an English or a Polish word inside a
% macro argument that the other edition also uses.
\input{lang/\booklang}

% ============================================================
%  PROGRAM NUMBERING
%  ---------------------------------------------------------
%  Part F programs are F1..F13; the main programs restart at 1. That is
%  Stroud's scheme and it is worth reproducing, because the Foundation part is
%  explicitly optional and numbering it separately is what makes skipping it
%  feel permitted rather than delinquent.
%
%  \theHchapter is hyperref's *destination* name, and it must stay unique
%  across the whole document. Resetting the chapter counter without also
%  changing \theHchapter produces duplicate PDF destinations, a pile of
%  warnings, and bookmarks that jump to the wrong program.
% ============================================================
\newcommand{\foundationnumbering}{%
  \setcounter{chapter}{0}%
  \renewcommand{\thechapter}{F\arabic{chapter}}%
  \renewcommand{\theHchapter}{F\arabic{chapter}}%
}
\newcommand{\mainnumbering}{%
  \setcounter{chapter}{0}%
  \renewcommand{\thechapter}{\arabic{chapter}}%
  \renewcommand{\theHchapter}{M\arabic{chapter}}%
}

% \appendix resets the chapter counter and switches \thechapter to letters, but
% it knows nothing about \theHchapter, so appendix A would claim the same PDF
% destination as main program 1. Patch it once, here.
\apptocmd{\appendix}{\renewcommand{\theHchapter}{A\Alph{chapter}}}{}{%
  \PackageWarning{mfabook}{Could not patch \string\appendix\space for hyperref destinations}}

% A program is a chapter. \chaptername is set from the language file so the
% word above the title is "Program" in both editions (it happens to be the
% same word, which is a small piece of luck).
\AtBeginDocument{\renewcommand{\chaptername}{\lblProgram}}

% ---------- Headings ----------
% \raggedright on the title body is load-bearing, not cosmetic: at \Huge on a
% 13 cm text block a long title that cannot break overflows the margin badly.
\titleformat{\chapter}[display]
  {\normalfont\huge\bfseries\color{mblue}\raggedright}
  {\normalfont\normalsize\color{mgrey}\MakeUppercase{\chaptertitlename}\ \thechapter}
  {8pt}{\Huge\raggedright}
\titlespacing*{\chapter}{0pt}{10pt}{28pt}
\titleformat{\section}{\normalfont\Large\bfseries\color{mblue}}{\thesection}{0.8em}{}
\titleformat{\subsection}{\normalfont\large\bfseries\color{mgrey}}{\thesubsection}{0.7em}{}

% head=23pt in the geometry options above, not \setlength{\headheight}:
% fancyhdr computes the height it needs from the running-head font and the
% rule, wants 22.5 pt at this size, and silently overprints the top of the text
% block if it does not get it. Setting it through geometry keeps the text block
% where geometry put it instead of pushing it down.
\pagestyle{fancy}
\fancyhf{}
\fancyhead[LE]{\small\nouppercase{\leftmark}}
% The recto head carries the FRAME number, not the section. In a book of
% numbered frames "where am I" is a frame question: every Summary
% back-reference, every Quiz route and every cross-reference between programs
% is a frame number, and hunting for one by page is the single most annoying
% thing about reading a programmed text. \theframeno at shipout is the last
% frame begun on the page, which is what a reader flicking through wants.
% Guarded: the front matter and the appendices have no frames, and a head
% reading "Frame 0" over the introduction is worse than no head at all.
\fancyhead[RO]{\ifnum\value{frameno}>0\relax\small\lblFrame~\theframeno\fi}
\fancyfoot[LE,RO]{\small\thepage}
\renewcommand{\headrulewidth}{0.4pt}

% This MUST come after \pagestyle{fancy}: fancyhdr installs its own \chaptermark
% at that point and silently overwrites anything defined earlier. The symptom is
% a running-head change that appears to do nothing at all.
%
% The running head carries the *frame* number as well as the program, because
% in a book of numbered frames "where am I" is a frame question, not a page
% question, and the summary's back-references are frame numbers.
% \@chapapp, not \lblProgram. \appendix switches \@chapapp from \chaptername
% to \appendixname, and babel supplies both in the right language; hard-coding
% the word gives an appendix a running head reading "Program A. Answers" over a
% page whose own chapter head reads "APPENDIX A".
\makeatletter
\renewcommand{\chaptermark}[1]{\markboth{\@chapapp\ \thechapter.\ #1}{}}
\makeatother
\renewcommand{\sectionmark}[1]{\markright{\thesection\ #1}}

% ============================================================
%  FRAMES — the unit the whole method is built from
%  ---------------------------------------------------------
%  A program is a stream of numbered frames. Nearly every frame ends by asking
%  the reader for something, and the answer opens the *next* frame, so that the
%  reader can cover what is below and commit to an answer before seeing it.
%
%  The environment is called `fr` rather than `frame` for two reasons. First,
%  \frame is already defined by LaTeX2e (it draws a box in a picture
%  environment) and an environment named `frame` would silently redefine it.
%  Second, this is typed forty to seventy times per program, so it wants to be
%  short in the source.
%
%    \begin{fr}
%    \ans{$x = 3$}            % the answer to the previous frame, optional
%    Teaching text, ending in a question.
%    \end{fr}
% ============================================================

\newcounter{frameno}[chapter]
\renewcommand{\theframeno}{\arabic{frameno}}

% Frames are referenced constantly -- by the Summary, by the Quiz, by the
% "Can you?" checklist and by other programs. \label inside a frame must
% therefore capture the frame number, not the section number.
\makeatletter
\newcommand{\framelabelfix}{%
  \@namedef{theHframeno}{\theHchapter.\arabic{frameno}}%
}
\makeatother

\newcommand{\framerule}{%
  \par\penalty-200\vspace{9pt}%
  \noindent\textcolor{codeframe}{\rule{\linewidth}{0.5pt}}%
  \par\vspace{5pt}%
}

% The frame number as a filled badge, set in the text flow rather than in the
% margin. A margin number would have to alternate sides under `twoside` and
% would be the first thing to overflow on a 17 cm page; a badge cannot.
\newcommand{\framebadge}{%
  \begingroup
  \setlength{\fboxsep}{2.5pt}%
  \colorbox{mblue}{\textcolor{white}{\bfseries\scriptsize\theframeno}}%
  \endgroup\hspace{0.7em}%
}

\newenvironment{fr}{%
  \framerule
  \refstepcounter{frameno}%
  \nopagebreak
  \noindent\framebadge\ignorespaces
}{\par}

% The answer to the previous frame. Set apart and centred so the eye can find
% the edge to cover, and so a reader who has answered can check in one glance
% without reading on.
\newcommand{\ans}[1]{%
  \par\nobreak\vspace{-4pt}%
  \begin{tcolorbox}[
    enhanced, sharp corners, boxrule=0pt, leftrule=2.5pt,
    colback=mfaint, colframe=mgreen,
    left=8pt, right=8pt, top=4pt, bottom=4pt,
    width=0.86\linewidth, center, halign=center,
    before skip=2pt, after skip=7pt]
    #1
  \end{tcolorbox}%
  \nobreak\noindent\ignorespaces
}

% A multi-paragraph or worked answer, where a centred box would be wrong.
\newenvironment{ansblock}{%
  \par\nopagebreak
  \begin{tcolorbox}[
    enhanced, breakable, sharp corners, boxrule=0pt, leftrule=2.5pt,
    colback=mfaint, colframe=mgreen,
    left=8pt, right=8pt, top=5pt, bottom=5pt]
}{%
  \end{tcolorbox}%
  \nopagebreak\par\vspace{2pt}\noindent\ignorespaces
}

% The row of dots. Stroud's signal that the reader is to supply a word, a
% number or an expression. \blank is inline; \dotline occupies its own line,
% which is what a longer derivation step wants.
% \penalty0 offers TeX a break opportunity immediately before the dots. Without
% it the run of dots is glued to the last word of the question, and a long
% question in the Polish edition -- where the words are longer and hyphenate
% less freely -- runs into the margin.
\newcommand{\blank}[1][4em]{\penalty0\makebox[#1]{\dotfill}}
\newcommand{\dotline}{\par\vspace{2pt}\noindent\makebox[\linewidth]{\dotfill}\par\vspace{2pt}}

% "Now one for you to do." The third rung of the scaffolding gradient, and the
% one most often skipped by authors who think they have written an exercise
% when they have written another worked example.
\newcommand{\yourturn}{%
  \par\vspace{4pt}\noindent
  {\bfseries\color{mgreen}\lblYourTurn}\par\nobreak\vspace{2pt}\noindent\ignorespaces
}

% ============================================================
%  THE PROGRAM SKELETON
%  ---------------------------------------------------------
%  Every program carries the same seven parts, in the same order. They are
%  macros rather than a convention so that a missing one is a build-time
%  absence rather than something a reader notices.
% ============================================================

% ---- Learning outcomes, on entry ----------------------------------------
% Stored as well as printed, because "Can you?" at the end of the program must
% be the same list, one for one. Storing it is what makes that structural
% rather than a promise: the checklist is generated from the outcomes, so the
% two cannot drift.
\newcounter{outcomeno}[chapter]

\makeatletter
% The key under which everything belonging to the current program is stored.
% It must expand to the printed program number (F1, F2, ..., 1, 2, ...), which
% is what makes an answer findable from the program that owns it.
\newcommand{\mfa@progkey}{\thechapter}
\newcommand{\outcome}[2]{% #1 = frames that teach it, #2 = the outcome
  \stepcounter{outcomeno}%
  \csgdef{mfa@out@\mfa@progkey @\theoutcomeno}{#2}%
  \csgdef{mfa@outfr@\mfa@progkey @\theoutcomeno}{#1}%
  % The "Can you?" row is built here rather than looped over later. A \loop
  % inside a tabularx body does not survive alignment scanning: TeX needs to
  % see the & and the row break while it reads the cell, and a conditional
  % hides them. Accumulating the rows as tokens at declaration time sidesteps
  % that entirely, and has the better property anyway -- the checklist is
  % literally made of the outcomes, so the two cannot disagree.
  \csgappto{mfa@rows@\mfa@progkey}{%
    #2 & {\footnotesize\color{mgrey}#1}
      & \ratingcell & \ratingcell & \ratingcell & \ratingcell & \ratingcell
    \tabularnewline}%
  \item #2\hfill{\footnotesize\color{mgrey}[\lblFrames~#1]}%
}
\makeatother

\makeatletter
\newenvironment{outcomes}{%
  \setcounter{outcomeno}{0}%
  \csgdef{mfa@rows@\mfa@progkey}{}%
  \begin{tcolorbox}[
    enhanced, breakable, sharp corners,
    colback=mlight, colframe=mblue, boxrule=0pt, leftrule=3pt,
    left=8pt, right=8pt, top=6pt, bottom=6pt,
    fonttitle=\bfseries\small, title={\lblOutcomes}]
  \begin{itemize}[leftmargin=1.3em, itemsep=3pt, topsep=2pt]
}{%
  \end{itemize}
  \end{tcolorbox}
}
\makeatother

% ---- "Can you?" -----------------------------------------------------------
% Generated from the stored outcomes, not written out again by hand. Five
% columns rather than a tick box, because a self-assessment with a midpoint is
% one people actually use and a yes/no one is one they lie to.
\newcommand{\ratingcell}{\textcolor{codeframe}{\rule[-0.15em]{0.8em}{0.8em}}}

\makeatletter
\newcommand{\canyou}{%
  \section*{\lblCanYou}%
  \phantomsection\addcontentsline{toc}{section}{\lblCanYou}%
  \noindent\lblCanYouIntro\par\vspace{8pt}
  \noindent
  \begingroup\small
  \ifcsvoid{mfa@rows@\mfa@progkey}{%
    {\color{mred}\lblNoOutcomes}\par
  }{%
    \begin{tabularx}{\linewidth}{@{}Xcccccc@{}}
    \toprule
    \textbf{\lblOnEntryYouCould} & \textbf{\lblFrames} & 1 & 2 & 3 & 4 & 5 \\
    \midrule
    \csuse{mfa@rows@\mfa@progkey}
    \bottomrule
    \end{tabularx}%
  }%
  \endgroup
  \par\vspace{6pt}\noindent{\footnotesize\color{mgrey}\lblCanYouFooter}\par
}
\makeatother

% ============================================================
%  QUIZ, SUMMARY, EXERCISES
% ============================================================

% ---- Quiz -----------------------------------------------------------------
% Foundation programs only. It is a diagnostic on the way in and the same test
% on the way out, which is what lets one book serve a reader who needs the
% whole program and a reader who needs four frames of it. Each question names
% the frames that teach it, so a wrong answer routes the reader somewhere
% specific rather than back to the beginning.
\newlist{quizlist}{enumerate}{1}
\setlist[quizlist]{label=\arabic*., leftmargin=2.0em, itemsep=5pt, topsep=3pt}

\makeatletter
\newenvironment{quiz}{%
  \section*{\lblQuiz}%
  \phantomsection\addcontentsline{toc}{section}{\lblQuiz}%
  \def\mfa@exkey{Q\arabic{quizlisti}}%
  \noindent\lblQuizIntro\par\vspace{5pt}
  \begin{quizlist}
}{%
  \end{quizlist}
}
\makeatother

% The frames that teach the question just asked. Two forms, because a Quiz
% question that routes to a single frame reading "Frames 22" is the kind of
% small wrongness a reader notices and an author never does.
\newcommand{\teachesat}[1]{\hfill{\footnotesize\color{mgrey}(\lblFrames~#1)}}
\newcommand{\teachesatone}[1]{\hfill{\footnotesize\color{mgrey}(\lblFrame~#1)}}

% ---- Summary --------------------------------------------------------------
% Every item carries the frame number it came from, in square brackets. That
% is the eighth-edition improvement and it is the single cheapest thing in the
% format: it turns a summary into a return index, so a reader who fails one
% line of "Can you?" knows exactly which four frames to reread.
\newenvironment{summarybox}{%
  \section*{\lblSummary}%
  \phantomsection\addcontentsline{toc}{section}{\lblSummary}%
  \begin{tcolorbox}[
    enhanced, breakable, sharp corners,
    colback=mlight, colframe=mblue, boxrule=0pt, leftrule=3pt,
    left=8pt, right=8pt, top=6pt, bottom=6pt]
  \begin{itemize}[leftmargin=1.3em, itemsep=5pt, topsep=2pt]
}{%
  \end{itemize}
  \end{tcolorbox}
}

% \sumitem{12}{text} -- the frame reference is not decoration, it is the point.
\newcommand{\sumitem}[2]{\item #2~{\footnotesize\color{mgrey}[#1]}}

% ---- Answers ---------------------------------------------------------------
%  Answers are authored at the point of use and printed at the back.
%
%  They are carried there by a global macro store rather than by an auxiliary
%  file. The sibling books collect their figure and screenshot manifests with
%  \@starttoc and a custom file extension, which works because those entries
%  are one line of text; it is the wrong mechanism for an answer, because an
%  answer contains displayed maths, and material written to an .aux-style file
%  and read back is expanded at shipout, where fragile commands break in ways
%  whose error messages name neither the answer nor the program.
%
%  \include is sequential within a run, so a macro defined globally while
%  typesetting program F1 is still defined when the back matter is typeset.
%  That is all this needs.
\makeatletter
\newcommand{\mfa@storeanswer}[3]{% #1 program key, #2 item key, #3 body
  \csgdef{mfa@a@#1@#2}{#3}%
  % \listcsxadd, not \listcsgadd: the item must be EXPANDED as it is stored.
  % Storing the token \mfa@exkey means every entry in the list expands to
  % whatever the key happens to be at the time the back matter is typeset,
  % which is the last one.
  \listcsxadd{mfa@akeys@#1}{#2}%
}
% Used inside an exercise item. Prints nothing here; the body reappears at the
% back of the book under this program's heading.
\newcommand{\answerto}[1]{\mfa@storeanswer{\mfa@progkey}{\mfa@exkey}{#1}}
\providecommand{\mfa@exkey}{?}
\makeatother

\newlist{exlist}{enumerate}{1}
\setlist[exlist]{label=\arabic*., leftmargin=2.0em, itemsep=6pt, topsep=3pt}
\newlist{fplist}{enumerate}{1}
\setlist[fplist]{label=\arabic*., leftmargin=2.0em, itemsep=5pt, topsep=3pt}

% These MUST be defined inside \makeatletter. Outside it, \def\mfa@exkey{...}
% parses as \def\mfa @exkey{...} -- a definition of \mfa with a delimited
% parameter text -- which raises no error at all and quietly leaves every
% answer filed under the same key. That failure is invisible until the answers
% section at the back of the book prints one program's answers twice.
\makeatletter
\newenvironment{testexercises}{%
  \section*{\lblTestExercises}%
  \phantomsection\addcontentsline{toc}{section}{\lblTestExercises}%
  \def\mfa@exkey{T\arabic{exlisti}}%
  \noindent\lblTestIntro\par\vspace{5pt}
  \begin{exlist}
}{%
  \end{exlist}
}

\newenvironment{furtherproblems}{%
  \section*{\lblFurther}%
  \phantomsection\addcontentsline{toc}{section}{\lblFurther}%
  \def\mfa@exkey{P\arabic{fplisti}}%
  \noindent\lblFurtherIntro\par\vspace{5pt}
  \begin{fplist}
}{%
  \end{fplist}
}
\makeatother

% ---- The programs register -------------------------------------------------
% \program{Title} opens a program: it is \chapter plus the bookkeeping that
% lets the back matter reproduce this program's answers under its own name.
\makeatletter
\newcommand{\program}[1]{%
  \chapter{#1}%
  \csgdef{mfa@title@\thechapter}{#1}%
  % Expanded, for the same reason as the answer keys above.
  \listxadd{\mfaprograms}{\thechapter}%
}
\newcommand{\mfa@printoneanswer}[2]{% #1 program key, #2 item key
  \item[\textbf{#2}] \csuse{mfa@a@#1@#2}%
}
\newcommand{\mfa@answerprogram}[1]{%
  \subsection*{\lblProgram~#1\quad\textmd{\itshape\csuse{mfa@title@#1}}}%
  \ifcsdef{mfa@akeys@#1}{%
    \begin{list}{}{%
      \setlength{\leftmargin}{3.2em}\setlength{\labelwidth}{2.7em}%
      \setlength{\labelsep}{0.5em}\setlength{\itemsep}{3pt}%
      \setlength{\topsep}{3pt}\setlength{\parsep}{0pt}%
      \renewcommand{\makelabel}[1]{\hfill##1}}
      \forlistcsloop{\mfa@printoneanswer{#1}}{mfa@akeys@#1}
    \end{list}%
  }{%
    \par\noindent{\color{mred}\lblNoAnswers}\par
  }%
}
% The answers body, without a heading of its own, so the appendix that carries
% it can be a numbered appendix like any other and appear in the running heads
% and the ToC in the ordinary way.
\newcommand{\answersbody}{%
  \noindent\lblAnswersIntro\par\vspace{8pt}
  \ifdef{\mfaprograms}{\forlistloop{\mfa@answerprogram}{\mfaprograms}}{%
    \noindent{\color{mred}\lblNoAnswers}\par}%
}
% Unnumbered form, for a draft build that has no appendix wiring yet.
\newcommand{\printanswers}{%
  \chapter*{\lblAnswers}%
  \phantomsection\addcontentsline{toc}{chapter}{\lblAnswers}%
  \markboth{\lblAnswers}{\lblAnswers}%
  \answersbody
}
\makeatother

% ============================================================
%  ADMONITIONS
%  ---------------------------------------------------------
%  A deliberately small, fixed vocabulary. Each one earns its place by doing
%  something no other box does; a seventh would be a smell.
% ============================================================

\newtcolorbox{note}[1][]{
  enhanced, breakable, sharp corners,
  colback=mlight, colframe=mblue, boxrule=0pt, leftrule=3pt,
  left=8pt, right=8pt, top=6pt, bottom=6pt,
  fonttitle=\bfseries\small, title={\lblNote}, #1
}

\newtcolorbox{warning}[1][]{
  enhanced, breakable, sharp corners,
  colback=mlight, colframe=mred, boxrule=0pt, leftrule=3pt,
  left=8pt, right=8pt, top=6pt, bottom=6pt,
  fonttitle=\bfseries\small, title={\lblWarning}, #1
}

% The misconception, stated as a summary after it has been walked into. The
% trap itself belongs in the frames -- elicit the wrong answer, then say
% flatly that it is wrong -- and this box is the thing the reader marks with a
% finger and comes back to.
\newtcolorbox{trapbox}[1][]{
  enhanced, breakable, sharp corners,
  colback=white, colframe=mred, boxrule=0.8pt,
  left=8pt, right=8pt, top=6pt, bottom=6pt,
  fonttitle=\bfseries\small, title={\lblTrap}, #1
}

% Where the mathematics just done shows up in a system the reader has actually
% deployed. This is the box that separates the book from a maths textbook, and
% it is also the one most easily filled with waffle: if it cannot name a
% specific line of a specific system, it should not be written.
\newtcolorbox{aibox}[1][]{
  enhanced, breakable, sharp corners,
  colback=mlight, colframe=mteal, boxrule=0pt, leftrule=3pt,
  left=8pt, right=8pt, top=6pt, bottom=6pt,
  fonttitle=\bfseries\small, title={\lblInAI}, #1
}

% What is being asserted rather than proved, and where to read the proof.
% Stroud's method buys fluency at the cost of rigour, and the honest thing is
% to say so at each point rather than once in an introduction nobody rereads.
\newtcolorbox{rigourbox}[1][]{
  enhanced, breakable, sharp corners,
  colback=white, colframe=mgrey, boxrule=0.6pt,
  left=8pt, right=8pt, top=6pt, bottom=6pt,
  fonttitle=\bfseries\small, title={\lblRigour}, #1
}

% Notation that genuinely differs -- between the Polish and English editions,
% and between disciplines. Matrix calculus alone has two incompatible layout
% conventions, and half the confusion about backpropagation is a transpose
% that came from the other one.
\newtcolorbox{notationbox}[1][]{
  enhanced, breakable, sharp corners,
  colback=mlight, colframe=mamber, boxrule=0pt, leftrule=3pt,
  left=8pt, right=8pt, top=6pt, bottom=6pt,
  fonttitle=\bfseries\small, title={\lblNotation}, #1
}

% The maths analogue of the sibling books' "run this before you trust it".
% A number in this book is either produced by a script under code/ and pulled
% in with \val, or it carries this box. Nothing ships to a reader with one
% still attached.
\newtcolorbox{verifybox}[1][]{
  enhanced, breakable, sharp corners,
  colback=white, colframe=mamber, boxrule=0.8pt,
  left=8pt, right=8pt, top=6pt, bottom=6pt,
  fonttitle=\bfseries\small, title={\lblVerify}, #1
}

\newtcolorbox{exercisebox}[1][]{
  enhanced, breakable, sharp corners,
  colback=white, colframe=mgreen, boxrule=0pt, leftrule=3pt,
  left=8pt, right=8pt, top=6pt, bottom=6pt,
  fonttitle=\bfseries\small, title={\lblExercise}, #1
}

% ============================================================
%  CODE LISTINGS
%  ---------------------------------------------------------
%  Code appears here for one reason only: to check a number. Every listing in
%  this book is runnable and every number it prints is the number in the text.
% ============================================================
\lstdefinelanguage{pythonx}{
  language=Python,
  morekeywords={as, with, assert, match, case},
  morekeywords=[2]{
    float64,float32,float16,bfloat16,int8,uint8,
    np,numpy,array,asarray,frombuffer,nextafter,finfo,iinfo,
    exp,log,log1p,expm1,logaddexp,sum,max,abs,sqrt,mean,var
  },
  sensitive=true
}

\lstdefinestyle{mfacode}{
  basicstyle=\ttfamily\footnotesize,
  keywordstyle=\color{kw}\bfseries,
  keywordstyle=[2]\color{typ},
  stringstyle=\color{str},
  commentstyle=\color{cmt}\itshape,
  numbers=left,
  numberstyle=\ttfamily\tiny\color{mgrey},
  numbersep=8pt,
  backgroundcolor=\color{codebg},
  frame=single,
  rulecolor=\color{codeframe},
  framesep=6pt,
  breaklines=true,
  breakatwhitespace=false,
  postbreak=\mbox{\textcolor{mgrey}{$\hookrightarrow$}\space},
  showstringspaces=false,
  tabsize=4,
  columns=fullflexible,
  keepspaces=true,
  captionpos=b,
  aboveskip=1.2em,
  belowskip=1.2em,
  % Maps the characters most likely to be pasted into a listing by accident.
  %
  % Do NOT add a mapping for U+00A0 (non-breaking space). It looks like the
  % obvious next entry and it makes `listings` abort the run with "Improper
  % alphabetic constant" -- fatal, no PDF, and the message names neither the
  % character nor this line. The ASCII-in-listings convention is the guard
  % against it instead. This trap has now been hit in two of the three books
  % in this series.
  literate={ł}{{\l}}1 {Ł}{{\L}}1 {ą}{{\k a}}1 {ę}{{\k e}}1
           {ś}{{\'s}}1 {ć}{{\'c}}1 {ż}{{\.z}}1 {ź}{{\'z}}1 {ó}{{\'o}}1 {ń}{{\'n}}1
           {Ą}{{\k A}}1 {Ę}{{\k E}}1 {Ś}{{\'S}}1 {Ć}{{\'C}}1
           {Ż}{{\.Z}}1 {Ź}{{\'Z}}1 {Ó}{{\'O}}1 {Ń}{{\'N}}1
           {—}{{-{}-}}2 {–}{{-}}1 {‘}{{'}}1 {’}{{'}}1
           {“}{{"}}1 {”}{{"}}1 {°}{{\textdegree}}1
}

\lstset{style=mfacode}

\lstnewenvironment{python}[1][]
  {\lstset{language=pythonx,style=mfacode,#1}}
  {}

\lstnewenvironment{shellcmd}[1][]
  {\lstset{language=bash,style=mfacode,numbers=none,keywordstyle=\color{mteal}\bfseries,#1}}
  {}

% Terminal transcripts: no line numbers, no keyword colouring. Everything in
% one of these was copied from a real run.
\lstnewenvironment{console}[1][]
  {\lstset{style=mfacode,numbers=none,basicstyle=\ttfamily\scriptsize,
           keywordstyle=\color{black},backgroundcolor=\color{white},#1}}
  {}

\newcommand{\pyfile}[3]{%
  \lstinputlisting[language=pythonx,style=mfacode,caption={#2},label={#3}]{#1}%
}
\newcommand{\pyregion}[4]{%
  \lstinputlisting[
    language=pythonx, style=mfacode,
    linerange={--8<--\ [start:#2]--- 8<--\ [end:#2]},
    includerangemarker=false,
    caption={#3}, label={#4}]{#1}%
}

% Inline literals
\newcommand{\code}[1]{\texttt{\small #1}}
\newcommand{\pkg}[1]{\texttt{\small\color{mblue}#1}}
\newcommand{\api}[1]{\texttt{\small\color{mteal}#1}}

% ============================================================
%  COMPUTED VALUES
%  ---------------------------------------------------------
%  The house rule in the sibling books is "run the listing, or mark it". The
%  analogue here is stronger, because a wrong digit in a maths book is not a
%  broken example, it is a lie the reader will carry.
%
%  Every number that came out of a computation is written by a script under
%  code/ into figures/values/<program>.tex as
%
%      \mfaval{f01.eps64}{2.220446049250313e-16}
%
%  and pulled into the text with \val{f01.eps64}. The script is the source of
%  truth; the book cannot disagree with it, because the book does not contain
%  the digits. A value the scripts no longer produce prints a visible marker
%  and is counted by `make debt`.
%
%  \val also applies the locale's decimal separator, which is how the Polish
%  edition gets its decimal comma without a translator having to remember.
% ============================================================
\IfFileExists{siunitx.sty}{%
  \usepackage{siunitx}%
  \sisetup{
    group-digits = integer,
    group-minimum-digits = 5,
    group-separator = {\,},
    exponent-product = \cdot,
    output-decimal-marker = {\mfadecimalmarker}
  }%
  \newcommand{\mfanum}[1]{\num{##1}}%
}{%
  \newcommand{\mfanum}[1]{##1}%
}

\makeatletter
\newcommand{\mfaval}[2]{\csgdef{mfaval@#1}{#2}\listxadd{\mfavalues}{#1}}
% \num on a non-numeric body is a FATAL siunitx error -- "Invalid number" --
% and under -halt-on-error that is the whole build. A script emitting a word
% where the book expected a number is a plausible mistake.
%
% The check is NOT done here. Deciding in TeX whether a token list is a number
% means parsing it character by character, which is fragile and was got wrong
% once already. The generator knows the answer for free, so it emits \mfaval
% for a number and \mfavaltext for a word, and this end only has to enforce
% which macro reads which.
\newcommand{\val}[1]{%
  \ifcsdef{mfaval@#1}{%
    \ifcsdef{mfavaltext@#1}{%
      \PackageError{mfabook}{Value `#1' is not a number}%
        {\string\val\space passes its body to siunitx, which refuses anything
         that is not a number. Use \string\valtext{#1} instead, or fix the
         script in code/ that emits this key.}%
    }{}%
    \mfanum{\csuse{mfaval@#1}}%
  }{\textcolor{mred}{\textbf{[?\,#1]}}}%
}
% A value that is deliberately a word rather than a number -- a format name, a
% verdict, a unit. Rare, and a visibly different macro so nobody reaches for
% \val and gets an error they cannot place.
\newcommand{\valtext}[1]{%
  \ifcsdef{mfaval@#1}{\csuse{mfaval@#1}}%
    {\textcolor{mred}{\textbf{[?\,#1]}}}%
}
\newcommand{\mfavaltext}[2]{%
  \csgdef{mfaval@#1}{#2}\csgdef{mfavaltext@#1}{}\listxadd{\mfavalues}{#1}%
}

% The raw stored string, for the cases where a number must appear exactly as
% the machine printed it -- inside a console block, or where the digits
% themselves are the point and grouping them would obscure the pattern.
\newcommand{\rawval}[1]{%
  \ifcsdef{mfaval@#1}{\csuse{mfaval@#1}}%
    {\textcolor{mred}{\textbf{[?\,#1]}}}%
}
\makeatother

% figures/values/all.tex is generated by `make numbers` and does nothing but
% \input each program's value file. A build with no values yet still compiles;
% every \val{} in it prints a visible marker instead of a number, which is the
% correct behaviour for a draft and an obvious failure for a release.
\IfFileExists{figures/values/all.tex}{% Generated by `make numbers` --- do not edit.
\input{figures/values/f01}
}{%
  \AtBeginDocument{\PackageWarning{mfabook}{No computed values found. Run `make numbers`.}}}

% ============================================================
%  DIAGRAMS — Mermaid pipeline
%  ---------------------------------------------------------
%  \mermaidfig{key}{caption}{one-line manifest description}
%
%  Source of truth is figures/mermaid/<lang>/<key>.mmd, committed. `make
%  diagrams` renders each to figures/diagrams/<lang>/<key>.pdf, which is build
%  output and is not committed, so a diagram change reviews as a text diff.
%
%  The source is per-language because a diagram with English labels in the
%  Polish edition is exactly the kind of half-translation that makes a
%  bilingual book feel machine-made. CI checks that both languages carry the
%  same set of keys.
%
%  If the rendered PDF is absent the macro draws a placeholder AND typesets the
%  Mermaid source, so a build without mermaid-cli still shows the structure.
% ============================================================
\makeatletter
\newcommand{\listofdiagrams}{%
  \section*{\lblDiagramManifest}%
  \phantomsection\addcontentsline{toc}{section}{\lblDiagramManifest}%
  \@starttoc{dgm}%
}
\makeatother

\newcommand{\mermaidfig}[3]{%
  \addcontentsline{dgm}{subsection}{\protect\numberline{}\texttt{#1.mmd} --- #3}%
  \IfFileExists{figures/diagrams/\booklang/#1.pdf}{%
    \begin{figure}[htbp]
    \centering
    \includegraphics[width=0.95\linewidth,height=0.42\textheight,keepaspectratio]{figures/diagrams/\booklang/#1.pdf}%
    \caption{#2}\label{fig:#1}
    \end{figure}%
  }{%
    % The unrendered case deliberately does NOT float. A Mermaid source
    % typeset verbatim is often taller than a page, and a float cannot break,
    % so wrapping it in `figure` produces an overfull box per diagram and a
    % pile of tcolorbox warnings. In the flow it breaks like any other text.
    \par\medskip
    \begin{tcolorbox}[
      enhanced, breakable, width=\linewidth, sharp corners,
      colback=white, colframe=mgrey, boxrule=0.8pt,
      left=8pt, right=8pt, top=8pt, bottom=8pt]
      {\bfseries\color{mgrey}\small \lblNotRendered}\\[4pt]
      {\ttfamily\scriptsize\color{mgrey} figures/mermaid/\booklang/#1.mmd}\\[6pt]
      \IfFileExists{figures/mermaid/\booklang/#1.mmd}{%
        % ASCII only, like every other listing in this book: the fallback runs
        % the Mermaid source through `listings`, which aborts on a multi-byte
        % character it has no literate mapping for.
        \lstinputlisting[style=mfacode,numbers=none,basicstyle=\ttfamily\scriptsize,
                         backgroundcolor=\color{mlight},frame=none,
                         keywordstyle=\color{black}]{figures/mermaid/\booklang/#1.mmd}%
      }{%
        {\small\color{mred} \lblSourceMissing: \texttt{figures/mermaid/\booklang/#1.mmd}}%
      }%
    \end{tcolorbox}%
    \captionof{figure}{#2}\label{fig:#1}
    \par\medskip
  }%
}

% ============================================================
%  PROGRAM STUBS
%  ---------------------------------------------------------
%  Prints a visible marker so an unwritten program cannot be mistaken for a
%  written one and the page count never flatters the draft. `make debt` counts
%  them, and CI publishes the count on every build.
% ============================================================
\newcommand{\programstub}[1]{%
  \begin{tcolorbox}[
    enhanced, breakable, width=\linewidth, sharp corners,
    colback=mlight, colframe=mred, boxrule=1pt,
    borderline={0.6pt}{3pt}{mred, dashed},
    left=12pt, right=12pt, top=12pt, bottom=12pt]
    {\bfseries\color{mred}\large \lblNotWritten}\\[8pt]
    \small\raggedright #1
  \end{tcolorbox}%
}

% ============================================================
%  MATHEMATICAL OPERATORS
%  ---------------------------------------------------------
%  Language-invariant names live here. The ones that genuinely differ between
%  Polish and English usage -- tan/tg, Var/D^2, gcd/NWD, and the interval
%  brackets -- are defined in lang/en.tex and lang/pl.tex instead, so the
%  notation contract is executed by the build rather than remembered by a
%  translator.
% ============================================================
\DeclareMathOperator*{\argmin}{arg\,min}
\DeclareMathOperator*{\argmax}{arg\,max}
\DeclareMathOperator{\tr}{tr}
\DeclareMathOperator{\rank}{rank}
\DeclareMathOperator{\diag}{diag}
\DeclareMathOperator{\sgn}{sgn}
\DeclareMathOperator{\relu}{ReLU}
\DeclareMathOperator{\softmax}{softmax}
\DeclareMathOperator{\logsumexp}{logsumexp}
\DeclareMathOperator{\fl}{fl}
\DeclareMathOperator{\ulp}{ulp}
\DeclareMathOperator{\round}{round}

% A bare \log is FORBIDDEN in this book, and tools/check_notation.py fails the
% build on one. In Polish textbooks \log means base ten; in machine-learning
% writing it means base e. The same three characters carry opposite meanings to
% the two readerships this book serves, and they collide hardest in exactly the
% programs where it matters most -- entropy in bits and cross-entropy in nats
% are two programs apart. Write \ln for nats and \logb{2} for bits. It costs
% two characters and removes the ambiguity entirely.
\makeatletter
\let\mfa@origlog\log
\renewcommand{\log}{%
  \PackageError{mfabook}{A bare \string\log\space is forbidden in this book}%
    {Write \string\ln\space for natural logarithms or \string\logb{2}\space for
     bits. The base is never left to a convention here; see Appendix B.}%
  \mfa@origlog}
% The one legitimate form: an explicit base.
\newcommand{\logb}[1]{\mathop{\mfa@origlog}\nolimits_{#1}}
% ...and a way to write the forbidden thing when the subject IS the ban, which
% is Appendix B and nowhere else.
\newcommand{\mfalogplain}{\mathop{\mfa@origlog}\nolimits}
\makeatother

\newcommand{\R}{\mathbb{R}}
\newcommand{\N}{\mathbb{N}}
\newcommand{\Z}{\mathbb{Z}}
\newcommand{\Q}{\mathbb{Q}}
\newcommand{\C}{\mathbb{C}}
\newcommand{\Fset}{\mathbb{F}}          % the floating-point numbers, F1's subject
\newcommand{\vect}[1]{\mathbf{#1}}
\newcommand{\mat}[1]{\mathbf{#1}}
\newcommand{\T}{^{\mathsf{T}}}
\newcommand{\norm}[1]{\left\lVert #1 \right\rVert}
\newcommand{\abs}[1]{\left\lvert #1 \right\rvert}
\newcommand{\inner}[2]{\left\langle #1, #2 \right\rangle}
\newcommand{\eps}{\varepsilon}

% ============================================================
%  INDEX
%  ---------------------------------------------------------
%  Two-column, and its entries include \texttt{} names that do not hyphenate.
%  Ragged right keeps multi-word entries off the margin; it cannot help a
%  single token wider than the column, so keep verbatim index entries under
%  about 29 characters and index longer names by concept instead. Both sibling
%  books learned this the same way.
% ============================================================
\apptocmd{\theindex}{\raggedright}{}{%
  \PackageWarning{mfabook}{Could not patch theindex for ragged-right setting}}

\makeindex

% ============================================================
%  PINNED FACTS — single source of truth
%  Change these here; they propagate through both editions.
% ============================================================
\newcommand{\bookedition}{0.1}
% \pinnedon is a user-visible string and lives in lang/<lang>.tex, not here.
% It said "August 2026" on the Polish title page for exactly as long as it
% took somebody to read the Polish title page.
\newcommand{\pyver}{Python 3.13}
\newcommand{\numpyver}{2.3}
 still compiles and says what it did.
\providecommand{\booklang}{en}

\usepackage{etoolbox}   % loaded first: the language switch below needs it

% ---------- Page geometry ----------
% Same 17 x 24 cm block as the other two books in this series, so a reader who
% owns one recognises the second. Frames are short, so the measure matters
% more here than in a prose book: a frame that runs past a page turn defeats
% the cover-the-next-frame discipline the whole method rests on.
\usepackage[
  paperwidth=17cm, paperheight=24cm,
  inner=2.2cm, outer=1.8cm, top=2.2cm, bottom=2.4cm,
  head=23pt, headsep=0.7cm, footskip=1.2cm
]{geometry}

% ---------- Encoding & fonts ----------
\usepackage[T1]{fontenc}
\usepackage[utf8]{inputenc}
\IfFileExists{lmodern.sty}{\usepackage{lmodern}}{}
\IfFileExists{newtxtext.sty}{\usepackage{newtxtext}}{}
% amsmath goes here, not down with the other packages, because newtxmath
% patches amsmath's internals and must be loaded AFTER it. The wrong order does
% not fail on a machine without newtxmath -- which is most CI images -- and
% fails on the author's full TeX Live.
%
% amssymb only when newtxmath is absent: newtxmath supplies the AMS symbols
% itself (its amssymbols option is on by default) and loading both clashes.
\usepackage{amsmath}
\IfFileExists{newtxmath.sty}{\usepackage{newtxmath}}{\usepackage{amssymb}}
\IfFileExists{inconsolata.sty}{\usepackage[varqu,scaled=0.95]{inconsolata}}{}
\IfFileExists{lmodern.sty}{\usepackage{microtype}}{\usepackage[expansion=false]{microtype}}

% ---------- Language ----------
% Loading babel with a language whose .ldf is absent is a *fatal* error, not a
% warning, so both the package and each language file are probed before either
% is requested. This is inherited from the sibling books, where a minimal TeX
% installation without texlive-lang-polish aborted the run with no PDF and an
% error message that named neither babel nor the language.
%
% The Polish edition needs Polish as the main language for hyphenation; the
% English edition needs British. Both editions load the other language too,
% because the glossary appendix is bilingual in both.
\IfFileExists{babel.sty}{%
  \IfFileExists{polish.ldf}{%
    \IfFileExists{british.ldf}{%
      \ifdefstring{\booklang}{pl}%
        {\usepackage[british,main=polish]{babel}}%
        {\usepackage[polish,main=british]{babel}}%
    }{\usepackage[polish]{babel}}%
  }{%
    \IfFileExists{british.ldf}{\usepackage[british]{babel}}{}%
  }%
}{}

% babel-polish makes " an active shorthand. listings reads its body verbatim
% and an active character in a verbatim context is a category-code accident
% waiting to happen, so the shorthand is turned off globally. Polish quotation
% marks are produced with \enquote-style macros from lang/pl.tex instead.
\ifdefstring{\booklang}{pl}{\AtBeginDocument{\shorthandoff{"}}}{}

% ---------- Core packages ----------
\usepackage{graphicx}
\usepackage{xcolor}
\usepackage{booktabs}
\usepackage{tabularx}
\usepackage{longtable}
\usepackage{array}
\usepackage{enumitem}
\IfFileExists{mathtools.sty}{\usepackage{mathtools}}{}
\usepackage{listings}
\usepackage[most]{tcolorbox}
\usepackage{fancyhdr}
\usepackage{titlesec}
\usepackage{caption}
\usepackage{imakeidx}
% csquotes with autostyle follows babel's active language, so the identical
% source \enquote{x} sets 'x' in the English edition and the Polish quotation
% marks in the Polish one. This is not a nicety: babel-polish makes " an active
% shorthand, which is turned off above precisely because an active character
% near a listings body is a category-code accident waiting to happen -- so a
% quotation mark has to come from a macro rather than from a keyboard key.
\IfFileExists{csquotes.sty}{\usepackage[autostyle=true]{csquotes}}{%
  \providecommand{\enquote}[1]{``##1''}}
\usepackage[hidelinks]{hyperref}

% ---------- Palette ----------
% Deliberately the same family as the sibling books, shifted green so the three
% spines are distinguishable on a shelf.
\definecolor{mblue}{HTML}{1F4E79}
\definecolor{mgreen}{HTML}{2E6B4F}
\definecolor{mteal}{HTML}{0E7C7B}
\definecolor{mamber}{HTML}{B26A00}
\definecolor{mred}{HTML}{9B2C2C}
\definecolor{mgrey}{HTML}{5A5A5A}
\definecolor{mlight}{HTML}{F4F6F8}
\definecolor{mfaint}{HTML}{FBFCFD}
\definecolor{codebg}{HTML}{FAFAFA}
\definecolor{codeframe}{HTML}{D8DEE4}
\definecolor{kw}{HTML}{0B5394}
\definecolor{str}{HTML}{9B2C2C}
\definecolor{cmt}{HTML}{4F7A4F}
\definecolor{typ}{HTML}{0E7C7B}

\hypersetup{
  colorlinks=true,
  linkcolor=mblue,
  urlcolor=mteal,
  citecolor=mblue,
  pdfauthor={Konrad Cinkusz}
}

% \ifpl{Polish text}{English text} -- for the handful of places where a string
% is structural rather than content and belongs in the shared files:
% structure.tex's part titles, and the main files' front matter wiring.
% Anything longer than a few words belongs in lang/ or in the edition's own
% source, not here.
% \DeclareRobustCommand, not \newcommand. A fragile \ifpl inside a \part title
% is expanded one level as it is written to the .toc, which lands
% "\ifdefstring {en}{pl}{...}{...}" in the file -- and that fails on the next
% run with "Extra }, or forgotten \endgroup", an error that names neither the
% part nor this macro.
\DeclareRobustCommand{\ifpl}[2]{\ifdefstring{\booklang}{pl}{#1}{#2}}

% A robust macro cannot go into a PDF bookmark string either -- hyperref strips
% \ifdefstring and warns once per part title. Tell it which branch to take
% instead, so the bookmarks carry the right language rather than nothing.
\makeatletter
\pdfstringdefDisableCommands{%
  \ifdefstring{\booklang}{pl}{\let\ifpl\@firstoftwo}{\let\ifpl\@secondoftwo}%
}
\makeatother

% ---------- Language strings ----------
% Every user-visible word the macros below typeset comes from here. Nothing in
% programs/ or appendices/ may hard-code an English or a Polish word inside a
% macro argument that the other edition also uses.
\input{lang/\booklang}

% ============================================================
%  PROGRAM NUMBERING
%  ---------------------------------------------------------
%  Part F programs are F1..F13; the main programs restart at 1. That is
%  Stroud's scheme and it is worth reproducing, because the Foundation part is
%  explicitly optional and numbering it separately is what makes skipping it
%  feel permitted rather than delinquent.
%
%  \theHchapter is hyperref's *destination* name, and it must stay unique
%  across the whole document. Resetting the chapter counter without also
%  changing \theHchapter produces duplicate PDF destinations, a pile of
%  warnings, and bookmarks that jump to the wrong program.
% ============================================================
\newcommand{\foundationnumbering}{%
  \setcounter{chapter}{0}%
  \renewcommand{\thechapter}{F\arabic{chapter}}%
  \renewcommand{\theHchapter}{F\arabic{chapter}}%
}
\newcommand{\mainnumbering}{%
  \setcounter{chapter}{0}%
  \renewcommand{\thechapter}{\arabic{chapter}}%
  \renewcommand{\theHchapter}{M\arabic{chapter}}%
}

% \appendix resets the chapter counter and switches \thechapter to letters, but
% it knows nothing about \theHchapter, so appendix A would claim the same PDF
% destination as main program 1. Patch it once, here.
\apptocmd{\appendix}{\renewcommand{\theHchapter}{A\Alph{chapter}}}{}{%
  \PackageWarning{mfabook}{Could not patch \string\appendix\space for hyperref destinations}}

% A program is a chapter. \chaptername is set from the language file so the
% word above the title is "Program" in both editions (it happens to be the
% same word, which is a small piece of luck).
\AtBeginDocument{\renewcommand{\chaptername}{\lblProgram}}

% ---------- Headings ----------
% \raggedright on the title body is load-bearing, not cosmetic: at \Huge on a
% 13 cm text block a long title that cannot break overflows the margin badly.
\titleformat{\chapter}[display]
  {\normalfont\huge\bfseries\color{mblue}\raggedright}
  {\normalfont\normalsize\color{mgrey}\MakeUppercase{\chaptertitlename}\ \thechapter}
  {8pt}{\Huge\raggedright}
\titlespacing*{\chapter}{0pt}{10pt}{28pt}
\titleformat{\section}{\normalfont\Large\bfseries\color{mblue}}{\thesection}{0.8em}{}
\titleformat{\subsection}{\normalfont\large\bfseries\color{mgrey}}{\thesubsection}{0.7em}{}

% head=23pt in the geometry options above, not \setlength{\headheight}:
% fancyhdr computes the height it needs from the running-head font and the
% rule, wants 22.5 pt at this size, and silently overprints the top of the text
% block if it does not get it. Setting it through geometry keeps the text block
% where geometry put it instead of pushing it down.
\pagestyle{fancy}
\fancyhf{}
\fancyhead[LE]{\small\nouppercase{\leftmark}}
% The recto head carries the FRAME number, not the section. In a book of
% numbered frames "where am I" is a frame question: every Summary
% back-reference, every Quiz route and every cross-reference between programs
% is a frame number, and hunting for one by page is the single most annoying
% thing about reading a programmed text. \theframeno at shipout is the last
% frame begun on the page, which is what a reader flicking through wants.
% Guarded: the front matter and the appendices have no frames, and a head
% reading "Frame 0" over the introduction is worse than no head at all.
\fancyhead[RO]{\ifnum\value{frameno}>0\relax\small\lblFrame~\theframeno\fi}
\fancyfoot[LE,RO]{\small\thepage}
\renewcommand{\headrulewidth}{0.4pt}

% This MUST come after \pagestyle{fancy}: fancyhdr installs its own \chaptermark
% at that point and silently overwrites anything defined earlier. The symptom is
% a running-head change that appears to do nothing at all.
%
% The running head carries the *frame* number as well as the program, because
% in a book of numbered frames "where am I" is a frame question, not a page
% question, and the summary's back-references are frame numbers.
% \@chapapp, not \lblProgram. \appendix switches \@chapapp from \chaptername
% to \appendixname, and babel supplies both in the right language; hard-coding
% the word gives an appendix a running head reading "Program A. Answers" over a
% page whose own chapter head reads "APPENDIX A".
\makeatletter
\renewcommand{\chaptermark}[1]{\markboth{\@chapapp\ \thechapter.\ #1}{}}
\makeatother
\renewcommand{\sectionmark}[1]{\markright{\thesection\ #1}}

% ============================================================
%  FRAMES — the unit the whole method is built from
%  ---------------------------------------------------------
%  A program is a stream of numbered frames. Nearly every frame ends by asking
%  the reader for something, and the answer opens the *next* frame, so that the
%  reader can cover what is below and commit to an answer before seeing it.
%
%  The environment is called `fr` rather than `frame` for two reasons. First,
%  \frame is already defined by LaTeX2e (it draws a box in a picture
%  environment) and an environment named `frame` would silently redefine it.
%  Second, this is typed forty to seventy times per program, so it wants to be
%  short in the source.
%
%    \begin{fr}
%    \ans{$x = 3$}            % the answer to the previous frame, optional
%    Teaching text, ending in a question.
%    \end{fr}
% ============================================================

\newcounter{frameno}[chapter]
\renewcommand{\theframeno}{\arabic{frameno}}

% Frames are referenced constantly -- by the Summary, by the Quiz, by the
% "Can you?" checklist and by other programs. \label inside a frame must
% therefore capture the frame number, not the section number.
\makeatletter
\newcommand{\framelabelfix}{%
  \@namedef{theHframeno}{\theHchapter.\arabic{frameno}}%
}
\makeatother

\newcommand{\framerule}{%
  \par\penalty-200\vspace{9pt}%
  \noindent\textcolor{codeframe}{\rule{\linewidth}{0.5pt}}%
  \par\vspace{5pt}%
}

% The frame number as a filled badge, set in the text flow rather than in the
% margin. A margin number would have to alternate sides under `twoside` and
% would be the first thing to overflow on a 17 cm page; a badge cannot.
\newcommand{\framebadge}{%
  \begingroup
  \setlength{\fboxsep}{2.5pt}%
  \colorbox{mblue}{\textcolor{white}{\bfseries\scriptsize\theframeno}}%
  \endgroup\hspace{0.7em}%
}

\newenvironment{fr}{%
  \framerule
  \refstepcounter{frameno}%
  \nopagebreak
  \noindent\framebadge\ignorespaces
}{\par}

% The answer to the previous frame. Set apart and centred so the eye can find
% the edge to cover, and so a reader who has answered can check in one glance
% without reading on.
\newcommand{\ans}[1]{%
  \par\nobreak\vspace{-4pt}%
  \begin{tcolorbox}[
    enhanced, sharp corners, boxrule=0pt, leftrule=2.5pt,
    colback=mfaint, colframe=mgreen,
    left=8pt, right=8pt, top=4pt, bottom=4pt,
    width=0.86\linewidth, center, halign=center,
    before skip=2pt, after skip=7pt]
    #1
  \end{tcolorbox}%
  \nobreak\noindent\ignorespaces
}

% A multi-paragraph or worked answer, where a centred box would be wrong.
\newenvironment{ansblock}{%
  \par\nopagebreak
  \begin{tcolorbox}[
    enhanced, breakable, sharp corners, boxrule=0pt, leftrule=2.5pt,
    colback=mfaint, colframe=mgreen,
    left=8pt, right=8pt, top=5pt, bottom=5pt]
}{%
  \end{tcolorbox}%
  \nopagebreak\par\vspace{2pt}\noindent\ignorespaces
}

% The row of dots. Stroud's signal that the reader is to supply a word, a
% number or an expression. \blank is inline; \dotline occupies its own line,
% which is what a longer derivation step wants.
% \penalty0 offers TeX a break opportunity immediately before the dots. Without
% it the run of dots is glued to the last word of the question, and a long
% question in the Polish edition -- where the words are longer and hyphenate
% less freely -- runs into the margin.
\newcommand{\blank}[1][4em]{\penalty0\makebox[#1]{\dotfill}}
\newcommand{\dotline}{\par\vspace{2pt}\noindent\makebox[\linewidth]{\dotfill}\par\vspace{2pt}}

% "Now one for you to do." The third rung of the scaffolding gradient, and the
% one most often skipped by authors who think they have written an exercise
% when they have written another worked example.
\newcommand{\yourturn}{%
  \par\vspace{4pt}\noindent
  {\bfseries\color{mgreen}\lblYourTurn}\par\nobreak\vspace{2pt}\noindent\ignorespaces
}

% ============================================================
%  THE PROGRAM SKELETON
%  ---------------------------------------------------------
%  Every program carries the same seven parts, in the same order. They are
%  macros rather than a convention so that a missing one is a build-time
%  absence rather than something a reader notices.
% ============================================================

% ---- Learning outcomes, on entry ----------------------------------------
% Stored as well as printed, because "Can you?" at the end of the program must
% be the same list, one for one. Storing it is what makes that structural
% rather than a promise: the checklist is generated from the outcomes, so the
% two cannot drift.
\newcounter{outcomeno}[chapter]

\makeatletter
% The key under which everything belonging to the current program is stored.
% It must expand to the printed program number (F1, F2, ..., 1, 2, ...), which
% is what makes an answer findable from the program that owns it.
\newcommand{\mfa@progkey}{\thechapter}
\newcommand{\outcome}[2]{% #1 = frames that teach it, #2 = the outcome
  \stepcounter{outcomeno}%
  \csgdef{mfa@out@\mfa@progkey @\theoutcomeno}{#2}%
  \csgdef{mfa@outfr@\mfa@progkey @\theoutcomeno}{#1}%
  % The "Can you?" row is built here rather than looped over later. A \loop
  % inside a tabularx body does not survive alignment scanning: TeX needs to
  % see the & and the row break while it reads the cell, and a conditional
  % hides them. Accumulating the rows as tokens at declaration time sidesteps
  % that entirely, and has the better property anyway -- the checklist is
  % literally made of the outcomes, so the two cannot disagree.
  \csgappto{mfa@rows@\mfa@progkey}{%
    #2 & {\footnotesize\color{mgrey}#1}
      & \ratingcell & \ratingcell & \ratingcell & \ratingcell & \ratingcell
    \tabularnewline}%
  \item #2\hfill{\footnotesize\color{mgrey}[\lblFrames~#1]}%
}
\makeatother

\makeatletter
\newenvironment{outcomes}{%
  \setcounter{outcomeno}{0}%
  \csgdef{mfa@rows@\mfa@progkey}{}%
  \begin{tcolorbox}[
    enhanced, breakable, sharp corners,
    colback=mlight, colframe=mblue, boxrule=0pt, leftrule=3pt,
    left=8pt, right=8pt, top=6pt, bottom=6pt,
    fonttitle=\bfseries\small, title={\lblOutcomes}]
  \begin{itemize}[leftmargin=1.3em, itemsep=3pt, topsep=2pt]
}{%
  \end{itemize}
  \end{tcolorbox}
}
\makeatother

% ---- "Can you?" -----------------------------------------------------------
% Generated from the stored outcomes, not written out again by hand. Five
% columns rather than a tick box, because a self-assessment with a midpoint is
% one people actually use and a yes/no one is one they lie to.
\newcommand{\ratingcell}{\textcolor{codeframe}{\rule[-0.15em]{0.8em}{0.8em}}}

\makeatletter
\newcommand{\canyou}{%
  \section*{\lblCanYou}%
  \phantomsection\addcontentsline{toc}{section}{\lblCanYou}%
  \noindent\lblCanYouIntro\par\vspace{8pt}
  \noindent
  \begingroup\small
  \ifcsvoid{mfa@rows@\mfa@progkey}{%
    {\color{mred}\lblNoOutcomes}\par
  }{%
    \begin{tabularx}{\linewidth}{@{}Xcccccc@{}}
    \toprule
    \textbf{\lblOnEntryYouCould} & \textbf{\lblFrames} & 1 & 2 & 3 & 4 & 5 \\
    \midrule
    \csuse{mfa@rows@\mfa@progkey}
    \bottomrule
    \end{tabularx}%
  }%
  \endgroup
  \par\vspace{6pt}\noindent{\footnotesize\color{mgrey}\lblCanYouFooter}\par
}
\makeatother

% ============================================================
%  QUIZ, SUMMARY, EXERCISES
% ============================================================

% ---- Quiz -----------------------------------------------------------------
% Foundation programs only. It is a diagnostic on the way in and the same test
% on the way out, which is what lets one book serve a reader who needs the
% whole program and a reader who needs four frames of it. Each question names
% the frames that teach it, so a wrong answer routes the reader somewhere
% specific rather than back to the beginning.
\newlist{quizlist}{enumerate}{1}
\setlist[quizlist]{label=\arabic*., leftmargin=2.0em, itemsep=5pt, topsep=3pt}

\makeatletter
\newenvironment{quiz}{%
  \section*{\lblQuiz}%
  \phantomsection\addcontentsline{toc}{section}{\lblQuiz}%
  \def\mfa@exkey{Q\arabic{quizlisti}}%
  \noindent\lblQuizIntro\par\vspace{5pt}
  \begin{quizlist}
}{%
  \end{quizlist}
}
\makeatother

% The frames that teach the question just asked. Two forms, because a Quiz
% question that routes to a single frame reading "Frames 22" is the kind of
% small wrongness a reader notices and an author never does.
\newcommand{\teachesat}[1]{\hfill{\footnotesize\color{mgrey}(\lblFrames~#1)}}
\newcommand{\teachesatone}[1]{\hfill{\footnotesize\color{mgrey}(\lblFrame~#1)}}

% ---- Summary --------------------------------------------------------------
% Every item carries the frame number it came from, in square brackets. That
% is the eighth-edition improvement and it is the single cheapest thing in the
% format: it turns a summary into a return index, so a reader who fails one
% line of "Can you?" knows exactly which four frames to reread.
\newenvironment{summarybox}{%
  \section*{\lblSummary}%
  \phantomsection\addcontentsline{toc}{section}{\lblSummary}%
  \begin{tcolorbox}[
    enhanced, breakable, sharp corners,
    colback=mlight, colframe=mblue, boxrule=0pt, leftrule=3pt,
    left=8pt, right=8pt, top=6pt, bottom=6pt]
  \begin{itemize}[leftmargin=1.3em, itemsep=5pt, topsep=2pt]
}{%
  \end{itemize}
  \end{tcolorbox}
}

% \sumitem{12}{text} -- the frame reference is not decoration, it is the point.
\newcommand{\sumitem}[2]{\item #2~{\footnotesize\color{mgrey}[#1]}}

% ---- Answers ---------------------------------------------------------------
%  Answers are authored at the point of use and printed at the back.
%
%  They are carried there by a global macro store rather than by an auxiliary
%  file. The sibling books collect their figure and screenshot manifests with
%  \@starttoc and a custom file extension, which works because those entries
%  are one line of text; it is the wrong mechanism for an answer, because an
%  answer contains displayed maths, and material written to an .aux-style file
%  and read back is expanded at shipout, where fragile commands break in ways
%  whose error messages name neither the answer nor the program.
%
%  \include is sequential within a run, so a macro defined globally while
%  typesetting program F1 is still defined when the back matter is typeset.
%  That is all this needs.
\makeatletter
\newcommand{\mfa@storeanswer}[3]{% #1 program key, #2 item key, #3 body
  \csgdef{mfa@a@#1@#2}{#3}%
  % \listcsxadd, not \listcsgadd: the item must be EXPANDED as it is stored.
  % Storing the token \mfa@exkey means every entry in the list expands to
  % whatever the key happens to be at the time the back matter is typeset,
  % which is the last one.
  \listcsxadd{mfa@akeys@#1}{#2}%
}
% Used inside an exercise item. Prints nothing here; the body reappears at the
% back of the book under this program's heading.
\newcommand{\answerto}[1]{\mfa@storeanswer{\mfa@progkey}{\mfa@exkey}{#1}}
\providecommand{\mfa@exkey}{?}
\makeatother

\newlist{exlist}{enumerate}{1}
\setlist[exlist]{label=\arabic*., leftmargin=2.0em, itemsep=6pt, topsep=3pt}
\newlist{fplist}{enumerate}{1}
\setlist[fplist]{label=\arabic*., leftmargin=2.0em, itemsep=5pt, topsep=3pt}

% These MUST be defined inside \makeatletter. Outside it, \def\mfa@exkey{...}
% parses as \def\mfa @exkey{...} -- a definition of \mfa with a delimited
% parameter text -- which raises no error at all and quietly leaves every
% answer filed under the same key. That failure is invisible until the answers
% section at the back of the book prints one program's answers twice.
\makeatletter
\newenvironment{testexercises}{%
  \section*{\lblTestExercises}%
  \phantomsection\addcontentsline{toc}{section}{\lblTestExercises}%
  \def\mfa@exkey{T\arabic{exlisti}}%
  \noindent\lblTestIntro\par\vspace{5pt}
  \begin{exlist}
}{%
  \end{exlist}
}

\newenvironment{furtherproblems}{%
  \section*{\lblFurther}%
  \phantomsection\addcontentsline{toc}{section}{\lblFurther}%
  \def\mfa@exkey{P\arabic{fplisti}}%
  \noindent\lblFurtherIntro\par\vspace{5pt}
  \begin{fplist}
}{%
  \end{fplist}
}
\makeatother

% ---- The programs register -------------------------------------------------
% \program{Title} opens a program: it is \chapter plus the bookkeeping that
% lets the back matter reproduce this program's answers under its own name.
\makeatletter
\newcommand{\program}[1]{%
  \chapter{#1}%
  \csgdef{mfa@title@\thechapter}{#1}%
  % Expanded, for the same reason as the answer keys above.
  \listxadd{\mfaprograms}{\thechapter}%
}
\newcommand{\mfa@printoneanswer}[2]{% #1 program key, #2 item key
  \item[\textbf{#2}] \csuse{mfa@a@#1@#2}%
}
\newcommand{\mfa@answerprogram}[1]{%
  \subsection*{\lblProgram~#1\quad\textmd{\itshape\csuse{mfa@title@#1}}}%
  \ifcsdef{mfa@akeys@#1}{%
    \begin{list}{}{%
      \setlength{\leftmargin}{3.2em}\setlength{\labelwidth}{2.7em}%
      \setlength{\labelsep}{0.5em}\setlength{\itemsep}{3pt}%
      \setlength{\topsep}{3pt}\setlength{\parsep}{0pt}%
      \renewcommand{\makelabel}[1]{\hfill##1}}
      \forlistcsloop{\mfa@printoneanswer{#1}}{mfa@akeys@#1}
    \end{list}%
  }{%
    \par\noindent{\color{mred}\lblNoAnswers}\par
  }%
}
% The answers body, without a heading of its own, so the appendix that carries
% it can be a numbered appendix like any other and appear in the running heads
% and the ToC in the ordinary way.
\newcommand{\answersbody}{%
  \noindent\lblAnswersIntro\par\vspace{8pt}
  \ifdef{\mfaprograms}{\forlistloop{\mfa@answerprogram}{\mfaprograms}}{%
    \noindent{\color{mred}\lblNoAnswers}\par}%
}
% Unnumbered form, for a draft build that has no appendix wiring yet.
\newcommand{\printanswers}{%
  \chapter*{\lblAnswers}%
  \phantomsection\addcontentsline{toc}{chapter}{\lblAnswers}%
  \markboth{\lblAnswers}{\lblAnswers}%
  \answersbody
}
\makeatother

% ============================================================
%  ADMONITIONS
%  ---------------------------------------------------------
%  A deliberately small, fixed vocabulary. Each one earns its place by doing
%  something no other box does; a seventh would be a smell.
% ============================================================

\newtcolorbox{note}[1][]{
  enhanced, breakable, sharp corners,
  colback=mlight, colframe=mblue, boxrule=0pt, leftrule=3pt,
  left=8pt, right=8pt, top=6pt, bottom=6pt,
  fonttitle=\bfseries\small, title={\lblNote}, #1
}

\newtcolorbox{warning}[1][]{
  enhanced, breakable, sharp corners,
  colback=mlight, colframe=mred, boxrule=0pt, leftrule=3pt,
  left=8pt, right=8pt, top=6pt, bottom=6pt,
  fonttitle=\bfseries\small, title={\lblWarning}, #1
}

% The misconception, stated as a summary after it has been walked into. The
% trap itself belongs in the frames -- elicit the wrong answer, then say
% flatly that it is wrong -- and this box is the thing the reader marks with a
% finger and comes back to.
\newtcolorbox{trapbox}[1][]{
  enhanced, breakable, sharp corners,
  colback=white, colframe=mred, boxrule=0.8pt,
  left=8pt, right=8pt, top=6pt, bottom=6pt,
  fonttitle=\bfseries\small, title={\lblTrap}, #1
}

% Where the mathematics just done shows up in a system the reader has actually
% deployed. This is the box that separates the book from a maths textbook, and
% it is also the one most easily filled with waffle: if it cannot name a
% specific line of a specific system, it should not be written.
\newtcolorbox{aibox}[1][]{
  enhanced, breakable, sharp corners,
  colback=mlight, colframe=mteal, boxrule=0pt, leftrule=3pt,
  left=8pt, right=8pt, top=6pt, bottom=6pt,
  fonttitle=\bfseries\small, title={\lblInAI}, #1
}

% What is being asserted rather than proved, and where to read the proof.
% Stroud's method buys fluency at the cost of rigour, and the honest thing is
% to say so at each point rather than once in an introduction nobody rereads.
\newtcolorbox{rigourbox}[1][]{
  enhanced, breakable, sharp corners,
  colback=white, colframe=mgrey, boxrule=0.6pt,
  left=8pt, right=8pt, top=6pt, bottom=6pt,
  fonttitle=\bfseries\small, title={\lblRigour}, #1
}

% Notation that genuinely differs -- between the Polish and English editions,
% and between disciplines. Matrix calculus alone has two incompatible layout
% conventions, and half the confusion about backpropagation is a transpose
% that came from the other one.
\newtcolorbox{notationbox}[1][]{
  enhanced, breakable, sharp corners,
  colback=mlight, colframe=mamber, boxrule=0pt, leftrule=3pt,
  left=8pt, right=8pt, top=6pt, bottom=6pt,
  fonttitle=\bfseries\small, title={\lblNotation}, #1
}

% The maths analogue of the sibling books' "run this before you trust it".
% A number in this book is either produced by a script under code/ and pulled
% in with \val, or it carries this box. Nothing ships to a reader with one
% still attached.
\newtcolorbox{verifybox}[1][]{
  enhanced, breakable, sharp corners,
  colback=white, colframe=mamber, boxrule=0.8pt,
  left=8pt, right=8pt, top=6pt, bottom=6pt,
  fonttitle=\bfseries\small, title={\lblVerify}, #1
}

\newtcolorbox{exercisebox}[1][]{
  enhanced, breakable, sharp corners,
  colback=white, colframe=mgreen, boxrule=0pt, leftrule=3pt,
  left=8pt, right=8pt, top=6pt, bottom=6pt,
  fonttitle=\bfseries\small, title={\lblExercise}, #1
}

% ============================================================
%  CODE LISTINGS
%  ---------------------------------------------------------
%  Code appears here for one reason only: to check a number. Every listing in
%  this book is runnable and every number it prints is the number in the text.
% ============================================================
\lstdefinelanguage{pythonx}{
  language=Python,
  morekeywords={as, with, assert, match, case},
  morekeywords=[2]{
    float64,float32,float16,bfloat16,int8,uint8,
    np,numpy,array,asarray,frombuffer,nextafter,finfo,iinfo,
    exp,log,log1p,expm1,logaddexp,sum,max,abs,sqrt,mean,var
  },
  sensitive=true
}

\lstdefinestyle{mfacode}{
  basicstyle=\ttfamily\footnotesize,
  keywordstyle=\color{kw}\bfseries,
  keywordstyle=[2]\color{typ},
  stringstyle=\color{str},
  commentstyle=\color{cmt}\itshape,
  numbers=left,
  numberstyle=\ttfamily\tiny\color{mgrey},
  numbersep=8pt,
  backgroundcolor=\color{codebg},
  frame=single,
  rulecolor=\color{codeframe},
  framesep=6pt,
  breaklines=true,
  breakatwhitespace=false,
  postbreak=\mbox{\textcolor{mgrey}{$\hookrightarrow$}\space},
  showstringspaces=false,
  tabsize=4,
  columns=fullflexible,
  keepspaces=true,
  captionpos=b,
  aboveskip=1.2em,
  belowskip=1.2em,
  % Maps the characters most likely to be pasted into a listing by accident.
  %
  % Do NOT add a mapping for U+00A0 (non-breaking space). It looks like the
  % obvious next entry and it makes `listings` abort the run with "Improper
  % alphabetic constant" -- fatal, no PDF, and the message names neither the
  % character nor this line. The ASCII-in-listings convention is the guard
  % against it instead. This trap has now been hit in two of the three books
  % in this series.
  literate={ł}{{\l}}1 {Ł}{{\L}}1 {ą}{{\k a}}1 {ę}{{\k e}}1
           {ś}{{\'s}}1 {ć}{{\'c}}1 {ż}{{\.z}}1 {ź}{{\'z}}1 {ó}{{\'o}}1 {ń}{{\'n}}1
           {Ą}{{\k A}}1 {Ę}{{\k E}}1 {Ś}{{\'S}}1 {Ć}{{\'C}}1
           {Ż}{{\.Z}}1 {Ź}{{\'Z}}1 {Ó}{{\'O}}1 {Ń}{{\'N}}1
           {—}{{-{}-}}2 {–}{{-}}1 {‘}{{'}}1 {’}{{'}}1
           {“}{{"}}1 {”}{{"}}1 {°}{{\textdegree}}1
}

\lstset{style=mfacode}

\lstnewenvironment{python}[1][]
  {\lstset{language=pythonx,style=mfacode,#1}}
  {}

\lstnewenvironment{shellcmd}[1][]
  {\lstset{language=bash,style=mfacode,numbers=none,keywordstyle=\color{mteal}\bfseries,#1}}
  {}

% Terminal transcripts: no line numbers, no keyword colouring. Everything in
% one of these was copied from a real run.
\lstnewenvironment{console}[1][]
  {\lstset{style=mfacode,numbers=none,basicstyle=\ttfamily\scriptsize,
           keywordstyle=\color{black},backgroundcolor=\color{white},#1}}
  {}

\newcommand{\pyfile}[3]{%
  \lstinputlisting[language=pythonx,style=mfacode,caption={#2},label={#3}]{#1}%
}
\newcommand{\pyregion}[4]{%
  \lstinputlisting[
    language=pythonx, style=mfacode,
    linerange={--8<--\ [start:#2]--- 8<--\ [end:#2]},
    includerangemarker=false,
    caption={#3}, label={#4}]{#1}%
}

% Inline literals
\newcommand{\code}[1]{\texttt{\small #1}}
\newcommand{\pkg}[1]{\texttt{\small\color{mblue}#1}}
\newcommand{\api}[1]{\texttt{\small\color{mteal}#1}}

% ============================================================
%  COMPUTED VALUES
%  ---------------------------------------------------------
%  The house rule in the sibling books is "run the listing, or mark it". The
%  analogue here is stronger, because a wrong digit in a maths book is not a
%  broken example, it is a lie the reader will carry.
%
%  Every number that came out of a computation is written by a script under
%  code/ into figures/values/<program>.tex as
%
%      \mfaval{f01.eps64}{2.220446049250313e-16}
%
%  and pulled into the text with \val{f01.eps64}. The script is the source of
%  truth; the book cannot disagree with it, because the book does not contain
%  the digits. A value the scripts no longer produce prints a visible marker
%  and is counted by `make debt`.
%
%  \val also applies the locale's decimal separator, which is how the Polish
%  edition gets its decimal comma without a translator having to remember.
% ============================================================
\IfFileExists{siunitx.sty}{%
  \usepackage{siunitx}%
  \sisetup{
    group-digits = integer,
    group-minimum-digits = 5,
    group-separator = {\,},
    exponent-product = \cdot,
    output-decimal-marker = {\mfadecimalmarker}
  }%
  \newcommand{\mfanum}[1]{\num{##1}}%
}{%
  \newcommand{\mfanum}[1]{##1}%
}

\makeatletter
\newcommand{\mfaval}[2]{\csgdef{mfaval@#1}{#2}\listxadd{\mfavalues}{#1}}
% \num on a non-numeric body is a FATAL siunitx error -- "Invalid number" --
% and under -halt-on-error that is the whole build. A script emitting a word
% where the book expected a number is a plausible mistake.
%
% The check is NOT done here. Deciding in TeX whether a token list is a number
% means parsing it character by character, which is fragile and was got wrong
% once already. The generator knows the answer for free, so it emits \mfaval
% for a number and \mfavaltext for a word, and this end only has to enforce
% which macro reads which.
\newcommand{\val}[1]{%
  \ifcsdef{mfaval@#1}{%
    \ifcsdef{mfavaltext@#1}{%
      \PackageError{mfabook}{Value `#1' is not a number}%
        {\string\val\space passes its body to siunitx, which refuses anything
         that is not a number. Use \string\valtext{#1} instead, or fix the
         script in code/ that emits this key.}%
    }{}%
    \mfanum{\csuse{mfaval@#1}}%
  }{\textcolor{mred}{\textbf{[?\,#1]}}}%
}
% A value that is deliberately a word rather than a number -- a format name, a
% verdict, a unit. Rare, and a visibly different macro so nobody reaches for
% \val and gets an error they cannot place.
\newcommand{\valtext}[1]{%
  \ifcsdef{mfaval@#1}{\csuse{mfaval@#1}}%
    {\textcolor{mred}{\textbf{[?\,#1]}}}%
}
\newcommand{\mfavaltext}[2]{%
  \csgdef{mfaval@#1}{#2}\csgdef{mfavaltext@#1}{}\listxadd{\mfavalues}{#1}%
}

% The raw stored string, for the cases where a number must appear exactly as
% the machine printed it -- inside a console block, or where the digits
% themselves are the point and grouping them would obscure the pattern.
\newcommand{\rawval}[1]{%
  \ifcsdef{mfaval@#1}{\csuse{mfaval@#1}}%
    {\textcolor{mred}{\textbf{[?\,#1]}}}%
}
\makeatother

% figures/values/all.tex is generated by `make numbers` and does nothing but
% \input each program's value file. A build with no values yet still compiles;
% every \val{} in it prints a visible marker instead of a number, which is the
% correct behaviour for a draft and an obvious failure for a release.
\IfFileExists{figures/values/all.tex}{% Generated by `make numbers` --- do not edit.
% Generated by code/f01_numbers.py --- do not edit.
% Regenerate with `make numbers`; `make verify` fails if this file and
% the script disagree, which is what stops a number in the book drifting
% away from the computation that justifies it.

\mfaval{f01.params}{7000000000}
\mfaval{f01.weights.bytes}{14000000000}
\mfaval{f01.weights.gb}{14}
\mfaval{f01.weights.gib}{13.04}
\mfaval{f01.weights.fp32.gb}{28}
\mfaval{f01.gib.per.gb}{0.9313}
\mfaval{f01.gib.gap.pct}{6.9}
\mfaval{f01.two.ten}{1024}
\mfaval{f01.two.ten.err.pct}{2.40}
\mfaval{f01.two.eighty}{1208925819614629174706176}
\mfaval{f01.two.eighty.err.pct}{20.89}
\mfaval{f01.tokens}{2000000000000}
\mfaval{f01.train.flops}{8.40e22}
\mfaval{f01.train.flops.exp}{22}
\mfaval{f01.device.flops}{4e14}
\mfaval{f01.device.util.pct}{40}
\mfaval{f01.device.days}{6076}
\mfaval{f01.device.years}{16.6}
\mfaval{f01.devices.for.30.days}{203}
\mfaval{f01.infer.flops.per.token}{14000000000}
\mfaval{f01.base.ms}{200}
\mfaval{f01.fifty.pct.less.ms}{100}
\mfaval{f01.fifty.pct.more.rate.ms}{133.3}
\mfaval{f01.speedup.discrepancy.ms}{33.3}

}{%
  \AtBeginDocument{\PackageWarning{mfabook}{No computed values found. Run `make numbers`.}}}

% ============================================================
%  DIAGRAMS — Mermaid pipeline
%  ---------------------------------------------------------
%  \mermaidfig{key}{caption}{one-line manifest description}
%
%  Source of truth is figures/mermaid/<lang>/<key>.mmd, committed. `make
%  diagrams` renders each to figures/diagrams/<lang>/<key>.pdf, which is build
%  output and is not committed, so a diagram change reviews as a text diff.
%
%  The source is per-language because a diagram with English labels in the
%  Polish edition is exactly the kind of half-translation that makes a
%  bilingual book feel machine-made. CI checks that both languages carry the
%  same set of keys.
%
%  If the rendered PDF is absent the macro draws a placeholder AND typesets the
%  Mermaid source, so a build without mermaid-cli still shows the structure.
% ============================================================
\makeatletter
\newcommand{\listofdiagrams}{%
  \section*{\lblDiagramManifest}%
  \phantomsection\addcontentsline{toc}{section}{\lblDiagramManifest}%
  \@starttoc{dgm}%
}
\makeatother

\newcommand{\mermaidfig}[3]{%
  \addcontentsline{dgm}{subsection}{\protect\numberline{}\texttt{#1.mmd} --- #3}%
  \IfFileExists{figures/diagrams/\booklang/#1.pdf}{%
    \begin{figure}[htbp]
    \centering
    \includegraphics[width=0.95\linewidth,height=0.42\textheight,keepaspectratio]{figures/diagrams/\booklang/#1.pdf}%
    \caption{#2}\label{fig:#1}
    \end{figure}%
  }{%
    % The unrendered case deliberately does NOT float. A Mermaid source
    % typeset verbatim is often taller than a page, and a float cannot break,
    % so wrapping it in `figure` produces an overfull box per diagram and a
    % pile of tcolorbox warnings. In the flow it breaks like any other text.
    \par\medskip
    \begin{tcolorbox}[
      enhanced, breakable, width=\linewidth, sharp corners,
      colback=white, colframe=mgrey, boxrule=0.8pt,
      left=8pt, right=8pt, top=8pt, bottom=8pt]
      {\bfseries\color{mgrey}\small \lblNotRendered}\\[4pt]
      {\ttfamily\scriptsize\color{mgrey} figures/mermaid/\booklang/#1.mmd}\\[6pt]
      \IfFileExists{figures/mermaid/\booklang/#1.mmd}{%
        % ASCII only, like every other listing in this book: the fallback runs
        % the Mermaid source through `listings`, which aborts on a multi-byte
        % character it has no literate mapping for.
        \lstinputlisting[style=mfacode,numbers=none,basicstyle=\ttfamily\scriptsize,
                         backgroundcolor=\color{mlight},frame=none,
                         keywordstyle=\color{black}]{figures/mermaid/\booklang/#1.mmd}%
      }{%
        {\small\color{mred} \lblSourceMissing: \texttt{figures/mermaid/\booklang/#1.mmd}}%
      }%
    \end{tcolorbox}%
    \captionof{figure}{#2}\label{fig:#1}
    \par\medskip
  }%
}

% ============================================================
%  PROGRAM STUBS
%  ---------------------------------------------------------
%  Prints a visible marker so an unwritten program cannot be mistaken for a
%  written one and the page count never flatters the draft. `make debt` counts
%  them, and CI publishes the count on every build.
% ============================================================
\newcommand{\programstub}[1]{%
  \begin{tcolorbox}[
    enhanced, breakable, width=\linewidth, sharp corners,
    colback=mlight, colframe=mred, boxrule=1pt,
    borderline={0.6pt}{3pt}{mred, dashed},
    left=12pt, right=12pt, top=12pt, bottom=12pt]
    {\bfseries\color{mred}\large \lblNotWritten}\\[8pt]
    \small\raggedright #1
  \end{tcolorbox}%
}

% ============================================================
%  MATHEMATICAL OPERATORS
%  ---------------------------------------------------------
%  Language-invariant names live here. The ones that genuinely differ between
%  Polish and English usage -- tan/tg, Var/D^2, gcd/NWD, and the interval
%  brackets -- are defined in lang/en.tex and lang/pl.tex instead, so the
%  notation contract is executed by the build rather than remembered by a
%  translator.
% ============================================================
\DeclareMathOperator*{\argmin}{arg\,min}
\DeclareMathOperator*{\argmax}{arg\,max}
\DeclareMathOperator{\tr}{tr}
\DeclareMathOperator{\rank}{rank}
\DeclareMathOperator{\diag}{diag}
\DeclareMathOperator{\sgn}{sgn}
\DeclareMathOperator{\relu}{ReLU}
\DeclareMathOperator{\softmax}{softmax}
\DeclareMathOperator{\logsumexp}{logsumexp}
\DeclareMathOperator{\fl}{fl}
\DeclareMathOperator{\ulp}{ulp}
\DeclareMathOperator{\round}{round}

% A bare \log is FORBIDDEN in this book, and tools/check_notation.py fails the
% build on one. In Polish textbooks \log means base ten; in machine-learning
% writing it means base e. The same three characters carry opposite meanings to
% the two readerships this book serves, and they collide hardest in exactly the
% programs where it matters most -- entropy in bits and cross-entropy in nats
% are two programs apart. Write \ln for nats and \logb{2} for bits. It costs
% two characters and removes the ambiguity entirely.
\makeatletter
\let\mfa@origlog\log
\renewcommand{\log}{%
  \PackageError{mfabook}{A bare \string\log\space is forbidden in this book}%
    {Write \string\ln\space for natural logarithms or \string\logb{2}\space for
     bits. The base is never left to a convention here; see Appendix B.}%
  \mfa@origlog}
% The one legitimate form: an explicit base.
\newcommand{\logb}[1]{\mathop{\mfa@origlog}\nolimits_{#1}}
% ...and a way to write the forbidden thing when the subject IS the ban, which
% is Appendix B and nowhere else.
\newcommand{\mfalogplain}{\mathop{\mfa@origlog}\nolimits}
\makeatother

\newcommand{\R}{\mathbb{R}}
\newcommand{\N}{\mathbb{N}}
\newcommand{\Z}{\mathbb{Z}}
\newcommand{\Q}{\mathbb{Q}}
\newcommand{\C}{\mathbb{C}}
\newcommand{\Fset}{\mathbb{F}}          % the floating-point numbers, F1's subject
\newcommand{\vect}[1]{\mathbf{#1}}
\newcommand{\mat}[1]{\mathbf{#1}}
\newcommand{\T}{^{\mathsf{T}}}
\newcommand{\norm}[1]{\left\lVert #1 \right\rVert}
\newcommand{\abs}[1]{\left\lvert #1 \right\rvert}
\newcommand{\inner}[2]{\left\langle #1, #2 \right\rangle}
\newcommand{\eps}{\varepsilon}

% ============================================================
%  INDEX
%  ---------------------------------------------------------
%  Two-column, and its entries include \texttt{} names that do not hyphenate.
%  Ragged right keeps multi-word entries off the margin; it cannot help a
%  single token wider than the column, so keep verbatim index entries under
%  about 29 characters and index longer names by concept instead. Both sibling
%  books learned this the same way.
% ============================================================
\apptocmd{\theindex}{\raggedright}{}{%
  \PackageWarning{mfabook}{Could not patch theindex for ragged-right setting}}

\makeindex

% ============================================================
%  PINNED FACTS — single source of truth
%  Change these here; they propagate through both editions.
% ============================================================
\newcommand{\bookedition}{0.1}
% \pinnedon is a user-visible string and lives in lang/<lang>.tex, not here.
% It said "August 2026" on the Polish title page for exactly as long as it
% took somebody to read the Polish title page.
\newcommand{\pyver}{Python 3.13}
\newcommand{\numpyver}{2.3}
; nothing else differs between them at
%  this level. Every user-visible string lives in lang/en.tex or lang/pl.tex,
%  so a label that appears in one edition and not the other is a missing macro
%  rather than a silent divergence.
% ============================================================

% \booklang must be set by the main file. Default to English rather than
% failing, so a stray % ============================================================
%  preamble.tex — styling, environments, macros
%
%  Shared by BOTH editions. main-en.tex and main-pl.tex each define
%  \booklang before % ============================================================
%  preamble.tex — styling, environments, macros
%
%  Shared by BOTH editions. main-en.tex and main-pl.tex each define
%  \booklang before % ============================================================
%  preamble.tex — styling, environments, macros
%
%  Shared by BOTH editions. main-en.tex and main-pl.tex each define
%  \booklang before \input{preamble}; nothing else differs between them at
%  this level. Every user-visible string lives in lang/en.tex or lang/pl.tex,
%  so a label that appears in one edition and not the other is a missing macro
%  rather than a silent divergence.
% ============================================================

% \booklang must be set by the main file. Default to English rather than
% failing, so a stray \input{preamble} still compiles and says what it did.
\providecommand{\booklang}{en}

\usepackage{etoolbox}   % loaded first: the language switch below needs it

% ---------- Page geometry ----------
% Same 17 x 24 cm block as the other two books in this series, so a reader who
% owns one recognises the second. Frames are short, so the measure matters
% more here than in a prose book: a frame that runs past a page turn defeats
% the cover-the-next-frame discipline the whole method rests on.
\usepackage[
  paperwidth=17cm, paperheight=24cm,
  inner=2.2cm, outer=1.8cm, top=2.2cm, bottom=2.4cm,
  head=23pt, headsep=0.7cm, footskip=1.2cm
]{geometry}

% ---------- Encoding & fonts ----------
\usepackage[T1]{fontenc}
\usepackage[utf8]{inputenc}
\IfFileExists{lmodern.sty}{\usepackage{lmodern}}{}
\IfFileExists{newtxtext.sty}{\usepackage{newtxtext}}{}
% amsmath goes here, not down with the other packages, because newtxmath
% patches amsmath's internals and must be loaded AFTER it. The wrong order does
% not fail on a machine without newtxmath -- which is most CI images -- and
% fails on the author's full TeX Live.
%
% amssymb only when newtxmath is absent: newtxmath supplies the AMS symbols
% itself (its amssymbols option is on by default) and loading both clashes.
\usepackage{amsmath}
\IfFileExists{newtxmath.sty}{\usepackage{newtxmath}}{\usepackage{amssymb}}
\IfFileExists{inconsolata.sty}{\usepackage[varqu,scaled=0.95]{inconsolata}}{}
\IfFileExists{lmodern.sty}{\usepackage{microtype}}{\usepackage[expansion=false]{microtype}}

% ---------- Language ----------
% Loading babel with a language whose .ldf is absent is a *fatal* error, not a
% warning, so both the package and each language file are probed before either
% is requested. This is inherited from the sibling books, where a minimal TeX
% installation without texlive-lang-polish aborted the run with no PDF and an
% error message that named neither babel nor the language.
%
% The Polish edition needs Polish as the main language for hyphenation; the
% English edition needs British. Both editions load the other language too,
% because the glossary appendix is bilingual in both.
\IfFileExists{babel.sty}{%
  \IfFileExists{polish.ldf}{%
    \IfFileExists{british.ldf}{%
      \ifdefstring{\booklang}{pl}%
        {\usepackage[british,main=polish]{babel}}%
        {\usepackage[polish,main=british]{babel}}%
    }{\usepackage[polish]{babel}}%
  }{%
    \IfFileExists{british.ldf}{\usepackage[british]{babel}}{}%
  }%
}{}

% babel-polish makes " an active shorthand. listings reads its body verbatim
% and an active character in a verbatim context is a category-code accident
% waiting to happen, so the shorthand is turned off globally. Polish quotation
% marks are produced with \enquote-style macros from lang/pl.tex instead.
\ifdefstring{\booklang}{pl}{\AtBeginDocument{\shorthandoff{"}}}{}

% ---------- Core packages ----------
\usepackage{graphicx}
\usepackage{xcolor}
\usepackage{booktabs}
\usepackage{tabularx}
\usepackage{longtable}
\usepackage{array}
\usepackage{enumitem}
\IfFileExists{mathtools.sty}{\usepackage{mathtools}}{}
\usepackage{listings}
\usepackage[most]{tcolorbox}
\usepackage{fancyhdr}
\usepackage{titlesec}
\usepackage{caption}
\usepackage{imakeidx}
% csquotes with autostyle follows babel's active language, so the identical
% source \enquote{x} sets 'x' in the English edition and the Polish quotation
% marks in the Polish one. This is not a nicety: babel-polish makes " an active
% shorthand, which is turned off above precisely because an active character
% near a listings body is a category-code accident waiting to happen -- so a
% quotation mark has to come from a macro rather than from a keyboard key.
\IfFileExists{csquotes.sty}{\usepackage[autostyle=true]{csquotes}}{%
  \providecommand{\enquote}[1]{``##1''}}
\usepackage[hidelinks]{hyperref}

% ---------- Palette ----------
% Deliberately the same family as the sibling books, shifted green so the three
% spines are distinguishable on a shelf.
\definecolor{mblue}{HTML}{1F4E79}
\definecolor{mgreen}{HTML}{2E6B4F}
\definecolor{mteal}{HTML}{0E7C7B}
\definecolor{mamber}{HTML}{B26A00}
\definecolor{mred}{HTML}{9B2C2C}
\definecolor{mgrey}{HTML}{5A5A5A}
\definecolor{mlight}{HTML}{F4F6F8}
\definecolor{mfaint}{HTML}{FBFCFD}
\definecolor{codebg}{HTML}{FAFAFA}
\definecolor{codeframe}{HTML}{D8DEE4}
\definecolor{kw}{HTML}{0B5394}
\definecolor{str}{HTML}{9B2C2C}
\definecolor{cmt}{HTML}{4F7A4F}
\definecolor{typ}{HTML}{0E7C7B}

\hypersetup{
  colorlinks=true,
  linkcolor=mblue,
  urlcolor=mteal,
  citecolor=mblue,
  pdfauthor={Konrad Cinkusz}
}

% \ifpl{Polish text}{English text} -- for the handful of places where a string
% is structural rather than content and belongs in the shared files:
% structure.tex's part titles, and the main files' front matter wiring.
% Anything longer than a few words belongs in lang/ or in the edition's own
% source, not here.
% \DeclareRobustCommand, not \newcommand. A fragile \ifpl inside a \part title
% is expanded one level as it is written to the .toc, which lands
% "\ifdefstring {en}{pl}{...}{...}" in the file -- and that fails on the next
% run with "Extra }, or forgotten \endgroup", an error that names neither the
% part nor this macro.
\DeclareRobustCommand{\ifpl}[2]{\ifdefstring{\booklang}{pl}{#1}{#2}}

% A robust macro cannot go into a PDF bookmark string either -- hyperref strips
% \ifdefstring and warns once per part title. Tell it which branch to take
% instead, so the bookmarks carry the right language rather than nothing.
\makeatletter
\pdfstringdefDisableCommands{%
  \ifdefstring{\booklang}{pl}{\let\ifpl\@firstoftwo}{\let\ifpl\@secondoftwo}%
}
\makeatother

% ---------- Language strings ----------
% Every user-visible word the macros below typeset comes from here. Nothing in
% programs/ or appendices/ may hard-code an English or a Polish word inside a
% macro argument that the other edition also uses.
\input{lang/\booklang}

% ============================================================
%  PROGRAM NUMBERING
%  ---------------------------------------------------------
%  Part F programs are F1..F13; the main programs restart at 1. That is
%  Stroud's scheme and it is worth reproducing, because the Foundation part is
%  explicitly optional and numbering it separately is what makes skipping it
%  feel permitted rather than delinquent.
%
%  \theHchapter is hyperref's *destination* name, and it must stay unique
%  across the whole document. Resetting the chapter counter without also
%  changing \theHchapter produces duplicate PDF destinations, a pile of
%  warnings, and bookmarks that jump to the wrong program.
% ============================================================
\newcommand{\foundationnumbering}{%
  \setcounter{chapter}{0}%
  \renewcommand{\thechapter}{F\arabic{chapter}}%
  \renewcommand{\theHchapter}{F\arabic{chapter}}%
}
\newcommand{\mainnumbering}{%
  \setcounter{chapter}{0}%
  \renewcommand{\thechapter}{\arabic{chapter}}%
  \renewcommand{\theHchapter}{M\arabic{chapter}}%
}

% \appendix resets the chapter counter and switches \thechapter to letters, but
% it knows nothing about \theHchapter, so appendix A would claim the same PDF
% destination as main program 1. Patch it once, here.
\apptocmd{\appendix}{\renewcommand{\theHchapter}{A\Alph{chapter}}}{}{%
  \PackageWarning{mfabook}{Could not patch \string\appendix\space for hyperref destinations}}

% A program is a chapter. \chaptername is set from the language file so the
% word above the title is "Program" in both editions (it happens to be the
% same word, which is a small piece of luck).
\AtBeginDocument{\renewcommand{\chaptername}{\lblProgram}}

% ---------- Headings ----------
% \raggedright on the title body is load-bearing, not cosmetic: at \Huge on a
% 13 cm text block a long title that cannot break overflows the margin badly.
\titleformat{\chapter}[display]
  {\normalfont\huge\bfseries\color{mblue}\raggedright}
  {\normalfont\normalsize\color{mgrey}\MakeUppercase{\chaptertitlename}\ \thechapter}
  {8pt}{\Huge\raggedright}
\titlespacing*{\chapter}{0pt}{10pt}{28pt}
\titleformat{\section}{\normalfont\Large\bfseries\color{mblue}}{\thesection}{0.8em}{}
\titleformat{\subsection}{\normalfont\large\bfseries\color{mgrey}}{\thesubsection}{0.7em}{}

% head=23pt in the geometry options above, not \setlength{\headheight}:
% fancyhdr computes the height it needs from the running-head font and the
% rule, wants 22.5 pt at this size, and silently overprints the top of the text
% block if it does not get it. Setting it through geometry keeps the text block
% where geometry put it instead of pushing it down.
\pagestyle{fancy}
\fancyhf{}
\fancyhead[LE]{\small\nouppercase{\leftmark}}
% The recto head carries the FRAME number, not the section. In a book of
% numbered frames "where am I" is a frame question: every Summary
% back-reference, every Quiz route and every cross-reference between programs
% is a frame number, and hunting for one by page is the single most annoying
% thing about reading a programmed text. \theframeno at shipout is the last
% frame begun on the page, which is what a reader flicking through wants.
% Guarded: the front matter and the appendices have no frames, and a head
% reading "Frame 0" over the introduction is worse than no head at all.
\fancyhead[RO]{\ifnum\value{frameno}>0\relax\small\lblFrame~\theframeno\fi}
\fancyfoot[LE,RO]{\small\thepage}
\renewcommand{\headrulewidth}{0.4pt}

% This MUST come after \pagestyle{fancy}: fancyhdr installs its own \chaptermark
% at that point and silently overwrites anything defined earlier. The symptom is
% a running-head change that appears to do nothing at all.
%
% The running head carries the *frame* number as well as the program, because
% in a book of numbered frames "where am I" is a frame question, not a page
% question, and the summary's back-references are frame numbers.
% \@chapapp, not \lblProgram. \appendix switches \@chapapp from \chaptername
% to \appendixname, and babel supplies both in the right language; hard-coding
% the word gives an appendix a running head reading "Program A. Answers" over a
% page whose own chapter head reads "APPENDIX A".
\makeatletter
\renewcommand{\chaptermark}[1]{\markboth{\@chapapp\ \thechapter.\ #1}{}}
\makeatother
\renewcommand{\sectionmark}[1]{\markright{\thesection\ #1}}

% ============================================================
%  FRAMES — the unit the whole method is built from
%  ---------------------------------------------------------
%  A program is a stream of numbered frames. Nearly every frame ends by asking
%  the reader for something, and the answer opens the *next* frame, so that the
%  reader can cover what is below and commit to an answer before seeing it.
%
%  The environment is called `fr` rather than `frame` for two reasons. First,
%  \frame is already defined by LaTeX2e (it draws a box in a picture
%  environment) and an environment named `frame` would silently redefine it.
%  Second, this is typed forty to seventy times per program, so it wants to be
%  short in the source.
%
%    \begin{fr}
%    \ans{$x = 3$}            % the answer to the previous frame, optional
%    Teaching text, ending in a question.
%    \end{fr}
% ============================================================

\newcounter{frameno}[chapter]
\renewcommand{\theframeno}{\arabic{frameno}}

% Frames are referenced constantly -- by the Summary, by the Quiz, by the
% "Can you?" checklist and by other programs. \label inside a frame must
% therefore capture the frame number, not the section number.
\makeatletter
\newcommand{\framelabelfix}{%
  \@namedef{theHframeno}{\theHchapter.\arabic{frameno}}%
}
\makeatother

\newcommand{\framerule}{%
  \par\penalty-200\vspace{9pt}%
  \noindent\textcolor{codeframe}{\rule{\linewidth}{0.5pt}}%
  \par\vspace{5pt}%
}

% The frame number as a filled badge, set in the text flow rather than in the
% margin. A margin number would have to alternate sides under `twoside` and
% would be the first thing to overflow on a 17 cm page; a badge cannot.
\newcommand{\framebadge}{%
  \begingroup
  \setlength{\fboxsep}{2.5pt}%
  \colorbox{mblue}{\textcolor{white}{\bfseries\scriptsize\theframeno}}%
  \endgroup\hspace{0.7em}%
}

\newenvironment{fr}{%
  \framerule
  \refstepcounter{frameno}%
  \nopagebreak
  \noindent\framebadge\ignorespaces
}{\par}

% The answer to the previous frame. Set apart and centred so the eye can find
% the edge to cover, and so a reader who has answered can check in one glance
% without reading on.
\newcommand{\ans}[1]{%
  \par\nobreak\vspace{-4pt}%
  \begin{tcolorbox}[
    enhanced, sharp corners, boxrule=0pt, leftrule=2.5pt,
    colback=mfaint, colframe=mgreen,
    left=8pt, right=8pt, top=4pt, bottom=4pt,
    width=0.86\linewidth, center, halign=center,
    before skip=2pt, after skip=7pt]
    #1
  \end{tcolorbox}%
  \nobreak\noindent\ignorespaces
}

% A multi-paragraph or worked answer, where a centred box would be wrong.
\newenvironment{ansblock}{%
  \par\nopagebreak
  \begin{tcolorbox}[
    enhanced, breakable, sharp corners, boxrule=0pt, leftrule=2.5pt,
    colback=mfaint, colframe=mgreen,
    left=8pt, right=8pt, top=5pt, bottom=5pt]
}{%
  \end{tcolorbox}%
  \nopagebreak\par\vspace{2pt}\noindent\ignorespaces
}

% The row of dots. Stroud's signal that the reader is to supply a word, a
% number or an expression. \blank is inline; \dotline occupies its own line,
% which is what a longer derivation step wants.
% \penalty0 offers TeX a break opportunity immediately before the dots. Without
% it the run of dots is glued to the last word of the question, and a long
% question in the Polish edition -- where the words are longer and hyphenate
% less freely -- runs into the margin.
\newcommand{\blank}[1][4em]{\penalty0\makebox[#1]{\dotfill}}
\newcommand{\dotline}{\par\vspace{2pt}\noindent\makebox[\linewidth]{\dotfill}\par\vspace{2pt}}

% "Now one for you to do." The third rung of the scaffolding gradient, and the
% one most often skipped by authors who think they have written an exercise
% when they have written another worked example.
\newcommand{\yourturn}{%
  \par\vspace{4pt}\noindent
  {\bfseries\color{mgreen}\lblYourTurn}\par\nobreak\vspace{2pt}\noindent\ignorespaces
}

% ============================================================
%  THE PROGRAM SKELETON
%  ---------------------------------------------------------
%  Every program carries the same seven parts, in the same order. They are
%  macros rather than a convention so that a missing one is a build-time
%  absence rather than something a reader notices.
% ============================================================

% ---- Learning outcomes, on entry ----------------------------------------
% Stored as well as printed, because "Can you?" at the end of the program must
% be the same list, one for one. Storing it is what makes that structural
% rather than a promise: the checklist is generated from the outcomes, so the
% two cannot drift.
\newcounter{outcomeno}[chapter]

\makeatletter
% The key under which everything belonging to the current program is stored.
% It must expand to the printed program number (F1, F2, ..., 1, 2, ...), which
% is what makes an answer findable from the program that owns it.
\newcommand{\mfa@progkey}{\thechapter}
\newcommand{\outcome}[2]{% #1 = frames that teach it, #2 = the outcome
  \stepcounter{outcomeno}%
  \csgdef{mfa@out@\mfa@progkey @\theoutcomeno}{#2}%
  \csgdef{mfa@outfr@\mfa@progkey @\theoutcomeno}{#1}%
  % The "Can you?" row is built here rather than looped over later. A \loop
  % inside a tabularx body does not survive alignment scanning: TeX needs to
  % see the & and the row break while it reads the cell, and a conditional
  % hides them. Accumulating the rows as tokens at declaration time sidesteps
  % that entirely, and has the better property anyway -- the checklist is
  % literally made of the outcomes, so the two cannot disagree.
  \csgappto{mfa@rows@\mfa@progkey}{%
    #2 & {\footnotesize\color{mgrey}#1}
      & \ratingcell & \ratingcell & \ratingcell & \ratingcell & \ratingcell
    \tabularnewline}%
  \item #2\hfill{\footnotesize\color{mgrey}[\lblFrames~#1]}%
}
\makeatother

\makeatletter
\newenvironment{outcomes}{%
  \setcounter{outcomeno}{0}%
  \csgdef{mfa@rows@\mfa@progkey}{}%
  \begin{tcolorbox}[
    enhanced, breakable, sharp corners,
    colback=mlight, colframe=mblue, boxrule=0pt, leftrule=3pt,
    left=8pt, right=8pt, top=6pt, bottom=6pt,
    fonttitle=\bfseries\small, title={\lblOutcomes}]
  \begin{itemize}[leftmargin=1.3em, itemsep=3pt, topsep=2pt]
}{%
  \end{itemize}
  \end{tcolorbox}
}
\makeatother

% ---- "Can you?" -----------------------------------------------------------
% Generated from the stored outcomes, not written out again by hand. Five
% columns rather than a tick box, because a self-assessment with a midpoint is
% one people actually use and a yes/no one is one they lie to.
\newcommand{\ratingcell}{\textcolor{codeframe}{\rule[-0.15em]{0.8em}{0.8em}}}

\makeatletter
\newcommand{\canyou}{%
  \section*{\lblCanYou}%
  \phantomsection\addcontentsline{toc}{section}{\lblCanYou}%
  \noindent\lblCanYouIntro\par\vspace{8pt}
  \noindent
  \begingroup\small
  \ifcsvoid{mfa@rows@\mfa@progkey}{%
    {\color{mred}\lblNoOutcomes}\par
  }{%
    \begin{tabularx}{\linewidth}{@{}Xcccccc@{}}
    \toprule
    \textbf{\lblOnEntryYouCould} & \textbf{\lblFrames} & 1 & 2 & 3 & 4 & 5 \\
    \midrule
    \csuse{mfa@rows@\mfa@progkey}
    \bottomrule
    \end{tabularx}%
  }%
  \endgroup
  \par\vspace{6pt}\noindent{\footnotesize\color{mgrey}\lblCanYouFooter}\par
}
\makeatother

% ============================================================
%  QUIZ, SUMMARY, EXERCISES
% ============================================================

% ---- Quiz -----------------------------------------------------------------
% Foundation programs only. It is a diagnostic on the way in and the same test
% on the way out, which is what lets one book serve a reader who needs the
% whole program and a reader who needs four frames of it. Each question names
% the frames that teach it, so a wrong answer routes the reader somewhere
% specific rather than back to the beginning.
\newlist{quizlist}{enumerate}{1}
\setlist[quizlist]{label=\arabic*., leftmargin=2.0em, itemsep=5pt, topsep=3pt}

\makeatletter
\newenvironment{quiz}{%
  \section*{\lblQuiz}%
  \phantomsection\addcontentsline{toc}{section}{\lblQuiz}%
  \def\mfa@exkey{Q\arabic{quizlisti}}%
  \noindent\lblQuizIntro\par\vspace{5pt}
  \begin{quizlist}
}{%
  \end{quizlist}
}
\makeatother

% The frames that teach the question just asked. Two forms, because a Quiz
% question that routes to a single frame reading "Frames 22" is the kind of
% small wrongness a reader notices and an author never does.
\newcommand{\teachesat}[1]{\hfill{\footnotesize\color{mgrey}(\lblFrames~#1)}}
\newcommand{\teachesatone}[1]{\hfill{\footnotesize\color{mgrey}(\lblFrame~#1)}}

% ---- Summary --------------------------------------------------------------
% Every item carries the frame number it came from, in square brackets. That
% is the eighth-edition improvement and it is the single cheapest thing in the
% format: it turns a summary into a return index, so a reader who fails one
% line of "Can you?" knows exactly which four frames to reread.
\newenvironment{summarybox}{%
  \section*{\lblSummary}%
  \phantomsection\addcontentsline{toc}{section}{\lblSummary}%
  \begin{tcolorbox}[
    enhanced, breakable, sharp corners,
    colback=mlight, colframe=mblue, boxrule=0pt, leftrule=3pt,
    left=8pt, right=8pt, top=6pt, bottom=6pt]
  \begin{itemize}[leftmargin=1.3em, itemsep=5pt, topsep=2pt]
}{%
  \end{itemize}
  \end{tcolorbox}
}

% \sumitem{12}{text} -- the frame reference is not decoration, it is the point.
\newcommand{\sumitem}[2]{\item #2~{\footnotesize\color{mgrey}[#1]}}

% ---- Answers ---------------------------------------------------------------
%  Answers are authored at the point of use and printed at the back.
%
%  They are carried there by a global macro store rather than by an auxiliary
%  file. The sibling books collect their figure and screenshot manifests with
%  \@starttoc and a custom file extension, which works because those entries
%  are one line of text; it is the wrong mechanism for an answer, because an
%  answer contains displayed maths, and material written to an .aux-style file
%  and read back is expanded at shipout, where fragile commands break in ways
%  whose error messages name neither the answer nor the program.
%
%  \include is sequential within a run, so a macro defined globally while
%  typesetting program F1 is still defined when the back matter is typeset.
%  That is all this needs.
\makeatletter
\newcommand{\mfa@storeanswer}[3]{% #1 program key, #2 item key, #3 body
  \csgdef{mfa@a@#1@#2}{#3}%
  % \listcsxadd, not \listcsgadd: the item must be EXPANDED as it is stored.
  % Storing the token \mfa@exkey means every entry in the list expands to
  % whatever the key happens to be at the time the back matter is typeset,
  % which is the last one.
  \listcsxadd{mfa@akeys@#1}{#2}%
}
% Used inside an exercise item. Prints nothing here; the body reappears at the
% back of the book under this program's heading.
\newcommand{\answerto}[1]{\mfa@storeanswer{\mfa@progkey}{\mfa@exkey}{#1}}
\providecommand{\mfa@exkey}{?}
\makeatother

\newlist{exlist}{enumerate}{1}
\setlist[exlist]{label=\arabic*., leftmargin=2.0em, itemsep=6pt, topsep=3pt}
\newlist{fplist}{enumerate}{1}
\setlist[fplist]{label=\arabic*., leftmargin=2.0em, itemsep=5pt, topsep=3pt}

% These MUST be defined inside \makeatletter. Outside it, \def\mfa@exkey{...}
% parses as \def\mfa @exkey{...} -- a definition of \mfa with a delimited
% parameter text -- which raises no error at all and quietly leaves every
% answer filed under the same key. That failure is invisible until the answers
% section at the back of the book prints one program's answers twice.
\makeatletter
\newenvironment{testexercises}{%
  \section*{\lblTestExercises}%
  \phantomsection\addcontentsline{toc}{section}{\lblTestExercises}%
  \def\mfa@exkey{T\arabic{exlisti}}%
  \noindent\lblTestIntro\par\vspace{5pt}
  \begin{exlist}
}{%
  \end{exlist}
}

\newenvironment{furtherproblems}{%
  \section*{\lblFurther}%
  \phantomsection\addcontentsline{toc}{section}{\lblFurther}%
  \def\mfa@exkey{P\arabic{fplisti}}%
  \noindent\lblFurtherIntro\par\vspace{5pt}
  \begin{fplist}
}{%
  \end{fplist}
}
\makeatother

% ---- The programs register -------------------------------------------------
% \program{Title} opens a program: it is \chapter plus the bookkeeping that
% lets the back matter reproduce this program's answers under its own name.
\makeatletter
\newcommand{\program}[1]{%
  \chapter{#1}%
  \csgdef{mfa@title@\thechapter}{#1}%
  % Expanded, for the same reason as the answer keys above.
  \listxadd{\mfaprograms}{\thechapter}%
}
\newcommand{\mfa@printoneanswer}[2]{% #1 program key, #2 item key
  \item[\textbf{#2}] \csuse{mfa@a@#1@#2}%
}
\newcommand{\mfa@answerprogram}[1]{%
  \subsection*{\lblProgram~#1\quad\textmd{\itshape\csuse{mfa@title@#1}}}%
  \ifcsdef{mfa@akeys@#1}{%
    \begin{list}{}{%
      \setlength{\leftmargin}{3.2em}\setlength{\labelwidth}{2.7em}%
      \setlength{\labelsep}{0.5em}\setlength{\itemsep}{3pt}%
      \setlength{\topsep}{3pt}\setlength{\parsep}{0pt}%
      \renewcommand{\makelabel}[1]{\hfill##1}}
      \forlistcsloop{\mfa@printoneanswer{#1}}{mfa@akeys@#1}
    \end{list}%
  }{%
    \par\noindent{\color{mred}\lblNoAnswers}\par
  }%
}
% The answers body, without a heading of its own, so the appendix that carries
% it can be a numbered appendix like any other and appear in the running heads
% and the ToC in the ordinary way.
\newcommand{\answersbody}{%
  \noindent\lblAnswersIntro\par\vspace{8pt}
  \ifdef{\mfaprograms}{\forlistloop{\mfa@answerprogram}{\mfaprograms}}{%
    \noindent{\color{mred}\lblNoAnswers}\par}%
}
% Unnumbered form, for a draft build that has no appendix wiring yet.
\newcommand{\printanswers}{%
  \chapter*{\lblAnswers}%
  \phantomsection\addcontentsline{toc}{chapter}{\lblAnswers}%
  \markboth{\lblAnswers}{\lblAnswers}%
  \answersbody
}
\makeatother

% ============================================================
%  ADMONITIONS
%  ---------------------------------------------------------
%  A deliberately small, fixed vocabulary. Each one earns its place by doing
%  something no other box does; a seventh would be a smell.
% ============================================================

\newtcolorbox{note}[1][]{
  enhanced, breakable, sharp corners,
  colback=mlight, colframe=mblue, boxrule=0pt, leftrule=3pt,
  left=8pt, right=8pt, top=6pt, bottom=6pt,
  fonttitle=\bfseries\small, title={\lblNote}, #1
}

\newtcolorbox{warning}[1][]{
  enhanced, breakable, sharp corners,
  colback=mlight, colframe=mred, boxrule=0pt, leftrule=3pt,
  left=8pt, right=8pt, top=6pt, bottom=6pt,
  fonttitle=\bfseries\small, title={\lblWarning}, #1
}

% The misconception, stated as a summary after it has been walked into. The
% trap itself belongs in the frames -- elicit the wrong answer, then say
% flatly that it is wrong -- and this box is the thing the reader marks with a
% finger and comes back to.
\newtcolorbox{trapbox}[1][]{
  enhanced, breakable, sharp corners,
  colback=white, colframe=mred, boxrule=0.8pt,
  left=8pt, right=8pt, top=6pt, bottom=6pt,
  fonttitle=\bfseries\small, title={\lblTrap}, #1
}

% Where the mathematics just done shows up in a system the reader has actually
% deployed. This is the box that separates the book from a maths textbook, and
% it is also the one most easily filled with waffle: if it cannot name a
% specific line of a specific system, it should not be written.
\newtcolorbox{aibox}[1][]{
  enhanced, breakable, sharp corners,
  colback=mlight, colframe=mteal, boxrule=0pt, leftrule=3pt,
  left=8pt, right=8pt, top=6pt, bottom=6pt,
  fonttitle=\bfseries\small, title={\lblInAI}, #1
}

% What is being asserted rather than proved, and where to read the proof.
% Stroud's method buys fluency at the cost of rigour, and the honest thing is
% to say so at each point rather than once in an introduction nobody rereads.
\newtcolorbox{rigourbox}[1][]{
  enhanced, breakable, sharp corners,
  colback=white, colframe=mgrey, boxrule=0.6pt,
  left=8pt, right=8pt, top=6pt, bottom=6pt,
  fonttitle=\bfseries\small, title={\lblRigour}, #1
}

% Notation that genuinely differs -- between the Polish and English editions,
% and between disciplines. Matrix calculus alone has two incompatible layout
% conventions, and half the confusion about backpropagation is a transpose
% that came from the other one.
\newtcolorbox{notationbox}[1][]{
  enhanced, breakable, sharp corners,
  colback=mlight, colframe=mamber, boxrule=0pt, leftrule=3pt,
  left=8pt, right=8pt, top=6pt, bottom=6pt,
  fonttitle=\bfseries\small, title={\lblNotation}, #1
}

% The maths analogue of the sibling books' "run this before you trust it".
% A number in this book is either produced by a script under code/ and pulled
% in with \val, or it carries this box. Nothing ships to a reader with one
% still attached.
\newtcolorbox{verifybox}[1][]{
  enhanced, breakable, sharp corners,
  colback=white, colframe=mamber, boxrule=0.8pt,
  left=8pt, right=8pt, top=6pt, bottom=6pt,
  fonttitle=\bfseries\small, title={\lblVerify}, #1
}

\newtcolorbox{exercisebox}[1][]{
  enhanced, breakable, sharp corners,
  colback=white, colframe=mgreen, boxrule=0pt, leftrule=3pt,
  left=8pt, right=8pt, top=6pt, bottom=6pt,
  fonttitle=\bfseries\small, title={\lblExercise}, #1
}

% ============================================================
%  CODE LISTINGS
%  ---------------------------------------------------------
%  Code appears here for one reason only: to check a number. Every listing in
%  this book is runnable and every number it prints is the number in the text.
% ============================================================
\lstdefinelanguage{pythonx}{
  language=Python,
  morekeywords={as, with, assert, match, case},
  morekeywords=[2]{
    float64,float32,float16,bfloat16,int8,uint8,
    np,numpy,array,asarray,frombuffer,nextafter,finfo,iinfo,
    exp,log,log1p,expm1,logaddexp,sum,max,abs,sqrt,mean,var
  },
  sensitive=true
}

\lstdefinestyle{mfacode}{
  basicstyle=\ttfamily\footnotesize,
  keywordstyle=\color{kw}\bfseries,
  keywordstyle=[2]\color{typ},
  stringstyle=\color{str},
  commentstyle=\color{cmt}\itshape,
  numbers=left,
  numberstyle=\ttfamily\tiny\color{mgrey},
  numbersep=8pt,
  backgroundcolor=\color{codebg},
  frame=single,
  rulecolor=\color{codeframe},
  framesep=6pt,
  breaklines=true,
  breakatwhitespace=false,
  postbreak=\mbox{\textcolor{mgrey}{$\hookrightarrow$}\space},
  showstringspaces=false,
  tabsize=4,
  columns=fullflexible,
  keepspaces=true,
  captionpos=b,
  aboveskip=1.2em,
  belowskip=1.2em,
  % Maps the characters most likely to be pasted into a listing by accident.
  %
  % Do NOT add a mapping for U+00A0 (non-breaking space). It looks like the
  % obvious next entry and it makes `listings` abort the run with "Improper
  % alphabetic constant" -- fatal, no PDF, and the message names neither the
  % character nor this line. The ASCII-in-listings convention is the guard
  % against it instead. This trap has now been hit in two of the three books
  % in this series.
  literate={ł}{{\l}}1 {Ł}{{\L}}1 {ą}{{\k a}}1 {ę}{{\k e}}1
           {ś}{{\'s}}1 {ć}{{\'c}}1 {ż}{{\.z}}1 {ź}{{\'z}}1 {ó}{{\'o}}1 {ń}{{\'n}}1
           {Ą}{{\k A}}1 {Ę}{{\k E}}1 {Ś}{{\'S}}1 {Ć}{{\'C}}1
           {Ż}{{\.Z}}1 {Ź}{{\'Z}}1 {Ó}{{\'O}}1 {Ń}{{\'N}}1
           {—}{{-{}-}}2 {–}{{-}}1 {‘}{{'}}1 {’}{{'}}1
           {“}{{"}}1 {”}{{"}}1 {°}{{\textdegree}}1
}

\lstset{style=mfacode}

\lstnewenvironment{python}[1][]
  {\lstset{language=pythonx,style=mfacode,#1}}
  {}

\lstnewenvironment{shellcmd}[1][]
  {\lstset{language=bash,style=mfacode,numbers=none,keywordstyle=\color{mteal}\bfseries,#1}}
  {}

% Terminal transcripts: no line numbers, no keyword colouring. Everything in
% one of these was copied from a real run.
\lstnewenvironment{console}[1][]
  {\lstset{style=mfacode,numbers=none,basicstyle=\ttfamily\scriptsize,
           keywordstyle=\color{black},backgroundcolor=\color{white},#1}}
  {}

\newcommand{\pyfile}[3]{%
  \lstinputlisting[language=pythonx,style=mfacode,caption={#2},label={#3}]{#1}%
}
\newcommand{\pyregion}[4]{%
  \lstinputlisting[
    language=pythonx, style=mfacode,
    linerange={--8<--\ [start:#2]--- 8<--\ [end:#2]},
    includerangemarker=false,
    caption={#3}, label={#4}]{#1}%
}

% Inline literals
\newcommand{\code}[1]{\texttt{\small #1}}
\newcommand{\pkg}[1]{\texttt{\small\color{mblue}#1}}
\newcommand{\api}[1]{\texttt{\small\color{mteal}#1}}

% ============================================================
%  COMPUTED VALUES
%  ---------------------------------------------------------
%  The house rule in the sibling books is "run the listing, or mark it". The
%  analogue here is stronger, because a wrong digit in a maths book is not a
%  broken example, it is a lie the reader will carry.
%
%  Every number that came out of a computation is written by a script under
%  code/ into figures/values/<program>.tex as
%
%      \mfaval{f01.eps64}{2.220446049250313e-16}
%
%  and pulled into the text with \val{f01.eps64}. The script is the source of
%  truth; the book cannot disagree with it, because the book does not contain
%  the digits. A value the scripts no longer produce prints a visible marker
%  and is counted by `make debt`.
%
%  \val also applies the locale's decimal separator, which is how the Polish
%  edition gets its decimal comma without a translator having to remember.
% ============================================================
\IfFileExists{siunitx.sty}{%
  \usepackage{siunitx}%
  \sisetup{
    group-digits = integer,
    group-minimum-digits = 5,
    group-separator = {\,},
    exponent-product = \cdot,
    output-decimal-marker = {\mfadecimalmarker}
  }%
  \newcommand{\mfanum}[1]{\num{##1}}%
}{%
  \newcommand{\mfanum}[1]{##1}%
}

\makeatletter
\newcommand{\mfaval}[2]{\csgdef{mfaval@#1}{#2}\listxadd{\mfavalues}{#1}}
% \num on a non-numeric body is a FATAL siunitx error -- "Invalid number" --
% and under -halt-on-error that is the whole build. A script emitting a word
% where the book expected a number is a plausible mistake.
%
% The check is NOT done here. Deciding in TeX whether a token list is a number
% means parsing it character by character, which is fragile and was got wrong
% once already. The generator knows the answer for free, so it emits \mfaval
% for a number and \mfavaltext for a word, and this end only has to enforce
% which macro reads which.
\newcommand{\val}[1]{%
  \ifcsdef{mfaval@#1}{%
    \ifcsdef{mfavaltext@#1}{%
      \PackageError{mfabook}{Value `#1' is not a number}%
        {\string\val\space passes its body to siunitx, which refuses anything
         that is not a number. Use \string\valtext{#1} instead, or fix the
         script in code/ that emits this key.}%
    }{}%
    \mfanum{\csuse{mfaval@#1}}%
  }{\textcolor{mred}{\textbf{[?\,#1]}}}%
}
% A value that is deliberately a word rather than a number -- a format name, a
% verdict, a unit. Rare, and a visibly different macro so nobody reaches for
% \val and gets an error they cannot place.
\newcommand{\valtext}[1]{%
  \ifcsdef{mfaval@#1}{\csuse{mfaval@#1}}%
    {\textcolor{mred}{\textbf{[?\,#1]}}}%
}
\newcommand{\mfavaltext}[2]{%
  \csgdef{mfaval@#1}{#2}\csgdef{mfavaltext@#1}{}\listxadd{\mfavalues}{#1}%
}

% The raw stored string, for the cases where a number must appear exactly as
% the machine printed it -- inside a console block, or where the digits
% themselves are the point and grouping them would obscure the pattern.
\newcommand{\rawval}[1]{%
  \ifcsdef{mfaval@#1}{\csuse{mfaval@#1}}%
    {\textcolor{mred}{\textbf{[?\,#1]}}}%
}
\makeatother

% figures/values/all.tex is generated by `make numbers` and does nothing but
% \input each program's value file. A build with no values yet still compiles;
% every \val{} in it prints a visible marker instead of a number, which is the
% correct behaviour for a draft and an obvious failure for a release.
\IfFileExists{figures/values/all.tex}{\input{figures/values/all}}{%
  \AtBeginDocument{\PackageWarning{mfabook}{No computed values found. Run `make numbers`.}}}

% ============================================================
%  DIAGRAMS — Mermaid pipeline
%  ---------------------------------------------------------
%  \mermaidfig{key}{caption}{one-line manifest description}
%
%  Source of truth is figures/mermaid/<lang>/<key>.mmd, committed. `make
%  diagrams` renders each to figures/diagrams/<lang>/<key>.pdf, which is build
%  output and is not committed, so a diagram change reviews as a text diff.
%
%  The source is per-language because a diagram with English labels in the
%  Polish edition is exactly the kind of half-translation that makes a
%  bilingual book feel machine-made. CI checks that both languages carry the
%  same set of keys.
%
%  If the rendered PDF is absent the macro draws a placeholder AND typesets the
%  Mermaid source, so a build without mermaid-cli still shows the structure.
% ============================================================
\makeatletter
\newcommand{\listofdiagrams}{%
  \section*{\lblDiagramManifest}%
  \phantomsection\addcontentsline{toc}{section}{\lblDiagramManifest}%
  \@starttoc{dgm}%
}
\makeatother

\newcommand{\mermaidfig}[3]{%
  \addcontentsline{dgm}{subsection}{\protect\numberline{}\texttt{#1.mmd} --- #3}%
  \IfFileExists{figures/diagrams/\booklang/#1.pdf}{%
    \begin{figure}[htbp]
    \centering
    \includegraphics[width=0.95\linewidth,height=0.42\textheight,keepaspectratio]{figures/diagrams/\booklang/#1.pdf}%
    \caption{#2}\label{fig:#1}
    \end{figure}%
  }{%
    % The unrendered case deliberately does NOT float. A Mermaid source
    % typeset verbatim is often taller than a page, and a float cannot break,
    % so wrapping it in `figure` produces an overfull box per diagram and a
    % pile of tcolorbox warnings. In the flow it breaks like any other text.
    \par\medskip
    \begin{tcolorbox}[
      enhanced, breakable, width=\linewidth, sharp corners,
      colback=white, colframe=mgrey, boxrule=0.8pt,
      left=8pt, right=8pt, top=8pt, bottom=8pt]
      {\bfseries\color{mgrey}\small \lblNotRendered}\\[4pt]
      {\ttfamily\scriptsize\color{mgrey} figures/mermaid/\booklang/#1.mmd}\\[6pt]
      \IfFileExists{figures/mermaid/\booklang/#1.mmd}{%
        % ASCII only, like every other listing in this book: the fallback runs
        % the Mermaid source through `listings`, which aborts on a multi-byte
        % character it has no literate mapping for.
        \lstinputlisting[style=mfacode,numbers=none,basicstyle=\ttfamily\scriptsize,
                         backgroundcolor=\color{mlight},frame=none,
                         keywordstyle=\color{black}]{figures/mermaid/\booklang/#1.mmd}%
      }{%
        {\small\color{mred} \lblSourceMissing: \texttt{figures/mermaid/\booklang/#1.mmd}}%
      }%
    \end{tcolorbox}%
    \captionof{figure}{#2}\label{fig:#1}
    \par\medskip
  }%
}

% ============================================================
%  PROGRAM STUBS
%  ---------------------------------------------------------
%  Prints a visible marker so an unwritten program cannot be mistaken for a
%  written one and the page count never flatters the draft. `make debt` counts
%  them, and CI publishes the count on every build.
% ============================================================
\newcommand{\programstub}[1]{%
  \begin{tcolorbox}[
    enhanced, breakable, width=\linewidth, sharp corners,
    colback=mlight, colframe=mred, boxrule=1pt,
    borderline={0.6pt}{3pt}{mred, dashed},
    left=12pt, right=12pt, top=12pt, bottom=12pt]
    {\bfseries\color{mred}\large \lblNotWritten}\\[8pt]
    \small\raggedright #1
  \end{tcolorbox}%
}

% ============================================================
%  MATHEMATICAL OPERATORS
%  ---------------------------------------------------------
%  Language-invariant names live here. The ones that genuinely differ between
%  Polish and English usage -- tan/tg, Var/D^2, gcd/NWD, and the interval
%  brackets -- are defined in lang/en.tex and lang/pl.tex instead, so the
%  notation contract is executed by the build rather than remembered by a
%  translator.
% ============================================================
\DeclareMathOperator*{\argmin}{arg\,min}
\DeclareMathOperator*{\argmax}{arg\,max}
\DeclareMathOperator{\tr}{tr}
\DeclareMathOperator{\rank}{rank}
\DeclareMathOperator{\diag}{diag}
\DeclareMathOperator{\sgn}{sgn}
\DeclareMathOperator{\relu}{ReLU}
\DeclareMathOperator{\softmax}{softmax}
\DeclareMathOperator{\logsumexp}{logsumexp}
\DeclareMathOperator{\fl}{fl}
\DeclareMathOperator{\ulp}{ulp}
\DeclareMathOperator{\round}{round}

% A bare \log is FORBIDDEN in this book, and tools/check_notation.py fails the
% build on one. In Polish textbooks \log means base ten; in machine-learning
% writing it means base e. The same three characters carry opposite meanings to
% the two readerships this book serves, and they collide hardest in exactly the
% programs where it matters most -- entropy in bits and cross-entropy in nats
% are two programs apart. Write \ln for nats and \logb{2} for bits. It costs
% two characters and removes the ambiguity entirely.
\makeatletter
\let\mfa@origlog\log
\renewcommand{\log}{%
  \PackageError{mfabook}{A bare \string\log\space is forbidden in this book}%
    {Write \string\ln\space for natural logarithms or \string\logb{2}\space for
     bits. The base is never left to a convention here; see Appendix B.}%
  \mfa@origlog}
% The one legitimate form: an explicit base.
\newcommand{\logb}[1]{\mathop{\mfa@origlog}\nolimits_{#1}}
% ...and a way to write the forbidden thing when the subject IS the ban, which
% is Appendix B and nowhere else.
\newcommand{\mfalogplain}{\mathop{\mfa@origlog}\nolimits}
\makeatother

\newcommand{\R}{\mathbb{R}}
\newcommand{\N}{\mathbb{N}}
\newcommand{\Z}{\mathbb{Z}}
\newcommand{\Q}{\mathbb{Q}}
\newcommand{\C}{\mathbb{C}}
\newcommand{\Fset}{\mathbb{F}}          % the floating-point numbers, F1's subject
\newcommand{\vect}[1]{\mathbf{#1}}
\newcommand{\mat}[1]{\mathbf{#1}}
\newcommand{\T}{^{\mathsf{T}}}
\newcommand{\norm}[1]{\left\lVert #1 \right\rVert}
\newcommand{\abs}[1]{\left\lvert #1 \right\rvert}
\newcommand{\inner}[2]{\left\langle #1, #2 \right\rangle}
\newcommand{\eps}{\varepsilon}

% ============================================================
%  INDEX
%  ---------------------------------------------------------
%  Two-column, and its entries include \texttt{} names that do not hyphenate.
%  Ragged right keeps multi-word entries off the margin; it cannot help a
%  single token wider than the column, so keep verbatim index entries under
%  about 29 characters and index longer names by concept instead. Both sibling
%  books learned this the same way.
% ============================================================
\apptocmd{\theindex}{\raggedright}{}{%
  \PackageWarning{mfabook}{Could not patch theindex for ragged-right setting}}

\makeindex

% ============================================================
%  PINNED FACTS — single source of truth
%  Change these here; they propagate through both editions.
% ============================================================
\newcommand{\bookedition}{0.1}
% \pinnedon is a user-visible string and lives in lang/<lang>.tex, not here.
% It said "August 2026" on the Polish title page for exactly as long as it
% took somebody to read the Polish title page.
\newcommand{\pyver}{Python 3.13}
\newcommand{\numpyver}{2.3}
; nothing else differs between them at
%  this level. Every user-visible string lives in lang/en.tex or lang/pl.tex,
%  so a label that appears in one edition and not the other is a missing macro
%  rather than a silent divergence.
% ============================================================

% \booklang must be set by the main file. Default to English rather than
% failing, so a stray % ============================================================
%  preamble.tex — styling, environments, macros
%
%  Shared by BOTH editions. main-en.tex and main-pl.tex each define
%  \booklang before \input{preamble}; nothing else differs between them at
%  this level. Every user-visible string lives in lang/en.tex or lang/pl.tex,
%  so a label that appears in one edition and not the other is a missing macro
%  rather than a silent divergence.
% ============================================================

% \booklang must be set by the main file. Default to English rather than
% failing, so a stray \input{preamble} still compiles and says what it did.
\providecommand{\booklang}{en}

\usepackage{etoolbox}   % loaded first: the language switch below needs it

% ---------- Page geometry ----------
% Same 17 x 24 cm block as the other two books in this series, so a reader who
% owns one recognises the second. Frames are short, so the measure matters
% more here than in a prose book: a frame that runs past a page turn defeats
% the cover-the-next-frame discipline the whole method rests on.
\usepackage[
  paperwidth=17cm, paperheight=24cm,
  inner=2.2cm, outer=1.8cm, top=2.2cm, bottom=2.4cm,
  head=23pt, headsep=0.7cm, footskip=1.2cm
]{geometry}

% ---------- Encoding & fonts ----------
\usepackage[T1]{fontenc}
\usepackage[utf8]{inputenc}
\IfFileExists{lmodern.sty}{\usepackage{lmodern}}{}
\IfFileExists{newtxtext.sty}{\usepackage{newtxtext}}{}
% amsmath goes here, not down with the other packages, because newtxmath
% patches amsmath's internals and must be loaded AFTER it. The wrong order does
% not fail on a machine without newtxmath -- which is most CI images -- and
% fails on the author's full TeX Live.
%
% amssymb only when newtxmath is absent: newtxmath supplies the AMS symbols
% itself (its amssymbols option is on by default) and loading both clashes.
\usepackage{amsmath}
\IfFileExists{newtxmath.sty}{\usepackage{newtxmath}}{\usepackage{amssymb}}
\IfFileExists{inconsolata.sty}{\usepackage[varqu,scaled=0.95]{inconsolata}}{}
\IfFileExists{lmodern.sty}{\usepackage{microtype}}{\usepackage[expansion=false]{microtype}}

% ---------- Language ----------
% Loading babel with a language whose .ldf is absent is a *fatal* error, not a
% warning, so both the package and each language file are probed before either
% is requested. This is inherited from the sibling books, where a minimal TeX
% installation without texlive-lang-polish aborted the run with no PDF and an
% error message that named neither babel nor the language.
%
% The Polish edition needs Polish as the main language for hyphenation; the
% English edition needs British. Both editions load the other language too,
% because the glossary appendix is bilingual in both.
\IfFileExists{babel.sty}{%
  \IfFileExists{polish.ldf}{%
    \IfFileExists{british.ldf}{%
      \ifdefstring{\booklang}{pl}%
        {\usepackage[british,main=polish]{babel}}%
        {\usepackage[polish,main=british]{babel}}%
    }{\usepackage[polish]{babel}}%
  }{%
    \IfFileExists{british.ldf}{\usepackage[british]{babel}}{}%
  }%
}{}

% babel-polish makes " an active shorthand. listings reads its body verbatim
% and an active character in a verbatim context is a category-code accident
% waiting to happen, so the shorthand is turned off globally. Polish quotation
% marks are produced with \enquote-style macros from lang/pl.tex instead.
\ifdefstring{\booklang}{pl}{\AtBeginDocument{\shorthandoff{"}}}{}

% ---------- Core packages ----------
\usepackage{graphicx}
\usepackage{xcolor}
\usepackage{booktabs}
\usepackage{tabularx}
\usepackage{longtable}
\usepackage{array}
\usepackage{enumitem}
\IfFileExists{mathtools.sty}{\usepackage{mathtools}}{}
\usepackage{listings}
\usepackage[most]{tcolorbox}
\usepackage{fancyhdr}
\usepackage{titlesec}
\usepackage{caption}
\usepackage{imakeidx}
% csquotes with autostyle follows babel's active language, so the identical
% source \enquote{x} sets 'x' in the English edition and the Polish quotation
% marks in the Polish one. This is not a nicety: babel-polish makes " an active
% shorthand, which is turned off above precisely because an active character
% near a listings body is a category-code accident waiting to happen -- so a
% quotation mark has to come from a macro rather than from a keyboard key.
\IfFileExists{csquotes.sty}{\usepackage[autostyle=true]{csquotes}}{%
  \providecommand{\enquote}[1]{``##1''}}
\usepackage[hidelinks]{hyperref}

% ---------- Palette ----------
% Deliberately the same family as the sibling books, shifted green so the three
% spines are distinguishable on a shelf.
\definecolor{mblue}{HTML}{1F4E79}
\definecolor{mgreen}{HTML}{2E6B4F}
\definecolor{mteal}{HTML}{0E7C7B}
\definecolor{mamber}{HTML}{B26A00}
\definecolor{mred}{HTML}{9B2C2C}
\definecolor{mgrey}{HTML}{5A5A5A}
\definecolor{mlight}{HTML}{F4F6F8}
\definecolor{mfaint}{HTML}{FBFCFD}
\definecolor{codebg}{HTML}{FAFAFA}
\definecolor{codeframe}{HTML}{D8DEE4}
\definecolor{kw}{HTML}{0B5394}
\definecolor{str}{HTML}{9B2C2C}
\definecolor{cmt}{HTML}{4F7A4F}
\definecolor{typ}{HTML}{0E7C7B}

\hypersetup{
  colorlinks=true,
  linkcolor=mblue,
  urlcolor=mteal,
  citecolor=mblue,
  pdfauthor={Konrad Cinkusz}
}

% \ifpl{Polish text}{English text} -- for the handful of places where a string
% is structural rather than content and belongs in the shared files:
% structure.tex's part titles, and the main files' front matter wiring.
% Anything longer than a few words belongs in lang/ or in the edition's own
% source, not here.
% \DeclareRobustCommand, not \newcommand. A fragile \ifpl inside a \part title
% is expanded one level as it is written to the .toc, which lands
% "\ifdefstring {en}{pl}{...}{...}" in the file -- and that fails on the next
% run with "Extra }, or forgotten \endgroup", an error that names neither the
% part nor this macro.
\DeclareRobustCommand{\ifpl}[2]{\ifdefstring{\booklang}{pl}{#1}{#2}}

% A robust macro cannot go into a PDF bookmark string either -- hyperref strips
% \ifdefstring and warns once per part title. Tell it which branch to take
% instead, so the bookmarks carry the right language rather than nothing.
\makeatletter
\pdfstringdefDisableCommands{%
  \ifdefstring{\booklang}{pl}{\let\ifpl\@firstoftwo}{\let\ifpl\@secondoftwo}%
}
\makeatother

% ---------- Language strings ----------
% Every user-visible word the macros below typeset comes from here. Nothing in
% programs/ or appendices/ may hard-code an English or a Polish word inside a
% macro argument that the other edition also uses.
\input{lang/\booklang}

% ============================================================
%  PROGRAM NUMBERING
%  ---------------------------------------------------------
%  Part F programs are F1..F13; the main programs restart at 1. That is
%  Stroud's scheme and it is worth reproducing, because the Foundation part is
%  explicitly optional and numbering it separately is what makes skipping it
%  feel permitted rather than delinquent.
%
%  \theHchapter is hyperref's *destination* name, and it must stay unique
%  across the whole document. Resetting the chapter counter without also
%  changing \theHchapter produces duplicate PDF destinations, a pile of
%  warnings, and bookmarks that jump to the wrong program.
% ============================================================
\newcommand{\foundationnumbering}{%
  \setcounter{chapter}{0}%
  \renewcommand{\thechapter}{F\arabic{chapter}}%
  \renewcommand{\theHchapter}{F\arabic{chapter}}%
}
\newcommand{\mainnumbering}{%
  \setcounter{chapter}{0}%
  \renewcommand{\thechapter}{\arabic{chapter}}%
  \renewcommand{\theHchapter}{M\arabic{chapter}}%
}

% \appendix resets the chapter counter and switches \thechapter to letters, but
% it knows nothing about \theHchapter, so appendix A would claim the same PDF
% destination as main program 1. Patch it once, here.
\apptocmd{\appendix}{\renewcommand{\theHchapter}{A\Alph{chapter}}}{}{%
  \PackageWarning{mfabook}{Could not patch \string\appendix\space for hyperref destinations}}

% A program is a chapter. \chaptername is set from the language file so the
% word above the title is "Program" in both editions (it happens to be the
% same word, which is a small piece of luck).
\AtBeginDocument{\renewcommand{\chaptername}{\lblProgram}}

% ---------- Headings ----------
% \raggedright on the title body is load-bearing, not cosmetic: at \Huge on a
% 13 cm text block a long title that cannot break overflows the margin badly.
\titleformat{\chapter}[display]
  {\normalfont\huge\bfseries\color{mblue}\raggedright}
  {\normalfont\normalsize\color{mgrey}\MakeUppercase{\chaptertitlename}\ \thechapter}
  {8pt}{\Huge\raggedright}
\titlespacing*{\chapter}{0pt}{10pt}{28pt}
\titleformat{\section}{\normalfont\Large\bfseries\color{mblue}}{\thesection}{0.8em}{}
\titleformat{\subsection}{\normalfont\large\bfseries\color{mgrey}}{\thesubsection}{0.7em}{}

% head=23pt in the geometry options above, not \setlength{\headheight}:
% fancyhdr computes the height it needs from the running-head font and the
% rule, wants 22.5 pt at this size, and silently overprints the top of the text
% block if it does not get it. Setting it through geometry keeps the text block
% where geometry put it instead of pushing it down.
\pagestyle{fancy}
\fancyhf{}
\fancyhead[LE]{\small\nouppercase{\leftmark}}
% The recto head carries the FRAME number, not the section. In a book of
% numbered frames "where am I" is a frame question: every Summary
% back-reference, every Quiz route and every cross-reference between programs
% is a frame number, and hunting for one by page is the single most annoying
% thing about reading a programmed text. \theframeno at shipout is the last
% frame begun on the page, which is what a reader flicking through wants.
% Guarded: the front matter and the appendices have no frames, and a head
% reading "Frame 0" over the introduction is worse than no head at all.
\fancyhead[RO]{\ifnum\value{frameno}>0\relax\small\lblFrame~\theframeno\fi}
\fancyfoot[LE,RO]{\small\thepage}
\renewcommand{\headrulewidth}{0.4pt}

% This MUST come after \pagestyle{fancy}: fancyhdr installs its own \chaptermark
% at that point and silently overwrites anything defined earlier. The symptom is
% a running-head change that appears to do nothing at all.
%
% The running head carries the *frame* number as well as the program, because
% in a book of numbered frames "where am I" is a frame question, not a page
% question, and the summary's back-references are frame numbers.
% \@chapapp, not \lblProgram. \appendix switches \@chapapp from \chaptername
% to \appendixname, and babel supplies both in the right language; hard-coding
% the word gives an appendix a running head reading "Program A. Answers" over a
% page whose own chapter head reads "APPENDIX A".
\makeatletter
\renewcommand{\chaptermark}[1]{\markboth{\@chapapp\ \thechapter.\ #1}{}}
\makeatother
\renewcommand{\sectionmark}[1]{\markright{\thesection\ #1}}

% ============================================================
%  FRAMES — the unit the whole method is built from
%  ---------------------------------------------------------
%  A program is a stream of numbered frames. Nearly every frame ends by asking
%  the reader for something, and the answer opens the *next* frame, so that the
%  reader can cover what is below and commit to an answer before seeing it.
%
%  The environment is called `fr` rather than `frame` for two reasons. First,
%  \frame is already defined by LaTeX2e (it draws a box in a picture
%  environment) and an environment named `frame` would silently redefine it.
%  Second, this is typed forty to seventy times per program, so it wants to be
%  short in the source.
%
%    \begin{fr}
%    \ans{$x = 3$}            % the answer to the previous frame, optional
%    Teaching text, ending in a question.
%    \end{fr}
% ============================================================

\newcounter{frameno}[chapter]
\renewcommand{\theframeno}{\arabic{frameno}}

% Frames are referenced constantly -- by the Summary, by the Quiz, by the
% "Can you?" checklist and by other programs. \label inside a frame must
% therefore capture the frame number, not the section number.
\makeatletter
\newcommand{\framelabelfix}{%
  \@namedef{theHframeno}{\theHchapter.\arabic{frameno}}%
}
\makeatother

\newcommand{\framerule}{%
  \par\penalty-200\vspace{9pt}%
  \noindent\textcolor{codeframe}{\rule{\linewidth}{0.5pt}}%
  \par\vspace{5pt}%
}

% The frame number as a filled badge, set in the text flow rather than in the
% margin. A margin number would have to alternate sides under `twoside` and
% would be the first thing to overflow on a 17 cm page; a badge cannot.
\newcommand{\framebadge}{%
  \begingroup
  \setlength{\fboxsep}{2.5pt}%
  \colorbox{mblue}{\textcolor{white}{\bfseries\scriptsize\theframeno}}%
  \endgroup\hspace{0.7em}%
}

\newenvironment{fr}{%
  \framerule
  \refstepcounter{frameno}%
  \nopagebreak
  \noindent\framebadge\ignorespaces
}{\par}

% The answer to the previous frame. Set apart and centred so the eye can find
% the edge to cover, and so a reader who has answered can check in one glance
% without reading on.
\newcommand{\ans}[1]{%
  \par\nobreak\vspace{-4pt}%
  \begin{tcolorbox}[
    enhanced, sharp corners, boxrule=0pt, leftrule=2.5pt,
    colback=mfaint, colframe=mgreen,
    left=8pt, right=8pt, top=4pt, bottom=4pt,
    width=0.86\linewidth, center, halign=center,
    before skip=2pt, after skip=7pt]
    #1
  \end{tcolorbox}%
  \nobreak\noindent\ignorespaces
}

% A multi-paragraph or worked answer, where a centred box would be wrong.
\newenvironment{ansblock}{%
  \par\nopagebreak
  \begin{tcolorbox}[
    enhanced, breakable, sharp corners, boxrule=0pt, leftrule=2.5pt,
    colback=mfaint, colframe=mgreen,
    left=8pt, right=8pt, top=5pt, bottom=5pt]
}{%
  \end{tcolorbox}%
  \nopagebreak\par\vspace{2pt}\noindent\ignorespaces
}

% The row of dots. Stroud's signal that the reader is to supply a word, a
% number or an expression. \blank is inline; \dotline occupies its own line,
% which is what a longer derivation step wants.
% \penalty0 offers TeX a break opportunity immediately before the dots. Without
% it the run of dots is glued to the last word of the question, and a long
% question in the Polish edition -- where the words are longer and hyphenate
% less freely -- runs into the margin.
\newcommand{\blank}[1][4em]{\penalty0\makebox[#1]{\dotfill}}
\newcommand{\dotline}{\par\vspace{2pt}\noindent\makebox[\linewidth]{\dotfill}\par\vspace{2pt}}

% "Now one for you to do." The third rung of the scaffolding gradient, and the
% one most often skipped by authors who think they have written an exercise
% when they have written another worked example.
\newcommand{\yourturn}{%
  \par\vspace{4pt}\noindent
  {\bfseries\color{mgreen}\lblYourTurn}\par\nobreak\vspace{2pt}\noindent\ignorespaces
}

% ============================================================
%  THE PROGRAM SKELETON
%  ---------------------------------------------------------
%  Every program carries the same seven parts, in the same order. They are
%  macros rather than a convention so that a missing one is a build-time
%  absence rather than something a reader notices.
% ============================================================

% ---- Learning outcomes, on entry ----------------------------------------
% Stored as well as printed, because "Can you?" at the end of the program must
% be the same list, one for one. Storing it is what makes that structural
% rather than a promise: the checklist is generated from the outcomes, so the
% two cannot drift.
\newcounter{outcomeno}[chapter]

\makeatletter
% The key under which everything belonging to the current program is stored.
% It must expand to the printed program number (F1, F2, ..., 1, 2, ...), which
% is what makes an answer findable from the program that owns it.
\newcommand{\mfa@progkey}{\thechapter}
\newcommand{\outcome}[2]{% #1 = frames that teach it, #2 = the outcome
  \stepcounter{outcomeno}%
  \csgdef{mfa@out@\mfa@progkey @\theoutcomeno}{#2}%
  \csgdef{mfa@outfr@\mfa@progkey @\theoutcomeno}{#1}%
  % The "Can you?" row is built here rather than looped over later. A \loop
  % inside a tabularx body does not survive alignment scanning: TeX needs to
  % see the & and the row break while it reads the cell, and a conditional
  % hides them. Accumulating the rows as tokens at declaration time sidesteps
  % that entirely, and has the better property anyway -- the checklist is
  % literally made of the outcomes, so the two cannot disagree.
  \csgappto{mfa@rows@\mfa@progkey}{%
    #2 & {\footnotesize\color{mgrey}#1}
      & \ratingcell & \ratingcell & \ratingcell & \ratingcell & \ratingcell
    \tabularnewline}%
  \item #2\hfill{\footnotesize\color{mgrey}[\lblFrames~#1]}%
}
\makeatother

\makeatletter
\newenvironment{outcomes}{%
  \setcounter{outcomeno}{0}%
  \csgdef{mfa@rows@\mfa@progkey}{}%
  \begin{tcolorbox}[
    enhanced, breakable, sharp corners,
    colback=mlight, colframe=mblue, boxrule=0pt, leftrule=3pt,
    left=8pt, right=8pt, top=6pt, bottom=6pt,
    fonttitle=\bfseries\small, title={\lblOutcomes}]
  \begin{itemize}[leftmargin=1.3em, itemsep=3pt, topsep=2pt]
}{%
  \end{itemize}
  \end{tcolorbox}
}
\makeatother

% ---- "Can you?" -----------------------------------------------------------
% Generated from the stored outcomes, not written out again by hand. Five
% columns rather than a tick box, because a self-assessment with a midpoint is
% one people actually use and a yes/no one is one they lie to.
\newcommand{\ratingcell}{\textcolor{codeframe}{\rule[-0.15em]{0.8em}{0.8em}}}

\makeatletter
\newcommand{\canyou}{%
  \section*{\lblCanYou}%
  \phantomsection\addcontentsline{toc}{section}{\lblCanYou}%
  \noindent\lblCanYouIntro\par\vspace{8pt}
  \noindent
  \begingroup\small
  \ifcsvoid{mfa@rows@\mfa@progkey}{%
    {\color{mred}\lblNoOutcomes}\par
  }{%
    \begin{tabularx}{\linewidth}{@{}Xcccccc@{}}
    \toprule
    \textbf{\lblOnEntryYouCould} & \textbf{\lblFrames} & 1 & 2 & 3 & 4 & 5 \\
    \midrule
    \csuse{mfa@rows@\mfa@progkey}
    \bottomrule
    \end{tabularx}%
  }%
  \endgroup
  \par\vspace{6pt}\noindent{\footnotesize\color{mgrey}\lblCanYouFooter}\par
}
\makeatother

% ============================================================
%  QUIZ, SUMMARY, EXERCISES
% ============================================================

% ---- Quiz -----------------------------------------------------------------
% Foundation programs only. It is a diagnostic on the way in and the same test
% on the way out, which is what lets one book serve a reader who needs the
% whole program and a reader who needs four frames of it. Each question names
% the frames that teach it, so a wrong answer routes the reader somewhere
% specific rather than back to the beginning.
\newlist{quizlist}{enumerate}{1}
\setlist[quizlist]{label=\arabic*., leftmargin=2.0em, itemsep=5pt, topsep=3pt}

\makeatletter
\newenvironment{quiz}{%
  \section*{\lblQuiz}%
  \phantomsection\addcontentsline{toc}{section}{\lblQuiz}%
  \def\mfa@exkey{Q\arabic{quizlisti}}%
  \noindent\lblQuizIntro\par\vspace{5pt}
  \begin{quizlist}
}{%
  \end{quizlist}
}
\makeatother

% The frames that teach the question just asked. Two forms, because a Quiz
% question that routes to a single frame reading "Frames 22" is the kind of
% small wrongness a reader notices and an author never does.
\newcommand{\teachesat}[1]{\hfill{\footnotesize\color{mgrey}(\lblFrames~#1)}}
\newcommand{\teachesatone}[1]{\hfill{\footnotesize\color{mgrey}(\lblFrame~#1)}}

% ---- Summary --------------------------------------------------------------
% Every item carries the frame number it came from, in square brackets. That
% is the eighth-edition improvement and it is the single cheapest thing in the
% format: it turns a summary into a return index, so a reader who fails one
% line of "Can you?" knows exactly which four frames to reread.
\newenvironment{summarybox}{%
  \section*{\lblSummary}%
  \phantomsection\addcontentsline{toc}{section}{\lblSummary}%
  \begin{tcolorbox}[
    enhanced, breakable, sharp corners,
    colback=mlight, colframe=mblue, boxrule=0pt, leftrule=3pt,
    left=8pt, right=8pt, top=6pt, bottom=6pt]
  \begin{itemize}[leftmargin=1.3em, itemsep=5pt, topsep=2pt]
}{%
  \end{itemize}
  \end{tcolorbox}
}

% \sumitem{12}{text} -- the frame reference is not decoration, it is the point.
\newcommand{\sumitem}[2]{\item #2~{\footnotesize\color{mgrey}[#1]}}

% ---- Answers ---------------------------------------------------------------
%  Answers are authored at the point of use and printed at the back.
%
%  They are carried there by a global macro store rather than by an auxiliary
%  file. The sibling books collect their figure and screenshot manifests with
%  \@starttoc and a custom file extension, which works because those entries
%  are one line of text; it is the wrong mechanism for an answer, because an
%  answer contains displayed maths, and material written to an .aux-style file
%  and read back is expanded at shipout, where fragile commands break in ways
%  whose error messages name neither the answer nor the program.
%
%  \include is sequential within a run, so a macro defined globally while
%  typesetting program F1 is still defined when the back matter is typeset.
%  That is all this needs.
\makeatletter
\newcommand{\mfa@storeanswer}[3]{% #1 program key, #2 item key, #3 body
  \csgdef{mfa@a@#1@#2}{#3}%
  % \listcsxadd, not \listcsgadd: the item must be EXPANDED as it is stored.
  % Storing the token \mfa@exkey means every entry in the list expands to
  % whatever the key happens to be at the time the back matter is typeset,
  % which is the last one.
  \listcsxadd{mfa@akeys@#1}{#2}%
}
% Used inside an exercise item. Prints nothing here; the body reappears at the
% back of the book under this program's heading.
\newcommand{\answerto}[1]{\mfa@storeanswer{\mfa@progkey}{\mfa@exkey}{#1}}
\providecommand{\mfa@exkey}{?}
\makeatother

\newlist{exlist}{enumerate}{1}
\setlist[exlist]{label=\arabic*., leftmargin=2.0em, itemsep=6pt, topsep=3pt}
\newlist{fplist}{enumerate}{1}
\setlist[fplist]{label=\arabic*., leftmargin=2.0em, itemsep=5pt, topsep=3pt}

% These MUST be defined inside \makeatletter. Outside it, \def\mfa@exkey{...}
% parses as \def\mfa @exkey{...} -- a definition of \mfa with a delimited
% parameter text -- which raises no error at all and quietly leaves every
% answer filed under the same key. That failure is invisible until the answers
% section at the back of the book prints one program's answers twice.
\makeatletter
\newenvironment{testexercises}{%
  \section*{\lblTestExercises}%
  \phantomsection\addcontentsline{toc}{section}{\lblTestExercises}%
  \def\mfa@exkey{T\arabic{exlisti}}%
  \noindent\lblTestIntro\par\vspace{5pt}
  \begin{exlist}
}{%
  \end{exlist}
}

\newenvironment{furtherproblems}{%
  \section*{\lblFurther}%
  \phantomsection\addcontentsline{toc}{section}{\lblFurther}%
  \def\mfa@exkey{P\arabic{fplisti}}%
  \noindent\lblFurtherIntro\par\vspace{5pt}
  \begin{fplist}
}{%
  \end{fplist}
}
\makeatother

% ---- The programs register -------------------------------------------------
% \program{Title} opens a program: it is \chapter plus the bookkeeping that
% lets the back matter reproduce this program's answers under its own name.
\makeatletter
\newcommand{\program}[1]{%
  \chapter{#1}%
  \csgdef{mfa@title@\thechapter}{#1}%
  % Expanded, for the same reason as the answer keys above.
  \listxadd{\mfaprograms}{\thechapter}%
}
\newcommand{\mfa@printoneanswer}[2]{% #1 program key, #2 item key
  \item[\textbf{#2}] \csuse{mfa@a@#1@#2}%
}
\newcommand{\mfa@answerprogram}[1]{%
  \subsection*{\lblProgram~#1\quad\textmd{\itshape\csuse{mfa@title@#1}}}%
  \ifcsdef{mfa@akeys@#1}{%
    \begin{list}{}{%
      \setlength{\leftmargin}{3.2em}\setlength{\labelwidth}{2.7em}%
      \setlength{\labelsep}{0.5em}\setlength{\itemsep}{3pt}%
      \setlength{\topsep}{3pt}\setlength{\parsep}{0pt}%
      \renewcommand{\makelabel}[1]{\hfill##1}}
      \forlistcsloop{\mfa@printoneanswer{#1}}{mfa@akeys@#1}
    \end{list}%
  }{%
    \par\noindent{\color{mred}\lblNoAnswers}\par
  }%
}
% The answers body, without a heading of its own, so the appendix that carries
% it can be a numbered appendix like any other and appear in the running heads
% and the ToC in the ordinary way.
\newcommand{\answersbody}{%
  \noindent\lblAnswersIntro\par\vspace{8pt}
  \ifdef{\mfaprograms}{\forlistloop{\mfa@answerprogram}{\mfaprograms}}{%
    \noindent{\color{mred}\lblNoAnswers}\par}%
}
% Unnumbered form, for a draft build that has no appendix wiring yet.
\newcommand{\printanswers}{%
  \chapter*{\lblAnswers}%
  \phantomsection\addcontentsline{toc}{chapter}{\lblAnswers}%
  \markboth{\lblAnswers}{\lblAnswers}%
  \answersbody
}
\makeatother

% ============================================================
%  ADMONITIONS
%  ---------------------------------------------------------
%  A deliberately small, fixed vocabulary. Each one earns its place by doing
%  something no other box does; a seventh would be a smell.
% ============================================================

\newtcolorbox{note}[1][]{
  enhanced, breakable, sharp corners,
  colback=mlight, colframe=mblue, boxrule=0pt, leftrule=3pt,
  left=8pt, right=8pt, top=6pt, bottom=6pt,
  fonttitle=\bfseries\small, title={\lblNote}, #1
}

\newtcolorbox{warning}[1][]{
  enhanced, breakable, sharp corners,
  colback=mlight, colframe=mred, boxrule=0pt, leftrule=3pt,
  left=8pt, right=8pt, top=6pt, bottom=6pt,
  fonttitle=\bfseries\small, title={\lblWarning}, #1
}

% The misconception, stated as a summary after it has been walked into. The
% trap itself belongs in the frames -- elicit the wrong answer, then say
% flatly that it is wrong -- and this box is the thing the reader marks with a
% finger and comes back to.
\newtcolorbox{trapbox}[1][]{
  enhanced, breakable, sharp corners,
  colback=white, colframe=mred, boxrule=0.8pt,
  left=8pt, right=8pt, top=6pt, bottom=6pt,
  fonttitle=\bfseries\small, title={\lblTrap}, #1
}

% Where the mathematics just done shows up in a system the reader has actually
% deployed. This is the box that separates the book from a maths textbook, and
% it is also the one most easily filled with waffle: if it cannot name a
% specific line of a specific system, it should not be written.
\newtcolorbox{aibox}[1][]{
  enhanced, breakable, sharp corners,
  colback=mlight, colframe=mteal, boxrule=0pt, leftrule=3pt,
  left=8pt, right=8pt, top=6pt, bottom=6pt,
  fonttitle=\bfseries\small, title={\lblInAI}, #1
}

% What is being asserted rather than proved, and where to read the proof.
% Stroud's method buys fluency at the cost of rigour, and the honest thing is
% to say so at each point rather than once in an introduction nobody rereads.
\newtcolorbox{rigourbox}[1][]{
  enhanced, breakable, sharp corners,
  colback=white, colframe=mgrey, boxrule=0.6pt,
  left=8pt, right=8pt, top=6pt, bottom=6pt,
  fonttitle=\bfseries\small, title={\lblRigour}, #1
}

% Notation that genuinely differs -- between the Polish and English editions,
% and between disciplines. Matrix calculus alone has two incompatible layout
% conventions, and half the confusion about backpropagation is a transpose
% that came from the other one.
\newtcolorbox{notationbox}[1][]{
  enhanced, breakable, sharp corners,
  colback=mlight, colframe=mamber, boxrule=0pt, leftrule=3pt,
  left=8pt, right=8pt, top=6pt, bottom=6pt,
  fonttitle=\bfseries\small, title={\lblNotation}, #1
}

% The maths analogue of the sibling books' "run this before you trust it".
% A number in this book is either produced by a script under code/ and pulled
% in with \val, or it carries this box. Nothing ships to a reader with one
% still attached.
\newtcolorbox{verifybox}[1][]{
  enhanced, breakable, sharp corners,
  colback=white, colframe=mamber, boxrule=0.8pt,
  left=8pt, right=8pt, top=6pt, bottom=6pt,
  fonttitle=\bfseries\small, title={\lblVerify}, #1
}

\newtcolorbox{exercisebox}[1][]{
  enhanced, breakable, sharp corners,
  colback=white, colframe=mgreen, boxrule=0pt, leftrule=3pt,
  left=8pt, right=8pt, top=6pt, bottom=6pt,
  fonttitle=\bfseries\small, title={\lblExercise}, #1
}

% ============================================================
%  CODE LISTINGS
%  ---------------------------------------------------------
%  Code appears here for one reason only: to check a number. Every listing in
%  this book is runnable and every number it prints is the number in the text.
% ============================================================
\lstdefinelanguage{pythonx}{
  language=Python,
  morekeywords={as, with, assert, match, case},
  morekeywords=[2]{
    float64,float32,float16,bfloat16,int8,uint8,
    np,numpy,array,asarray,frombuffer,nextafter,finfo,iinfo,
    exp,log,log1p,expm1,logaddexp,sum,max,abs,sqrt,mean,var
  },
  sensitive=true
}

\lstdefinestyle{mfacode}{
  basicstyle=\ttfamily\footnotesize,
  keywordstyle=\color{kw}\bfseries,
  keywordstyle=[2]\color{typ},
  stringstyle=\color{str},
  commentstyle=\color{cmt}\itshape,
  numbers=left,
  numberstyle=\ttfamily\tiny\color{mgrey},
  numbersep=8pt,
  backgroundcolor=\color{codebg},
  frame=single,
  rulecolor=\color{codeframe},
  framesep=6pt,
  breaklines=true,
  breakatwhitespace=false,
  postbreak=\mbox{\textcolor{mgrey}{$\hookrightarrow$}\space},
  showstringspaces=false,
  tabsize=4,
  columns=fullflexible,
  keepspaces=true,
  captionpos=b,
  aboveskip=1.2em,
  belowskip=1.2em,
  % Maps the characters most likely to be pasted into a listing by accident.
  %
  % Do NOT add a mapping for U+00A0 (non-breaking space). It looks like the
  % obvious next entry and it makes `listings` abort the run with "Improper
  % alphabetic constant" -- fatal, no PDF, and the message names neither the
  % character nor this line. The ASCII-in-listings convention is the guard
  % against it instead. This trap has now been hit in two of the three books
  % in this series.
  literate={ł}{{\l}}1 {Ł}{{\L}}1 {ą}{{\k a}}1 {ę}{{\k e}}1
           {ś}{{\'s}}1 {ć}{{\'c}}1 {ż}{{\.z}}1 {ź}{{\'z}}1 {ó}{{\'o}}1 {ń}{{\'n}}1
           {Ą}{{\k A}}1 {Ę}{{\k E}}1 {Ś}{{\'S}}1 {Ć}{{\'C}}1
           {Ż}{{\.Z}}1 {Ź}{{\'Z}}1 {Ó}{{\'O}}1 {Ń}{{\'N}}1
           {—}{{-{}-}}2 {–}{{-}}1 {‘}{{'}}1 {’}{{'}}1
           {“}{{"}}1 {”}{{"}}1 {°}{{\textdegree}}1
}

\lstset{style=mfacode}

\lstnewenvironment{python}[1][]
  {\lstset{language=pythonx,style=mfacode,#1}}
  {}

\lstnewenvironment{shellcmd}[1][]
  {\lstset{language=bash,style=mfacode,numbers=none,keywordstyle=\color{mteal}\bfseries,#1}}
  {}

% Terminal transcripts: no line numbers, no keyword colouring. Everything in
% one of these was copied from a real run.
\lstnewenvironment{console}[1][]
  {\lstset{style=mfacode,numbers=none,basicstyle=\ttfamily\scriptsize,
           keywordstyle=\color{black},backgroundcolor=\color{white},#1}}
  {}

\newcommand{\pyfile}[3]{%
  \lstinputlisting[language=pythonx,style=mfacode,caption={#2},label={#3}]{#1}%
}
\newcommand{\pyregion}[4]{%
  \lstinputlisting[
    language=pythonx, style=mfacode,
    linerange={--8<--\ [start:#2]--- 8<--\ [end:#2]},
    includerangemarker=false,
    caption={#3}, label={#4}]{#1}%
}

% Inline literals
\newcommand{\code}[1]{\texttt{\small #1}}
\newcommand{\pkg}[1]{\texttt{\small\color{mblue}#1}}
\newcommand{\api}[1]{\texttt{\small\color{mteal}#1}}

% ============================================================
%  COMPUTED VALUES
%  ---------------------------------------------------------
%  The house rule in the sibling books is "run the listing, or mark it". The
%  analogue here is stronger, because a wrong digit in a maths book is not a
%  broken example, it is a lie the reader will carry.
%
%  Every number that came out of a computation is written by a script under
%  code/ into figures/values/<program>.tex as
%
%      \mfaval{f01.eps64}{2.220446049250313e-16}
%
%  and pulled into the text with \val{f01.eps64}. The script is the source of
%  truth; the book cannot disagree with it, because the book does not contain
%  the digits. A value the scripts no longer produce prints a visible marker
%  and is counted by `make debt`.
%
%  \val also applies the locale's decimal separator, which is how the Polish
%  edition gets its decimal comma without a translator having to remember.
% ============================================================
\IfFileExists{siunitx.sty}{%
  \usepackage{siunitx}%
  \sisetup{
    group-digits = integer,
    group-minimum-digits = 5,
    group-separator = {\,},
    exponent-product = \cdot,
    output-decimal-marker = {\mfadecimalmarker}
  }%
  \newcommand{\mfanum}[1]{\num{##1}}%
}{%
  \newcommand{\mfanum}[1]{##1}%
}

\makeatletter
\newcommand{\mfaval}[2]{\csgdef{mfaval@#1}{#2}\listxadd{\mfavalues}{#1}}
% \num on a non-numeric body is a FATAL siunitx error -- "Invalid number" --
% and under -halt-on-error that is the whole build. A script emitting a word
% where the book expected a number is a plausible mistake.
%
% The check is NOT done here. Deciding in TeX whether a token list is a number
% means parsing it character by character, which is fragile and was got wrong
% once already. The generator knows the answer for free, so it emits \mfaval
% for a number and \mfavaltext for a word, and this end only has to enforce
% which macro reads which.
\newcommand{\val}[1]{%
  \ifcsdef{mfaval@#1}{%
    \ifcsdef{mfavaltext@#1}{%
      \PackageError{mfabook}{Value `#1' is not a number}%
        {\string\val\space passes its body to siunitx, which refuses anything
         that is not a number. Use \string\valtext{#1} instead, or fix the
         script in code/ that emits this key.}%
    }{}%
    \mfanum{\csuse{mfaval@#1}}%
  }{\textcolor{mred}{\textbf{[?\,#1]}}}%
}
% A value that is deliberately a word rather than a number -- a format name, a
% verdict, a unit. Rare, and a visibly different macro so nobody reaches for
% \val and gets an error they cannot place.
\newcommand{\valtext}[1]{%
  \ifcsdef{mfaval@#1}{\csuse{mfaval@#1}}%
    {\textcolor{mred}{\textbf{[?\,#1]}}}%
}
\newcommand{\mfavaltext}[2]{%
  \csgdef{mfaval@#1}{#2}\csgdef{mfavaltext@#1}{}\listxadd{\mfavalues}{#1}%
}

% The raw stored string, for the cases where a number must appear exactly as
% the machine printed it -- inside a console block, or where the digits
% themselves are the point and grouping them would obscure the pattern.
\newcommand{\rawval}[1]{%
  \ifcsdef{mfaval@#1}{\csuse{mfaval@#1}}%
    {\textcolor{mred}{\textbf{[?\,#1]}}}%
}
\makeatother

% figures/values/all.tex is generated by `make numbers` and does nothing but
% \input each program's value file. A build with no values yet still compiles;
% every \val{} in it prints a visible marker instead of a number, which is the
% correct behaviour for a draft and an obvious failure for a release.
\IfFileExists{figures/values/all.tex}{\input{figures/values/all}}{%
  \AtBeginDocument{\PackageWarning{mfabook}{No computed values found. Run `make numbers`.}}}

% ============================================================
%  DIAGRAMS — Mermaid pipeline
%  ---------------------------------------------------------
%  \mermaidfig{key}{caption}{one-line manifest description}
%
%  Source of truth is figures/mermaid/<lang>/<key>.mmd, committed. `make
%  diagrams` renders each to figures/diagrams/<lang>/<key>.pdf, which is build
%  output and is not committed, so a diagram change reviews as a text diff.
%
%  The source is per-language because a diagram with English labels in the
%  Polish edition is exactly the kind of half-translation that makes a
%  bilingual book feel machine-made. CI checks that both languages carry the
%  same set of keys.
%
%  If the rendered PDF is absent the macro draws a placeholder AND typesets the
%  Mermaid source, so a build without mermaid-cli still shows the structure.
% ============================================================
\makeatletter
\newcommand{\listofdiagrams}{%
  \section*{\lblDiagramManifest}%
  \phantomsection\addcontentsline{toc}{section}{\lblDiagramManifest}%
  \@starttoc{dgm}%
}
\makeatother

\newcommand{\mermaidfig}[3]{%
  \addcontentsline{dgm}{subsection}{\protect\numberline{}\texttt{#1.mmd} --- #3}%
  \IfFileExists{figures/diagrams/\booklang/#1.pdf}{%
    \begin{figure}[htbp]
    \centering
    \includegraphics[width=0.95\linewidth,height=0.42\textheight,keepaspectratio]{figures/diagrams/\booklang/#1.pdf}%
    \caption{#2}\label{fig:#1}
    \end{figure}%
  }{%
    % The unrendered case deliberately does NOT float. A Mermaid source
    % typeset verbatim is often taller than a page, and a float cannot break,
    % so wrapping it in `figure` produces an overfull box per diagram and a
    % pile of tcolorbox warnings. In the flow it breaks like any other text.
    \par\medskip
    \begin{tcolorbox}[
      enhanced, breakable, width=\linewidth, sharp corners,
      colback=white, colframe=mgrey, boxrule=0.8pt,
      left=8pt, right=8pt, top=8pt, bottom=8pt]
      {\bfseries\color{mgrey}\small \lblNotRendered}\\[4pt]
      {\ttfamily\scriptsize\color{mgrey} figures/mermaid/\booklang/#1.mmd}\\[6pt]
      \IfFileExists{figures/mermaid/\booklang/#1.mmd}{%
        % ASCII only, like every other listing in this book: the fallback runs
        % the Mermaid source through `listings`, which aborts on a multi-byte
        % character it has no literate mapping for.
        \lstinputlisting[style=mfacode,numbers=none,basicstyle=\ttfamily\scriptsize,
                         backgroundcolor=\color{mlight},frame=none,
                         keywordstyle=\color{black}]{figures/mermaid/\booklang/#1.mmd}%
      }{%
        {\small\color{mred} \lblSourceMissing: \texttt{figures/mermaid/\booklang/#1.mmd}}%
      }%
    \end{tcolorbox}%
    \captionof{figure}{#2}\label{fig:#1}
    \par\medskip
  }%
}

% ============================================================
%  PROGRAM STUBS
%  ---------------------------------------------------------
%  Prints a visible marker so an unwritten program cannot be mistaken for a
%  written one and the page count never flatters the draft. `make debt` counts
%  them, and CI publishes the count on every build.
% ============================================================
\newcommand{\programstub}[1]{%
  \begin{tcolorbox}[
    enhanced, breakable, width=\linewidth, sharp corners,
    colback=mlight, colframe=mred, boxrule=1pt,
    borderline={0.6pt}{3pt}{mred, dashed},
    left=12pt, right=12pt, top=12pt, bottom=12pt]
    {\bfseries\color{mred}\large \lblNotWritten}\\[8pt]
    \small\raggedright #1
  \end{tcolorbox}%
}

% ============================================================
%  MATHEMATICAL OPERATORS
%  ---------------------------------------------------------
%  Language-invariant names live here. The ones that genuinely differ between
%  Polish and English usage -- tan/tg, Var/D^2, gcd/NWD, and the interval
%  brackets -- are defined in lang/en.tex and lang/pl.tex instead, so the
%  notation contract is executed by the build rather than remembered by a
%  translator.
% ============================================================
\DeclareMathOperator*{\argmin}{arg\,min}
\DeclareMathOperator*{\argmax}{arg\,max}
\DeclareMathOperator{\tr}{tr}
\DeclareMathOperator{\rank}{rank}
\DeclareMathOperator{\diag}{diag}
\DeclareMathOperator{\sgn}{sgn}
\DeclareMathOperator{\relu}{ReLU}
\DeclareMathOperator{\softmax}{softmax}
\DeclareMathOperator{\logsumexp}{logsumexp}
\DeclareMathOperator{\fl}{fl}
\DeclareMathOperator{\ulp}{ulp}
\DeclareMathOperator{\round}{round}

% A bare \log is FORBIDDEN in this book, and tools/check_notation.py fails the
% build on one. In Polish textbooks \log means base ten; in machine-learning
% writing it means base e. The same three characters carry opposite meanings to
% the two readerships this book serves, and they collide hardest in exactly the
% programs where it matters most -- entropy in bits and cross-entropy in nats
% are two programs apart. Write \ln for nats and \logb{2} for bits. It costs
% two characters and removes the ambiguity entirely.
\makeatletter
\let\mfa@origlog\log
\renewcommand{\log}{%
  \PackageError{mfabook}{A bare \string\log\space is forbidden in this book}%
    {Write \string\ln\space for natural logarithms or \string\logb{2}\space for
     bits. The base is never left to a convention here; see Appendix B.}%
  \mfa@origlog}
% The one legitimate form: an explicit base.
\newcommand{\logb}[1]{\mathop{\mfa@origlog}\nolimits_{#1}}
% ...and a way to write the forbidden thing when the subject IS the ban, which
% is Appendix B and nowhere else.
\newcommand{\mfalogplain}{\mathop{\mfa@origlog}\nolimits}
\makeatother

\newcommand{\R}{\mathbb{R}}
\newcommand{\N}{\mathbb{N}}
\newcommand{\Z}{\mathbb{Z}}
\newcommand{\Q}{\mathbb{Q}}
\newcommand{\C}{\mathbb{C}}
\newcommand{\Fset}{\mathbb{F}}          % the floating-point numbers, F1's subject
\newcommand{\vect}[1]{\mathbf{#1}}
\newcommand{\mat}[1]{\mathbf{#1}}
\newcommand{\T}{^{\mathsf{T}}}
\newcommand{\norm}[1]{\left\lVert #1 \right\rVert}
\newcommand{\abs}[1]{\left\lvert #1 \right\rvert}
\newcommand{\inner}[2]{\left\langle #1, #2 \right\rangle}
\newcommand{\eps}{\varepsilon}

% ============================================================
%  INDEX
%  ---------------------------------------------------------
%  Two-column, and its entries include \texttt{} names that do not hyphenate.
%  Ragged right keeps multi-word entries off the margin; it cannot help a
%  single token wider than the column, so keep verbatim index entries under
%  about 29 characters and index longer names by concept instead. Both sibling
%  books learned this the same way.
% ============================================================
\apptocmd{\theindex}{\raggedright}{}{%
  \PackageWarning{mfabook}{Could not patch theindex for ragged-right setting}}

\makeindex

% ============================================================
%  PINNED FACTS — single source of truth
%  Change these here; they propagate through both editions.
% ============================================================
\newcommand{\bookedition}{0.1}
% \pinnedon is a user-visible string and lives in lang/<lang>.tex, not here.
% It said "August 2026" on the Polish title page for exactly as long as it
% took somebody to read the Polish title page.
\newcommand{\pyver}{Python 3.13}
\newcommand{\numpyver}{2.3}
 still compiles and says what it did.
\providecommand{\booklang}{en}

\usepackage{etoolbox}   % loaded first: the language switch below needs it

% ---------- Page geometry ----------
% Same 17 x 24 cm block as the other two books in this series, so a reader who
% owns one recognises the second. Frames are short, so the measure matters
% more here than in a prose book: a frame that runs past a page turn defeats
% the cover-the-next-frame discipline the whole method rests on.
\usepackage[
  paperwidth=17cm, paperheight=24cm,
  inner=2.2cm, outer=1.8cm, top=2.2cm, bottom=2.4cm,
  head=23pt, headsep=0.7cm, footskip=1.2cm
]{geometry}

% ---------- Encoding & fonts ----------
\usepackage[T1]{fontenc}
\usepackage[utf8]{inputenc}
\IfFileExists{lmodern.sty}{\usepackage{lmodern}}{}
\IfFileExists{newtxtext.sty}{\usepackage{newtxtext}}{}
% amsmath goes here, not down with the other packages, because newtxmath
% patches amsmath's internals and must be loaded AFTER it. The wrong order does
% not fail on a machine without newtxmath -- which is most CI images -- and
% fails on the author's full TeX Live.
%
% amssymb only when newtxmath is absent: newtxmath supplies the AMS symbols
% itself (its amssymbols option is on by default) and loading both clashes.
\usepackage{amsmath}
\IfFileExists{newtxmath.sty}{\usepackage{newtxmath}}{\usepackage{amssymb}}
\IfFileExists{inconsolata.sty}{\usepackage[varqu,scaled=0.95]{inconsolata}}{}
\IfFileExists{lmodern.sty}{\usepackage{microtype}}{\usepackage[expansion=false]{microtype}}

% ---------- Language ----------
% Loading babel with a language whose .ldf is absent is a *fatal* error, not a
% warning, so both the package and each language file are probed before either
% is requested. This is inherited from the sibling books, where a minimal TeX
% installation without texlive-lang-polish aborted the run with no PDF and an
% error message that named neither babel nor the language.
%
% The Polish edition needs Polish as the main language for hyphenation; the
% English edition needs British. Both editions load the other language too,
% because the glossary appendix is bilingual in both.
\IfFileExists{babel.sty}{%
  \IfFileExists{polish.ldf}{%
    \IfFileExists{british.ldf}{%
      \ifdefstring{\booklang}{pl}%
        {\usepackage[british,main=polish]{babel}}%
        {\usepackage[polish,main=british]{babel}}%
    }{\usepackage[polish]{babel}}%
  }{%
    \IfFileExists{british.ldf}{\usepackage[british]{babel}}{}%
  }%
}{}

% babel-polish makes " an active shorthand. listings reads its body verbatim
% and an active character in a verbatim context is a category-code accident
% waiting to happen, so the shorthand is turned off globally. Polish quotation
% marks are produced with \enquote-style macros from lang/pl.tex instead.
\ifdefstring{\booklang}{pl}{\AtBeginDocument{\shorthandoff{"}}}{}

% ---------- Core packages ----------
\usepackage{graphicx}
\usepackage{xcolor}
\usepackage{booktabs}
\usepackage{tabularx}
\usepackage{longtable}
\usepackage{array}
\usepackage{enumitem}
\IfFileExists{mathtools.sty}{\usepackage{mathtools}}{}
\usepackage{listings}
\usepackage[most]{tcolorbox}
\usepackage{fancyhdr}
\usepackage{titlesec}
\usepackage{caption}
\usepackage{imakeidx}
% csquotes with autostyle follows babel's active language, so the identical
% source \enquote{x} sets 'x' in the English edition and the Polish quotation
% marks in the Polish one. This is not a nicety: babel-polish makes " an active
% shorthand, which is turned off above precisely because an active character
% near a listings body is a category-code accident waiting to happen -- so a
% quotation mark has to come from a macro rather than from a keyboard key.
\IfFileExists{csquotes.sty}{\usepackage[autostyle=true]{csquotes}}{%
  \providecommand{\enquote}[1]{``##1''}}
\usepackage[hidelinks]{hyperref}

% ---------- Palette ----------
% Deliberately the same family as the sibling books, shifted green so the three
% spines are distinguishable on a shelf.
\definecolor{mblue}{HTML}{1F4E79}
\definecolor{mgreen}{HTML}{2E6B4F}
\definecolor{mteal}{HTML}{0E7C7B}
\definecolor{mamber}{HTML}{B26A00}
\definecolor{mred}{HTML}{9B2C2C}
\definecolor{mgrey}{HTML}{5A5A5A}
\definecolor{mlight}{HTML}{F4F6F8}
\definecolor{mfaint}{HTML}{FBFCFD}
\definecolor{codebg}{HTML}{FAFAFA}
\definecolor{codeframe}{HTML}{D8DEE4}
\definecolor{kw}{HTML}{0B5394}
\definecolor{str}{HTML}{9B2C2C}
\definecolor{cmt}{HTML}{4F7A4F}
\definecolor{typ}{HTML}{0E7C7B}

\hypersetup{
  colorlinks=true,
  linkcolor=mblue,
  urlcolor=mteal,
  citecolor=mblue,
  pdfauthor={Konrad Cinkusz}
}

% \ifpl{Polish text}{English text} -- for the handful of places where a string
% is structural rather than content and belongs in the shared files:
% structure.tex's part titles, and the main files' front matter wiring.
% Anything longer than a few words belongs in lang/ or in the edition's own
% source, not here.
% \DeclareRobustCommand, not \newcommand. A fragile \ifpl inside a \part title
% is expanded one level as it is written to the .toc, which lands
% "\ifdefstring {en}{pl}{...}{...}" in the file -- and that fails on the next
% run with "Extra }, or forgotten \endgroup", an error that names neither the
% part nor this macro.
\DeclareRobustCommand{\ifpl}[2]{\ifdefstring{\booklang}{pl}{#1}{#2}}

% A robust macro cannot go into a PDF bookmark string either -- hyperref strips
% \ifdefstring and warns once per part title. Tell it which branch to take
% instead, so the bookmarks carry the right language rather than nothing.
\makeatletter
\pdfstringdefDisableCommands{%
  \ifdefstring{\booklang}{pl}{\let\ifpl\@firstoftwo}{\let\ifpl\@secondoftwo}%
}
\makeatother

% ---------- Language strings ----------
% Every user-visible word the macros below typeset comes from here. Nothing in
% programs/ or appendices/ may hard-code an English or a Polish word inside a
% macro argument that the other edition also uses.
\input{lang/\booklang}

% ============================================================
%  PROGRAM NUMBERING
%  ---------------------------------------------------------
%  Part F programs are F1..F13; the main programs restart at 1. That is
%  Stroud's scheme and it is worth reproducing, because the Foundation part is
%  explicitly optional and numbering it separately is what makes skipping it
%  feel permitted rather than delinquent.
%
%  \theHchapter is hyperref's *destination* name, and it must stay unique
%  across the whole document. Resetting the chapter counter without also
%  changing \theHchapter produces duplicate PDF destinations, a pile of
%  warnings, and bookmarks that jump to the wrong program.
% ============================================================
\newcommand{\foundationnumbering}{%
  \setcounter{chapter}{0}%
  \renewcommand{\thechapter}{F\arabic{chapter}}%
  \renewcommand{\theHchapter}{F\arabic{chapter}}%
}
\newcommand{\mainnumbering}{%
  \setcounter{chapter}{0}%
  \renewcommand{\thechapter}{\arabic{chapter}}%
  \renewcommand{\theHchapter}{M\arabic{chapter}}%
}

% \appendix resets the chapter counter and switches \thechapter to letters, but
% it knows nothing about \theHchapter, so appendix A would claim the same PDF
% destination as main program 1. Patch it once, here.
\apptocmd{\appendix}{\renewcommand{\theHchapter}{A\Alph{chapter}}}{}{%
  \PackageWarning{mfabook}{Could not patch \string\appendix\space for hyperref destinations}}

% A program is a chapter. \chaptername is set from the language file so the
% word above the title is "Program" in both editions (it happens to be the
% same word, which is a small piece of luck).
\AtBeginDocument{\renewcommand{\chaptername}{\lblProgram}}

% ---------- Headings ----------
% \raggedright on the title body is load-bearing, not cosmetic: at \Huge on a
% 13 cm text block a long title that cannot break overflows the margin badly.
\titleformat{\chapter}[display]
  {\normalfont\huge\bfseries\color{mblue}\raggedright}
  {\normalfont\normalsize\color{mgrey}\MakeUppercase{\chaptertitlename}\ \thechapter}
  {8pt}{\Huge\raggedright}
\titlespacing*{\chapter}{0pt}{10pt}{28pt}
\titleformat{\section}{\normalfont\Large\bfseries\color{mblue}}{\thesection}{0.8em}{}
\titleformat{\subsection}{\normalfont\large\bfseries\color{mgrey}}{\thesubsection}{0.7em}{}

% head=23pt in the geometry options above, not \setlength{\headheight}:
% fancyhdr computes the height it needs from the running-head font and the
% rule, wants 22.5 pt at this size, and silently overprints the top of the text
% block if it does not get it. Setting it through geometry keeps the text block
% where geometry put it instead of pushing it down.
\pagestyle{fancy}
\fancyhf{}
\fancyhead[LE]{\small\nouppercase{\leftmark}}
% The recto head carries the FRAME number, not the section. In a book of
% numbered frames "where am I" is a frame question: every Summary
% back-reference, every Quiz route and every cross-reference between programs
% is a frame number, and hunting for one by page is the single most annoying
% thing about reading a programmed text. \theframeno at shipout is the last
% frame begun on the page, which is what a reader flicking through wants.
% Guarded: the front matter and the appendices have no frames, and a head
% reading "Frame 0" over the introduction is worse than no head at all.
\fancyhead[RO]{\ifnum\value{frameno}>0\relax\small\lblFrame~\theframeno\fi}
\fancyfoot[LE,RO]{\small\thepage}
\renewcommand{\headrulewidth}{0.4pt}

% This MUST come after \pagestyle{fancy}: fancyhdr installs its own \chaptermark
% at that point and silently overwrites anything defined earlier. The symptom is
% a running-head change that appears to do nothing at all.
%
% The running head carries the *frame* number as well as the program, because
% in a book of numbered frames "where am I" is a frame question, not a page
% question, and the summary's back-references are frame numbers.
% \@chapapp, not \lblProgram. \appendix switches \@chapapp from \chaptername
% to \appendixname, and babel supplies both in the right language; hard-coding
% the word gives an appendix a running head reading "Program A. Answers" over a
% page whose own chapter head reads "APPENDIX A".
\makeatletter
\renewcommand{\chaptermark}[1]{\markboth{\@chapapp\ \thechapter.\ #1}{}}
\makeatother
\renewcommand{\sectionmark}[1]{\markright{\thesection\ #1}}

% ============================================================
%  FRAMES — the unit the whole method is built from
%  ---------------------------------------------------------
%  A program is a stream of numbered frames. Nearly every frame ends by asking
%  the reader for something, and the answer opens the *next* frame, so that the
%  reader can cover what is below and commit to an answer before seeing it.
%
%  The environment is called `fr` rather than `frame` for two reasons. First,
%  \frame is already defined by LaTeX2e (it draws a box in a picture
%  environment) and an environment named `frame` would silently redefine it.
%  Second, this is typed forty to seventy times per program, so it wants to be
%  short in the source.
%
%    \begin{fr}
%    \ans{$x = 3$}            % the answer to the previous frame, optional
%    Teaching text, ending in a question.
%    \end{fr}
% ============================================================

\newcounter{frameno}[chapter]
\renewcommand{\theframeno}{\arabic{frameno}}

% Frames are referenced constantly -- by the Summary, by the Quiz, by the
% "Can you?" checklist and by other programs. \label inside a frame must
% therefore capture the frame number, not the section number.
\makeatletter
\newcommand{\framelabelfix}{%
  \@namedef{theHframeno}{\theHchapter.\arabic{frameno}}%
}
\makeatother

\newcommand{\framerule}{%
  \par\penalty-200\vspace{9pt}%
  \noindent\textcolor{codeframe}{\rule{\linewidth}{0.5pt}}%
  \par\vspace{5pt}%
}

% The frame number as a filled badge, set in the text flow rather than in the
% margin. A margin number would have to alternate sides under `twoside` and
% would be the first thing to overflow on a 17 cm page; a badge cannot.
\newcommand{\framebadge}{%
  \begingroup
  \setlength{\fboxsep}{2.5pt}%
  \colorbox{mblue}{\textcolor{white}{\bfseries\scriptsize\theframeno}}%
  \endgroup\hspace{0.7em}%
}

\newenvironment{fr}{%
  \framerule
  \refstepcounter{frameno}%
  \nopagebreak
  \noindent\framebadge\ignorespaces
}{\par}

% The answer to the previous frame. Set apart and centred so the eye can find
% the edge to cover, and so a reader who has answered can check in one glance
% without reading on.
\newcommand{\ans}[1]{%
  \par\nobreak\vspace{-4pt}%
  \begin{tcolorbox}[
    enhanced, sharp corners, boxrule=0pt, leftrule=2.5pt,
    colback=mfaint, colframe=mgreen,
    left=8pt, right=8pt, top=4pt, bottom=4pt,
    width=0.86\linewidth, center, halign=center,
    before skip=2pt, after skip=7pt]
    #1
  \end{tcolorbox}%
  \nobreak\noindent\ignorespaces
}

% A multi-paragraph or worked answer, where a centred box would be wrong.
\newenvironment{ansblock}{%
  \par\nopagebreak
  \begin{tcolorbox}[
    enhanced, breakable, sharp corners, boxrule=0pt, leftrule=2.5pt,
    colback=mfaint, colframe=mgreen,
    left=8pt, right=8pt, top=5pt, bottom=5pt]
}{%
  \end{tcolorbox}%
  \nopagebreak\par\vspace{2pt}\noindent\ignorespaces
}

% The row of dots. Stroud's signal that the reader is to supply a word, a
% number or an expression. \blank is inline; \dotline occupies its own line,
% which is what a longer derivation step wants.
% \penalty0 offers TeX a break opportunity immediately before the dots. Without
% it the run of dots is glued to the last word of the question, and a long
% question in the Polish edition -- where the words are longer and hyphenate
% less freely -- runs into the margin.
\newcommand{\blank}[1][4em]{\penalty0\makebox[#1]{\dotfill}}
\newcommand{\dotline}{\par\vspace{2pt}\noindent\makebox[\linewidth]{\dotfill}\par\vspace{2pt}}

% "Now one for you to do." The third rung of the scaffolding gradient, and the
% one most often skipped by authors who think they have written an exercise
% when they have written another worked example.
\newcommand{\yourturn}{%
  \par\vspace{4pt}\noindent
  {\bfseries\color{mgreen}\lblYourTurn}\par\nobreak\vspace{2pt}\noindent\ignorespaces
}

% ============================================================
%  THE PROGRAM SKELETON
%  ---------------------------------------------------------
%  Every program carries the same seven parts, in the same order. They are
%  macros rather than a convention so that a missing one is a build-time
%  absence rather than something a reader notices.
% ============================================================

% ---- Learning outcomes, on entry ----------------------------------------
% Stored as well as printed, because "Can you?" at the end of the program must
% be the same list, one for one. Storing it is what makes that structural
% rather than a promise: the checklist is generated from the outcomes, so the
% two cannot drift.
\newcounter{outcomeno}[chapter]

\makeatletter
% The key under which everything belonging to the current program is stored.
% It must expand to the printed program number (F1, F2, ..., 1, 2, ...), which
% is what makes an answer findable from the program that owns it.
\newcommand{\mfa@progkey}{\thechapter}
\newcommand{\outcome}[2]{% #1 = frames that teach it, #2 = the outcome
  \stepcounter{outcomeno}%
  \csgdef{mfa@out@\mfa@progkey @\theoutcomeno}{#2}%
  \csgdef{mfa@outfr@\mfa@progkey @\theoutcomeno}{#1}%
  % The "Can you?" row is built here rather than looped over later. A \loop
  % inside a tabularx body does not survive alignment scanning: TeX needs to
  % see the & and the row break while it reads the cell, and a conditional
  % hides them. Accumulating the rows as tokens at declaration time sidesteps
  % that entirely, and has the better property anyway -- the checklist is
  % literally made of the outcomes, so the two cannot disagree.
  \csgappto{mfa@rows@\mfa@progkey}{%
    #2 & {\footnotesize\color{mgrey}#1}
      & \ratingcell & \ratingcell & \ratingcell & \ratingcell & \ratingcell
    \tabularnewline}%
  \item #2\hfill{\footnotesize\color{mgrey}[\lblFrames~#1]}%
}
\makeatother

\makeatletter
\newenvironment{outcomes}{%
  \setcounter{outcomeno}{0}%
  \csgdef{mfa@rows@\mfa@progkey}{}%
  \begin{tcolorbox}[
    enhanced, breakable, sharp corners,
    colback=mlight, colframe=mblue, boxrule=0pt, leftrule=3pt,
    left=8pt, right=8pt, top=6pt, bottom=6pt,
    fonttitle=\bfseries\small, title={\lblOutcomes}]
  \begin{itemize}[leftmargin=1.3em, itemsep=3pt, topsep=2pt]
}{%
  \end{itemize}
  \end{tcolorbox}
}
\makeatother

% ---- "Can you?" -----------------------------------------------------------
% Generated from the stored outcomes, not written out again by hand. Five
% columns rather than a tick box, because a self-assessment with a midpoint is
% one people actually use and a yes/no one is one they lie to.
\newcommand{\ratingcell}{\textcolor{codeframe}{\rule[-0.15em]{0.8em}{0.8em}}}

\makeatletter
\newcommand{\canyou}{%
  \section*{\lblCanYou}%
  \phantomsection\addcontentsline{toc}{section}{\lblCanYou}%
  \noindent\lblCanYouIntro\par\vspace{8pt}
  \noindent
  \begingroup\small
  \ifcsvoid{mfa@rows@\mfa@progkey}{%
    {\color{mred}\lblNoOutcomes}\par
  }{%
    \begin{tabularx}{\linewidth}{@{}Xcccccc@{}}
    \toprule
    \textbf{\lblOnEntryYouCould} & \textbf{\lblFrames} & 1 & 2 & 3 & 4 & 5 \\
    \midrule
    \csuse{mfa@rows@\mfa@progkey}
    \bottomrule
    \end{tabularx}%
  }%
  \endgroup
  \par\vspace{6pt}\noindent{\footnotesize\color{mgrey}\lblCanYouFooter}\par
}
\makeatother

% ============================================================
%  QUIZ, SUMMARY, EXERCISES
% ============================================================

% ---- Quiz -----------------------------------------------------------------
% Foundation programs only. It is a diagnostic on the way in and the same test
% on the way out, which is what lets one book serve a reader who needs the
% whole program and a reader who needs four frames of it. Each question names
% the frames that teach it, so a wrong answer routes the reader somewhere
% specific rather than back to the beginning.
\newlist{quizlist}{enumerate}{1}
\setlist[quizlist]{label=\arabic*., leftmargin=2.0em, itemsep=5pt, topsep=3pt}

\makeatletter
\newenvironment{quiz}{%
  \section*{\lblQuiz}%
  \phantomsection\addcontentsline{toc}{section}{\lblQuiz}%
  \def\mfa@exkey{Q\arabic{quizlisti}}%
  \noindent\lblQuizIntro\par\vspace{5pt}
  \begin{quizlist}
}{%
  \end{quizlist}
}
\makeatother

% The frames that teach the question just asked. Two forms, because a Quiz
% question that routes to a single frame reading "Frames 22" is the kind of
% small wrongness a reader notices and an author never does.
\newcommand{\teachesat}[1]{\hfill{\footnotesize\color{mgrey}(\lblFrames~#1)}}
\newcommand{\teachesatone}[1]{\hfill{\footnotesize\color{mgrey}(\lblFrame~#1)}}

% ---- Summary --------------------------------------------------------------
% Every item carries the frame number it came from, in square brackets. That
% is the eighth-edition improvement and it is the single cheapest thing in the
% format: it turns a summary into a return index, so a reader who fails one
% line of "Can you?" knows exactly which four frames to reread.
\newenvironment{summarybox}{%
  \section*{\lblSummary}%
  \phantomsection\addcontentsline{toc}{section}{\lblSummary}%
  \begin{tcolorbox}[
    enhanced, breakable, sharp corners,
    colback=mlight, colframe=mblue, boxrule=0pt, leftrule=3pt,
    left=8pt, right=8pt, top=6pt, bottom=6pt]
  \begin{itemize}[leftmargin=1.3em, itemsep=5pt, topsep=2pt]
}{%
  \end{itemize}
  \end{tcolorbox}
}

% \sumitem{12}{text} -- the frame reference is not decoration, it is the point.
\newcommand{\sumitem}[2]{\item #2~{\footnotesize\color{mgrey}[#1]}}

% ---- Answers ---------------------------------------------------------------
%  Answers are authored at the point of use and printed at the back.
%
%  They are carried there by a global macro store rather than by an auxiliary
%  file. The sibling books collect their figure and screenshot manifests with
%  \@starttoc and a custom file extension, which works because those entries
%  are one line of text; it is the wrong mechanism for an answer, because an
%  answer contains displayed maths, and material written to an .aux-style file
%  and read back is expanded at shipout, where fragile commands break in ways
%  whose error messages name neither the answer nor the program.
%
%  \include is sequential within a run, so a macro defined globally while
%  typesetting program F1 is still defined when the back matter is typeset.
%  That is all this needs.
\makeatletter
\newcommand{\mfa@storeanswer}[3]{% #1 program key, #2 item key, #3 body
  \csgdef{mfa@a@#1@#2}{#3}%
  % \listcsxadd, not \listcsgadd: the item must be EXPANDED as it is stored.
  % Storing the token \mfa@exkey means every entry in the list expands to
  % whatever the key happens to be at the time the back matter is typeset,
  % which is the last one.
  \listcsxadd{mfa@akeys@#1}{#2}%
}
% Used inside an exercise item. Prints nothing here; the body reappears at the
% back of the book under this program's heading.
\newcommand{\answerto}[1]{\mfa@storeanswer{\mfa@progkey}{\mfa@exkey}{#1}}
\providecommand{\mfa@exkey}{?}
\makeatother

\newlist{exlist}{enumerate}{1}
\setlist[exlist]{label=\arabic*., leftmargin=2.0em, itemsep=6pt, topsep=3pt}
\newlist{fplist}{enumerate}{1}
\setlist[fplist]{label=\arabic*., leftmargin=2.0em, itemsep=5pt, topsep=3pt}

% These MUST be defined inside \makeatletter. Outside it, \def\mfa@exkey{...}
% parses as \def\mfa @exkey{...} -- a definition of \mfa with a delimited
% parameter text -- which raises no error at all and quietly leaves every
% answer filed under the same key. That failure is invisible until the answers
% section at the back of the book prints one program's answers twice.
\makeatletter
\newenvironment{testexercises}{%
  \section*{\lblTestExercises}%
  \phantomsection\addcontentsline{toc}{section}{\lblTestExercises}%
  \def\mfa@exkey{T\arabic{exlisti}}%
  \noindent\lblTestIntro\par\vspace{5pt}
  \begin{exlist}
}{%
  \end{exlist}
}

\newenvironment{furtherproblems}{%
  \section*{\lblFurther}%
  \phantomsection\addcontentsline{toc}{section}{\lblFurther}%
  \def\mfa@exkey{P\arabic{fplisti}}%
  \noindent\lblFurtherIntro\par\vspace{5pt}
  \begin{fplist}
}{%
  \end{fplist}
}
\makeatother

% ---- The programs register -------------------------------------------------
% \program{Title} opens a program: it is \chapter plus the bookkeeping that
% lets the back matter reproduce this program's answers under its own name.
\makeatletter
\newcommand{\program}[1]{%
  \chapter{#1}%
  \csgdef{mfa@title@\thechapter}{#1}%
  % Expanded, for the same reason as the answer keys above.
  \listxadd{\mfaprograms}{\thechapter}%
}
\newcommand{\mfa@printoneanswer}[2]{% #1 program key, #2 item key
  \item[\textbf{#2}] \csuse{mfa@a@#1@#2}%
}
\newcommand{\mfa@answerprogram}[1]{%
  \subsection*{\lblProgram~#1\quad\textmd{\itshape\csuse{mfa@title@#1}}}%
  \ifcsdef{mfa@akeys@#1}{%
    \begin{list}{}{%
      \setlength{\leftmargin}{3.2em}\setlength{\labelwidth}{2.7em}%
      \setlength{\labelsep}{0.5em}\setlength{\itemsep}{3pt}%
      \setlength{\topsep}{3pt}\setlength{\parsep}{0pt}%
      \renewcommand{\makelabel}[1]{\hfill##1}}
      \forlistcsloop{\mfa@printoneanswer{#1}}{mfa@akeys@#1}
    \end{list}%
  }{%
    \par\noindent{\color{mred}\lblNoAnswers}\par
  }%
}
% The answers body, without a heading of its own, so the appendix that carries
% it can be a numbered appendix like any other and appear in the running heads
% and the ToC in the ordinary way.
\newcommand{\answersbody}{%
  \noindent\lblAnswersIntro\par\vspace{8pt}
  \ifdef{\mfaprograms}{\forlistloop{\mfa@answerprogram}{\mfaprograms}}{%
    \noindent{\color{mred}\lblNoAnswers}\par}%
}
% Unnumbered form, for a draft build that has no appendix wiring yet.
\newcommand{\printanswers}{%
  \chapter*{\lblAnswers}%
  \phantomsection\addcontentsline{toc}{chapter}{\lblAnswers}%
  \markboth{\lblAnswers}{\lblAnswers}%
  \answersbody
}
\makeatother

% ============================================================
%  ADMONITIONS
%  ---------------------------------------------------------
%  A deliberately small, fixed vocabulary. Each one earns its place by doing
%  something no other box does; a seventh would be a smell.
% ============================================================

\newtcolorbox{note}[1][]{
  enhanced, breakable, sharp corners,
  colback=mlight, colframe=mblue, boxrule=0pt, leftrule=3pt,
  left=8pt, right=8pt, top=6pt, bottom=6pt,
  fonttitle=\bfseries\small, title={\lblNote}, #1
}

\newtcolorbox{warning}[1][]{
  enhanced, breakable, sharp corners,
  colback=mlight, colframe=mred, boxrule=0pt, leftrule=3pt,
  left=8pt, right=8pt, top=6pt, bottom=6pt,
  fonttitle=\bfseries\small, title={\lblWarning}, #1
}

% The misconception, stated as a summary after it has been walked into. The
% trap itself belongs in the frames -- elicit the wrong answer, then say
% flatly that it is wrong -- and this box is the thing the reader marks with a
% finger and comes back to.
\newtcolorbox{trapbox}[1][]{
  enhanced, breakable, sharp corners,
  colback=white, colframe=mred, boxrule=0.8pt,
  left=8pt, right=8pt, top=6pt, bottom=6pt,
  fonttitle=\bfseries\small, title={\lblTrap}, #1
}

% Where the mathematics just done shows up in a system the reader has actually
% deployed. This is the box that separates the book from a maths textbook, and
% it is also the one most easily filled with waffle: if it cannot name a
% specific line of a specific system, it should not be written.
\newtcolorbox{aibox}[1][]{
  enhanced, breakable, sharp corners,
  colback=mlight, colframe=mteal, boxrule=0pt, leftrule=3pt,
  left=8pt, right=8pt, top=6pt, bottom=6pt,
  fonttitle=\bfseries\small, title={\lblInAI}, #1
}

% What is being asserted rather than proved, and where to read the proof.
% Stroud's method buys fluency at the cost of rigour, and the honest thing is
% to say so at each point rather than once in an introduction nobody rereads.
\newtcolorbox{rigourbox}[1][]{
  enhanced, breakable, sharp corners,
  colback=white, colframe=mgrey, boxrule=0.6pt,
  left=8pt, right=8pt, top=6pt, bottom=6pt,
  fonttitle=\bfseries\small, title={\lblRigour}, #1
}

% Notation that genuinely differs -- between the Polish and English editions,
% and between disciplines. Matrix calculus alone has two incompatible layout
% conventions, and half the confusion about backpropagation is a transpose
% that came from the other one.
\newtcolorbox{notationbox}[1][]{
  enhanced, breakable, sharp corners,
  colback=mlight, colframe=mamber, boxrule=0pt, leftrule=3pt,
  left=8pt, right=8pt, top=6pt, bottom=6pt,
  fonttitle=\bfseries\small, title={\lblNotation}, #1
}

% The maths analogue of the sibling books' "run this before you trust it".
% A number in this book is either produced by a script under code/ and pulled
% in with \val, or it carries this box. Nothing ships to a reader with one
% still attached.
\newtcolorbox{verifybox}[1][]{
  enhanced, breakable, sharp corners,
  colback=white, colframe=mamber, boxrule=0.8pt,
  left=8pt, right=8pt, top=6pt, bottom=6pt,
  fonttitle=\bfseries\small, title={\lblVerify}, #1
}

\newtcolorbox{exercisebox}[1][]{
  enhanced, breakable, sharp corners,
  colback=white, colframe=mgreen, boxrule=0pt, leftrule=3pt,
  left=8pt, right=8pt, top=6pt, bottom=6pt,
  fonttitle=\bfseries\small, title={\lblExercise}, #1
}

% ============================================================
%  CODE LISTINGS
%  ---------------------------------------------------------
%  Code appears here for one reason only: to check a number. Every listing in
%  this book is runnable and every number it prints is the number in the text.
% ============================================================
\lstdefinelanguage{pythonx}{
  language=Python,
  morekeywords={as, with, assert, match, case},
  morekeywords=[2]{
    float64,float32,float16,bfloat16,int8,uint8,
    np,numpy,array,asarray,frombuffer,nextafter,finfo,iinfo,
    exp,log,log1p,expm1,logaddexp,sum,max,abs,sqrt,mean,var
  },
  sensitive=true
}

\lstdefinestyle{mfacode}{
  basicstyle=\ttfamily\footnotesize,
  keywordstyle=\color{kw}\bfseries,
  keywordstyle=[2]\color{typ},
  stringstyle=\color{str},
  commentstyle=\color{cmt}\itshape,
  numbers=left,
  numberstyle=\ttfamily\tiny\color{mgrey},
  numbersep=8pt,
  backgroundcolor=\color{codebg},
  frame=single,
  rulecolor=\color{codeframe},
  framesep=6pt,
  breaklines=true,
  breakatwhitespace=false,
  postbreak=\mbox{\textcolor{mgrey}{$\hookrightarrow$}\space},
  showstringspaces=false,
  tabsize=4,
  columns=fullflexible,
  keepspaces=true,
  captionpos=b,
  aboveskip=1.2em,
  belowskip=1.2em,
  % Maps the characters most likely to be pasted into a listing by accident.
  %
  % Do NOT add a mapping for U+00A0 (non-breaking space). It looks like the
  % obvious next entry and it makes `listings` abort the run with "Improper
  % alphabetic constant" -- fatal, no PDF, and the message names neither the
  % character nor this line. The ASCII-in-listings convention is the guard
  % against it instead. This trap has now been hit in two of the three books
  % in this series.
  literate={ł}{{\l}}1 {Ł}{{\L}}1 {ą}{{\k a}}1 {ę}{{\k e}}1
           {ś}{{\'s}}1 {ć}{{\'c}}1 {ż}{{\.z}}1 {ź}{{\'z}}1 {ó}{{\'o}}1 {ń}{{\'n}}1
           {Ą}{{\k A}}1 {Ę}{{\k E}}1 {Ś}{{\'S}}1 {Ć}{{\'C}}1
           {Ż}{{\.Z}}1 {Ź}{{\'Z}}1 {Ó}{{\'O}}1 {Ń}{{\'N}}1
           {—}{{-{}-}}2 {–}{{-}}1 {‘}{{'}}1 {’}{{'}}1
           {“}{{"}}1 {”}{{"}}1 {°}{{\textdegree}}1
}

\lstset{style=mfacode}

\lstnewenvironment{python}[1][]
  {\lstset{language=pythonx,style=mfacode,#1}}
  {}

\lstnewenvironment{shellcmd}[1][]
  {\lstset{language=bash,style=mfacode,numbers=none,keywordstyle=\color{mteal}\bfseries,#1}}
  {}

% Terminal transcripts: no line numbers, no keyword colouring. Everything in
% one of these was copied from a real run.
\lstnewenvironment{console}[1][]
  {\lstset{style=mfacode,numbers=none,basicstyle=\ttfamily\scriptsize,
           keywordstyle=\color{black},backgroundcolor=\color{white},#1}}
  {}

\newcommand{\pyfile}[3]{%
  \lstinputlisting[language=pythonx,style=mfacode,caption={#2},label={#3}]{#1}%
}
\newcommand{\pyregion}[4]{%
  \lstinputlisting[
    language=pythonx, style=mfacode,
    linerange={--8<--\ [start:#2]--- 8<--\ [end:#2]},
    includerangemarker=false,
    caption={#3}, label={#4}]{#1}%
}

% Inline literals
\newcommand{\code}[1]{\texttt{\small #1}}
\newcommand{\pkg}[1]{\texttt{\small\color{mblue}#1}}
\newcommand{\api}[1]{\texttt{\small\color{mteal}#1}}

% ============================================================
%  COMPUTED VALUES
%  ---------------------------------------------------------
%  The house rule in the sibling books is "run the listing, or mark it". The
%  analogue here is stronger, because a wrong digit in a maths book is not a
%  broken example, it is a lie the reader will carry.
%
%  Every number that came out of a computation is written by a script under
%  code/ into figures/values/<program>.tex as
%
%      \mfaval{f01.eps64}{2.220446049250313e-16}
%
%  and pulled into the text with \val{f01.eps64}. The script is the source of
%  truth; the book cannot disagree with it, because the book does not contain
%  the digits. A value the scripts no longer produce prints a visible marker
%  and is counted by `make debt`.
%
%  \val also applies the locale's decimal separator, which is how the Polish
%  edition gets its decimal comma without a translator having to remember.
% ============================================================
\IfFileExists{siunitx.sty}{%
  \usepackage{siunitx}%
  \sisetup{
    group-digits = integer,
    group-minimum-digits = 5,
    group-separator = {\,},
    exponent-product = \cdot,
    output-decimal-marker = {\mfadecimalmarker}
  }%
  \newcommand{\mfanum}[1]{\num{##1}}%
}{%
  \newcommand{\mfanum}[1]{##1}%
}

\makeatletter
\newcommand{\mfaval}[2]{\csgdef{mfaval@#1}{#2}\listxadd{\mfavalues}{#1}}
% \num on a non-numeric body is a FATAL siunitx error -- "Invalid number" --
% and under -halt-on-error that is the whole build. A script emitting a word
% where the book expected a number is a plausible mistake.
%
% The check is NOT done here. Deciding in TeX whether a token list is a number
% means parsing it character by character, which is fragile and was got wrong
% once already. The generator knows the answer for free, so it emits \mfaval
% for a number and \mfavaltext for a word, and this end only has to enforce
% which macro reads which.
\newcommand{\val}[1]{%
  \ifcsdef{mfaval@#1}{%
    \ifcsdef{mfavaltext@#1}{%
      \PackageError{mfabook}{Value `#1' is not a number}%
        {\string\val\space passes its body to siunitx, which refuses anything
         that is not a number. Use \string\valtext{#1} instead, or fix the
         script in code/ that emits this key.}%
    }{}%
    \mfanum{\csuse{mfaval@#1}}%
  }{\textcolor{mred}{\textbf{[?\,#1]}}}%
}
% A value that is deliberately a word rather than a number -- a format name, a
% verdict, a unit. Rare, and a visibly different macro so nobody reaches for
% \val and gets an error they cannot place.
\newcommand{\valtext}[1]{%
  \ifcsdef{mfaval@#1}{\csuse{mfaval@#1}}%
    {\textcolor{mred}{\textbf{[?\,#1]}}}%
}
\newcommand{\mfavaltext}[2]{%
  \csgdef{mfaval@#1}{#2}\csgdef{mfavaltext@#1}{}\listxadd{\mfavalues}{#1}%
}

% The raw stored string, for the cases where a number must appear exactly as
% the machine printed it -- inside a console block, or where the digits
% themselves are the point and grouping them would obscure the pattern.
\newcommand{\rawval}[1]{%
  \ifcsdef{mfaval@#1}{\csuse{mfaval@#1}}%
    {\textcolor{mred}{\textbf{[?\,#1]}}}%
}
\makeatother

% figures/values/all.tex is generated by `make numbers` and does nothing but
% \input each program's value file. A build with no values yet still compiles;
% every \val{} in it prints a visible marker instead of a number, which is the
% correct behaviour for a draft and an obvious failure for a release.
\IfFileExists{figures/values/all.tex}{% Generated by `make numbers` --- do not edit.
\input{figures/values/f01}
}{%
  \AtBeginDocument{\PackageWarning{mfabook}{No computed values found. Run `make numbers`.}}}

% ============================================================
%  DIAGRAMS — Mermaid pipeline
%  ---------------------------------------------------------
%  \mermaidfig{key}{caption}{one-line manifest description}
%
%  Source of truth is figures/mermaid/<lang>/<key>.mmd, committed. `make
%  diagrams` renders each to figures/diagrams/<lang>/<key>.pdf, which is build
%  output and is not committed, so a diagram change reviews as a text diff.
%
%  The source is per-language because a diagram with English labels in the
%  Polish edition is exactly the kind of half-translation that makes a
%  bilingual book feel machine-made. CI checks that both languages carry the
%  same set of keys.
%
%  If the rendered PDF is absent the macro draws a placeholder AND typesets the
%  Mermaid source, so a build without mermaid-cli still shows the structure.
% ============================================================
\makeatletter
\newcommand{\listofdiagrams}{%
  \section*{\lblDiagramManifest}%
  \phantomsection\addcontentsline{toc}{section}{\lblDiagramManifest}%
  \@starttoc{dgm}%
}
\makeatother

\newcommand{\mermaidfig}[3]{%
  \addcontentsline{dgm}{subsection}{\protect\numberline{}\texttt{#1.mmd} --- #3}%
  \IfFileExists{figures/diagrams/\booklang/#1.pdf}{%
    \begin{figure}[htbp]
    \centering
    \includegraphics[width=0.95\linewidth,height=0.42\textheight,keepaspectratio]{figures/diagrams/\booklang/#1.pdf}%
    \caption{#2}\label{fig:#1}
    \end{figure}%
  }{%
    % The unrendered case deliberately does NOT float. A Mermaid source
    % typeset verbatim is often taller than a page, and a float cannot break,
    % so wrapping it in `figure` produces an overfull box per diagram and a
    % pile of tcolorbox warnings. In the flow it breaks like any other text.
    \par\medskip
    \begin{tcolorbox}[
      enhanced, breakable, width=\linewidth, sharp corners,
      colback=white, colframe=mgrey, boxrule=0.8pt,
      left=8pt, right=8pt, top=8pt, bottom=8pt]
      {\bfseries\color{mgrey}\small \lblNotRendered}\\[4pt]
      {\ttfamily\scriptsize\color{mgrey} figures/mermaid/\booklang/#1.mmd}\\[6pt]
      \IfFileExists{figures/mermaid/\booklang/#1.mmd}{%
        % ASCII only, like every other listing in this book: the fallback runs
        % the Mermaid source through `listings`, which aborts on a multi-byte
        % character it has no literate mapping for.
        \lstinputlisting[style=mfacode,numbers=none,basicstyle=\ttfamily\scriptsize,
                         backgroundcolor=\color{mlight},frame=none,
                         keywordstyle=\color{black}]{figures/mermaid/\booklang/#1.mmd}%
      }{%
        {\small\color{mred} \lblSourceMissing: \texttt{figures/mermaid/\booklang/#1.mmd}}%
      }%
    \end{tcolorbox}%
    \captionof{figure}{#2}\label{fig:#1}
    \par\medskip
  }%
}

% ============================================================
%  PROGRAM STUBS
%  ---------------------------------------------------------
%  Prints a visible marker so an unwritten program cannot be mistaken for a
%  written one and the page count never flatters the draft. `make debt` counts
%  them, and CI publishes the count on every build.
% ============================================================
\newcommand{\programstub}[1]{%
  \begin{tcolorbox}[
    enhanced, breakable, width=\linewidth, sharp corners,
    colback=mlight, colframe=mred, boxrule=1pt,
    borderline={0.6pt}{3pt}{mred, dashed},
    left=12pt, right=12pt, top=12pt, bottom=12pt]
    {\bfseries\color{mred}\large \lblNotWritten}\\[8pt]
    \small\raggedright #1
  \end{tcolorbox}%
}

% ============================================================
%  MATHEMATICAL OPERATORS
%  ---------------------------------------------------------
%  Language-invariant names live here. The ones that genuinely differ between
%  Polish and English usage -- tan/tg, Var/D^2, gcd/NWD, and the interval
%  brackets -- are defined in lang/en.tex and lang/pl.tex instead, so the
%  notation contract is executed by the build rather than remembered by a
%  translator.
% ============================================================
\DeclareMathOperator*{\argmin}{arg\,min}
\DeclareMathOperator*{\argmax}{arg\,max}
\DeclareMathOperator{\tr}{tr}
\DeclareMathOperator{\rank}{rank}
\DeclareMathOperator{\diag}{diag}
\DeclareMathOperator{\sgn}{sgn}
\DeclareMathOperator{\relu}{ReLU}
\DeclareMathOperator{\softmax}{softmax}
\DeclareMathOperator{\logsumexp}{logsumexp}
\DeclareMathOperator{\fl}{fl}
\DeclareMathOperator{\ulp}{ulp}
\DeclareMathOperator{\round}{round}

% A bare \log is FORBIDDEN in this book, and tools/check_notation.py fails the
% build on one. In Polish textbooks \log means base ten; in machine-learning
% writing it means base e. The same three characters carry opposite meanings to
% the two readerships this book serves, and they collide hardest in exactly the
% programs where it matters most -- entropy in bits and cross-entropy in nats
% are two programs apart. Write \ln for nats and \logb{2} for bits. It costs
% two characters and removes the ambiguity entirely.
\makeatletter
\let\mfa@origlog\log
\renewcommand{\log}{%
  \PackageError{mfabook}{A bare \string\log\space is forbidden in this book}%
    {Write \string\ln\space for natural logarithms or \string\logb{2}\space for
     bits. The base is never left to a convention here; see Appendix B.}%
  \mfa@origlog}
% The one legitimate form: an explicit base.
\newcommand{\logb}[1]{\mathop{\mfa@origlog}\nolimits_{#1}}
% ...and a way to write the forbidden thing when the subject IS the ban, which
% is Appendix B and nowhere else.
\newcommand{\mfalogplain}{\mathop{\mfa@origlog}\nolimits}
\makeatother

\newcommand{\R}{\mathbb{R}}
\newcommand{\N}{\mathbb{N}}
\newcommand{\Z}{\mathbb{Z}}
\newcommand{\Q}{\mathbb{Q}}
\newcommand{\C}{\mathbb{C}}
\newcommand{\Fset}{\mathbb{F}}          % the floating-point numbers, F1's subject
\newcommand{\vect}[1]{\mathbf{#1}}
\newcommand{\mat}[1]{\mathbf{#1}}
\newcommand{\T}{^{\mathsf{T}}}
\newcommand{\norm}[1]{\left\lVert #1 \right\rVert}
\newcommand{\abs}[1]{\left\lvert #1 \right\rvert}
\newcommand{\inner}[2]{\left\langle #1, #2 \right\rangle}
\newcommand{\eps}{\varepsilon}

% ============================================================
%  INDEX
%  ---------------------------------------------------------
%  Two-column, and its entries include \texttt{} names that do not hyphenate.
%  Ragged right keeps multi-word entries off the margin; it cannot help a
%  single token wider than the column, so keep verbatim index entries under
%  about 29 characters and index longer names by concept instead. Both sibling
%  books learned this the same way.
% ============================================================
\apptocmd{\theindex}{\raggedright}{}{%
  \PackageWarning{mfabook}{Could not patch theindex for ragged-right setting}}

\makeindex

% ============================================================
%  PINNED FACTS — single source of truth
%  Change these here; they propagate through both editions.
% ============================================================
\newcommand{\bookedition}{0.1}
% \pinnedon is a user-visible string and lives in lang/<lang>.tex, not here.
% It said "August 2026" on the Polish title page for exactly as long as it
% took somebody to read the Polish title page.
\newcommand{\pyver}{Python 3.13}
\newcommand{\numpyver}{2.3}
; nothing else differs between them at
%  this level. Every user-visible string lives in lang/en.tex or lang/pl.tex,
%  so a label that appears in one edition and not the other is a missing macro
%  rather than a silent divergence.
% ============================================================

% \booklang must be set by the main file. Default to English rather than
% failing, so a stray % ============================================================
%  preamble.tex — styling, environments, macros
%
%  Shared by BOTH editions. main-en.tex and main-pl.tex each define
%  \booklang before % ============================================================
%  preamble.tex — styling, environments, macros
%
%  Shared by BOTH editions. main-en.tex and main-pl.tex each define
%  \booklang before \input{preamble}; nothing else differs between them at
%  this level. Every user-visible string lives in lang/en.tex or lang/pl.tex,
%  so a label that appears in one edition and not the other is a missing macro
%  rather than a silent divergence.
% ============================================================

% \booklang must be set by the main file. Default to English rather than
% failing, so a stray \input{preamble} still compiles and says what it did.
\providecommand{\booklang}{en}

\usepackage{etoolbox}   % loaded first: the language switch below needs it

% ---------- Page geometry ----------
% Same 17 x 24 cm block as the other two books in this series, so a reader who
% owns one recognises the second. Frames are short, so the measure matters
% more here than in a prose book: a frame that runs past a page turn defeats
% the cover-the-next-frame discipline the whole method rests on.
\usepackage[
  paperwidth=17cm, paperheight=24cm,
  inner=2.2cm, outer=1.8cm, top=2.2cm, bottom=2.4cm,
  head=23pt, headsep=0.7cm, footskip=1.2cm
]{geometry}

% ---------- Encoding & fonts ----------
\usepackage[T1]{fontenc}
\usepackage[utf8]{inputenc}
\IfFileExists{lmodern.sty}{\usepackage{lmodern}}{}
\IfFileExists{newtxtext.sty}{\usepackage{newtxtext}}{}
% amsmath goes here, not down with the other packages, because newtxmath
% patches amsmath's internals and must be loaded AFTER it. The wrong order does
% not fail on a machine without newtxmath -- which is most CI images -- and
% fails on the author's full TeX Live.
%
% amssymb only when newtxmath is absent: newtxmath supplies the AMS symbols
% itself (its amssymbols option is on by default) and loading both clashes.
\usepackage{amsmath}
\IfFileExists{newtxmath.sty}{\usepackage{newtxmath}}{\usepackage{amssymb}}
\IfFileExists{inconsolata.sty}{\usepackage[varqu,scaled=0.95]{inconsolata}}{}
\IfFileExists{lmodern.sty}{\usepackage{microtype}}{\usepackage[expansion=false]{microtype}}

% ---------- Language ----------
% Loading babel with a language whose .ldf is absent is a *fatal* error, not a
% warning, so both the package and each language file are probed before either
% is requested. This is inherited from the sibling books, where a minimal TeX
% installation without texlive-lang-polish aborted the run with no PDF and an
% error message that named neither babel nor the language.
%
% The Polish edition needs Polish as the main language for hyphenation; the
% English edition needs British. Both editions load the other language too,
% because the glossary appendix is bilingual in both.
\IfFileExists{babel.sty}{%
  \IfFileExists{polish.ldf}{%
    \IfFileExists{british.ldf}{%
      \ifdefstring{\booklang}{pl}%
        {\usepackage[british,main=polish]{babel}}%
        {\usepackage[polish,main=british]{babel}}%
    }{\usepackage[polish]{babel}}%
  }{%
    \IfFileExists{british.ldf}{\usepackage[british]{babel}}{}%
  }%
}{}

% babel-polish makes " an active shorthand. listings reads its body verbatim
% and an active character in a verbatim context is a category-code accident
% waiting to happen, so the shorthand is turned off globally. Polish quotation
% marks are produced with \enquote-style macros from lang/pl.tex instead.
\ifdefstring{\booklang}{pl}{\AtBeginDocument{\shorthandoff{"}}}{}

% ---------- Core packages ----------
\usepackage{graphicx}
\usepackage{xcolor}
\usepackage{booktabs}
\usepackage{tabularx}
\usepackage{longtable}
\usepackage{array}
\usepackage{enumitem}
\IfFileExists{mathtools.sty}{\usepackage{mathtools}}{}
\usepackage{listings}
\usepackage[most]{tcolorbox}
\usepackage{fancyhdr}
\usepackage{titlesec}
\usepackage{caption}
\usepackage{imakeidx}
% csquotes with autostyle follows babel's active language, so the identical
% source \enquote{x} sets 'x' in the English edition and the Polish quotation
% marks in the Polish one. This is not a nicety: babel-polish makes " an active
% shorthand, which is turned off above precisely because an active character
% near a listings body is a category-code accident waiting to happen -- so a
% quotation mark has to come from a macro rather than from a keyboard key.
\IfFileExists{csquotes.sty}{\usepackage[autostyle=true]{csquotes}}{%
  \providecommand{\enquote}[1]{``##1''}}
\usepackage[hidelinks]{hyperref}

% ---------- Palette ----------
% Deliberately the same family as the sibling books, shifted green so the three
% spines are distinguishable on a shelf.
\definecolor{mblue}{HTML}{1F4E79}
\definecolor{mgreen}{HTML}{2E6B4F}
\definecolor{mteal}{HTML}{0E7C7B}
\definecolor{mamber}{HTML}{B26A00}
\definecolor{mred}{HTML}{9B2C2C}
\definecolor{mgrey}{HTML}{5A5A5A}
\definecolor{mlight}{HTML}{F4F6F8}
\definecolor{mfaint}{HTML}{FBFCFD}
\definecolor{codebg}{HTML}{FAFAFA}
\definecolor{codeframe}{HTML}{D8DEE4}
\definecolor{kw}{HTML}{0B5394}
\definecolor{str}{HTML}{9B2C2C}
\definecolor{cmt}{HTML}{4F7A4F}
\definecolor{typ}{HTML}{0E7C7B}

\hypersetup{
  colorlinks=true,
  linkcolor=mblue,
  urlcolor=mteal,
  citecolor=mblue,
  pdfauthor={Konrad Cinkusz}
}

% \ifpl{Polish text}{English text} -- for the handful of places where a string
% is structural rather than content and belongs in the shared files:
% structure.tex's part titles, and the main files' front matter wiring.
% Anything longer than a few words belongs in lang/ or in the edition's own
% source, not here.
% \DeclareRobustCommand, not \newcommand. A fragile \ifpl inside a \part title
% is expanded one level as it is written to the .toc, which lands
% "\ifdefstring {en}{pl}{...}{...}" in the file -- and that fails on the next
% run with "Extra }, or forgotten \endgroup", an error that names neither the
% part nor this macro.
\DeclareRobustCommand{\ifpl}[2]{\ifdefstring{\booklang}{pl}{#1}{#2}}

% A robust macro cannot go into a PDF bookmark string either -- hyperref strips
% \ifdefstring and warns once per part title. Tell it which branch to take
% instead, so the bookmarks carry the right language rather than nothing.
\makeatletter
\pdfstringdefDisableCommands{%
  \ifdefstring{\booklang}{pl}{\let\ifpl\@firstoftwo}{\let\ifpl\@secondoftwo}%
}
\makeatother

% ---------- Language strings ----------
% Every user-visible word the macros below typeset comes from here. Nothing in
% programs/ or appendices/ may hard-code an English or a Polish word inside a
% macro argument that the other edition also uses.
\input{lang/\booklang}

% ============================================================
%  PROGRAM NUMBERING
%  ---------------------------------------------------------
%  Part F programs are F1..F13; the main programs restart at 1. That is
%  Stroud's scheme and it is worth reproducing, because the Foundation part is
%  explicitly optional and numbering it separately is what makes skipping it
%  feel permitted rather than delinquent.
%
%  \theHchapter is hyperref's *destination* name, and it must stay unique
%  across the whole document. Resetting the chapter counter without also
%  changing \theHchapter produces duplicate PDF destinations, a pile of
%  warnings, and bookmarks that jump to the wrong program.
% ============================================================
\newcommand{\foundationnumbering}{%
  \setcounter{chapter}{0}%
  \renewcommand{\thechapter}{F\arabic{chapter}}%
  \renewcommand{\theHchapter}{F\arabic{chapter}}%
}
\newcommand{\mainnumbering}{%
  \setcounter{chapter}{0}%
  \renewcommand{\thechapter}{\arabic{chapter}}%
  \renewcommand{\theHchapter}{M\arabic{chapter}}%
}

% \appendix resets the chapter counter and switches \thechapter to letters, but
% it knows nothing about \theHchapter, so appendix A would claim the same PDF
% destination as main program 1. Patch it once, here.
\apptocmd{\appendix}{\renewcommand{\theHchapter}{A\Alph{chapter}}}{}{%
  \PackageWarning{mfabook}{Could not patch \string\appendix\space for hyperref destinations}}

% A program is a chapter. \chaptername is set from the language file so the
% word above the title is "Program" in both editions (it happens to be the
% same word, which is a small piece of luck).
\AtBeginDocument{\renewcommand{\chaptername}{\lblProgram}}

% ---------- Headings ----------
% \raggedright on the title body is load-bearing, not cosmetic: at \Huge on a
% 13 cm text block a long title that cannot break overflows the margin badly.
\titleformat{\chapter}[display]
  {\normalfont\huge\bfseries\color{mblue}\raggedright}
  {\normalfont\normalsize\color{mgrey}\MakeUppercase{\chaptertitlename}\ \thechapter}
  {8pt}{\Huge\raggedright}
\titlespacing*{\chapter}{0pt}{10pt}{28pt}
\titleformat{\section}{\normalfont\Large\bfseries\color{mblue}}{\thesection}{0.8em}{}
\titleformat{\subsection}{\normalfont\large\bfseries\color{mgrey}}{\thesubsection}{0.7em}{}

% head=23pt in the geometry options above, not \setlength{\headheight}:
% fancyhdr computes the height it needs from the running-head font and the
% rule, wants 22.5 pt at this size, and silently overprints the top of the text
% block if it does not get it. Setting it through geometry keeps the text block
% where geometry put it instead of pushing it down.
\pagestyle{fancy}
\fancyhf{}
\fancyhead[LE]{\small\nouppercase{\leftmark}}
% The recto head carries the FRAME number, not the section. In a book of
% numbered frames "where am I" is a frame question: every Summary
% back-reference, every Quiz route and every cross-reference between programs
% is a frame number, and hunting for one by page is the single most annoying
% thing about reading a programmed text. \theframeno at shipout is the last
% frame begun on the page, which is what a reader flicking through wants.
% Guarded: the front matter and the appendices have no frames, and a head
% reading "Frame 0" over the introduction is worse than no head at all.
\fancyhead[RO]{\ifnum\value{frameno}>0\relax\small\lblFrame~\theframeno\fi}
\fancyfoot[LE,RO]{\small\thepage}
\renewcommand{\headrulewidth}{0.4pt}

% This MUST come after \pagestyle{fancy}: fancyhdr installs its own \chaptermark
% at that point and silently overwrites anything defined earlier. The symptom is
% a running-head change that appears to do nothing at all.
%
% The running head carries the *frame* number as well as the program, because
% in a book of numbered frames "where am I" is a frame question, not a page
% question, and the summary's back-references are frame numbers.
% \@chapapp, not \lblProgram. \appendix switches \@chapapp from \chaptername
% to \appendixname, and babel supplies both in the right language; hard-coding
% the word gives an appendix a running head reading "Program A. Answers" over a
% page whose own chapter head reads "APPENDIX A".
\makeatletter
\renewcommand{\chaptermark}[1]{\markboth{\@chapapp\ \thechapter.\ #1}{}}
\makeatother
\renewcommand{\sectionmark}[1]{\markright{\thesection\ #1}}

% ============================================================
%  FRAMES — the unit the whole method is built from
%  ---------------------------------------------------------
%  A program is a stream of numbered frames. Nearly every frame ends by asking
%  the reader for something, and the answer opens the *next* frame, so that the
%  reader can cover what is below and commit to an answer before seeing it.
%
%  The environment is called `fr` rather than `frame` for two reasons. First,
%  \frame is already defined by LaTeX2e (it draws a box in a picture
%  environment) and an environment named `frame` would silently redefine it.
%  Second, this is typed forty to seventy times per program, so it wants to be
%  short in the source.
%
%    \begin{fr}
%    \ans{$x = 3$}            % the answer to the previous frame, optional
%    Teaching text, ending in a question.
%    \end{fr}
% ============================================================

\newcounter{frameno}[chapter]
\renewcommand{\theframeno}{\arabic{frameno}}

% Frames are referenced constantly -- by the Summary, by the Quiz, by the
% "Can you?" checklist and by other programs. \label inside a frame must
% therefore capture the frame number, not the section number.
\makeatletter
\newcommand{\framelabelfix}{%
  \@namedef{theHframeno}{\theHchapter.\arabic{frameno}}%
}
\makeatother

\newcommand{\framerule}{%
  \par\penalty-200\vspace{9pt}%
  \noindent\textcolor{codeframe}{\rule{\linewidth}{0.5pt}}%
  \par\vspace{5pt}%
}

% The frame number as a filled badge, set in the text flow rather than in the
% margin. A margin number would have to alternate sides under `twoside` and
% would be the first thing to overflow on a 17 cm page; a badge cannot.
\newcommand{\framebadge}{%
  \begingroup
  \setlength{\fboxsep}{2.5pt}%
  \colorbox{mblue}{\textcolor{white}{\bfseries\scriptsize\theframeno}}%
  \endgroup\hspace{0.7em}%
}

\newenvironment{fr}{%
  \framerule
  \refstepcounter{frameno}%
  \nopagebreak
  \noindent\framebadge\ignorespaces
}{\par}

% The answer to the previous frame. Set apart and centred so the eye can find
% the edge to cover, and so a reader who has answered can check in one glance
% without reading on.
\newcommand{\ans}[1]{%
  \par\nobreak\vspace{-4pt}%
  \begin{tcolorbox}[
    enhanced, sharp corners, boxrule=0pt, leftrule=2.5pt,
    colback=mfaint, colframe=mgreen,
    left=8pt, right=8pt, top=4pt, bottom=4pt,
    width=0.86\linewidth, center, halign=center,
    before skip=2pt, after skip=7pt]
    #1
  \end{tcolorbox}%
  \nobreak\noindent\ignorespaces
}

% A multi-paragraph or worked answer, where a centred box would be wrong.
\newenvironment{ansblock}{%
  \par\nopagebreak
  \begin{tcolorbox}[
    enhanced, breakable, sharp corners, boxrule=0pt, leftrule=2.5pt,
    colback=mfaint, colframe=mgreen,
    left=8pt, right=8pt, top=5pt, bottom=5pt]
}{%
  \end{tcolorbox}%
  \nopagebreak\par\vspace{2pt}\noindent\ignorespaces
}

% The row of dots. Stroud's signal that the reader is to supply a word, a
% number or an expression. \blank is inline; \dotline occupies its own line,
% which is what a longer derivation step wants.
% \penalty0 offers TeX a break opportunity immediately before the dots. Without
% it the run of dots is glued to the last word of the question, and a long
% question in the Polish edition -- where the words are longer and hyphenate
% less freely -- runs into the margin.
\newcommand{\blank}[1][4em]{\penalty0\makebox[#1]{\dotfill}}
\newcommand{\dotline}{\par\vspace{2pt}\noindent\makebox[\linewidth]{\dotfill}\par\vspace{2pt}}

% "Now one for you to do." The third rung of the scaffolding gradient, and the
% one most often skipped by authors who think they have written an exercise
% when they have written another worked example.
\newcommand{\yourturn}{%
  \par\vspace{4pt}\noindent
  {\bfseries\color{mgreen}\lblYourTurn}\par\nobreak\vspace{2pt}\noindent\ignorespaces
}

% ============================================================
%  THE PROGRAM SKELETON
%  ---------------------------------------------------------
%  Every program carries the same seven parts, in the same order. They are
%  macros rather than a convention so that a missing one is a build-time
%  absence rather than something a reader notices.
% ============================================================

% ---- Learning outcomes, on entry ----------------------------------------
% Stored as well as printed, because "Can you?" at the end of the program must
% be the same list, one for one. Storing it is what makes that structural
% rather than a promise: the checklist is generated from the outcomes, so the
% two cannot drift.
\newcounter{outcomeno}[chapter]

\makeatletter
% The key under which everything belonging to the current program is stored.
% It must expand to the printed program number (F1, F2, ..., 1, 2, ...), which
% is what makes an answer findable from the program that owns it.
\newcommand{\mfa@progkey}{\thechapter}
\newcommand{\outcome}[2]{% #1 = frames that teach it, #2 = the outcome
  \stepcounter{outcomeno}%
  \csgdef{mfa@out@\mfa@progkey @\theoutcomeno}{#2}%
  \csgdef{mfa@outfr@\mfa@progkey @\theoutcomeno}{#1}%
  % The "Can you?" row is built here rather than looped over later. A \loop
  % inside a tabularx body does not survive alignment scanning: TeX needs to
  % see the & and the row break while it reads the cell, and a conditional
  % hides them. Accumulating the rows as tokens at declaration time sidesteps
  % that entirely, and has the better property anyway -- the checklist is
  % literally made of the outcomes, so the two cannot disagree.
  \csgappto{mfa@rows@\mfa@progkey}{%
    #2 & {\footnotesize\color{mgrey}#1}
      & \ratingcell & \ratingcell & \ratingcell & \ratingcell & \ratingcell
    \tabularnewline}%
  \item #2\hfill{\footnotesize\color{mgrey}[\lblFrames~#1]}%
}
\makeatother

\makeatletter
\newenvironment{outcomes}{%
  \setcounter{outcomeno}{0}%
  \csgdef{mfa@rows@\mfa@progkey}{}%
  \begin{tcolorbox}[
    enhanced, breakable, sharp corners,
    colback=mlight, colframe=mblue, boxrule=0pt, leftrule=3pt,
    left=8pt, right=8pt, top=6pt, bottom=6pt,
    fonttitle=\bfseries\small, title={\lblOutcomes}]
  \begin{itemize}[leftmargin=1.3em, itemsep=3pt, topsep=2pt]
}{%
  \end{itemize}
  \end{tcolorbox}
}
\makeatother

% ---- "Can you?" -----------------------------------------------------------
% Generated from the stored outcomes, not written out again by hand. Five
% columns rather than a tick box, because a self-assessment with a midpoint is
% one people actually use and a yes/no one is one they lie to.
\newcommand{\ratingcell}{\textcolor{codeframe}{\rule[-0.15em]{0.8em}{0.8em}}}

\makeatletter
\newcommand{\canyou}{%
  \section*{\lblCanYou}%
  \phantomsection\addcontentsline{toc}{section}{\lblCanYou}%
  \noindent\lblCanYouIntro\par\vspace{8pt}
  \noindent
  \begingroup\small
  \ifcsvoid{mfa@rows@\mfa@progkey}{%
    {\color{mred}\lblNoOutcomes}\par
  }{%
    \begin{tabularx}{\linewidth}{@{}Xcccccc@{}}
    \toprule
    \textbf{\lblOnEntryYouCould} & \textbf{\lblFrames} & 1 & 2 & 3 & 4 & 5 \\
    \midrule
    \csuse{mfa@rows@\mfa@progkey}
    \bottomrule
    \end{tabularx}%
  }%
  \endgroup
  \par\vspace{6pt}\noindent{\footnotesize\color{mgrey}\lblCanYouFooter}\par
}
\makeatother

% ============================================================
%  QUIZ, SUMMARY, EXERCISES
% ============================================================

% ---- Quiz -----------------------------------------------------------------
% Foundation programs only. It is a diagnostic on the way in and the same test
% on the way out, which is what lets one book serve a reader who needs the
% whole program and a reader who needs four frames of it. Each question names
% the frames that teach it, so a wrong answer routes the reader somewhere
% specific rather than back to the beginning.
\newlist{quizlist}{enumerate}{1}
\setlist[quizlist]{label=\arabic*., leftmargin=2.0em, itemsep=5pt, topsep=3pt}

\makeatletter
\newenvironment{quiz}{%
  \section*{\lblQuiz}%
  \phantomsection\addcontentsline{toc}{section}{\lblQuiz}%
  \def\mfa@exkey{Q\arabic{quizlisti}}%
  \noindent\lblQuizIntro\par\vspace{5pt}
  \begin{quizlist}
}{%
  \end{quizlist}
}
\makeatother

% The frames that teach the question just asked. Two forms, because a Quiz
% question that routes to a single frame reading "Frames 22" is the kind of
% small wrongness a reader notices and an author never does.
\newcommand{\teachesat}[1]{\hfill{\footnotesize\color{mgrey}(\lblFrames~#1)}}
\newcommand{\teachesatone}[1]{\hfill{\footnotesize\color{mgrey}(\lblFrame~#1)}}

% ---- Summary --------------------------------------------------------------
% Every item carries the frame number it came from, in square brackets. That
% is the eighth-edition improvement and it is the single cheapest thing in the
% format: it turns a summary into a return index, so a reader who fails one
% line of "Can you?" knows exactly which four frames to reread.
\newenvironment{summarybox}{%
  \section*{\lblSummary}%
  \phantomsection\addcontentsline{toc}{section}{\lblSummary}%
  \begin{tcolorbox}[
    enhanced, breakable, sharp corners,
    colback=mlight, colframe=mblue, boxrule=0pt, leftrule=3pt,
    left=8pt, right=8pt, top=6pt, bottom=6pt]
  \begin{itemize}[leftmargin=1.3em, itemsep=5pt, topsep=2pt]
}{%
  \end{itemize}
  \end{tcolorbox}
}

% \sumitem{12}{text} -- the frame reference is not decoration, it is the point.
\newcommand{\sumitem}[2]{\item #2~{\footnotesize\color{mgrey}[#1]}}

% ---- Answers ---------------------------------------------------------------
%  Answers are authored at the point of use and printed at the back.
%
%  They are carried there by a global macro store rather than by an auxiliary
%  file. The sibling books collect their figure and screenshot manifests with
%  \@starttoc and a custom file extension, which works because those entries
%  are one line of text; it is the wrong mechanism for an answer, because an
%  answer contains displayed maths, and material written to an .aux-style file
%  and read back is expanded at shipout, where fragile commands break in ways
%  whose error messages name neither the answer nor the program.
%
%  \include is sequential within a run, so a macro defined globally while
%  typesetting program F1 is still defined when the back matter is typeset.
%  That is all this needs.
\makeatletter
\newcommand{\mfa@storeanswer}[3]{% #1 program key, #2 item key, #3 body
  \csgdef{mfa@a@#1@#2}{#3}%
  % \listcsxadd, not \listcsgadd: the item must be EXPANDED as it is stored.
  % Storing the token \mfa@exkey means every entry in the list expands to
  % whatever the key happens to be at the time the back matter is typeset,
  % which is the last one.
  \listcsxadd{mfa@akeys@#1}{#2}%
}
% Used inside an exercise item. Prints nothing here; the body reappears at the
% back of the book under this program's heading.
\newcommand{\answerto}[1]{\mfa@storeanswer{\mfa@progkey}{\mfa@exkey}{#1}}
\providecommand{\mfa@exkey}{?}
\makeatother

\newlist{exlist}{enumerate}{1}
\setlist[exlist]{label=\arabic*., leftmargin=2.0em, itemsep=6pt, topsep=3pt}
\newlist{fplist}{enumerate}{1}
\setlist[fplist]{label=\arabic*., leftmargin=2.0em, itemsep=5pt, topsep=3pt}

% These MUST be defined inside \makeatletter. Outside it, \def\mfa@exkey{...}
% parses as \def\mfa @exkey{...} -- a definition of \mfa with a delimited
% parameter text -- which raises no error at all and quietly leaves every
% answer filed under the same key. That failure is invisible until the answers
% section at the back of the book prints one program's answers twice.
\makeatletter
\newenvironment{testexercises}{%
  \section*{\lblTestExercises}%
  \phantomsection\addcontentsline{toc}{section}{\lblTestExercises}%
  \def\mfa@exkey{T\arabic{exlisti}}%
  \noindent\lblTestIntro\par\vspace{5pt}
  \begin{exlist}
}{%
  \end{exlist}
}

\newenvironment{furtherproblems}{%
  \section*{\lblFurther}%
  \phantomsection\addcontentsline{toc}{section}{\lblFurther}%
  \def\mfa@exkey{P\arabic{fplisti}}%
  \noindent\lblFurtherIntro\par\vspace{5pt}
  \begin{fplist}
}{%
  \end{fplist}
}
\makeatother

% ---- The programs register -------------------------------------------------
% \program{Title} opens a program: it is \chapter plus the bookkeeping that
% lets the back matter reproduce this program's answers under its own name.
\makeatletter
\newcommand{\program}[1]{%
  \chapter{#1}%
  \csgdef{mfa@title@\thechapter}{#1}%
  % Expanded, for the same reason as the answer keys above.
  \listxadd{\mfaprograms}{\thechapter}%
}
\newcommand{\mfa@printoneanswer}[2]{% #1 program key, #2 item key
  \item[\textbf{#2}] \csuse{mfa@a@#1@#2}%
}
\newcommand{\mfa@answerprogram}[1]{%
  \subsection*{\lblProgram~#1\quad\textmd{\itshape\csuse{mfa@title@#1}}}%
  \ifcsdef{mfa@akeys@#1}{%
    \begin{list}{}{%
      \setlength{\leftmargin}{3.2em}\setlength{\labelwidth}{2.7em}%
      \setlength{\labelsep}{0.5em}\setlength{\itemsep}{3pt}%
      \setlength{\topsep}{3pt}\setlength{\parsep}{0pt}%
      \renewcommand{\makelabel}[1]{\hfill##1}}
      \forlistcsloop{\mfa@printoneanswer{#1}}{mfa@akeys@#1}
    \end{list}%
  }{%
    \par\noindent{\color{mred}\lblNoAnswers}\par
  }%
}
% The answers body, without a heading of its own, so the appendix that carries
% it can be a numbered appendix like any other and appear in the running heads
% and the ToC in the ordinary way.
\newcommand{\answersbody}{%
  \noindent\lblAnswersIntro\par\vspace{8pt}
  \ifdef{\mfaprograms}{\forlistloop{\mfa@answerprogram}{\mfaprograms}}{%
    \noindent{\color{mred}\lblNoAnswers}\par}%
}
% Unnumbered form, for a draft build that has no appendix wiring yet.
\newcommand{\printanswers}{%
  \chapter*{\lblAnswers}%
  \phantomsection\addcontentsline{toc}{chapter}{\lblAnswers}%
  \markboth{\lblAnswers}{\lblAnswers}%
  \answersbody
}
\makeatother

% ============================================================
%  ADMONITIONS
%  ---------------------------------------------------------
%  A deliberately small, fixed vocabulary. Each one earns its place by doing
%  something no other box does; a seventh would be a smell.
% ============================================================

\newtcolorbox{note}[1][]{
  enhanced, breakable, sharp corners,
  colback=mlight, colframe=mblue, boxrule=0pt, leftrule=3pt,
  left=8pt, right=8pt, top=6pt, bottom=6pt,
  fonttitle=\bfseries\small, title={\lblNote}, #1
}

\newtcolorbox{warning}[1][]{
  enhanced, breakable, sharp corners,
  colback=mlight, colframe=mred, boxrule=0pt, leftrule=3pt,
  left=8pt, right=8pt, top=6pt, bottom=6pt,
  fonttitle=\bfseries\small, title={\lblWarning}, #1
}

% The misconception, stated as a summary after it has been walked into. The
% trap itself belongs in the frames -- elicit the wrong answer, then say
% flatly that it is wrong -- and this box is the thing the reader marks with a
% finger and comes back to.
\newtcolorbox{trapbox}[1][]{
  enhanced, breakable, sharp corners,
  colback=white, colframe=mred, boxrule=0.8pt,
  left=8pt, right=8pt, top=6pt, bottom=6pt,
  fonttitle=\bfseries\small, title={\lblTrap}, #1
}

% Where the mathematics just done shows up in a system the reader has actually
% deployed. This is the box that separates the book from a maths textbook, and
% it is also the one most easily filled with waffle: if it cannot name a
% specific line of a specific system, it should not be written.
\newtcolorbox{aibox}[1][]{
  enhanced, breakable, sharp corners,
  colback=mlight, colframe=mteal, boxrule=0pt, leftrule=3pt,
  left=8pt, right=8pt, top=6pt, bottom=6pt,
  fonttitle=\bfseries\small, title={\lblInAI}, #1
}

% What is being asserted rather than proved, and where to read the proof.
% Stroud's method buys fluency at the cost of rigour, and the honest thing is
% to say so at each point rather than once in an introduction nobody rereads.
\newtcolorbox{rigourbox}[1][]{
  enhanced, breakable, sharp corners,
  colback=white, colframe=mgrey, boxrule=0.6pt,
  left=8pt, right=8pt, top=6pt, bottom=6pt,
  fonttitle=\bfseries\small, title={\lblRigour}, #1
}

% Notation that genuinely differs -- between the Polish and English editions,
% and between disciplines. Matrix calculus alone has two incompatible layout
% conventions, and half the confusion about backpropagation is a transpose
% that came from the other one.
\newtcolorbox{notationbox}[1][]{
  enhanced, breakable, sharp corners,
  colback=mlight, colframe=mamber, boxrule=0pt, leftrule=3pt,
  left=8pt, right=8pt, top=6pt, bottom=6pt,
  fonttitle=\bfseries\small, title={\lblNotation}, #1
}

% The maths analogue of the sibling books' "run this before you trust it".
% A number in this book is either produced by a script under code/ and pulled
% in with \val, or it carries this box. Nothing ships to a reader with one
% still attached.
\newtcolorbox{verifybox}[1][]{
  enhanced, breakable, sharp corners,
  colback=white, colframe=mamber, boxrule=0.8pt,
  left=8pt, right=8pt, top=6pt, bottom=6pt,
  fonttitle=\bfseries\small, title={\lblVerify}, #1
}

\newtcolorbox{exercisebox}[1][]{
  enhanced, breakable, sharp corners,
  colback=white, colframe=mgreen, boxrule=0pt, leftrule=3pt,
  left=8pt, right=8pt, top=6pt, bottom=6pt,
  fonttitle=\bfseries\small, title={\lblExercise}, #1
}

% ============================================================
%  CODE LISTINGS
%  ---------------------------------------------------------
%  Code appears here for one reason only: to check a number. Every listing in
%  this book is runnable and every number it prints is the number in the text.
% ============================================================
\lstdefinelanguage{pythonx}{
  language=Python,
  morekeywords={as, with, assert, match, case},
  morekeywords=[2]{
    float64,float32,float16,bfloat16,int8,uint8,
    np,numpy,array,asarray,frombuffer,nextafter,finfo,iinfo,
    exp,log,log1p,expm1,logaddexp,sum,max,abs,sqrt,mean,var
  },
  sensitive=true
}

\lstdefinestyle{mfacode}{
  basicstyle=\ttfamily\footnotesize,
  keywordstyle=\color{kw}\bfseries,
  keywordstyle=[2]\color{typ},
  stringstyle=\color{str},
  commentstyle=\color{cmt}\itshape,
  numbers=left,
  numberstyle=\ttfamily\tiny\color{mgrey},
  numbersep=8pt,
  backgroundcolor=\color{codebg},
  frame=single,
  rulecolor=\color{codeframe},
  framesep=6pt,
  breaklines=true,
  breakatwhitespace=false,
  postbreak=\mbox{\textcolor{mgrey}{$\hookrightarrow$}\space},
  showstringspaces=false,
  tabsize=4,
  columns=fullflexible,
  keepspaces=true,
  captionpos=b,
  aboveskip=1.2em,
  belowskip=1.2em,
  % Maps the characters most likely to be pasted into a listing by accident.
  %
  % Do NOT add a mapping for U+00A0 (non-breaking space). It looks like the
  % obvious next entry and it makes `listings` abort the run with "Improper
  % alphabetic constant" -- fatal, no PDF, and the message names neither the
  % character nor this line. The ASCII-in-listings convention is the guard
  % against it instead. This trap has now been hit in two of the three books
  % in this series.
  literate={ł}{{\l}}1 {Ł}{{\L}}1 {ą}{{\k a}}1 {ę}{{\k e}}1
           {ś}{{\'s}}1 {ć}{{\'c}}1 {ż}{{\.z}}1 {ź}{{\'z}}1 {ó}{{\'o}}1 {ń}{{\'n}}1
           {Ą}{{\k A}}1 {Ę}{{\k E}}1 {Ś}{{\'S}}1 {Ć}{{\'C}}1
           {Ż}{{\.Z}}1 {Ź}{{\'Z}}1 {Ó}{{\'O}}1 {Ń}{{\'N}}1
           {—}{{-{}-}}2 {–}{{-}}1 {‘}{{'}}1 {’}{{'}}1
           {“}{{"}}1 {”}{{"}}1 {°}{{\textdegree}}1
}

\lstset{style=mfacode}

\lstnewenvironment{python}[1][]
  {\lstset{language=pythonx,style=mfacode,#1}}
  {}

\lstnewenvironment{shellcmd}[1][]
  {\lstset{language=bash,style=mfacode,numbers=none,keywordstyle=\color{mteal}\bfseries,#1}}
  {}

% Terminal transcripts: no line numbers, no keyword colouring. Everything in
% one of these was copied from a real run.
\lstnewenvironment{console}[1][]
  {\lstset{style=mfacode,numbers=none,basicstyle=\ttfamily\scriptsize,
           keywordstyle=\color{black},backgroundcolor=\color{white},#1}}
  {}

\newcommand{\pyfile}[3]{%
  \lstinputlisting[language=pythonx,style=mfacode,caption={#2},label={#3}]{#1}%
}
\newcommand{\pyregion}[4]{%
  \lstinputlisting[
    language=pythonx, style=mfacode,
    linerange={--8<--\ [start:#2]--- 8<--\ [end:#2]},
    includerangemarker=false,
    caption={#3}, label={#4}]{#1}%
}

% Inline literals
\newcommand{\code}[1]{\texttt{\small #1}}
\newcommand{\pkg}[1]{\texttt{\small\color{mblue}#1}}
\newcommand{\api}[1]{\texttt{\small\color{mteal}#1}}

% ============================================================
%  COMPUTED VALUES
%  ---------------------------------------------------------
%  The house rule in the sibling books is "run the listing, or mark it". The
%  analogue here is stronger, because a wrong digit in a maths book is not a
%  broken example, it is a lie the reader will carry.
%
%  Every number that came out of a computation is written by a script under
%  code/ into figures/values/<program>.tex as
%
%      \mfaval{f01.eps64}{2.220446049250313e-16}
%
%  and pulled into the text with \val{f01.eps64}. The script is the source of
%  truth; the book cannot disagree with it, because the book does not contain
%  the digits. A value the scripts no longer produce prints a visible marker
%  and is counted by `make debt`.
%
%  \val also applies the locale's decimal separator, which is how the Polish
%  edition gets its decimal comma without a translator having to remember.
% ============================================================
\IfFileExists{siunitx.sty}{%
  \usepackage{siunitx}%
  \sisetup{
    group-digits = integer,
    group-minimum-digits = 5,
    group-separator = {\,},
    exponent-product = \cdot,
    output-decimal-marker = {\mfadecimalmarker}
  }%
  \newcommand{\mfanum}[1]{\num{##1}}%
}{%
  \newcommand{\mfanum}[1]{##1}%
}

\makeatletter
\newcommand{\mfaval}[2]{\csgdef{mfaval@#1}{#2}\listxadd{\mfavalues}{#1}}
% \num on a non-numeric body is a FATAL siunitx error -- "Invalid number" --
% and under -halt-on-error that is the whole build. A script emitting a word
% where the book expected a number is a plausible mistake.
%
% The check is NOT done here. Deciding in TeX whether a token list is a number
% means parsing it character by character, which is fragile and was got wrong
% once already. The generator knows the answer for free, so it emits \mfaval
% for a number and \mfavaltext for a word, and this end only has to enforce
% which macro reads which.
\newcommand{\val}[1]{%
  \ifcsdef{mfaval@#1}{%
    \ifcsdef{mfavaltext@#1}{%
      \PackageError{mfabook}{Value `#1' is not a number}%
        {\string\val\space passes its body to siunitx, which refuses anything
         that is not a number. Use \string\valtext{#1} instead, or fix the
         script in code/ that emits this key.}%
    }{}%
    \mfanum{\csuse{mfaval@#1}}%
  }{\textcolor{mred}{\textbf{[?\,#1]}}}%
}
% A value that is deliberately a word rather than a number -- a format name, a
% verdict, a unit. Rare, and a visibly different macro so nobody reaches for
% \val and gets an error they cannot place.
\newcommand{\valtext}[1]{%
  \ifcsdef{mfaval@#1}{\csuse{mfaval@#1}}%
    {\textcolor{mred}{\textbf{[?\,#1]}}}%
}
\newcommand{\mfavaltext}[2]{%
  \csgdef{mfaval@#1}{#2}\csgdef{mfavaltext@#1}{}\listxadd{\mfavalues}{#1}%
}

% The raw stored string, for the cases where a number must appear exactly as
% the machine printed it -- inside a console block, or where the digits
% themselves are the point and grouping them would obscure the pattern.
\newcommand{\rawval}[1]{%
  \ifcsdef{mfaval@#1}{\csuse{mfaval@#1}}%
    {\textcolor{mred}{\textbf{[?\,#1]}}}%
}
\makeatother

% figures/values/all.tex is generated by `make numbers` and does nothing but
% \input each program's value file. A build with no values yet still compiles;
% every \val{} in it prints a visible marker instead of a number, which is the
% correct behaviour for a draft and an obvious failure for a release.
\IfFileExists{figures/values/all.tex}{\input{figures/values/all}}{%
  \AtBeginDocument{\PackageWarning{mfabook}{No computed values found. Run `make numbers`.}}}

% ============================================================
%  DIAGRAMS — Mermaid pipeline
%  ---------------------------------------------------------
%  \mermaidfig{key}{caption}{one-line manifest description}
%
%  Source of truth is figures/mermaid/<lang>/<key>.mmd, committed. `make
%  diagrams` renders each to figures/diagrams/<lang>/<key>.pdf, which is build
%  output and is not committed, so a diagram change reviews as a text diff.
%
%  The source is per-language because a diagram with English labels in the
%  Polish edition is exactly the kind of half-translation that makes a
%  bilingual book feel machine-made. CI checks that both languages carry the
%  same set of keys.
%
%  If the rendered PDF is absent the macro draws a placeholder AND typesets the
%  Mermaid source, so a build without mermaid-cli still shows the structure.
% ============================================================
\makeatletter
\newcommand{\listofdiagrams}{%
  \section*{\lblDiagramManifest}%
  \phantomsection\addcontentsline{toc}{section}{\lblDiagramManifest}%
  \@starttoc{dgm}%
}
\makeatother

\newcommand{\mermaidfig}[3]{%
  \addcontentsline{dgm}{subsection}{\protect\numberline{}\texttt{#1.mmd} --- #3}%
  \IfFileExists{figures/diagrams/\booklang/#1.pdf}{%
    \begin{figure}[htbp]
    \centering
    \includegraphics[width=0.95\linewidth,height=0.42\textheight,keepaspectratio]{figures/diagrams/\booklang/#1.pdf}%
    \caption{#2}\label{fig:#1}
    \end{figure}%
  }{%
    % The unrendered case deliberately does NOT float. A Mermaid source
    % typeset verbatim is often taller than a page, and a float cannot break,
    % so wrapping it in `figure` produces an overfull box per diagram and a
    % pile of tcolorbox warnings. In the flow it breaks like any other text.
    \par\medskip
    \begin{tcolorbox}[
      enhanced, breakable, width=\linewidth, sharp corners,
      colback=white, colframe=mgrey, boxrule=0.8pt,
      left=8pt, right=8pt, top=8pt, bottom=8pt]
      {\bfseries\color{mgrey}\small \lblNotRendered}\\[4pt]
      {\ttfamily\scriptsize\color{mgrey} figures/mermaid/\booklang/#1.mmd}\\[6pt]
      \IfFileExists{figures/mermaid/\booklang/#1.mmd}{%
        % ASCII only, like every other listing in this book: the fallback runs
        % the Mermaid source through `listings`, which aborts on a multi-byte
        % character it has no literate mapping for.
        \lstinputlisting[style=mfacode,numbers=none,basicstyle=\ttfamily\scriptsize,
                         backgroundcolor=\color{mlight},frame=none,
                         keywordstyle=\color{black}]{figures/mermaid/\booklang/#1.mmd}%
      }{%
        {\small\color{mred} \lblSourceMissing: \texttt{figures/mermaid/\booklang/#1.mmd}}%
      }%
    \end{tcolorbox}%
    \captionof{figure}{#2}\label{fig:#1}
    \par\medskip
  }%
}

% ============================================================
%  PROGRAM STUBS
%  ---------------------------------------------------------
%  Prints a visible marker so an unwritten program cannot be mistaken for a
%  written one and the page count never flatters the draft. `make debt` counts
%  them, and CI publishes the count on every build.
% ============================================================
\newcommand{\programstub}[1]{%
  \begin{tcolorbox}[
    enhanced, breakable, width=\linewidth, sharp corners,
    colback=mlight, colframe=mred, boxrule=1pt,
    borderline={0.6pt}{3pt}{mred, dashed},
    left=12pt, right=12pt, top=12pt, bottom=12pt]
    {\bfseries\color{mred}\large \lblNotWritten}\\[8pt]
    \small\raggedright #1
  \end{tcolorbox}%
}

% ============================================================
%  MATHEMATICAL OPERATORS
%  ---------------------------------------------------------
%  Language-invariant names live here. The ones that genuinely differ between
%  Polish and English usage -- tan/tg, Var/D^2, gcd/NWD, and the interval
%  brackets -- are defined in lang/en.tex and lang/pl.tex instead, so the
%  notation contract is executed by the build rather than remembered by a
%  translator.
% ============================================================
\DeclareMathOperator*{\argmin}{arg\,min}
\DeclareMathOperator*{\argmax}{arg\,max}
\DeclareMathOperator{\tr}{tr}
\DeclareMathOperator{\rank}{rank}
\DeclareMathOperator{\diag}{diag}
\DeclareMathOperator{\sgn}{sgn}
\DeclareMathOperator{\relu}{ReLU}
\DeclareMathOperator{\softmax}{softmax}
\DeclareMathOperator{\logsumexp}{logsumexp}
\DeclareMathOperator{\fl}{fl}
\DeclareMathOperator{\ulp}{ulp}
\DeclareMathOperator{\round}{round}

% A bare \log is FORBIDDEN in this book, and tools/check_notation.py fails the
% build on one. In Polish textbooks \log means base ten; in machine-learning
% writing it means base e. The same three characters carry opposite meanings to
% the two readerships this book serves, and they collide hardest in exactly the
% programs where it matters most -- entropy in bits and cross-entropy in nats
% are two programs apart. Write \ln for nats and \logb{2} for bits. It costs
% two characters and removes the ambiguity entirely.
\makeatletter
\let\mfa@origlog\log
\renewcommand{\log}{%
  \PackageError{mfabook}{A bare \string\log\space is forbidden in this book}%
    {Write \string\ln\space for natural logarithms or \string\logb{2}\space for
     bits. The base is never left to a convention here; see Appendix B.}%
  \mfa@origlog}
% The one legitimate form: an explicit base.
\newcommand{\logb}[1]{\mathop{\mfa@origlog}\nolimits_{#1}}
% ...and a way to write the forbidden thing when the subject IS the ban, which
% is Appendix B and nowhere else.
\newcommand{\mfalogplain}{\mathop{\mfa@origlog}\nolimits}
\makeatother

\newcommand{\R}{\mathbb{R}}
\newcommand{\N}{\mathbb{N}}
\newcommand{\Z}{\mathbb{Z}}
\newcommand{\Q}{\mathbb{Q}}
\newcommand{\C}{\mathbb{C}}
\newcommand{\Fset}{\mathbb{F}}          % the floating-point numbers, F1's subject
\newcommand{\vect}[1]{\mathbf{#1}}
\newcommand{\mat}[1]{\mathbf{#1}}
\newcommand{\T}{^{\mathsf{T}}}
\newcommand{\norm}[1]{\left\lVert #1 \right\rVert}
\newcommand{\abs}[1]{\left\lvert #1 \right\rvert}
\newcommand{\inner}[2]{\left\langle #1, #2 \right\rangle}
\newcommand{\eps}{\varepsilon}

% ============================================================
%  INDEX
%  ---------------------------------------------------------
%  Two-column, and its entries include \texttt{} names that do not hyphenate.
%  Ragged right keeps multi-word entries off the margin; it cannot help a
%  single token wider than the column, so keep verbatim index entries under
%  about 29 characters and index longer names by concept instead. Both sibling
%  books learned this the same way.
% ============================================================
\apptocmd{\theindex}{\raggedright}{}{%
  \PackageWarning{mfabook}{Could not patch theindex for ragged-right setting}}

\makeindex

% ============================================================
%  PINNED FACTS — single source of truth
%  Change these here; they propagate through both editions.
% ============================================================
\newcommand{\bookedition}{0.1}
% \pinnedon is a user-visible string and lives in lang/<lang>.tex, not here.
% It said "August 2026" on the Polish title page for exactly as long as it
% took somebody to read the Polish title page.
\newcommand{\pyver}{Python 3.13}
\newcommand{\numpyver}{2.3}
; nothing else differs between them at
%  this level. Every user-visible string lives in lang/en.tex or lang/pl.tex,
%  so a label that appears in one edition and not the other is a missing macro
%  rather than a silent divergence.
% ============================================================

% \booklang must be set by the main file. Default to English rather than
% failing, so a stray % ============================================================
%  preamble.tex — styling, environments, macros
%
%  Shared by BOTH editions. main-en.tex and main-pl.tex each define
%  \booklang before \input{preamble}; nothing else differs between them at
%  this level. Every user-visible string lives in lang/en.tex or lang/pl.tex,
%  so a label that appears in one edition and not the other is a missing macro
%  rather than a silent divergence.
% ============================================================

% \booklang must be set by the main file. Default to English rather than
% failing, so a stray \input{preamble} still compiles and says what it did.
\providecommand{\booklang}{en}

\usepackage{etoolbox}   % loaded first: the language switch below needs it

% ---------- Page geometry ----------
% Same 17 x 24 cm block as the other two books in this series, so a reader who
% owns one recognises the second. Frames are short, so the measure matters
% more here than in a prose book: a frame that runs past a page turn defeats
% the cover-the-next-frame discipline the whole method rests on.
\usepackage[
  paperwidth=17cm, paperheight=24cm,
  inner=2.2cm, outer=1.8cm, top=2.2cm, bottom=2.4cm,
  head=23pt, headsep=0.7cm, footskip=1.2cm
]{geometry}

% ---------- Encoding & fonts ----------
\usepackage[T1]{fontenc}
\usepackage[utf8]{inputenc}
\IfFileExists{lmodern.sty}{\usepackage{lmodern}}{}
\IfFileExists{newtxtext.sty}{\usepackage{newtxtext}}{}
% amsmath goes here, not down with the other packages, because newtxmath
% patches amsmath's internals and must be loaded AFTER it. The wrong order does
% not fail on a machine without newtxmath -- which is most CI images -- and
% fails on the author's full TeX Live.
%
% amssymb only when newtxmath is absent: newtxmath supplies the AMS symbols
% itself (its amssymbols option is on by default) and loading both clashes.
\usepackage{amsmath}
\IfFileExists{newtxmath.sty}{\usepackage{newtxmath}}{\usepackage{amssymb}}
\IfFileExists{inconsolata.sty}{\usepackage[varqu,scaled=0.95]{inconsolata}}{}
\IfFileExists{lmodern.sty}{\usepackage{microtype}}{\usepackage[expansion=false]{microtype}}

% ---------- Language ----------
% Loading babel with a language whose .ldf is absent is a *fatal* error, not a
% warning, so both the package and each language file are probed before either
% is requested. This is inherited from the sibling books, where a minimal TeX
% installation without texlive-lang-polish aborted the run with no PDF and an
% error message that named neither babel nor the language.
%
% The Polish edition needs Polish as the main language for hyphenation; the
% English edition needs British. Both editions load the other language too,
% because the glossary appendix is bilingual in both.
\IfFileExists{babel.sty}{%
  \IfFileExists{polish.ldf}{%
    \IfFileExists{british.ldf}{%
      \ifdefstring{\booklang}{pl}%
        {\usepackage[british,main=polish]{babel}}%
        {\usepackage[polish,main=british]{babel}}%
    }{\usepackage[polish]{babel}}%
  }{%
    \IfFileExists{british.ldf}{\usepackage[british]{babel}}{}%
  }%
}{}

% babel-polish makes " an active shorthand. listings reads its body verbatim
% and an active character in a verbatim context is a category-code accident
% waiting to happen, so the shorthand is turned off globally. Polish quotation
% marks are produced with \enquote-style macros from lang/pl.tex instead.
\ifdefstring{\booklang}{pl}{\AtBeginDocument{\shorthandoff{"}}}{}

% ---------- Core packages ----------
\usepackage{graphicx}
\usepackage{xcolor}
\usepackage{booktabs}
\usepackage{tabularx}
\usepackage{longtable}
\usepackage{array}
\usepackage{enumitem}
\IfFileExists{mathtools.sty}{\usepackage{mathtools}}{}
\usepackage{listings}
\usepackage[most]{tcolorbox}
\usepackage{fancyhdr}
\usepackage{titlesec}
\usepackage{caption}
\usepackage{imakeidx}
% csquotes with autostyle follows babel's active language, so the identical
% source \enquote{x} sets 'x' in the English edition and the Polish quotation
% marks in the Polish one. This is not a nicety: babel-polish makes " an active
% shorthand, which is turned off above precisely because an active character
% near a listings body is a category-code accident waiting to happen -- so a
% quotation mark has to come from a macro rather than from a keyboard key.
\IfFileExists{csquotes.sty}{\usepackage[autostyle=true]{csquotes}}{%
  \providecommand{\enquote}[1]{``##1''}}
\usepackage[hidelinks]{hyperref}

% ---------- Palette ----------
% Deliberately the same family as the sibling books, shifted green so the three
% spines are distinguishable on a shelf.
\definecolor{mblue}{HTML}{1F4E79}
\definecolor{mgreen}{HTML}{2E6B4F}
\definecolor{mteal}{HTML}{0E7C7B}
\definecolor{mamber}{HTML}{B26A00}
\definecolor{mred}{HTML}{9B2C2C}
\definecolor{mgrey}{HTML}{5A5A5A}
\definecolor{mlight}{HTML}{F4F6F8}
\definecolor{mfaint}{HTML}{FBFCFD}
\definecolor{codebg}{HTML}{FAFAFA}
\definecolor{codeframe}{HTML}{D8DEE4}
\definecolor{kw}{HTML}{0B5394}
\definecolor{str}{HTML}{9B2C2C}
\definecolor{cmt}{HTML}{4F7A4F}
\definecolor{typ}{HTML}{0E7C7B}

\hypersetup{
  colorlinks=true,
  linkcolor=mblue,
  urlcolor=mteal,
  citecolor=mblue,
  pdfauthor={Konrad Cinkusz}
}

% \ifpl{Polish text}{English text} -- for the handful of places where a string
% is structural rather than content and belongs in the shared files:
% structure.tex's part titles, and the main files' front matter wiring.
% Anything longer than a few words belongs in lang/ or in the edition's own
% source, not here.
% \DeclareRobustCommand, not \newcommand. A fragile \ifpl inside a \part title
% is expanded one level as it is written to the .toc, which lands
% "\ifdefstring {en}{pl}{...}{...}" in the file -- and that fails on the next
% run with "Extra }, or forgotten \endgroup", an error that names neither the
% part nor this macro.
\DeclareRobustCommand{\ifpl}[2]{\ifdefstring{\booklang}{pl}{#1}{#2}}

% A robust macro cannot go into a PDF bookmark string either -- hyperref strips
% \ifdefstring and warns once per part title. Tell it which branch to take
% instead, so the bookmarks carry the right language rather than nothing.
\makeatletter
\pdfstringdefDisableCommands{%
  \ifdefstring{\booklang}{pl}{\let\ifpl\@firstoftwo}{\let\ifpl\@secondoftwo}%
}
\makeatother

% ---------- Language strings ----------
% Every user-visible word the macros below typeset comes from here. Nothing in
% programs/ or appendices/ may hard-code an English or a Polish word inside a
% macro argument that the other edition also uses.
\input{lang/\booklang}

% ============================================================
%  PROGRAM NUMBERING
%  ---------------------------------------------------------
%  Part F programs are F1..F13; the main programs restart at 1. That is
%  Stroud's scheme and it is worth reproducing, because the Foundation part is
%  explicitly optional and numbering it separately is what makes skipping it
%  feel permitted rather than delinquent.
%
%  \theHchapter is hyperref's *destination* name, and it must stay unique
%  across the whole document. Resetting the chapter counter without also
%  changing \theHchapter produces duplicate PDF destinations, a pile of
%  warnings, and bookmarks that jump to the wrong program.
% ============================================================
\newcommand{\foundationnumbering}{%
  \setcounter{chapter}{0}%
  \renewcommand{\thechapter}{F\arabic{chapter}}%
  \renewcommand{\theHchapter}{F\arabic{chapter}}%
}
\newcommand{\mainnumbering}{%
  \setcounter{chapter}{0}%
  \renewcommand{\thechapter}{\arabic{chapter}}%
  \renewcommand{\theHchapter}{M\arabic{chapter}}%
}

% \appendix resets the chapter counter and switches \thechapter to letters, but
% it knows nothing about \theHchapter, so appendix A would claim the same PDF
% destination as main program 1. Patch it once, here.
\apptocmd{\appendix}{\renewcommand{\theHchapter}{A\Alph{chapter}}}{}{%
  \PackageWarning{mfabook}{Could not patch \string\appendix\space for hyperref destinations}}

% A program is a chapter. \chaptername is set from the language file so the
% word above the title is "Program" in both editions (it happens to be the
% same word, which is a small piece of luck).
\AtBeginDocument{\renewcommand{\chaptername}{\lblProgram}}

% ---------- Headings ----------
% \raggedright on the title body is load-bearing, not cosmetic: at \Huge on a
% 13 cm text block a long title that cannot break overflows the margin badly.
\titleformat{\chapter}[display]
  {\normalfont\huge\bfseries\color{mblue}\raggedright}
  {\normalfont\normalsize\color{mgrey}\MakeUppercase{\chaptertitlename}\ \thechapter}
  {8pt}{\Huge\raggedright}
\titlespacing*{\chapter}{0pt}{10pt}{28pt}
\titleformat{\section}{\normalfont\Large\bfseries\color{mblue}}{\thesection}{0.8em}{}
\titleformat{\subsection}{\normalfont\large\bfseries\color{mgrey}}{\thesubsection}{0.7em}{}

% head=23pt in the geometry options above, not \setlength{\headheight}:
% fancyhdr computes the height it needs from the running-head font and the
% rule, wants 22.5 pt at this size, and silently overprints the top of the text
% block if it does not get it. Setting it through geometry keeps the text block
% where geometry put it instead of pushing it down.
\pagestyle{fancy}
\fancyhf{}
\fancyhead[LE]{\small\nouppercase{\leftmark}}
% The recto head carries the FRAME number, not the section. In a book of
% numbered frames "where am I" is a frame question: every Summary
% back-reference, every Quiz route and every cross-reference between programs
% is a frame number, and hunting for one by page is the single most annoying
% thing about reading a programmed text. \theframeno at shipout is the last
% frame begun on the page, which is what a reader flicking through wants.
% Guarded: the front matter and the appendices have no frames, and a head
% reading "Frame 0" over the introduction is worse than no head at all.
\fancyhead[RO]{\ifnum\value{frameno}>0\relax\small\lblFrame~\theframeno\fi}
\fancyfoot[LE,RO]{\small\thepage}
\renewcommand{\headrulewidth}{0.4pt}

% This MUST come after \pagestyle{fancy}: fancyhdr installs its own \chaptermark
% at that point and silently overwrites anything defined earlier. The symptom is
% a running-head change that appears to do nothing at all.
%
% The running head carries the *frame* number as well as the program, because
% in a book of numbered frames "where am I" is a frame question, not a page
% question, and the summary's back-references are frame numbers.
% \@chapapp, not \lblProgram. \appendix switches \@chapapp from \chaptername
% to \appendixname, and babel supplies both in the right language; hard-coding
% the word gives an appendix a running head reading "Program A. Answers" over a
% page whose own chapter head reads "APPENDIX A".
\makeatletter
\renewcommand{\chaptermark}[1]{\markboth{\@chapapp\ \thechapter.\ #1}{}}
\makeatother
\renewcommand{\sectionmark}[1]{\markright{\thesection\ #1}}

% ============================================================
%  FRAMES — the unit the whole method is built from
%  ---------------------------------------------------------
%  A program is a stream of numbered frames. Nearly every frame ends by asking
%  the reader for something, and the answer opens the *next* frame, so that the
%  reader can cover what is below and commit to an answer before seeing it.
%
%  The environment is called `fr` rather than `frame` for two reasons. First,
%  \frame is already defined by LaTeX2e (it draws a box in a picture
%  environment) and an environment named `frame` would silently redefine it.
%  Second, this is typed forty to seventy times per program, so it wants to be
%  short in the source.
%
%    \begin{fr}
%    \ans{$x = 3$}            % the answer to the previous frame, optional
%    Teaching text, ending in a question.
%    \end{fr}
% ============================================================

\newcounter{frameno}[chapter]
\renewcommand{\theframeno}{\arabic{frameno}}

% Frames are referenced constantly -- by the Summary, by the Quiz, by the
% "Can you?" checklist and by other programs. \label inside a frame must
% therefore capture the frame number, not the section number.
\makeatletter
\newcommand{\framelabelfix}{%
  \@namedef{theHframeno}{\theHchapter.\arabic{frameno}}%
}
\makeatother

\newcommand{\framerule}{%
  \par\penalty-200\vspace{9pt}%
  \noindent\textcolor{codeframe}{\rule{\linewidth}{0.5pt}}%
  \par\vspace{5pt}%
}

% The frame number as a filled badge, set in the text flow rather than in the
% margin. A margin number would have to alternate sides under `twoside` and
% would be the first thing to overflow on a 17 cm page; a badge cannot.
\newcommand{\framebadge}{%
  \begingroup
  \setlength{\fboxsep}{2.5pt}%
  \colorbox{mblue}{\textcolor{white}{\bfseries\scriptsize\theframeno}}%
  \endgroup\hspace{0.7em}%
}

\newenvironment{fr}{%
  \framerule
  \refstepcounter{frameno}%
  \nopagebreak
  \noindent\framebadge\ignorespaces
}{\par}

% The answer to the previous frame. Set apart and centred so the eye can find
% the edge to cover, and so a reader who has answered can check in one glance
% without reading on.
\newcommand{\ans}[1]{%
  \par\nobreak\vspace{-4pt}%
  \begin{tcolorbox}[
    enhanced, sharp corners, boxrule=0pt, leftrule=2.5pt,
    colback=mfaint, colframe=mgreen,
    left=8pt, right=8pt, top=4pt, bottom=4pt,
    width=0.86\linewidth, center, halign=center,
    before skip=2pt, after skip=7pt]
    #1
  \end{tcolorbox}%
  \nobreak\noindent\ignorespaces
}

% A multi-paragraph or worked answer, where a centred box would be wrong.
\newenvironment{ansblock}{%
  \par\nopagebreak
  \begin{tcolorbox}[
    enhanced, breakable, sharp corners, boxrule=0pt, leftrule=2.5pt,
    colback=mfaint, colframe=mgreen,
    left=8pt, right=8pt, top=5pt, bottom=5pt]
}{%
  \end{tcolorbox}%
  \nopagebreak\par\vspace{2pt}\noindent\ignorespaces
}

% The row of dots. Stroud's signal that the reader is to supply a word, a
% number or an expression. \blank is inline; \dotline occupies its own line,
% which is what a longer derivation step wants.
% \penalty0 offers TeX a break opportunity immediately before the dots. Without
% it the run of dots is glued to the last word of the question, and a long
% question in the Polish edition -- where the words are longer and hyphenate
% less freely -- runs into the margin.
\newcommand{\blank}[1][4em]{\penalty0\makebox[#1]{\dotfill}}
\newcommand{\dotline}{\par\vspace{2pt}\noindent\makebox[\linewidth]{\dotfill}\par\vspace{2pt}}

% "Now one for you to do." The third rung of the scaffolding gradient, and the
% one most often skipped by authors who think they have written an exercise
% when they have written another worked example.
\newcommand{\yourturn}{%
  \par\vspace{4pt}\noindent
  {\bfseries\color{mgreen}\lblYourTurn}\par\nobreak\vspace{2pt}\noindent\ignorespaces
}

% ============================================================
%  THE PROGRAM SKELETON
%  ---------------------------------------------------------
%  Every program carries the same seven parts, in the same order. They are
%  macros rather than a convention so that a missing one is a build-time
%  absence rather than something a reader notices.
% ============================================================

% ---- Learning outcomes, on entry ----------------------------------------
% Stored as well as printed, because "Can you?" at the end of the program must
% be the same list, one for one. Storing it is what makes that structural
% rather than a promise: the checklist is generated from the outcomes, so the
% two cannot drift.
\newcounter{outcomeno}[chapter]

\makeatletter
% The key under which everything belonging to the current program is stored.
% It must expand to the printed program number (F1, F2, ..., 1, 2, ...), which
% is what makes an answer findable from the program that owns it.
\newcommand{\mfa@progkey}{\thechapter}
\newcommand{\outcome}[2]{% #1 = frames that teach it, #2 = the outcome
  \stepcounter{outcomeno}%
  \csgdef{mfa@out@\mfa@progkey @\theoutcomeno}{#2}%
  \csgdef{mfa@outfr@\mfa@progkey @\theoutcomeno}{#1}%
  % The "Can you?" row is built here rather than looped over later. A \loop
  % inside a tabularx body does not survive alignment scanning: TeX needs to
  % see the & and the row break while it reads the cell, and a conditional
  % hides them. Accumulating the rows as tokens at declaration time sidesteps
  % that entirely, and has the better property anyway -- the checklist is
  % literally made of the outcomes, so the two cannot disagree.
  \csgappto{mfa@rows@\mfa@progkey}{%
    #2 & {\footnotesize\color{mgrey}#1}
      & \ratingcell & \ratingcell & \ratingcell & \ratingcell & \ratingcell
    \tabularnewline}%
  \item #2\hfill{\footnotesize\color{mgrey}[\lblFrames~#1]}%
}
\makeatother

\makeatletter
\newenvironment{outcomes}{%
  \setcounter{outcomeno}{0}%
  \csgdef{mfa@rows@\mfa@progkey}{}%
  \begin{tcolorbox}[
    enhanced, breakable, sharp corners,
    colback=mlight, colframe=mblue, boxrule=0pt, leftrule=3pt,
    left=8pt, right=8pt, top=6pt, bottom=6pt,
    fonttitle=\bfseries\small, title={\lblOutcomes}]
  \begin{itemize}[leftmargin=1.3em, itemsep=3pt, topsep=2pt]
}{%
  \end{itemize}
  \end{tcolorbox}
}
\makeatother

% ---- "Can you?" -----------------------------------------------------------
% Generated from the stored outcomes, not written out again by hand. Five
% columns rather than a tick box, because a self-assessment with a midpoint is
% one people actually use and a yes/no one is one they lie to.
\newcommand{\ratingcell}{\textcolor{codeframe}{\rule[-0.15em]{0.8em}{0.8em}}}

\makeatletter
\newcommand{\canyou}{%
  \section*{\lblCanYou}%
  \phantomsection\addcontentsline{toc}{section}{\lblCanYou}%
  \noindent\lblCanYouIntro\par\vspace{8pt}
  \noindent
  \begingroup\small
  \ifcsvoid{mfa@rows@\mfa@progkey}{%
    {\color{mred}\lblNoOutcomes}\par
  }{%
    \begin{tabularx}{\linewidth}{@{}Xcccccc@{}}
    \toprule
    \textbf{\lblOnEntryYouCould} & \textbf{\lblFrames} & 1 & 2 & 3 & 4 & 5 \\
    \midrule
    \csuse{mfa@rows@\mfa@progkey}
    \bottomrule
    \end{tabularx}%
  }%
  \endgroup
  \par\vspace{6pt}\noindent{\footnotesize\color{mgrey}\lblCanYouFooter}\par
}
\makeatother

% ============================================================
%  QUIZ, SUMMARY, EXERCISES
% ============================================================

% ---- Quiz -----------------------------------------------------------------
% Foundation programs only. It is a diagnostic on the way in and the same test
% on the way out, which is what lets one book serve a reader who needs the
% whole program and a reader who needs four frames of it. Each question names
% the frames that teach it, so a wrong answer routes the reader somewhere
% specific rather than back to the beginning.
\newlist{quizlist}{enumerate}{1}
\setlist[quizlist]{label=\arabic*., leftmargin=2.0em, itemsep=5pt, topsep=3pt}

\makeatletter
\newenvironment{quiz}{%
  \section*{\lblQuiz}%
  \phantomsection\addcontentsline{toc}{section}{\lblQuiz}%
  \def\mfa@exkey{Q\arabic{quizlisti}}%
  \noindent\lblQuizIntro\par\vspace{5pt}
  \begin{quizlist}
}{%
  \end{quizlist}
}
\makeatother

% The frames that teach the question just asked. Two forms, because a Quiz
% question that routes to a single frame reading "Frames 22" is the kind of
% small wrongness a reader notices and an author never does.
\newcommand{\teachesat}[1]{\hfill{\footnotesize\color{mgrey}(\lblFrames~#1)}}
\newcommand{\teachesatone}[1]{\hfill{\footnotesize\color{mgrey}(\lblFrame~#1)}}

% ---- Summary --------------------------------------------------------------
% Every item carries the frame number it came from, in square brackets. That
% is the eighth-edition improvement and it is the single cheapest thing in the
% format: it turns a summary into a return index, so a reader who fails one
% line of "Can you?" knows exactly which four frames to reread.
\newenvironment{summarybox}{%
  \section*{\lblSummary}%
  \phantomsection\addcontentsline{toc}{section}{\lblSummary}%
  \begin{tcolorbox}[
    enhanced, breakable, sharp corners,
    colback=mlight, colframe=mblue, boxrule=0pt, leftrule=3pt,
    left=8pt, right=8pt, top=6pt, bottom=6pt]
  \begin{itemize}[leftmargin=1.3em, itemsep=5pt, topsep=2pt]
}{%
  \end{itemize}
  \end{tcolorbox}
}

% \sumitem{12}{text} -- the frame reference is not decoration, it is the point.
\newcommand{\sumitem}[2]{\item #2~{\footnotesize\color{mgrey}[#1]}}

% ---- Answers ---------------------------------------------------------------
%  Answers are authored at the point of use and printed at the back.
%
%  They are carried there by a global macro store rather than by an auxiliary
%  file. The sibling books collect their figure and screenshot manifests with
%  \@starttoc and a custom file extension, which works because those entries
%  are one line of text; it is the wrong mechanism for an answer, because an
%  answer contains displayed maths, and material written to an .aux-style file
%  and read back is expanded at shipout, where fragile commands break in ways
%  whose error messages name neither the answer nor the program.
%
%  \include is sequential within a run, so a macro defined globally while
%  typesetting program F1 is still defined when the back matter is typeset.
%  That is all this needs.
\makeatletter
\newcommand{\mfa@storeanswer}[3]{% #1 program key, #2 item key, #3 body
  \csgdef{mfa@a@#1@#2}{#3}%
  % \listcsxadd, not \listcsgadd: the item must be EXPANDED as it is stored.
  % Storing the token \mfa@exkey means every entry in the list expands to
  % whatever the key happens to be at the time the back matter is typeset,
  % which is the last one.
  \listcsxadd{mfa@akeys@#1}{#2}%
}
% Used inside an exercise item. Prints nothing here; the body reappears at the
% back of the book under this program's heading.
\newcommand{\answerto}[1]{\mfa@storeanswer{\mfa@progkey}{\mfa@exkey}{#1}}
\providecommand{\mfa@exkey}{?}
\makeatother

\newlist{exlist}{enumerate}{1}
\setlist[exlist]{label=\arabic*., leftmargin=2.0em, itemsep=6pt, topsep=3pt}
\newlist{fplist}{enumerate}{1}
\setlist[fplist]{label=\arabic*., leftmargin=2.0em, itemsep=5pt, topsep=3pt}

% These MUST be defined inside \makeatletter. Outside it, \def\mfa@exkey{...}
% parses as \def\mfa @exkey{...} -- a definition of \mfa with a delimited
% parameter text -- which raises no error at all and quietly leaves every
% answer filed under the same key. That failure is invisible until the answers
% section at the back of the book prints one program's answers twice.
\makeatletter
\newenvironment{testexercises}{%
  \section*{\lblTestExercises}%
  \phantomsection\addcontentsline{toc}{section}{\lblTestExercises}%
  \def\mfa@exkey{T\arabic{exlisti}}%
  \noindent\lblTestIntro\par\vspace{5pt}
  \begin{exlist}
}{%
  \end{exlist}
}

\newenvironment{furtherproblems}{%
  \section*{\lblFurther}%
  \phantomsection\addcontentsline{toc}{section}{\lblFurther}%
  \def\mfa@exkey{P\arabic{fplisti}}%
  \noindent\lblFurtherIntro\par\vspace{5pt}
  \begin{fplist}
}{%
  \end{fplist}
}
\makeatother

% ---- The programs register -------------------------------------------------
% \program{Title} opens a program: it is \chapter plus the bookkeeping that
% lets the back matter reproduce this program's answers under its own name.
\makeatletter
\newcommand{\program}[1]{%
  \chapter{#1}%
  \csgdef{mfa@title@\thechapter}{#1}%
  % Expanded, for the same reason as the answer keys above.
  \listxadd{\mfaprograms}{\thechapter}%
}
\newcommand{\mfa@printoneanswer}[2]{% #1 program key, #2 item key
  \item[\textbf{#2}] \csuse{mfa@a@#1@#2}%
}
\newcommand{\mfa@answerprogram}[1]{%
  \subsection*{\lblProgram~#1\quad\textmd{\itshape\csuse{mfa@title@#1}}}%
  \ifcsdef{mfa@akeys@#1}{%
    \begin{list}{}{%
      \setlength{\leftmargin}{3.2em}\setlength{\labelwidth}{2.7em}%
      \setlength{\labelsep}{0.5em}\setlength{\itemsep}{3pt}%
      \setlength{\topsep}{3pt}\setlength{\parsep}{0pt}%
      \renewcommand{\makelabel}[1]{\hfill##1}}
      \forlistcsloop{\mfa@printoneanswer{#1}}{mfa@akeys@#1}
    \end{list}%
  }{%
    \par\noindent{\color{mred}\lblNoAnswers}\par
  }%
}
% The answers body, without a heading of its own, so the appendix that carries
% it can be a numbered appendix like any other and appear in the running heads
% and the ToC in the ordinary way.
\newcommand{\answersbody}{%
  \noindent\lblAnswersIntro\par\vspace{8pt}
  \ifdef{\mfaprograms}{\forlistloop{\mfa@answerprogram}{\mfaprograms}}{%
    \noindent{\color{mred}\lblNoAnswers}\par}%
}
% Unnumbered form, for a draft build that has no appendix wiring yet.
\newcommand{\printanswers}{%
  \chapter*{\lblAnswers}%
  \phantomsection\addcontentsline{toc}{chapter}{\lblAnswers}%
  \markboth{\lblAnswers}{\lblAnswers}%
  \answersbody
}
\makeatother

% ============================================================
%  ADMONITIONS
%  ---------------------------------------------------------
%  A deliberately small, fixed vocabulary. Each one earns its place by doing
%  something no other box does; a seventh would be a smell.
% ============================================================

\newtcolorbox{note}[1][]{
  enhanced, breakable, sharp corners,
  colback=mlight, colframe=mblue, boxrule=0pt, leftrule=3pt,
  left=8pt, right=8pt, top=6pt, bottom=6pt,
  fonttitle=\bfseries\small, title={\lblNote}, #1
}

\newtcolorbox{warning}[1][]{
  enhanced, breakable, sharp corners,
  colback=mlight, colframe=mred, boxrule=0pt, leftrule=3pt,
  left=8pt, right=8pt, top=6pt, bottom=6pt,
  fonttitle=\bfseries\small, title={\lblWarning}, #1
}

% The misconception, stated as a summary after it has been walked into. The
% trap itself belongs in the frames -- elicit the wrong answer, then say
% flatly that it is wrong -- and this box is the thing the reader marks with a
% finger and comes back to.
\newtcolorbox{trapbox}[1][]{
  enhanced, breakable, sharp corners,
  colback=white, colframe=mred, boxrule=0.8pt,
  left=8pt, right=8pt, top=6pt, bottom=6pt,
  fonttitle=\bfseries\small, title={\lblTrap}, #1
}

% Where the mathematics just done shows up in a system the reader has actually
% deployed. This is the box that separates the book from a maths textbook, and
% it is also the one most easily filled with waffle: if it cannot name a
% specific line of a specific system, it should not be written.
\newtcolorbox{aibox}[1][]{
  enhanced, breakable, sharp corners,
  colback=mlight, colframe=mteal, boxrule=0pt, leftrule=3pt,
  left=8pt, right=8pt, top=6pt, bottom=6pt,
  fonttitle=\bfseries\small, title={\lblInAI}, #1
}

% What is being asserted rather than proved, and where to read the proof.
% Stroud's method buys fluency at the cost of rigour, and the honest thing is
% to say so at each point rather than once in an introduction nobody rereads.
\newtcolorbox{rigourbox}[1][]{
  enhanced, breakable, sharp corners,
  colback=white, colframe=mgrey, boxrule=0.6pt,
  left=8pt, right=8pt, top=6pt, bottom=6pt,
  fonttitle=\bfseries\small, title={\lblRigour}, #1
}

% Notation that genuinely differs -- between the Polish and English editions,
% and between disciplines. Matrix calculus alone has two incompatible layout
% conventions, and half the confusion about backpropagation is a transpose
% that came from the other one.
\newtcolorbox{notationbox}[1][]{
  enhanced, breakable, sharp corners,
  colback=mlight, colframe=mamber, boxrule=0pt, leftrule=3pt,
  left=8pt, right=8pt, top=6pt, bottom=6pt,
  fonttitle=\bfseries\small, title={\lblNotation}, #1
}

% The maths analogue of the sibling books' "run this before you trust it".
% A number in this book is either produced by a script under code/ and pulled
% in with \val, or it carries this box. Nothing ships to a reader with one
% still attached.
\newtcolorbox{verifybox}[1][]{
  enhanced, breakable, sharp corners,
  colback=white, colframe=mamber, boxrule=0.8pt,
  left=8pt, right=8pt, top=6pt, bottom=6pt,
  fonttitle=\bfseries\small, title={\lblVerify}, #1
}

\newtcolorbox{exercisebox}[1][]{
  enhanced, breakable, sharp corners,
  colback=white, colframe=mgreen, boxrule=0pt, leftrule=3pt,
  left=8pt, right=8pt, top=6pt, bottom=6pt,
  fonttitle=\bfseries\small, title={\lblExercise}, #1
}

% ============================================================
%  CODE LISTINGS
%  ---------------------------------------------------------
%  Code appears here for one reason only: to check a number. Every listing in
%  this book is runnable and every number it prints is the number in the text.
% ============================================================
\lstdefinelanguage{pythonx}{
  language=Python,
  morekeywords={as, with, assert, match, case},
  morekeywords=[2]{
    float64,float32,float16,bfloat16,int8,uint8,
    np,numpy,array,asarray,frombuffer,nextafter,finfo,iinfo,
    exp,log,log1p,expm1,logaddexp,sum,max,abs,sqrt,mean,var
  },
  sensitive=true
}

\lstdefinestyle{mfacode}{
  basicstyle=\ttfamily\footnotesize,
  keywordstyle=\color{kw}\bfseries,
  keywordstyle=[2]\color{typ},
  stringstyle=\color{str},
  commentstyle=\color{cmt}\itshape,
  numbers=left,
  numberstyle=\ttfamily\tiny\color{mgrey},
  numbersep=8pt,
  backgroundcolor=\color{codebg},
  frame=single,
  rulecolor=\color{codeframe},
  framesep=6pt,
  breaklines=true,
  breakatwhitespace=false,
  postbreak=\mbox{\textcolor{mgrey}{$\hookrightarrow$}\space},
  showstringspaces=false,
  tabsize=4,
  columns=fullflexible,
  keepspaces=true,
  captionpos=b,
  aboveskip=1.2em,
  belowskip=1.2em,
  % Maps the characters most likely to be pasted into a listing by accident.
  %
  % Do NOT add a mapping for U+00A0 (non-breaking space). It looks like the
  % obvious next entry and it makes `listings` abort the run with "Improper
  % alphabetic constant" -- fatal, no PDF, and the message names neither the
  % character nor this line. The ASCII-in-listings convention is the guard
  % against it instead. This trap has now been hit in two of the three books
  % in this series.
  literate={ł}{{\l}}1 {Ł}{{\L}}1 {ą}{{\k a}}1 {ę}{{\k e}}1
           {ś}{{\'s}}1 {ć}{{\'c}}1 {ż}{{\.z}}1 {ź}{{\'z}}1 {ó}{{\'o}}1 {ń}{{\'n}}1
           {Ą}{{\k A}}1 {Ę}{{\k E}}1 {Ś}{{\'S}}1 {Ć}{{\'C}}1
           {Ż}{{\.Z}}1 {Ź}{{\'Z}}1 {Ó}{{\'O}}1 {Ń}{{\'N}}1
           {—}{{-{}-}}2 {–}{{-}}1 {‘}{{'}}1 {’}{{'}}1
           {“}{{"}}1 {”}{{"}}1 {°}{{\textdegree}}1
}

\lstset{style=mfacode}

\lstnewenvironment{python}[1][]
  {\lstset{language=pythonx,style=mfacode,#1}}
  {}

\lstnewenvironment{shellcmd}[1][]
  {\lstset{language=bash,style=mfacode,numbers=none,keywordstyle=\color{mteal}\bfseries,#1}}
  {}

% Terminal transcripts: no line numbers, no keyword colouring. Everything in
% one of these was copied from a real run.
\lstnewenvironment{console}[1][]
  {\lstset{style=mfacode,numbers=none,basicstyle=\ttfamily\scriptsize,
           keywordstyle=\color{black},backgroundcolor=\color{white},#1}}
  {}

\newcommand{\pyfile}[3]{%
  \lstinputlisting[language=pythonx,style=mfacode,caption={#2},label={#3}]{#1}%
}
\newcommand{\pyregion}[4]{%
  \lstinputlisting[
    language=pythonx, style=mfacode,
    linerange={--8<--\ [start:#2]--- 8<--\ [end:#2]},
    includerangemarker=false,
    caption={#3}, label={#4}]{#1}%
}

% Inline literals
\newcommand{\code}[1]{\texttt{\small #1}}
\newcommand{\pkg}[1]{\texttt{\small\color{mblue}#1}}
\newcommand{\api}[1]{\texttt{\small\color{mteal}#1}}

% ============================================================
%  COMPUTED VALUES
%  ---------------------------------------------------------
%  The house rule in the sibling books is "run the listing, or mark it". The
%  analogue here is stronger, because a wrong digit in a maths book is not a
%  broken example, it is a lie the reader will carry.
%
%  Every number that came out of a computation is written by a script under
%  code/ into figures/values/<program>.tex as
%
%      \mfaval{f01.eps64}{2.220446049250313e-16}
%
%  and pulled into the text with \val{f01.eps64}. The script is the source of
%  truth; the book cannot disagree with it, because the book does not contain
%  the digits. A value the scripts no longer produce prints a visible marker
%  and is counted by `make debt`.
%
%  \val also applies the locale's decimal separator, which is how the Polish
%  edition gets its decimal comma without a translator having to remember.
% ============================================================
\IfFileExists{siunitx.sty}{%
  \usepackage{siunitx}%
  \sisetup{
    group-digits = integer,
    group-minimum-digits = 5,
    group-separator = {\,},
    exponent-product = \cdot,
    output-decimal-marker = {\mfadecimalmarker}
  }%
  \newcommand{\mfanum}[1]{\num{##1}}%
}{%
  \newcommand{\mfanum}[1]{##1}%
}

\makeatletter
\newcommand{\mfaval}[2]{\csgdef{mfaval@#1}{#2}\listxadd{\mfavalues}{#1}}
% \num on a non-numeric body is a FATAL siunitx error -- "Invalid number" --
% and under -halt-on-error that is the whole build. A script emitting a word
% where the book expected a number is a plausible mistake.
%
% The check is NOT done here. Deciding in TeX whether a token list is a number
% means parsing it character by character, which is fragile and was got wrong
% once already. The generator knows the answer for free, so it emits \mfaval
% for a number and \mfavaltext for a word, and this end only has to enforce
% which macro reads which.
\newcommand{\val}[1]{%
  \ifcsdef{mfaval@#1}{%
    \ifcsdef{mfavaltext@#1}{%
      \PackageError{mfabook}{Value `#1' is not a number}%
        {\string\val\space passes its body to siunitx, which refuses anything
         that is not a number. Use \string\valtext{#1} instead, or fix the
         script in code/ that emits this key.}%
    }{}%
    \mfanum{\csuse{mfaval@#1}}%
  }{\textcolor{mred}{\textbf{[?\,#1]}}}%
}
% A value that is deliberately a word rather than a number -- a format name, a
% verdict, a unit. Rare, and a visibly different macro so nobody reaches for
% \val and gets an error they cannot place.
\newcommand{\valtext}[1]{%
  \ifcsdef{mfaval@#1}{\csuse{mfaval@#1}}%
    {\textcolor{mred}{\textbf{[?\,#1]}}}%
}
\newcommand{\mfavaltext}[2]{%
  \csgdef{mfaval@#1}{#2}\csgdef{mfavaltext@#1}{}\listxadd{\mfavalues}{#1}%
}

% The raw stored string, for the cases where a number must appear exactly as
% the machine printed it -- inside a console block, or where the digits
% themselves are the point and grouping them would obscure the pattern.
\newcommand{\rawval}[1]{%
  \ifcsdef{mfaval@#1}{\csuse{mfaval@#1}}%
    {\textcolor{mred}{\textbf{[?\,#1]}}}%
}
\makeatother

% figures/values/all.tex is generated by `make numbers` and does nothing but
% \input each program's value file. A build with no values yet still compiles;
% every \val{} in it prints a visible marker instead of a number, which is the
% correct behaviour for a draft and an obvious failure for a release.
\IfFileExists{figures/values/all.tex}{\input{figures/values/all}}{%
  \AtBeginDocument{\PackageWarning{mfabook}{No computed values found. Run `make numbers`.}}}

% ============================================================
%  DIAGRAMS — Mermaid pipeline
%  ---------------------------------------------------------
%  \mermaidfig{key}{caption}{one-line manifest description}
%
%  Source of truth is figures/mermaid/<lang>/<key>.mmd, committed. `make
%  diagrams` renders each to figures/diagrams/<lang>/<key>.pdf, which is build
%  output and is not committed, so a diagram change reviews as a text diff.
%
%  The source is per-language because a diagram with English labels in the
%  Polish edition is exactly the kind of half-translation that makes a
%  bilingual book feel machine-made. CI checks that both languages carry the
%  same set of keys.
%
%  If the rendered PDF is absent the macro draws a placeholder AND typesets the
%  Mermaid source, so a build without mermaid-cli still shows the structure.
% ============================================================
\makeatletter
\newcommand{\listofdiagrams}{%
  \section*{\lblDiagramManifest}%
  \phantomsection\addcontentsline{toc}{section}{\lblDiagramManifest}%
  \@starttoc{dgm}%
}
\makeatother

\newcommand{\mermaidfig}[3]{%
  \addcontentsline{dgm}{subsection}{\protect\numberline{}\texttt{#1.mmd} --- #3}%
  \IfFileExists{figures/diagrams/\booklang/#1.pdf}{%
    \begin{figure}[htbp]
    \centering
    \includegraphics[width=0.95\linewidth,height=0.42\textheight,keepaspectratio]{figures/diagrams/\booklang/#1.pdf}%
    \caption{#2}\label{fig:#1}
    \end{figure}%
  }{%
    % The unrendered case deliberately does NOT float. A Mermaid source
    % typeset verbatim is often taller than a page, and a float cannot break,
    % so wrapping it in `figure` produces an overfull box per diagram and a
    % pile of tcolorbox warnings. In the flow it breaks like any other text.
    \par\medskip
    \begin{tcolorbox}[
      enhanced, breakable, width=\linewidth, sharp corners,
      colback=white, colframe=mgrey, boxrule=0.8pt,
      left=8pt, right=8pt, top=8pt, bottom=8pt]
      {\bfseries\color{mgrey}\small \lblNotRendered}\\[4pt]
      {\ttfamily\scriptsize\color{mgrey} figures/mermaid/\booklang/#1.mmd}\\[6pt]
      \IfFileExists{figures/mermaid/\booklang/#1.mmd}{%
        % ASCII only, like every other listing in this book: the fallback runs
        % the Mermaid source through `listings`, which aborts on a multi-byte
        % character it has no literate mapping for.
        \lstinputlisting[style=mfacode,numbers=none,basicstyle=\ttfamily\scriptsize,
                         backgroundcolor=\color{mlight},frame=none,
                         keywordstyle=\color{black}]{figures/mermaid/\booklang/#1.mmd}%
      }{%
        {\small\color{mred} \lblSourceMissing: \texttt{figures/mermaid/\booklang/#1.mmd}}%
      }%
    \end{tcolorbox}%
    \captionof{figure}{#2}\label{fig:#1}
    \par\medskip
  }%
}

% ============================================================
%  PROGRAM STUBS
%  ---------------------------------------------------------
%  Prints a visible marker so an unwritten program cannot be mistaken for a
%  written one and the page count never flatters the draft. `make debt` counts
%  them, and CI publishes the count on every build.
% ============================================================
\newcommand{\programstub}[1]{%
  \begin{tcolorbox}[
    enhanced, breakable, width=\linewidth, sharp corners,
    colback=mlight, colframe=mred, boxrule=1pt,
    borderline={0.6pt}{3pt}{mred, dashed},
    left=12pt, right=12pt, top=12pt, bottom=12pt]
    {\bfseries\color{mred}\large \lblNotWritten}\\[8pt]
    \small\raggedright #1
  \end{tcolorbox}%
}

% ============================================================
%  MATHEMATICAL OPERATORS
%  ---------------------------------------------------------
%  Language-invariant names live here. The ones that genuinely differ between
%  Polish and English usage -- tan/tg, Var/D^2, gcd/NWD, and the interval
%  brackets -- are defined in lang/en.tex and lang/pl.tex instead, so the
%  notation contract is executed by the build rather than remembered by a
%  translator.
% ============================================================
\DeclareMathOperator*{\argmin}{arg\,min}
\DeclareMathOperator*{\argmax}{arg\,max}
\DeclareMathOperator{\tr}{tr}
\DeclareMathOperator{\rank}{rank}
\DeclareMathOperator{\diag}{diag}
\DeclareMathOperator{\sgn}{sgn}
\DeclareMathOperator{\relu}{ReLU}
\DeclareMathOperator{\softmax}{softmax}
\DeclareMathOperator{\logsumexp}{logsumexp}
\DeclareMathOperator{\fl}{fl}
\DeclareMathOperator{\ulp}{ulp}
\DeclareMathOperator{\round}{round}

% A bare \log is FORBIDDEN in this book, and tools/check_notation.py fails the
% build on one. In Polish textbooks \log means base ten; in machine-learning
% writing it means base e. The same three characters carry opposite meanings to
% the two readerships this book serves, and they collide hardest in exactly the
% programs where it matters most -- entropy in bits and cross-entropy in nats
% are two programs apart. Write \ln for nats and \logb{2} for bits. It costs
% two characters and removes the ambiguity entirely.
\makeatletter
\let\mfa@origlog\log
\renewcommand{\log}{%
  \PackageError{mfabook}{A bare \string\log\space is forbidden in this book}%
    {Write \string\ln\space for natural logarithms or \string\logb{2}\space for
     bits. The base is never left to a convention here; see Appendix B.}%
  \mfa@origlog}
% The one legitimate form: an explicit base.
\newcommand{\logb}[1]{\mathop{\mfa@origlog}\nolimits_{#1}}
% ...and a way to write the forbidden thing when the subject IS the ban, which
% is Appendix B and nowhere else.
\newcommand{\mfalogplain}{\mathop{\mfa@origlog}\nolimits}
\makeatother

\newcommand{\R}{\mathbb{R}}
\newcommand{\N}{\mathbb{N}}
\newcommand{\Z}{\mathbb{Z}}
\newcommand{\Q}{\mathbb{Q}}
\newcommand{\C}{\mathbb{C}}
\newcommand{\Fset}{\mathbb{F}}          % the floating-point numbers, F1's subject
\newcommand{\vect}[1]{\mathbf{#1}}
\newcommand{\mat}[1]{\mathbf{#1}}
\newcommand{\T}{^{\mathsf{T}}}
\newcommand{\norm}[1]{\left\lVert #1 \right\rVert}
\newcommand{\abs}[1]{\left\lvert #1 \right\rvert}
\newcommand{\inner}[2]{\left\langle #1, #2 \right\rangle}
\newcommand{\eps}{\varepsilon}

% ============================================================
%  INDEX
%  ---------------------------------------------------------
%  Two-column, and its entries include \texttt{} names that do not hyphenate.
%  Ragged right keeps multi-word entries off the margin; it cannot help a
%  single token wider than the column, so keep verbatim index entries under
%  about 29 characters and index longer names by concept instead. Both sibling
%  books learned this the same way.
% ============================================================
\apptocmd{\theindex}{\raggedright}{}{%
  \PackageWarning{mfabook}{Could not patch theindex for ragged-right setting}}

\makeindex

% ============================================================
%  PINNED FACTS — single source of truth
%  Change these here; they propagate through both editions.
% ============================================================
\newcommand{\bookedition}{0.1}
% \pinnedon is a user-visible string and lives in lang/<lang>.tex, not here.
% It said "August 2026" on the Polish title page for exactly as long as it
% took somebody to read the Polish title page.
\newcommand{\pyver}{Python 3.13}
\newcommand{\numpyver}{2.3}
 still compiles and says what it did.
\providecommand{\booklang}{en}

\usepackage{etoolbox}   % loaded first: the language switch below needs it

% ---------- Page geometry ----------
% Same 17 x 24 cm block as the other two books in this series, so a reader who
% owns one recognises the second. Frames are short, so the measure matters
% more here than in a prose book: a frame that runs past a page turn defeats
% the cover-the-next-frame discipline the whole method rests on.
\usepackage[
  paperwidth=17cm, paperheight=24cm,
  inner=2.2cm, outer=1.8cm, top=2.2cm, bottom=2.4cm,
  head=23pt, headsep=0.7cm, footskip=1.2cm
]{geometry}

% ---------- Encoding & fonts ----------
\usepackage[T1]{fontenc}
\usepackage[utf8]{inputenc}
\IfFileExists{lmodern.sty}{\usepackage{lmodern}}{}
\IfFileExists{newtxtext.sty}{\usepackage{newtxtext}}{}
% amsmath goes here, not down with the other packages, because newtxmath
% patches amsmath's internals and must be loaded AFTER it. The wrong order does
% not fail on a machine without newtxmath -- which is most CI images -- and
% fails on the author's full TeX Live.
%
% amssymb only when newtxmath is absent: newtxmath supplies the AMS symbols
% itself (its amssymbols option is on by default) and loading both clashes.
\usepackage{amsmath}
\IfFileExists{newtxmath.sty}{\usepackage{newtxmath}}{\usepackage{amssymb}}
\IfFileExists{inconsolata.sty}{\usepackage[varqu,scaled=0.95]{inconsolata}}{}
\IfFileExists{lmodern.sty}{\usepackage{microtype}}{\usepackage[expansion=false]{microtype}}

% ---------- Language ----------
% Loading babel with a language whose .ldf is absent is a *fatal* error, not a
% warning, so both the package and each language file are probed before either
% is requested. This is inherited from the sibling books, where a minimal TeX
% installation without texlive-lang-polish aborted the run with no PDF and an
% error message that named neither babel nor the language.
%
% The Polish edition needs Polish as the main language for hyphenation; the
% English edition needs British. Both editions load the other language too,
% because the glossary appendix is bilingual in both.
\IfFileExists{babel.sty}{%
  \IfFileExists{polish.ldf}{%
    \IfFileExists{british.ldf}{%
      \ifdefstring{\booklang}{pl}%
        {\usepackage[british,main=polish]{babel}}%
        {\usepackage[polish,main=british]{babel}}%
    }{\usepackage[polish]{babel}}%
  }{%
    \IfFileExists{british.ldf}{\usepackage[british]{babel}}{}%
  }%
}{}

% babel-polish makes " an active shorthand. listings reads its body verbatim
% and an active character in a verbatim context is a category-code accident
% waiting to happen, so the shorthand is turned off globally. Polish quotation
% marks are produced with \enquote-style macros from lang/pl.tex instead.
\ifdefstring{\booklang}{pl}{\AtBeginDocument{\shorthandoff{"}}}{}

% ---------- Core packages ----------
\usepackage{graphicx}
\usepackage{xcolor}
\usepackage{booktabs}
\usepackage{tabularx}
\usepackage{longtable}
\usepackage{array}
\usepackage{enumitem}
\IfFileExists{mathtools.sty}{\usepackage{mathtools}}{}
\usepackage{listings}
\usepackage[most]{tcolorbox}
\usepackage{fancyhdr}
\usepackage{titlesec}
\usepackage{caption}
\usepackage{imakeidx}
% csquotes with autostyle follows babel's active language, so the identical
% source \enquote{x} sets 'x' in the English edition and the Polish quotation
% marks in the Polish one. This is not a nicety: babel-polish makes " an active
% shorthand, which is turned off above precisely because an active character
% near a listings body is a category-code accident waiting to happen -- so a
% quotation mark has to come from a macro rather than from a keyboard key.
\IfFileExists{csquotes.sty}{\usepackage[autostyle=true]{csquotes}}{%
  \providecommand{\enquote}[1]{``##1''}}
\usepackage[hidelinks]{hyperref}

% ---------- Palette ----------
% Deliberately the same family as the sibling books, shifted green so the three
% spines are distinguishable on a shelf.
\definecolor{mblue}{HTML}{1F4E79}
\definecolor{mgreen}{HTML}{2E6B4F}
\definecolor{mteal}{HTML}{0E7C7B}
\definecolor{mamber}{HTML}{B26A00}
\definecolor{mred}{HTML}{9B2C2C}
\definecolor{mgrey}{HTML}{5A5A5A}
\definecolor{mlight}{HTML}{F4F6F8}
\definecolor{mfaint}{HTML}{FBFCFD}
\definecolor{codebg}{HTML}{FAFAFA}
\definecolor{codeframe}{HTML}{D8DEE4}
\definecolor{kw}{HTML}{0B5394}
\definecolor{str}{HTML}{9B2C2C}
\definecolor{cmt}{HTML}{4F7A4F}
\definecolor{typ}{HTML}{0E7C7B}

\hypersetup{
  colorlinks=true,
  linkcolor=mblue,
  urlcolor=mteal,
  citecolor=mblue,
  pdfauthor={Konrad Cinkusz}
}

% \ifpl{Polish text}{English text} -- for the handful of places where a string
% is structural rather than content and belongs in the shared files:
% structure.tex's part titles, and the main files' front matter wiring.
% Anything longer than a few words belongs in lang/ or in the edition's own
% source, not here.
% \DeclareRobustCommand, not \newcommand. A fragile \ifpl inside a \part title
% is expanded one level as it is written to the .toc, which lands
% "\ifdefstring {en}{pl}{...}{...}" in the file -- and that fails on the next
% run with "Extra }, or forgotten \endgroup", an error that names neither the
% part nor this macro.
\DeclareRobustCommand{\ifpl}[2]{\ifdefstring{\booklang}{pl}{#1}{#2}}

% A robust macro cannot go into a PDF bookmark string either -- hyperref strips
% \ifdefstring and warns once per part title. Tell it which branch to take
% instead, so the bookmarks carry the right language rather than nothing.
\makeatletter
\pdfstringdefDisableCommands{%
  \ifdefstring{\booklang}{pl}{\let\ifpl\@firstoftwo}{\let\ifpl\@secondoftwo}%
}
\makeatother

% ---------- Language strings ----------
% Every user-visible word the macros below typeset comes from here. Nothing in
% programs/ or appendices/ may hard-code an English or a Polish word inside a
% macro argument that the other edition also uses.
\input{lang/\booklang}

% ============================================================
%  PROGRAM NUMBERING
%  ---------------------------------------------------------
%  Part F programs are F1..F13; the main programs restart at 1. That is
%  Stroud's scheme and it is worth reproducing, because the Foundation part is
%  explicitly optional and numbering it separately is what makes skipping it
%  feel permitted rather than delinquent.
%
%  \theHchapter is hyperref's *destination* name, and it must stay unique
%  across the whole document. Resetting the chapter counter without also
%  changing \theHchapter produces duplicate PDF destinations, a pile of
%  warnings, and bookmarks that jump to the wrong program.
% ============================================================
\newcommand{\foundationnumbering}{%
  \setcounter{chapter}{0}%
  \renewcommand{\thechapter}{F\arabic{chapter}}%
  \renewcommand{\theHchapter}{F\arabic{chapter}}%
}
\newcommand{\mainnumbering}{%
  \setcounter{chapter}{0}%
  \renewcommand{\thechapter}{\arabic{chapter}}%
  \renewcommand{\theHchapter}{M\arabic{chapter}}%
}

% \appendix resets the chapter counter and switches \thechapter to letters, but
% it knows nothing about \theHchapter, so appendix A would claim the same PDF
% destination as main program 1. Patch it once, here.
\apptocmd{\appendix}{\renewcommand{\theHchapter}{A\Alph{chapter}}}{}{%
  \PackageWarning{mfabook}{Could not patch \string\appendix\space for hyperref destinations}}

% A program is a chapter. \chaptername is set from the language file so the
% word above the title is "Program" in both editions (it happens to be the
% same word, which is a small piece of luck).
\AtBeginDocument{\renewcommand{\chaptername}{\lblProgram}}

% ---------- Headings ----------
% \raggedright on the title body is load-bearing, not cosmetic: at \Huge on a
% 13 cm text block a long title that cannot break overflows the margin badly.
\titleformat{\chapter}[display]
  {\normalfont\huge\bfseries\color{mblue}\raggedright}
  {\normalfont\normalsize\color{mgrey}\MakeUppercase{\chaptertitlename}\ \thechapter}
  {8pt}{\Huge\raggedright}
\titlespacing*{\chapter}{0pt}{10pt}{28pt}
\titleformat{\section}{\normalfont\Large\bfseries\color{mblue}}{\thesection}{0.8em}{}
\titleformat{\subsection}{\normalfont\large\bfseries\color{mgrey}}{\thesubsection}{0.7em}{}

% head=23pt in the geometry options above, not \setlength{\headheight}:
% fancyhdr computes the height it needs from the running-head font and the
% rule, wants 22.5 pt at this size, and silently overprints the top of the text
% block if it does not get it. Setting it through geometry keeps the text block
% where geometry put it instead of pushing it down.
\pagestyle{fancy}
\fancyhf{}
\fancyhead[LE]{\small\nouppercase{\leftmark}}
% The recto head carries the FRAME number, not the section. In a book of
% numbered frames "where am I" is a frame question: every Summary
% back-reference, every Quiz route and every cross-reference between programs
% is a frame number, and hunting for one by page is the single most annoying
% thing about reading a programmed text. \theframeno at shipout is the last
% frame begun on the page, which is what a reader flicking through wants.
% Guarded: the front matter and the appendices have no frames, and a head
% reading "Frame 0" over the introduction is worse than no head at all.
\fancyhead[RO]{\ifnum\value{frameno}>0\relax\small\lblFrame~\theframeno\fi}
\fancyfoot[LE,RO]{\small\thepage}
\renewcommand{\headrulewidth}{0.4pt}

% This MUST come after \pagestyle{fancy}: fancyhdr installs its own \chaptermark
% at that point and silently overwrites anything defined earlier. The symptom is
% a running-head change that appears to do nothing at all.
%
% The running head carries the *frame* number as well as the program, because
% in a book of numbered frames "where am I" is a frame question, not a page
% question, and the summary's back-references are frame numbers.
% \@chapapp, not \lblProgram. \appendix switches \@chapapp from \chaptername
% to \appendixname, and babel supplies both in the right language; hard-coding
% the word gives an appendix a running head reading "Program A. Answers" over a
% page whose own chapter head reads "APPENDIX A".
\makeatletter
\renewcommand{\chaptermark}[1]{\markboth{\@chapapp\ \thechapter.\ #1}{}}
\makeatother
\renewcommand{\sectionmark}[1]{\markright{\thesection\ #1}}

% ============================================================
%  FRAMES — the unit the whole method is built from
%  ---------------------------------------------------------
%  A program is a stream of numbered frames. Nearly every frame ends by asking
%  the reader for something, and the answer opens the *next* frame, so that the
%  reader can cover what is below and commit to an answer before seeing it.
%
%  The environment is called `fr` rather than `frame` for two reasons. First,
%  \frame is already defined by LaTeX2e (it draws a box in a picture
%  environment) and an environment named `frame` would silently redefine it.
%  Second, this is typed forty to seventy times per program, so it wants to be
%  short in the source.
%
%    \begin{fr}
%    \ans{$x = 3$}            % the answer to the previous frame, optional
%    Teaching text, ending in a question.
%    \end{fr}
% ============================================================

\newcounter{frameno}[chapter]
\renewcommand{\theframeno}{\arabic{frameno}}

% Frames are referenced constantly -- by the Summary, by the Quiz, by the
% "Can you?" checklist and by other programs. \label inside a frame must
% therefore capture the frame number, not the section number.
\makeatletter
\newcommand{\framelabelfix}{%
  \@namedef{theHframeno}{\theHchapter.\arabic{frameno}}%
}
\makeatother

\newcommand{\framerule}{%
  \par\penalty-200\vspace{9pt}%
  \noindent\textcolor{codeframe}{\rule{\linewidth}{0.5pt}}%
  \par\vspace{5pt}%
}

% The frame number as a filled badge, set in the text flow rather than in the
% margin. A margin number would have to alternate sides under `twoside` and
% would be the first thing to overflow on a 17 cm page; a badge cannot.
\newcommand{\framebadge}{%
  \begingroup
  \setlength{\fboxsep}{2.5pt}%
  \colorbox{mblue}{\textcolor{white}{\bfseries\scriptsize\theframeno}}%
  \endgroup\hspace{0.7em}%
}

\newenvironment{fr}{%
  \framerule
  \refstepcounter{frameno}%
  \nopagebreak
  \noindent\framebadge\ignorespaces
}{\par}

% The answer to the previous frame. Set apart and centred so the eye can find
% the edge to cover, and so a reader who has answered can check in one glance
% without reading on.
\newcommand{\ans}[1]{%
  \par\nobreak\vspace{-4pt}%
  \begin{tcolorbox}[
    enhanced, sharp corners, boxrule=0pt, leftrule=2.5pt,
    colback=mfaint, colframe=mgreen,
    left=8pt, right=8pt, top=4pt, bottom=4pt,
    width=0.86\linewidth, center, halign=center,
    before skip=2pt, after skip=7pt]
    #1
  \end{tcolorbox}%
  \nobreak\noindent\ignorespaces
}

% A multi-paragraph or worked answer, where a centred box would be wrong.
\newenvironment{ansblock}{%
  \par\nopagebreak
  \begin{tcolorbox}[
    enhanced, breakable, sharp corners, boxrule=0pt, leftrule=2.5pt,
    colback=mfaint, colframe=mgreen,
    left=8pt, right=8pt, top=5pt, bottom=5pt]
}{%
  \end{tcolorbox}%
  \nopagebreak\par\vspace{2pt}\noindent\ignorespaces
}

% The row of dots. Stroud's signal that the reader is to supply a word, a
% number or an expression. \blank is inline; \dotline occupies its own line,
% which is what a longer derivation step wants.
% \penalty0 offers TeX a break opportunity immediately before the dots. Without
% it the run of dots is glued to the last word of the question, and a long
% question in the Polish edition -- where the words are longer and hyphenate
% less freely -- runs into the margin.
\newcommand{\blank}[1][4em]{\penalty0\makebox[#1]{\dotfill}}
\newcommand{\dotline}{\par\vspace{2pt}\noindent\makebox[\linewidth]{\dotfill}\par\vspace{2pt}}

% "Now one for you to do." The third rung of the scaffolding gradient, and the
% one most often skipped by authors who think they have written an exercise
% when they have written another worked example.
\newcommand{\yourturn}{%
  \par\vspace{4pt}\noindent
  {\bfseries\color{mgreen}\lblYourTurn}\par\nobreak\vspace{2pt}\noindent\ignorespaces
}

% ============================================================
%  THE PROGRAM SKELETON
%  ---------------------------------------------------------
%  Every program carries the same seven parts, in the same order. They are
%  macros rather than a convention so that a missing one is a build-time
%  absence rather than something a reader notices.
% ============================================================

% ---- Learning outcomes, on entry ----------------------------------------
% Stored as well as printed, because "Can you?" at the end of the program must
% be the same list, one for one. Storing it is what makes that structural
% rather than a promise: the checklist is generated from the outcomes, so the
% two cannot drift.
\newcounter{outcomeno}[chapter]

\makeatletter
% The key under which everything belonging to the current program is stored.
% It must expand to the printed program number (F1, F2, ..., 1, 2, ...), which
% is what makes an answer findable from the program that owns it.
\newcommand{\mfa@progkey}{\thechapter}
\newcommand{\outcome}[2]{% #1 = frames that teach it, #2 = the outcome
  \stepcounter{outcomeno}%
  \csgdef{mfa@out@\mfa@progkey @\theoutcomeno}{#2}%
  \csgdef{mfa@outfr@\mfa@progkey @\theoutcomeno}{#1}%
  % The "Can you?" row is built here rather than looped over later. A \loop
  % inside a tabularx body does not survive alignment scanning: TeX needs to
  % see the & and the row break while it reads the cell, and a conditional
  % hides them. Accumulating the rows as tokens at declaration time sidesteps
  % that entirely, and has the better property anyway -- the checklist is
  % literally made of the outcomes, so the two cannot disagree.
  \csgappto{mfa@rows@\mfa@progkey}{%
    #2 & {\footnotesize\color{mgrey}#1}
      & \ratingcell & \ratingcell & \ratingcell & \ratingcell & \ratingcell
    \tabularnewline}%
  \item #2\hfill{\footnotesize\color{mgrey}[\lblFrames~#1]}%
}
\makeatother

\makeatletter
\newenvironment{outcomes}{%
  \setcounter{outcomeno}{0}%
  \csgdef{mfa@rows@\mfa@progkey}{}%
  \begin{tcolorbox}[
    enhanced, breakable, sharp corners,
    colback=mlight, colframe=mblue, boxrule=0pt, leftrule=3pt,
    left=8pt, right=8pt, top=6pt, bottom=6pt,
    fonttitle=\bfseries\small, title={\lblOutcomes}]
  \begin{itemize}[leftmargin=1.3em, itemsep=3pt, topsep=2pt]
}{%
  \end{itemize}
  \end{tcolorbox}
}
\makeatother

% ---- "Can you?" -----------------------------------------------------------
% Generated from the stored outcomes, not written out again by hand. Five
% columns rather than a tick box, because a self-assessment with a midpoint is
% one people actually use and a yes/no one is one they lie to.
\newcommand{\ratingcell}{\textcolor{codeframe}{\rule[-0.15em]{0.8em}{0.8em}}}

\makeatletter
\newcommand{\canyou}{%
  \section*{\lblCanYou}%
  \phantomsection\addcontentsline{toc}{section}{\lblCanYou}%
  \noindent\lblCanYouIntro\par\vspace{8pt}
  \noindent
  \begingroup\small
  \ifcsvoid{mfa@rows@\mfa@progkey}{%
    {\color{mred}\lblNoOutcomes}\par
  }{%
    \begin{tabularx}{\linewidth}{@{}Xcccccc@{}}
    \toprule
    \textbf{\lblOnEntryYouCould} & \textbf{\lblFrames} & 1 & 2 & 3 & 4 & 5 \\
    \midrule
    \csuse{mfa@rows@\mfa@progkey}
    \bottomrule
    \end{tabularx}%
  }%
  \endgroup
  \par\vspace{6pt}\noindent{\footnotesize\color{mgrey}\lblCanYouFooter}\par
}
\makeatother

% ============================================================
%  QUIZ, SUMMARY, EXERCISES
% ============================================================

% ---- Quiz -----------------------------------------------------------------
% Foundation programs only. It is a diagnostic on the way in and the same test
% on the way out, which is what lets one book serve a reader who needs the
% whole program and a reader who needs four frames of it. Each question names
% the frames that teach it, so a wrong answer routes the reader somewhere
% specific rather than back to the beginning.
\newlist{quizlist}{enumerate}{1}
\setlist[quizlist]{label=\arabic*., leftmargin=2.0em, itemsep=5pt, topsep=3pt}

\makeatletter
\newenvironment{quiz}{%
  \section*{\lblQuiz}%
  \phantomsection\addcontentsline{toc}{section}{\lblQuiz}%
  \def\mfa@exkey{Q\arabic{quizlisti}}%
  \noindent\lblQuizIntro\par\vspace{5pt}
  \begin{quizlist}
}{%
  \end{quizlist}
}
\makeatother

% The frames that teach the question just asked. Two forms, because a Quiz
% question that routes to a single frame reading "Frames 22" is the kind of
% small wrongness a reader notices and an author never does.
\newcommand{\teachesat}[1]{\hfill{\footnotesize\color{mgrey}(\lblFrames~#1)}}
\newcommand{\teachesatone}[1]{\hfill{\footnotesize\color{mgrey}(\lblFrame~#1)}}

% ---- Summary --------------------------------------------------------------
% Every item carries the frame number it came from, in square brackets. That
% is the eighth-edition improvement and it is the single cheapest thing in the
% format: it turns a summary into a return index, so a reader who fails one
% line of "Can you?" knows exactly which four frames to reread.
\newenvironment{summarybox}{%
  \section*{\lblSummary}%
  \phantomsection\addcontentsline{toc}{section}{\lblSummary}%
  \begin{tcolorbox}[
    enhanced, breakable, sharp corners,
    colback=mlight, colframe=mblue, boxrule=0pt, leftrule=3pt,
    left=8pt, right=8pt, top=6pt, bottom=6pt]
  \begin{itemize}[leftmargin=1.3em, itemsep=5pt, topsep=2pt]
}{%
  \end{itemize}
  \end{tcolorbox}
}

% \sumitem{12}{text} -- the frame reference is not decoration, it is the point.
\newcommand{\sumitem}[2]{\item #2~{\footnotesize\color{mgrey}[#1]}}

% ---- Answers ---------------------------------------------------------------
%  Answers are authored at the point of use and printed at the back.
%
%  They are carried there by a global macro store rather than by an auxiliary
%  file. The sibling books collect their figure and screenshot manifests with
%  \@starttoc and a custom file extension, which works because those entries
%  are one line of text; it is the wrong mechanism for an answer, because an
%  answer contains displayed maths, and material written to an .aux-style file
%  and read back is expanded at shipout, where fragile commands break in ways
%  whose error messages name neither the answer nor the program.
%
%  \include is sequential within a run, so a macro defined globally while
%  typesetting program F1 is still defined when the back matter is typeset.
%  That is all this needs.
\makeatletter
\newcommand{\mfa@storeanswer}[3]{% #1 program key, #2 item key, #3 body
  \csgdef{mfa@a@#1@#2}{#3}%
  % \listcsxadd, not \listcsgadd: the item must be EXPANDED as it is stored.
  % Storing the token \mfa@exkey means every entry in the list expands to
  % whatever the key happens to be at the time the back matter is typeset,
  % which is the last one.
  \listcsxadd{mfa@akeys@#1}{#2}%
}
% Used inside an exercise item. Prints nothing here; the body reappears at the
% back of the book under this program's heading.
\newcommand{\answerto}[1]{\mfa@storeanswer{\mfa@progkey}{\mfa@exkey}{#1}}
\providecommand{\mfa@exkey}{?}
\makeatother

\newlist{exlist}{enumerate}{1}
\setlist[exlist]{label=\arabic*., leftmargin=2.0em, itemsep=6pt, topsep=3pt}
\newlist{fplist}{enumerate}{1}
\setlist[fplist]{label=\arabic*., leftmargin=2.0em, itemsep=5pt, topsep=3pt}

% These MUST be defined inside \makeatletter. Outside it, \def\mfa@exkey{...}
% parses as \def\mfa @exkey{...} -- a definition of \mfa with a delimited
% parameter text -- which raises no error at all and quietly leaves every
% answer filed under the same key. That failure is invisible until the answers
% section at the back of the book prints one program's answers twice.
\makeatletter
\newenvironment{testexercises}{%
  \section*{\lblTestExercises}%
  \phantomsection\addcontentsline{toc}{section}{\lblTestExercises}%
  \def\mfa@exkey{T\arabic{exlisti}}%
  \noindent\lblTestIntro\par\vspace{5pt}
  \begin{exlist}
}{%
  \end{exlist}
}

\newenvironment{furtherproblems}{%
  \section*{\lblFurther}%
  \phantomsection\addcontentsline{toc}{section}{\lblFurther}%
  \def\mfa@exkey{P\arabic{fplisti}}%
  \noindent\lblFurtherIntro\par\vspace{5pt}
  \begin{fplist}
}{%
  \end{fplist}
}
\makeatother

% ---- The programs register -------------------------------------------------
% \program{Title} opens a program: it is \chapter plus the bookkeeping that
% lets the back matter reproduce this program's answers under its own name.
\makeatletter
\newcommand{\program}[1]{%
  \chapter{#1}%
  \csgdef{mfa@title@\thechapter}{#1}%
  % Expanded, for the same reason as the answer keys above.
  \listxadd{\mfaprograms}{\thechapter}%
}
\newcommand{\mfa@printoneanswer}[2]{% #1 program key, #2 item key
  \item[\textbf{#2}] \csuse{mfa@a@#1@#2}%
}
\newcommand{\mfa@answerprogram}[1]{%
  \subsection*{\lblProgram~#1\quad\textmd{\itshape\csuse{mfa@title@#1}}}%
  \ifcsdef{mfa@akeys@#1}{%
    \begin{list}{}{%
      \setlength{\leftmargin}{3.2em}\setlength{\labelwidth}{2.7em}%
      \setlength{\labelsep}{0.5em}\setlength{\itemsep}{3pt}%
      \setlength{\topsep}{3pt}\setlength{\parsep}{0pt}%
      \renewcommand{\makelabel}[1]{\hfill##1}}
      \forlistcsloop{\mfa@printoneanswer{#1}}{mfa@akeys@#1}
    \end{list}%
  }{%
    \par\noindent{\color{mred}\lblNoAnswers}\par
  }%
}
% The answers body, without a heading of its own, so the appendix that carries
% it can be a numbered appendix like any other and appear in the running heads
% and the ToC in the ordinary way.
\newcommand{\answersbody}{%
  \noindent\lblAnswersIntro\par\vspace{8pt}
  \ifdef{\mfaprograms}{\forlistloop{\mfa@answerprogram}{\mfaprograms}}{%
    \noindent{\color{mred}\lblNoAnswers}\par}%
}
% Unnumbered form, for a draft build that has no appendix wiring yet.
\newcommand{\printanswers}{%
  \chapter*{\lblAnswers}%
  \phantomsection\addcontentsline{toc}{chapter}{\lblAnswers}%
  \markboth{\lblAnswers}{\lblAnswers}%
  \answersbody
}
\makeatother

% ============================================================
%  ADMONITIONS
%  ---------------------------------------------------------
%  A deliberately small, fixed vocabulary. Each one earns its place by doing
%  something no other box does; a seventh would be a smell.
% ============================================================

\newtcolorbox{note}[1][]{
  enhanced, breakable, sharp corners,
  colback=mlight, colframe=mblue, boxrule=0pt, leftrule=3pt,
  left=8pt, right=8pt, top=6pt, bottom=6pt,
  fonttitle=\bfseries\small, title={\lblNote}, #1
}

\newtcolorbox{warning}[1][]{
  enhanced, breakable, sharp corners,
  colback=mlight, colframe=mred, boxrule=0pt, leftrule=3pt,
  left=8pt, right=8pt, top=6pt, bottom=6pt,
  fonttitle=\bfseries\small, title={\lblWarning}, #1
}

% The misconception, stated as a summary after it has been walked into. The
% trap itself belongs in the frames -- elicit the wrong answer, then say
% flatly that it is wrong -- and this box is the thing the reader marks with a
% finger and comes back to.
\newtcolorbox{trapbox}[1][]{
  enhanced, breakable, sharp corners,
  colback=white, colframe=mred, boxrule=0.8pt,
  left=8pt, right=8pt, top=6pt, bottom=6pt,
  fonttitle=\bfseries\small, title={\lblTrap}, #1
}

% Where the mathematics just done shows up in a system the reader has actually
% deployed. This is the box that separates the book from a maths textbook, and
% it is also the one most easily filled with waffle: if it cannot name a
% specific line of a specific system, it should not be written.
\newtcolorbox{aibox}[1][]{
  enhanced, breakable, sharp corners,
  colback=mlight, colframe=mteal, boxrule=0pt, leftrule=3pt,
  left=8pt, right=8pt, top=6pt, bottom=6pt,
  fonttitle=\bfseries\small, title={\lblInAI}, #1
}

% What is being asserted rather than proved, and where to read the proof.
% Stroud's method buys fluency at the cost of rigour, and the honest thing is
% to say so at each point rather than once in an introduction nobody rereads.
\newtcolorbox{rigourbox}[1][]{
  enhanced, breakable, sharp corners,
  colback=white, colframe=mgrey, boxrule=0.6pt,
  left=8pt, right=8pt, top=6pt, bottom=6pt,
  fonttitle=\bfseries\small, title={\lblRigour}, #1
}

% Notation that genuinely differs -- between the Polish and English editions,
% and between disciplines. Matrix calculus alone has two incompatible layout
% conventions, and half the confusion about backpropagation is a transpose
% that came from the other one.
\newtcolorbox{notationbox}[1][]{
  enhanced, breakable, sharp corners,
  colback=mlight, colframe=mamber, boxrule=0pt, leftrule=3pt,
  left=8pt, right=8pt, top=6pt, bottom=6pt,
  fonttitle=\bfseries\small, title={\lblNotation}, #1
}

% The maths analogue of the sibling books' "run this before you trust it".
% A number in this book is either produced by a script under code/ and pulled
% in with \val, or it carries this box. Nothing ships to a reader with one
% still attached.
\newtcolorbox{verifybox}[1][]{
  enhanced, breakable, sharp corners,
  colback=white, colframe=mamber, boxrule=0.8pt,
  left=8pt, right=8pt, top=6pt, bottom=6pt,
  fonttitle=\bfseries\small, title={\lblVerify}, #1
}

\newtcolorbox{exercisebox}[1][]{
  enhanced, breakable, sharp corners,
  colback=white, colframe=mgreen, boxrule=0pt, leftrule=3pt,
  left=8pt, right=8pt, top=6pt, bottom=6pt,
  fonttitle=\bfseries\small, title={\lblExercise}, #1
}

% ============================================================
%  CODE LISTINGS
%  ---------------------------------------------------------
%  Code appears here for one reason only: to check a number. Every listing in
%  this book is runnable and every number it prints is the number in the text.
% ============================================================
\lstdefinelanguage{pythonx}{
  language=Python,
  morekeywords={as, with, assert, match, case},
  morekeywords=[2]{
    float64,float32,float16,bfloat16,int8,uint8,
    np,numpy,array,asarray,frombuffer,nextafter,finfo,iinfo,
    exp,log,log1p,expm1,logaddexp,sum,max,abs,sqrt,mean,var
  },
  sensitive=true
}

\lstdefinestyle{mfacode}{
  basicstyle=\ttfamily\footnotesize,
  keywordstyle=\color{kw}\bfseries,
  keywordstyle=[2]\color{typ},
  stringstyle=\color{str},
  commentstyle=\color{cmt}\itshape,
  numbers=left,
  numberstyle=\ttfamily\tiny\color{mgrey},
  numbersep=8pt,
  backgroundcolor=\color{codebg},
  frame=single,
  rulecolor=\color{codeframe},
  framesep=6pt,
  breaklines=true,
  breakatwhitespace=false,
  postbreak=\mbox{\textcolor{mgrey}{$\hookrightarrow$}\space},
  showstringspaces=false,
  tabsize=4,
  columns=fullflexible,
  keepspaces=true,
  captionpos=b,
  aboveskip=1.2em,
  belowskip=1.2em,
  % Maps the characters most likely to be pasted into a listing by accident.
  %
  % Do NOT add a mapping for U+00A0 (non-breaking space). It looks like the
  % obvious next entry and it makes `listings` abort the run with "Improper
  % alphabetic constant" -- fatal, no PDF, and the message names neither the
  % character nor this line. The ASCII-in-listings convention is the guard
  % against it instead. This trap has now been hit in two of the three books
  % in this series.
  literate={ł}{{\l}}1 {Ł}{{\L}}1 {ą}{{\k a}}1 {ę}{{\k e}}1
           {ś}{{\'s}}1 {ć}{{\'c}}1 {ż}{{\.z}}1 {ź}{{\'z}}1 {ó}{{\'o}}1 {ń}{{\'n}}1
           {Ą}{{\k A}}1 {Ę}{{\k E}}1 {Ś}{{\'S}}1 {Ć}{{\'C}}1
           {Ż}{{\.Z}}1 {Ź}{{\'Z}}1 {Ó}{{\'O}}1 {Ń}{{\'N}}1
           {—}{{-{}-}}2 {–}{{-}}1 {‘}{{'}}1 {’}{{'}}1
           {“}{{"}}1 {”}{{"}}1 {°}{{\textdegree}}1
}

\lstset{style=mfacode}

\lstnewenvironment{python}[1][]
  {\lstset{language=pythonx,style=mfacode,#1}}
  {}

\lstnewenvironment{shellcmd}[1][]
  {\lstset{language=bash,style=mfacode,numbers=none,keywordstyle=\color{mteal}\bfseries,#1}}
  {}

% Terminal transcripts: no line numbers, no keyword colouring. Everything in
% one of these was copied from a real run.
\lstnewenvironment{console}[1][]
  {\lstset{style=mfacode,numbers=none,basicstyle=\ttfamily\scriptsize,
           keywordstyle=\color{black},backgroundcolor=\color{white},#1}}
  {}

\newcommand{\pyfile}[3]{%
  \lstinputlisting[language=pythonx,style=mfacode,caption={#2},label={#3}]{#1}%
}
\newcommand{\pyregion}[4]{%
  \lstinputlisting[
    language=pythonx, style=mfacode,
    linerange={--8<--\ [start:#2]--- 8<--\ [end:#2]},
    includerangemarker=false,
    caption={#3}, label={#4}]{#1}%
}

% Inline literals
\newcommand{\code}[1]{\texttt{\small #1}}
\newcommand{\pkg}[1]{\texttt{\small\color{mblue}#1}}
\newcommand{\api}[1]{\texttt{\small\color{mteal}#1}}

% ============================================================
%  COMPUTED VALUES
%  ---------------------------------------------------------
%  The house rule in the sibling books is "run the listing, or mark it". The
%  analogue here is stronger, because a wrong digit in a maths book is not a
%  broken example, it is a lie the reader will carry.
%
%  Every number that came out of a computation is written by a script under
%  code/ into figures/values/<program>.tex as
%
%      \mfaval{f01.eps64}{2.220446049250313e-16}
%
%  and pulled into the text with \val{f01.eps64}. The script is the source of
%  truth; the book cannot disagree with it, because the book does not contain
%  the digits. A value the scripts no longer produce prints a visible marker
%  and is counted by `make debt`.
%
%  \val also applies the locale's decimal separator, which is how the Polish
%  edition gets its decimal comma without a translator having to remember.
% ============================================================
\IfFileExists{siunitx.sty}{%
  \usepackage{siunitx}%
  \sisetup{
    group-digits = integer,
    group-minimum-digits = 5,
    group-separator = {\,},
    exponent-product = \cdot,
    output-decimal-marker = {\mfadecimalmarker}
  }%
  \newcommand{\mfanum}[1]{\num{##1}}%
}{%
  \newcommand{\mfanum}[1]{##1}%
}

\makeatletter
\newcommand{\mfaval}[2]{\csgdef{mfaval@#1}{#2}\listxadd{\mfavalues}{#1}}
% \num on a non-numeric body is a FATAL siunitx error -- "Invalid number" --
% and under -halt-on-error that is the whole build. A script emitting a word
% where the book expected a number is a plausible mistake.
%
% The check is NOT done here. Deciding in TeX whether a token list is a number
% means parsing it character by character, which is fragile and was got wrong
% once already. The generator knows the answer for free, so it emits \mfaval
% for a number and \mfavaltext for a word, and this end only has to enforce
% which macro reads which.
\newcommand{\val}[1]{%
  \ifcsdef{mfaval@#1}{%
    \ifcsdef{mfavaltext@#1}{%
      \PackageError{mfabook}{Value `#1' is not a number}%
        {\string\val\space passes its body to siunitx, which refuses anything
         that is not a number. Use \string\valtext{#1} instead, or fix the
         script in code/ that emits this key.}%
    }{}%
    \mfanum{\csuse{mfaval@#1}}%
  }{\textcolor{mred}{\textbf{[?\,#1]}}}%
}
% A value that is deliberately a word rather than a number -- a format name, a
% verdict, a unit. Rare, and a visibly different macro so nobody reaches for
% \val and gets an error they cannot place.
\newcommand{\valtext}[1]{%
  \ifcsdef{mfaval@#1}{\csuse{mfaval@#1}}%
    {\textcolor{mred}{\textbf{[?\,#1]}}}%
}
\newcommand{\mfavaltext}[2]{%
  \csgdef{mfaval@#1}{#2}\csgdef{mfavaltext@#1}{}\listxadd{\mfavalues}{#1}%
}

% The raw stored string, for the cases where a number must appear exactly as
% the machine printed it -- inside a console block, or where the digits
% themselves are the point and grouping them would obscure the pattern.
\newcommand{\rawval}[1]{%
  \ifcsdef{mfaval@#1}{\csuse{mfaval@#1}}%
    {\textcolor{mred}{\textbf{[?\,#1]}}}%
}
\makeatother

% figures/values/all.tex is generated by `make numbers` and does nothing but
% \input each program's value file. A build with no values yet still compiles;
% every \val{} in it prints a visible marker instead of a number, which is the
% correct behaviour for a draft and an obvious failure for a release.
\IfFileExists{figures/values/all.tex}{% Generated by `make numbers` --- do not edit.
\input{figures/values/f01}
}{%
  \AtBeginDocument{\PackageWarning{mfabook}{No computed values found. Run `make numbers`.}}}

% ============================================================
%  DIAGRAMS — Mermaid pipeline
%  ---------------------------------------------------------
%  \mermaidfig{key}{caption}{one-line manifest description}
%
%  Source of truth is figures/mermaid/<lang>/<key>.mmd, committed. `make
%  diagrams` renders each to figures/diagrams/<lang>/<key>.pdf, which is build
%  output and is not committed, so a diagram change reviews as a text diff.
%
%  The source is per-language because a diagram with English labels in the
%  Polish edition is exactly the kind of half-translation that makes a
%  bilingual book feel machine-made. CI checks that both languages carry the
%  same set of keys.
%
%  If the rendered PDF is absent the macro draws a placeholder AND typesets the
%  Mermaid source, so a build without mermaid-cli still shows the structure.
% ============================================================
\makeatletter
\newcommand{\listofdiagrams}{%
  \section*{\lblDiagramManifest}%
  \phantomsection\addcontentsline{toc}{section}{\lblDiagramManifest}%
  \@starttoc{dgm}%
}
\makeatother

\newcommand{\mermaidfig}[3]{%
  \addcontentsline{dgm}{subsection}{\protect\numberline{}\texttt{#1.mmd} --- #3}%
  \IfFileExists{figures/diagrams/\booklang/#1.pdf}{%
    \begin{figure}[htbp]
    \centering
    \includegraphics[width=0.95\linewidth,height=0.42\textheight,keepaspectratio]{figures/diagrams/\booklang/#1.pdf}%
    \caption{#2}\label{fig:#1}
    \end{figure}%
  }{%
    % The unrendered case deliberately does NOT float. A Mermaid source
    % typeset verbatim is often taller than a page, and a float cannot break,
    % so wrapping it in `figure` produces an overfull box per diagram and a
    % pile of tcolorbox warnings. In the flow it breaks like any other text.
    \par\medskip
    \begin{tcolorbox}[
      enhanced, breakable, width=\linewidth, sharp corners,
      colback=white, colframe=mgrey, boxrule=0.8pt,
      left=8pt, right=8pt, top=8pt, bottom=8pt]
      {\bfseries\color{mgrey}\small \lblNotRendered}\\[4pt]
      {\ttfamily\scriptsize\color{mgrey} figures/mermaid/\booklang/#1.mmd}\\[6pt]
      \IfFileExists{figures/mermaid/\booklang/#1.mmd}{%
        % ASCII only, like every other listing in this book: the fallback runs
        % the Mermaid source through `listings`, which aborts on a multi-byte
        % character it has no literate mapping for.
        \lstinputlisting[style=mfacode,numbers=none,basicstyle=\ttfamily\scriptsize,
                         backgroundcolor=\color{mlight},frame=none,
                         keywordstyle=\color{black}]{figures/mermaid/\booklang/#1.mmd}%
      }{%
        {\small\color{mred} \lblSourceMissing: \texttt{figures/mermaid/\booklang/#1.mmd}}%
      }%
    \end{tcolorbox}%
    \captionof{figure}{#2}\label{fig:#1}
    \par\medskip
  }%
}

% ============================================================
%  PROGRAM STUBS
%  ---------------------------------------------------------
%  Prints a visible marker so an unwritten program cannot be mistaken for a
%  written one and the page count never flatters the draft. `make debt` counts
%  them, and CI publishes the count on every build.
% ============================================================
\newcommand{\programstub}[1]{%
  \begin{tcolorbox}[
    enhanced, breakable, width=\linewidth, sharp corners,
    colback=mlight, colframe=mred, boxrule=1pt,
    borderline={0.6pt}{3pt}{mred, dashed},
    left=12pt, right=12pt, top=12pt, bottom=12pt]
    {\bfseries\color{mred}\large \lblNotWritten}\\[8pt]
    \small\raggedright #1
  \end{tcolorbox}%
}

% ============================================================
%  MATHEMATICAL OPERATORS
%  ---------------------------------------------------------
%  Language-invariant names live here. The ones that genuinely differ between
%  Polish and English usage -- tan/tg, Var/D^2, gcd/NWD, and the interval
%  brackets -- are defined in lang/en.tex and lang/pl.tex instead, so the
%  notation contract is executed by the build rather than remembered by a
%  translator.
% ============================================================
\DeclareMathOperator*{\argmin}{arg\,min}
\DeclareMathOperator*{\argmax}{arg\,max}
\DeclareMathOperator{\tr}{tr}
\DeclareMathOperator{\rank}{rank}
\DeclareMathOperator{\diag}{diag}
\DeclareMathOperator{\sgn}{sgn}
\DeclareMathOperator{\relu}{ReLU}
\DeclareMathOperator{\softmax}{softmax}
\DeclareMathOperator{\logsumexp}{logsumexp}
\DeclareMathOperator{\fl}{fl}
\DeclareMathOperator{\ulp}{ulp}
\DeclareMathOperator{\round}{round}

% A bare \log is FORBIDDEN in this book, and tools/check_notation.py fails the
% build on one. In Polish textbooks \log means base ten; in machine-learning
% writing it means base e. The same three characters carry opposite meanings to
% the two readerships this book serves, and they collide hardest in exactly the
% programs where it matters most -- entropy in bits and cross-entropy in nats
% are two programs apart. Write \ln for nats and \logb{2} for bits. It costs
% two characters and removes the ambiguity entirely.
\makeatletter
\let\mfa@origlog\log
\renewcommand{\log}{%
  \PackageError{mfabook}{A bare \string\log\space is forbidden in this book}%
    {Write \string\ln\space for natural logarithms or \string\logb{2}\space for
     bits. The base is never left to a convention here; see Appendix B.}%
  \mfa@origlog}
% The one legitimate form: an explicit base.
\newcommand{\logb}[1]{\mathop{\mfa@origlog}\nolimits_{#1}}
% ...and a way to write the forbidden thing when the subject IS the ban, which
% is Appendix B and nowhere else.
\newcommand{\mfalogplain}{\mathop{\mfa@origlog}\nolimits}
\makeatother

\newcommand{\R}{\mathbb{R}}
\newcommand{\N}{\mathbb{N}}
\newcommand{\Z}{\mathbb{Z}}
\newcommand{\Q}{\mathbb{Q}}
\newcommand{\C}{\mathbb{C}}
\newcommand{\Fset}{\mathbb{F}}          % the floating-point numbers, F1's subject
\newcommand{\vect}[1]{\mathbf{#1}}
\newcommand{\mat}[1]{\mathbf{#1}}
\newcommand{\T}{^{\mathsf{T}}}
\newcommand{\norm}[1]{\left\lVert #1 \right\rVert}
\newcommand{\abs}[1]{\left\lvert #1 \right\rvert}
\newcommand{\inner}[2]{\left\langle #1, #2 \right\rangle}
\newcommand{\eps}{\varepsilon}

% ============================================================
%  INDEX
%  ---------------------------------------------------------
%  Two-column, and its entries include \texttt{} names that do not hyphenate.
%  Ragged right keeps multi-word entries off the margin; it cannot help a
%  single token wider than the column, so keep verbatim index entries under
%  about 29 characters and index longer names by concept instead. Both sibling
%  books learned this the same way.
% ============================================================
\apptocmd{\theindex}{\raggedright}{}{%
  \PackageWarning{mfabook}{Could not patch theindex for ragged-right setting}}

\makeindex

% ============================================================
%  PINNED FACTS — single source of truth
%  Change these here; they propagate through both editions.
% ============================================================
\newcommand{\bookedition}{0.1}
% \pinnedon is a user-visible string and lives in lang/<lang>.tex, not here.
% It said "August 2026" on the Polish title page for exactly as long as it
% took somebody to read the Polish title page.
\newcommand{\pyver}{Python 3.13}
\newcommand{\numpyver}{2.3}
 still compiles and says what it did.
\providecommand{\booklang}{en}

\usepackage{etoolbox}   % loaded first: the language switch below needs it

% ---------- Page geometry ----------
% Same 17 x 24 cm block as the other two books in this series, so a reader who
% owns one recognises the second. Frames are short, so the measure matters
% more here than in a prose book: a frame that runs past a page turn defeats
% the cover-the-next-frame discipline the whole method rests on.
\usepackage[
  paperwidth=17cm, paperheight=24cm,
  inner=2.2cm, outer=1.8cm, top=2.2cm, bottom=2.4cm,
  head=23pt, headsep=0.7cm, footskip=1.2cm
]{geometry}

% ---------- Encoding & fonts ----------
\usepackage[T1]{fontenc}
\usepackage[utf8]{inputenc}
\IfFileExists{lmodern.sty}{\usepackage{lmodern}}{}
\IfFileExists{newtxtext.sty}{\usepackage{newtxtext}}{}
% amsmath goes here, not down with the other packages, because newtxmath
% patches amsmath's internals and must be loaded AFTER it. The wrong order does
% not fail on a machine without newtxmath -- which is most CI images -- and
% fails on the author's full TeX Live.
%
% amssymb only when newtxmath is absent: newtxmath supplies the AMS symbols
% itself (its amssymbols option is on by default) and loading both clashes.
\usepackage{amsmath}
\IfFileExists{newtxmath.sty}{\usepackage{newtxmath}}{\usepackage{amssymb}}
\IfFileExists{inconsolata.sty}{\usepackage[varqu,scaled=0.95]{inconsolata}}{}
\IfFileExists{lmodern.sty}{\usepackage{microtype}}{\usepackage[expansion=false]{microtype}}

% ---------- Language ----------
% Loading babel with a language whose .ldf is absent is a *fatal* error, not a
% warning, so both the package and each language file are probed before either
% is requested. This is inherited from the sibling books, where a minimal TeX
% installation without texlive-lang-polish aborted the run with no PDF and an
% error message that named neither babel nor the language.
%
% The Polish edition needs Polish as the main language for hyphenation; the
% English edition needs British. Both editions load the other language too,
% because the glossary appendix is bilingual in both.
\IfFileExists{babel.sty}{%
  \IfFileExists{polish.ldf}{%
    \IfFileExists{british.ldf}{%
      \ifdefstring{\booklang}{pl}%
        {\usepackage[british,main=polish]{babel}}%
        {\usepackage[polish,main=british]{babel}}%
    }{\usepackage[polish]{babel}}%
  }{%
    \IfFileExists{british.ldf}{\usepackage[british]{babel}}{}%
  }%
}{}

% babel-polish makes " an active shorthand. listings reads its body verbatim
% and an active character in a verbatim context is a category-code accident
% waiting to happen, so the shorthand is turned off globally. Polish quotation
% marks are produced with \enquote-style macros from lang/pl.tex instead.
\ifdefstring{\booklang}{pl}{\AtBeginDocument{\shorthandoff{"}}}{}

% ---------- Core packages ----------
\usepackage{graphicx}
\usepackage{xcolor}
\usepackage{booktabs}
\usepackage{tabularx}
\usepackage{longtable}
\usepackage{array}
\usepackage{enumitem}
\IfFileExists{mathtools.sty}{\usepackage{mathtools}}{}
\usepackage{listings}
\usepackage[most]{tcolorbox}
\usepackage{fancyhdr}
\usepackage{titlesec}
\usepackage{caption}
\usepackage{imakeidx}
% csquotes with autostyle follows babel's active language, so the identical
% source \enquote{x} sets 'x' in the English edition and the Polish quotation
% marks in the Polish one. This is not a nicety: babel-polish makes " an active
% shorthand, which is turned off above precisely because an active character
% near a listings body is a category-code accident waiting to happen -- so a
% quotation mark has to come from a macro rather than from a keyboard key.
\IfFileExists{csquotes.sty}{\usepackage[autostyle=true]{csquotes}}{%
  \providecommand{\enquote}[1]{``##1''}}
\usepackage[hidelinks]{hyperref}

% ---------- Palette ----------
% Deliberately the same family as the sibling books, shifted green so the three
% spines are distinguishable on a shelf.
\definecolor{mblue}{HTML}{1F4E79}
\definecolor{mgreen}{HTML}{2E6B4F}
\definecolor{mteal}{HTML}{0E7C7B}
\definecolor{mamber}{HTML}{B26A00}
\definecolor{mred}{HTML}{9B2C2C}
\definecolor{mgrey}{HTML}{5A5A5A}
\definecolor{mlight}{HTML}{F4F6F8}
\definecolor{mfaint}{HTML}{FBFCFD}
\definecolor{codebg}{HTML}{FAFAFA}
\definecolor{codeframe}{HTML}{D8DEE4}
\definecolor{kw}{HTML}{0B5394}
\definecolor{str}{HTML}{9B2C2C}
\definecolor{cmt}{HTML}{4F7A4F}
\definecolor{typ}{HTML}{0E7C7B}

\hypersetup{
  colorlinks=true,
  linkcolor=mblue,
  urlcolor=mteal,
  citecolor=mblue,
  pdfauthor={Konrad Cinkusz}
}

% \ifpl{Polish text}{English text} -- for the handful of places where a string
% is structural rather than content and belongs in the shared files:
% structure.tex's part titles, and the main files' front matter wiring.
% Anything longer than a few words belongs in lang/ or in the edition's own
% source, not here.
% \DeclareRobustCommand, not \newcommand. A fragile \ifpl inside a \part title
% is expanded one level as it is written to the .toc, which lands
% "\ifdefstring {en}{pl}{...}{...}" in the file -- and that fails on the next
% run with "Extra }, or forgotten \endgroup", an error that names neither the
% part nor this macro.
\DeclareRobustCommand{\ifpl}[2]{\ifdefstring{\booklang}{pl}{#1}{#2}}

% A robust macro cannot go into a PDF bookmark string either -- hyperref strips
% \ifdefstring and warns once per part title. Tell it which branch to take
% instead, so the bookmarks carry the right language rather than nothing.
\makeatletter
\pdfstringdefDisableCommands{%
  \ifdefstring{\booklang}{pl}{\let\ifpl\@firstoftwo}{\let\ifpl\@secondoftwo}%
}
\makeatother

% ---------- Language strings ----------
% Every user-visible word the macros below typeset comes from here. Nothing in
% programs/ or appendices/ may hard-code an English or a Polish word inside a
% macro argument that the other edition also uses.
\input{lang/\booklang}

% ============================================================
%  PROGRAM NUMBERING
%  ---------------------------------------------------------
%  Part F programs are F1..F13; the main programs restart at 1. That is
%  Stroud's scheme and it is worth reproducing, because the Foundation part is
%  explicitly optional and numbering it separately is what makes skipping it
%  feel permitted rather than delinquent.
%
%  \theHchapter is hyperref's *destination* name, and it must stay unique
%  across the whole document. Resetting the chapter counter without also
%  changing \theHchapter produces duplicate PDF destinations, a pile of
%  warnings, and bookmarks that jump to the wrong program.
% ============================================================
\newcommand{\foundationnumbering}{%
  \setcounter{chapter}{0}%
  \renewcommand{\thechapter}{F\arabic{chapter}}%
  \renewcommand{\theHchapter}{F\arabic{chapter}}%
}
\newcommand{\mainnumbering}{%
  \setcounter{chapter}{0}%
  \renewcommand{\thechapter}{\arabic{chapter}}%
  \renewcommand{\theHchapter}{M\arabic{chapter}}%
}

% \appendix resets the chapter counter and switches \thechapter to letters, but
% it knows nothing about \theHchapter, so appendix A would claim the same PDF
% destination as main program 1. Patch it once, here.
\apptocmd{\appendix}{\renewcommand{\theHchapter}{A\Alph{chapter}}}{}{%
  \PackageWarning{mfabook}{Could not patch \string\appendix\space for hyperref destinations}}

% A program is a chapter. \chaptername is set from the language file so the
% word above the title is "Program" in both editions (it happens to be the
% same word, which is a small piece of luck).
\AtBeginDocument{\renewcommand{\chaptername}{\lblProgram}}

% ---------- Headings ----------
% \raggedright on the title body is load-bearing, not cosmetic: at \Huge on a
% 13 cm text block a long title that cannot break overflows the margin badly.
\titleformat{\chapter}[display]
  {\normalfont\huge\bfseries\color{mblue}\raggedright}
  {\normalfont\normalsize\color{mgrey}\MakeUppercase{\chaptertitlename}\ \thechapter}
  {8pt}{\Huge\raggedright}
\titlespacing*{\chapter}{0pt}{10pt}{28pt}
\titleformat{\section}{\normalfont\Large\bfseries\color{mblue}}{\thesection}{0.8em}{}
\titleformat{\subsection}{\normalfont\large\bfseries\color{mgrey}}{\thesubsection}{0.7em}{}

% head=23pt in the geometry options above, not \setlength{\headheight}:
% fancyhdr computes the height it needs from the running-head font and the
% rule, wants 22.5 pt at this size, and silently overprints the top of the text
% block if it does not get it. Setting it through geometry keeps the text block
% where geometry put it instead of pushing it down.
\pagestyle{fancy}
\fancyhf{}
\fancyhead[LE]{\small\nouppercase{\leftmark}}
% The recto head carries the FRAME number, not the section. In a book of
% numbered frames "where am I" is a frame question: every Summary
% back-reference, every Quiz route and every cross-reference between programs
% is a frame number, and hunting for one by page is the single most annoying
% thing about reading a programmed text. \theframeno at shipout is the last
% frame begun on the page, which is what a reader flicking through wants.
% Guarded: the front matter and the appendices have no frames, and a head
% reading "Frame 0" over the introduction is worse than no head at all.
\fancyhead[RO]{\ifnum\value{frameno}>0\relax\small\lblFrame~\theframeno\fi}
\fancyfoot[LE,RO]{\small\thepage}
\renewcommand{\headrulewidth}{0.4pt}

% This MUST come after \pagestyle{fancy}: fancyhdr installs its own \chaptermark
% at that point and silently overwrites anything defined earlier. The symptom is
% a running-head change that appears to do nothing at all.
%
% The running head carries the *frame* number as well as the program, because
% in a book of numbered frames "where am I" is a frame question, not a page
% question, and the summary's back-references are frame numbers.
% \@chapapp, not \lblProgram. \appendix switches \@chapapp from \chaptername
% to \appendixname, and babel supplies both in the right language; hard-coding
% the word gives an appendix a running head reading "Program A. Answers" over a
% page whose own chapter head reads "APPENDIX A".
\makeatletter
\renewcommand{\chaptermark}[1]{\markboth{\@chapapp\ \thechapter.\ #1}{}}
\makeatother
\renewcommand{\sectionmark}[1]{\markright{\thesection\ #1}}

% ============================================================
%  FRAMES — the unit the whole method is built from
%  ---------------------------------------------------------
%  A program is a stream of numbered frames. Nearly every frame ends by asking
%  the reader for something, and the answer opens the *next* frame, so that the
%  reader can cover what is below and commit to an answer before seeing it.
%
%  The environment is called `fr` rather than `frame` for two reasons. First,
%  \frame is already defined by LaTeX2e (it draws a box in a picture
%  environment) and an environment named `frame` would silently redefine it.
%  Second, this is typed forty to seventy times per program, so it wants to be
%  short in the source.
%
%    \begin{fr}
%    \ans{$x = 3$}            % the answer to the previous frame, optional
%    Teaching text, ending in a question.
%    \end{fr}
% ============================================================

\newcounter{frameno}[chapter]
\renewcommand{\theframeno}{\arabic{frameno}}

% Frames are referenced constantly -- by the Summary, by the Quiz, by the
% "Can you?" checklist and by other programs. \label inside a frame must
% therefore capture the frame number, not the section number.
\makeatletter
\newcommand{\framelabelfix}{%
  \@namedef{theHframeno}{\theHchapter.\arabic{frameno}}%
}
\makeatother

\newcommand{\framerule}{%
  \par\penalty-200\vspace{9pt}%
  \noindent\textcolor{codeframe}{\rule{\linewidth}{0.5pt}}%
  \par\vspace{5pt}%
}

% The frame number as a filled badge, set in the text flow rather than in the
% margin. A margin number would have to alternate sides under `twoside` and
% would be the first thing to overflow on a 17 cm page; a badge cannot.
\newcommand{\framebadge}{%
  \begingroup
  \setlength{\fboxsep}{2.5pt}%
  \colorbox{mblue}{\textcolor{white}{\bfseries\scriptsize\theframeno}}%
  \endgroup\hspace{0.7em}%
}

\newenvironment{fr}{%
  \framerule
  \refstepcounter{frameno}%
  \nopagebreak
  \noindent\framebadge\ignorespaces
}{\par}

% The answer to the previous frame. Set apart and centred so the eye can find
% the edge to cover, and so a reader who has answered can check in one glance
% without reading on.
\newcommand{\ans}[1]{%
  \par\nobreak\vspace{-4pt}%
  \begin{tcolorbox}[
    enhanced, sharp corners, boxrule=0pt, leftrule=2.5pt,
    colback=mfaint, colframe=mgreen,
    left=8pt, right=8pt, top=4pt, bottom=4pt,
    width=0.86\linewidth, center, halign=center,
    before skip=2pt, after skip=7pt]
    #1
  \end{tcolorbox}%
  \nobreak\noindent\ignorespaces
}

% A multi-paragraph or worked answer, where a centred box would be wrong.
\newenvironment{ansblock}{%
  \par\nopagebreak
  \begin{tcolorbox}[
    enhanced, breakable, sharp corners, boxrule=0pt, leftrule=2.5pt,
    colback=mfaint, colframe=mgreen,
    left=8pt, right=8pt, top=5pt, bottom=5pt]
}{%
  \end{tcolorbox}%
  \nopagebreak\par\vspace{2pt}\noindent\ignorespaces
}

% The row of dots. Stroud's signal that the reader is to supply a word, a
% number or an expression. \blank is inline; \dotline occupies its own line,
% which is what a longer derivation step wants.
% \penalty0 offers TeX a break opportunity immediately before the dots. Without
% it the run of dots is glued to the last word of the question, and a long
% question in the Polish edition -- where the words are longer and hyphenate
% less freely -- runs into the margin.
\newcommand{\blank}[1][4em]{\penalty0\makebox[#1]{\dotfill}}
\newcommand{\dotline}{\par\vspace{2pt}\noindent\makebox[\linewidth]{\dotfill}\par\vspace{2pt}}

% "Now one for you to do." The third rung of the scaffolding gradient, and the
% one most often skipped by authors who think they have written an exercise
% when they have written another worked example.
\newcommand{\yourturn}{%
  \par\vspace{4pt}\noindent
  {\bfseries\color{mgreen}\lblYourTurn}\par\nobreak\vspace{2pt}\noindent\ignorespaces
}

% ============================================================
%  THE PROGRAM SKELETON
%  ---------------------------------------------------------
%  Every program carries the same seven parts, in the same order. They are
%  macros rather than a convention so that a missing one is a build-time
%  absence rather than something a reader notices.
% ============================================================

% ---- Learning outcomes, on entry ----------------------------------------
% Stored as well as printed, because "Can you?" at the end of the program must
% be the same list, one for one. Storing it is what makes that structural
% rather than a promise: the checklist is generated from the outcomes, so the
% two cannot drift.
\newcounter{outcomeno}[chapter]

\makeatletter
% The key under which everything belonging to the current program is stored.
% It must expand to the printed program number (F1, F2, ..., 1, 2, ...), which
% is what makes an answer findable from the program that owns it.
\newcommand{\mfa@progkey}{\thechapter}
\newcommand{\outcome}[2]{% #1 = frames that teach it, #2 = the outcome
  \stepcounter{outcomeno}%
  \csgdef{mfa@out@\mfa@progkey @\theoutcomeno}{#2}%
  \csgdef{mfa@outfr@\mfa@progkey @\theoutcomeno}{#1}%
  % The "Can you?" row is built here rather than looped over later. A \loop
  % inside a tabularx body does not survive alignment scanning: TeX needs to
  % see the & and the row break while it reads the cell, and a conditional
  % hides them. Accumulating the rows as tokens at declaration time sidesteps
  % that entirely, and has the better property anyway -- the checklist is
  % literally made of the outcomes, so the two cannot disagree.
  \csgappto{mfa@rows@\mfa@progkey}{%
    #2 & {\footnotesize\color{mgrey}#1}
      & \ratingcell & \ratingcell & \ratingcell & \ratingcell & \ratingcell
    \tabularnewline}%
  \item #2\hfill{\footnotesize\color{mgrey}[\lblFrames~#1]}%
}
\makeatother

\makeatletter
\newenvironment{outcomes}{%
  \setcounter{outcomeno}{0}%
  \csgdef{mfa@rows@\mfa@progkey}{}%
  \begin{tcolorbox}[
    enhanced, breakable, sharp corners,
    colback=mlight, colframe=mblue, boxrule=0pt, leftrule=3pt,
    left=8pt, right=8pt, top=6pt, bottom=6pt,
    fonttitle=\bfseries\small, title={\lblOutcomes}]
  \begin{itemize}[leftmargin=1.3em, itemsep=3pt, topsep=2pt]
}{%
  \end{itemize}
  \end{tcolorbox}
}
\makeatother

% ---- "Can you?" -----------------------------------------------------------
% Generated from the stored outcomes, not written out again by hand. Five
% columns rather than a tick box, because a self-assessment with a midpoint is
% one people actually use and a yes/no one is one they lie to.
\newcommand{\ratingcell}{\textcolor{codeframe}{\rule[-0.15em]{0.8em}{0.8em}}}

\makeatletter
\newcommand{\canyou}{%
  \section*{\lblCanYou}%
  \phantomsection\addcontentsline{toc}{section}{\lblCanYou}%
  \noindent\lblCanYouIntro\par\vspace{8pt}
  \noindent
  \begingroup\small
  \ifcsvoid{mfa@rows@\mfa@progkey}{%
    {\color{mred}\lblNoOutcomes}\par
  }{%
    \begin{tabularx}{\linewidth}{@{}Xcccccc@{}}
    \toprule
    \textbf{\lblOnEntryYouCould} & \textbf{\lblFrames} & 1 & 2 & 3 & 4 & 5 \\
    \midrule
    \csuse{mfa@rows@\mfa@progkey}
    \bottomrule
    \end{tabularx}%
  }%
  \endgroup
  \par\vspace{6pt}\noindent{\footnotesize\color{mgrey}\lblCanYouFooter}\par
}
\makeatother

% ============================================================
%  QUIZ, SUMMARY, EXERCISES
% ============================================================

% ---- Quiz -----------------------------------------------------------------
% Foundation programs only. It is a diagnostic on the way in and the same test
% on the way out, which is what lets one book serve a reader who needs the
% whole program and a reader who needs four frames of it. Each question names
% the frames that teach it, so a wrong answer routes the reader somewhere
% specific rather than back to the beginning.
\newlist{quizlist}{enumerate}{1}
\setlist[quizlist]{label=\arabic*., leftmargin=2.0em, itemsep=5pt, topsep=3pt}

\makeatletter
\newenvironment{quiz}{%
  \section*{\lblQuiz}%
  \phantomsection\addcontentsline{toc}{section}{\lblQuiz}%
  \def\mfa@exkey{Q\arabic{quizlisti}}%
  \noindent\lblQuizIntro\par\vspace{5pt}
  \begin{quizlist}
}{%
  \end{quizlist}
}
\makeatother

% The frames that teach the question just asked. Two forms, because a Quiz
% question that routes to a single frame reading "Frames 22" is the kind of
% small wrongness a reader notices and an author never does.
\newcommand{\teachesat}[1]{\hfill{\footnotesize\color{mgrey}(\lblFrames~#1)}}
\newcommand{\teachesatone}[1]{\hfill{\footnotesize\color{mgrey}(\lblFrame~#1)}}

% ---- Summary --------------------------------------------------------------
% Every item carries the frame number it came from, in square brackets. That
% is the eighth-edition improvement and it is the single cheapest thing in the
% format: it turns a summary into a return index, so a reader who fails one
% line of "Can you?" knows exactly which four frames to reread.
\newenvironment{summarybox}{%
  \section*{\lblSummary}%
  \phantomsection\addcontentsline{toc}{section}{\lblSummary}%
  \begin{tcolorbox}[
    enhanced, breakable, sharp corners,
    colback=mlight, colframe=mblue, boxrule=0pt, leftrule=3pt,
    left=8pt, right=8pt, top=6pt, bottom=6pt]
  \begin{itemize}[leftmargin=1.3em, itemsep=5pt, topsep=2pt]
}{%
  \end{itemize}
  \end{tcolorbox}
}

% \sumitem{12}{text} -- the frame reference is not decoration, it is the point.
\newcommand{\sumitem}[2]{\item #2~{\footnotesize\color{mgrey}[#1]}}

% ---- Answers ---------------------------------------------------------------
%  Answers are authored at the point of use and printed at the back.
%
%  They are carried there by a global macro store rather than by an auxiliary
%  file. The sibling books collect their figure and screenshot manifests with
%  \@starttoc and a custom file extension, which works because those entries
%  are one line of text; it is the wrong mechanism for an answer, because an
%  answer contains displayed maths, and material written to an .aux-style file
%  and read back is expanded at shipout, where fragile commands break in ways
%  whose error messages name neither the answer nor the program.
%
%  \include is sequential within a run, so a macro defined globally while
%  typesetting program F1 is still defined when the back matter is typeset.
%  That is all this needs.
\makeatletter
\newcommand{\mfa@storeanswer}[3]{% #1 program key, #2 item key, #3 body
  \csgdef{mfa@a@#1@#2}{#3}%
  % \listcsxadd, not \listcsgadd: the item must be EXPANDED as it is stored.
  % Storing the token \mfa@exkey means every entry in the list expands to
  % whatever the key happens to be at the time the back matter is typeset,
  % which is the last one.
  \listcsxadd{mfa@akeys@#1}{#2}%
}
% Used inside an exercise item. Prints nothing here; the body reappears at the
% back of the book under this program's heading.
\newcommand{\answerto}[1]{\mfa@storeanswer{\mfa@progkey}{\mfa@exkey}{#1}}
\providecommand{\mfa@exkey}{?}
\makeatother

\newlist{exlist}{enumerate}{1}
\setlist[exlist]{label=\arabic*., leftmargin=2.0em, itemsep=6pt, topsep=3pt}
\newlist{fplist}{enumerate}{1}
\setlist[fplist]{label=\arabic*., leftmargin=2.0em, itemsep=5pt, topsep=3pt}

% These MUST be defined inside \makeatletter. Outside it, \def\mfa@exkey{...}
% parses as \def\mfa @exkey{...} -- a definition of \mfa with a delimited
% parameter text -- which raises no error at all and quietly leaves every
% answer filed under the same key. That failure is invisible until the answers
% section at the back of the book prints one program's answers twice.
\makeatletter
\newenvironment{testexercises}{%
  \section*{\lblTestExercises}%
  \phantomsection\addcontentsline{toc}{section}{\lblTestExercises}%
  \def\mfa@exkey{T\arabic{exlisti}}%
  \noindent\lblTestIntro\par\vspace{5pt}
  \begin{exlist}
}{%
  \end{exlist}
}

\newenvironment{furtherproblems}{%
  \section*{\lblFurther}%
  \phantomsection\addcontentsline{toc}{section}{\lblFurther}%
  \def\mfa@exkey{P\arabic{fplisti}}%
  \noindent\lblFurtherIntro\par\vspace{5pt}
  \begin{fplist}
}{%
  \end{fplist}
}
\makeatother

% ---- The programs register -------------------------------------------------
% \program{Title} opens a program: it is \chapter plus the bookkeeping that
% lets the back matter reproduce this program's answers under its own name.
\makeatletter
\newcommand{\program}[1]{%
  \chapter{#1}%
  \csgdef{mfa@title@\thechapter}{#1}%
  % Expanded, for the same reason as the answer keys above.
  \listxadd{\mfaprograms}{\thechapter}%
}
\newcommand{\mfa@printoneanswer}[2]{% #1 program key, #2 item key
  \item[\textbf{#2}] \csuse{mfa@a@#1@#2}%
}
\newcommand{\mfa@answerprogram}[1]{%
  \subsection*{\lblProgram~#1\quad\textmd{\itshape\csuse{mfa@title@#1}}}%
  \ifcsdef{mfa@akeys@#1}{%
    \begin{list}{}{%
      \setlength{\leftmargin}{3.2em}\setlength{\labelwidth}{2.7em}%
      \setlength{\labelsep}{0.5em}\setlength{\itemsep}{3pt}%
      \setlength{\topsep}{3pt}\setlength{\parsep}{0pt}%
      \renewcommand{\makelabel}[1]{\hfill##1}}
      \forlistcsloop{\mfa@printoneanswer{#1}}{mfa@akeys@#1}
    \end{list}%
  }{%
    \par\noindent{\color{mred}\lblNoAnswers}\par
  }%
}
% The answers body, without a heading of its own, so the appendix that carries
% it can be a numbered appendix like any other and appear in the running heads
% and the ToC in the ordinary way.
\newcommand{\answersbody}{%
  \noindent\lblAnswersIntro\par\vspace{8pt}
  \ifdef{\mfaprograms}{\forlistloop{\mfa@answerprogram}{\mfaprograms}}{%
    \noindent{\color{mred}\lblNoAnswers}\par}%
}
% Unnumbered form, for a draft build that has no appendix wiring yet.
\newcommand{\printanswers}{%
  \chapter*{\lblAnswers}%
  \phantomsection\addcontentsline{toc}{chapter}{\lblAnswers}%
  \markboth{\lblAnswers}{\lblAnswers}%
  \answersbody
}
\makeatother

% ============================================================
%  ADMONITIONS
%  ---------------------------------------------------------
%  A deliberately small, fixed vocabulary. Each one earns its place by doing
%  something no other box does; a seventh would be a smell.
% ============================================================

\newtcolorbox{note}[1][]{
  enhanced, breakable, sharp corners,
  colback=mlight, colframe=mblue, boxrule=0pt, leftrule=3pt,
  left=8pt, right=8pt, top=6pt, bottom=6pt,
  fonttitle=\bfseries\small, title={\lblNote}, #1
}

\newtcolorbox{warning}[1][]{
  enhanced, breakable, sharp corners,
  colback=mlight, colframe=mred, boxrule=0pt, leftrule=3pt,
  left=8pt, right=8pt, top=6pt, bottom=6pt,
  fonttitle=\bfseries\small, title={\lblWarning}, #1
}

% The misconception, stated as a summary after it has been walked into. The
% trap itself belongs in the frames -- elicit the wrong answer, then say
% flatly that it is wrong -- and this box is the thing the reader marks with a
% finger and comes back to.
\newtcolorbox{trapbox}[1][]{
  enhanced, breakable, sharp corners,
  colback=white, colframe=mred, boxrule=0.8pt,
  left=8pt, right=8pt, top=6pt, bottom=6pt,
  fonttitle=\bfseries\small, title={\lblTrap}, #1
}

% Where the mathematics just done shows up in a system the reader has actually
% deployed. This is the box that separates the book from a maths textbook, and
% it is also the one most easily filled with waffle: if it cannot name a
% specific line of a specific system, it should not be written.
\newtcolorbox{aibox}[1][]{
  enhanced, breakable, sharp corners,
  colback=mlight, colframe=mteal, boxrule=0pt, leftrule=3pt,
  left=8pt, right=8pt, top=6pt, bottom=6pt,
  fonttitle=\bfseries\small, title={\lblInAI}, #1
}

% What is being asserted rather than proved, and where to read the proof.
% Stroud's method buys fluency at the cost of rigour, and the honest thing is
% to say so at each point rather than once in an introduction nobody rereads.
\newtcolorbox{rigourbox}[1][]{
  enhanced, breakable, sharp corners,
  colback=white, colframe=mgrey, boxrule=0.6pt,
  left=8pt, right=8pt, top=6pt, bottom=6pt,
  fonttitle=\bfseries\small, title={\lblRigour}, #1
}

% Notation that genuinely differs -- between the Polish and English editions,
% and between disciplines. Matrix calculus alone has two incompatible layout
% conventions, and half the confusion about backpropagation is a transpose
% that came from the other one.
\newtcolorbox{notationbox}[1][]{
  enhanced, breakable, sharp corners,
  colback=mlight, colframe=mamber, boxrule=0pt, leftrule=3pt,
  left=8pt, right=8pt, top=6pt, bottom=6pt,
  fonttitle=\bfseries\small, title={\lblNotation}, #1
}

% The maths analogue of the sibling books' "run this before you trust it".
% A number in this book is either produced by a script under code/ and pulled
% in with \val, or it carries this box. Nothing ships to a reader with one
% still attached.
\newtcolorbox{verifybox}[1][]{
  enhanced, breakable, sharp corners,
  colback=white, colframe=mamber, boxrule=0.8pt,
  left=8pt, right=8pt, top=6pt, bottom=6pt,
  fonttitle=\bfseries\small, title={\lblVerify}, #1
}

\newtcolorbox{exercisebox}[1][]{
  enhanced, breakable, sharp corners,
  colback=white, colframe=mgreen, boxrule=0pt, leftrule=3pt,
  left=8pt, right=8pt, top=6pt, bottom=6pt,
  fonttitle=\bfseries\small, title={\lblExercise}, #1
}

% ============================================================
%  CODE LISTINGS
%  ---------------------------------------------------------
%  Code appears here for one reason only: to check a number. Every listing in
%  this book is runnable and every number it prints is the number in the text.
% ============================================================
\lstdefinelanguage{pythonx}{
  language=Python,
  morekeywords={as, with, assert, match, case},
  morekeywords=[2]{
    float64,float32,float16,bfloat16,int8,uint8,
    np,numpy,array,asarray,frombuffer,nextafter,finfo,iinfo,
    exp,log,log1p,expm1,logaddexp,sum,max,abs,sqrt,mean,var
  },
  sensitive=true
}

\lstdefinestyle{mfacode}{
  basicstyle=\ttfamily\footnotesize,
  keywordstyle=\color{kw}\bfseries,
  keywordstyle=[2]\color{typ},
  stringstyle=\color{str},
  commentstyle=\color{cmt}\itshape,
  numbers=left,
  numberstyle=\ttfamily\tiny\color{mgrey},
  numbersep=8pt,
  backgroundcolor=\color{codebg},
  frame=single,
  rulecolor=\color{codeframe},
  framesep=6pt,
  breaklines=true,
  breakatwhitespace=false,
  postbreak=\mbox{\textcolor{mgrey}{$\hookrightarrow$}\space},
  showstringspaces=false,
  tabsize=4,
  columns=fullflexible,
  keepspaces=true,
  captionpos=b,
  aboveskip=1.2em,
  belowskip=1.2em,
  % Maps the characters most likely to be pasted into a listing by accident.
  %
  % Do NOT add a mapping for U+00A0 (non-breaking space). It looks like the
  % obvious next entry and it makes `listings` abort the run with "Improper
  % alphabetic constant" -- fatal, no PDF, and the message names neither the
  % character nor this line. The ASCII-in-listings convention is the guard
  % against it instead. This trap has now been hit in two of the three books
  % in this series.
  literate={ł}{{\l}}1 {Ł}{{\L}}1 {ą}{{\k a}}1 {ę}{{\k e}}1
           {ś}{{\'s}}1 {ć}{{\'c}}1 {ż}{{\.z}}1 {ź}{{\'z}}1 {ó}{{\'o}}1 {ń}{{\'n}}1
           {Ą}{{\k A}}1 {Ę}{{\k E}}1 {Ś}{{\'S}}1 {Ć}{{\'C}}1
           {Ż}{{\.Z}}1 {Ź}{{\'Z}}1 {Ó}{{\'O}}1 {Ń}{{\'N}}1
           {—}{{-{}-}}2 {–}{{-}}1 {‘}{{'}}1 {’}{{'}}1
           {“}{{"}}1 {”}{{"}}1 {°}{{\textdegree}}1
}

\lstset{style=mfacode}

\lstnewenvironment{python}[1][]
  {\lstset{language=pythonx,style=mfacode,#1}}
  {}

\lstnewenvironment{shellcmd}[1][]
  {\lstset{language=bash,style=mfacode,numbers=none,keywordstyle=\color{mteal}\bfseries,#1}}
  {}

% Terminal transcripts: no line numbers, no keyword colouring. Everything in
% one of these was copied from a real run.
\lstnewenvironment{console}[1][]
  {\lstset{style=mfacode,numbers=none,basicstyle=\ttfamily\scriptsize,
           keywordstyle=\color{black},backgroundcolor=\color{white},#1}}
  {}

\newcommand{\pyfile}[3]{%
  \lstinputlisting[language=pythonx,style=mfacode,caption={#2},label={#3}]{#1}%
}
\newcommand{\pyregion}[4]{%
  \lstinputlisting[
    language=pythonx, style=mfacode,
    linerange={--8<--\ [start:#2]--- 8<--\ [end:#2]},
    includerangemarker=false,
    caption={#3}, label={#4}]{#1}%
}

% Inline literals
\newcommand{\code}[1]{\texttt{\small #1}}
\newcommand{\pkg}[1]{\texttt{\small\color{mblue}#1}}
\newcommand{\api}[1]{\texttt{\small\color{mteal}#1}}

% ============================================================
%  COMPUTED VALUES
%  ---------------------------------------------------------
%  The house rule in the sibling books is "run the listing, or mark it". The
%  analogue here is stronger, because a wrong digit in a maths book is not a
%  broken example, it is a lie the reader will carry.
%
%  Every number that came out of a computation is written by a script under
%  code/ into figures/values/<program>.tex as
%
%      \mfaval{f01.eps64}{2.220446049250313e-16}
%
%  and pulled into the text with \val{f01.eps64}. The script is the source of
%  truth; the book cannot disagree with it, because the book does not contain
%  the digits. A value the scripts no longer produce prints a visible marker
%  and is counted by `make debt`.
%
%  \val also applies the locale's decimal separator, which is how the Polish
%  edition gets its decimal comma without a translator having to remember.
% ============================================================
\IfFileExists{siunitx.sty}{%
  \usepackage{siunitx}%
  \sisetup{
    group-digits = integer,
    group-minimum-digits = 5,
    group-separator = {\,},
    exponent-product = \cdot,
    output-decimal-marker = {\mfadecimalmarker}
  }%
  \newcommand{\mfanum}[1]{\num{##1}}%
}{%
  \newcommand{\mfanum}[1]{##1}%
}

\makeatletter
\newcommand{\mfaval}[2]{\csgdef{mfaval@#1}{#2}\listxadd{\mfavalues}{#1}}
% \num on a non-numeric body is a FATAL siunitx error -- "Invalid number" --
% and under -halt-on-error that is the whole build. A script emitting a word
% where the book expected a number is a plausible mistake.
%
% The check is NOT done here. Deciding in TeX whether a token list is a number
% means parsing it character by character, which is fragile and was got wrong
% once already. The generator knows the answer for free, so it emits \mfaval
% for a number and \mfavaltext for a word, and this end only has to enforce
% which macro reads which.
\newcommand{\val}[1]{%
  \ifcsdef{mfaval@#1}{%
    \ifcsdef{mfavaltext@#1}{%
      \PackageError{mfabook}{Value `#1' is not a number}%
        {\string\val\space passes its body to siunitx, which refuses anything
         that is not a number. Use \string\valtext{#1} instead, or fix the
         script in code/ that emits this key.}%
    }{}%
    \mfanum{\csuse{mfaval@#1}}%
  }{\textcolor{mred}{\textbf{[?\,#1]}}}%
}
% A value that is deliberately a word rather than a number -- a format name, a
% verdict, a unit. Rare, and a visibly different macro so nobody reaches for
% \val and gets an error they cannot place.
\newcommand{\valtext}[1]{%
  \ifcsdef{mfaval@#1}{\csuse{mfaval@#1}}%
    {\textcolor{mred}{\textbf{[?\,#1]}}}%
}
\newcommand{\mfavaltext}[2]{%
  \csgdef{mfaval@#1}{#2}\csgdef{mfavaltext@#1}{}\listxadd{\mfavalues}{#1}%
}

% The raw stored string, for the cases where a number must appear exactly as
% the machine printed it -- inside a console block, or where the digits
% themselves are the point and grouping them would obscure the pattern.
\newcommand{\rawval}[1]{%
  \ifcsdef{mfaval@#1}{\csuse{mfaval@#1}}%
    {\textcolor{mred}{\textbf{[?\,#1]}}}%
}
\makeatother

% figures/values/all.tex is generated by `make numbers` and does nothing but
% \input each program's value file. A build with no values yet still compiles;
% every \val{} in it prints a visible marker instead of a number, which is the
% correct behaviour for a draft and an obvious failure for a release.
\IfFileExists{figures/values/all.tex}{% Generated by `make numbers` --- do not edit.
% Generated by code/f01_numbers.py --- do not edit.
% Regenerate with `make numbers`; `make verify` fails if this file and
% the script disagree, which is what stops a number in the book drifting
% away from the computation that justifies it.

\mfaval{f01.params}{7000000000}
\mfaval{f01.weights.bytes}{14000000000}
\mfaval{f01.weights.gb}{14}
\mfaval{f01.weights.gib}{13.04}
\mfaval{f01.weights.fp32.gb}{28}
\mfaval{f01.gib.per.gb}{0.9313}
\mfaval{f01.gib.gap.pct}{6.9}
\mfaval{f01.two.ten}{1024}
\mfaval{f01.two.ten.err.pct}{2.40}
\mfaval{f01.two.eighty}{1208925819614629174706176}
\mfaval{f01.two.eighty.err.pct}{20.89}
\mfaval{f01.tokens}{2000000000000}
\mfaval{f01.train.flops}{8.40e22}
\mfaval{f01.train.flops.exp}{22}
\mfaval{f01.device.flops}{4e14}
\mfaval{f01.device.util.pct}{40}
\mfaval{f01.device.days}{6076}
\mfaval{f01.device.years}{16.6}
\mfaval{f01.devices.for.30.days}{203}
\mfaval{f01.infer.flops.per.token}{14000000000}
\mfaval{f01.base.ms}{200}
\mfaval{f01.fifty.pct.less.ms}{100}
\mfaval{f01.fifty.pct.more.rate.ms}{133.3}
\mfaval{f01.speedup.discrepancy.ms}{33.3}

}{%
  \AtBeginDocument{\PackageWarning{mfabook}{No computed values found. Run `make numbers`.}}}

% ============================================================
%  DIAGRAMS — Mermaid pipeline
%  ---------------------------------------------------------
%  \mermaidfig{key}{caption}{one-line manifest description}
%
%  Source of truth is figures/mermaid/<lang>/<key>.mmd, committed. `make
%  diagrams` renders each to figures/diagrams/<lang>/<key>.pdf, which is build
%  output and is not committed, so a diagram change reviews as a text diff.
%
%  The source is per-language because a diagram with English labels in the
%  Polish edition is exactly the kind of half-translation that makes a
%  bilingual book feel machine-made. CI checks that both languages carry the
%  same set of keys.
%
%  If the rendered PDF is absent the macro draws a placeholder AND typesets the
%  Mermaid source, so a build without mermaid-cli still shows the structure.
% ============================================================
\makeatletter
\newcommand{\listofdiagrams}{%
  \section*{\lblDiagramManifest}%
  \phantomsection\addcontentsline{toc}{section}{\lblDiagramManifest}%
  \@starttoc{dgm}%
}
\makeatother

\newcommand{\mermaidfig}[3]{%
  \addcontentsline{dgm}{subsection}{\protect\numberline{}\texttt{#1.mmd} --- #3}%
  \IfFileExists{figures/diagrams/\booklang/#1.pdf}{%
    \begin{figure}[htbp]
    \centering
    \includegraphics[width=0.95\linewidth,height=0.42\textheight,keepaspectratio]{figures/diagrams/\booklang/#1.pdf}%
    \caption{#2}\label{fig:#1}
    \end{figure}%
  }{%
    % The unrendered case deliberately does NOT float. A Mermaid source
    % typeset verbatim is often taller than a page, and a float cannot break,
    % so wrapping it in `figure` produces an overfull box per diagram and a
    % pile of tcolorbox warnings. In the flow it breaks like any other text.
    \par\medskip
    \begin{tcolorbox}[
      enhanced, breakable, width=\linewidth, sharp corners,
      colback=white, colframe=mgrey, boxrule=0.8pt,
      left=8pt, right=8pt, top=8pt, bottom=8pt]
      {\bfseries\color{mgrey}\small \lblNotRendered}\\[4pt]
      {\ttfamily\scriptsize\color{mgrey} figures/mermaid/\booklang/#1.mmd}\\[6pt]
      \IfFileExists{figures/mermaid/\booklang/#1.mmd}{%
        % ASCII only, like every other listing in this book: the fallback runs
        % the Mermaid source through `listings`, which aborts on a multi-byte
        % character it has no literate mapping for.
        \lstinputlisting[style=mfacode,numbers=none,basicstyle=\ttfamily\scriptsize,
                         backgroundcolor=\color{mlight},frame=none,
                         keywordstyle=\color{black}]{figures/mermaid/\booklang/#1.mmd}%
      }{%
        {\small\color{mred} \lblSourceMissing: \texttt{figures/mermaid/\booklang/#1.mmd}}%
      }%
    \end{tcolorbox}%
    \captionof{figure}{#2}\label{fig:#1}
    \par\medskip
  }%
}

% ============================================================
%  PROGRAM STUBS
%  ---------------------------------------------------------
%  Prints a visible marker so an unwritten program cannot be mistaken for a
%  written one and the page count never flatters the draft. `make debt` counts
%  them, and CI publishes the count on every build.
% ============================================================
\newcommand{\programstub}[1]{%
  \begin{tcolorbox}[
    enhanced, breakable, width=\linewidth, sharp corners,
    colback=mlight, colframe=mred, boxrule=1pt,
    borderline={0.6pt}{3pt}{mred, dashed},
    left=12pt, right=12pt, top=12pt, bottom=12pt]
    {\bfseries\color{mred}\large \lblNotWritten}\\[8pt]
    \small\raggedright #1
  \end{tcolorbox}%
}

% ============================================================
%  MATHEMATICAL OPERATORS
%  ---------------------------------------------------------
%  Language-invariant names live here. The ones that genuinely differ between
%  Polish and English usage -- tan/tg, Var/D^2, gcd/NWD, and the interval
%  brackets -- are defined in lang/en.tex and lang/pl.tex instead, so the
%  notation contract is executed by the build rather than remembered by a
%  translator.
% ============================================================
\DeclareMathOperator*{\argmin}{arg\,min}
\DeclareMathOperator*{\argmax}{arg\,max}
\DeclareMathOperator{\tr}{tr}
\DeclareMathOperator{\rank}{rank}
\DeclareMathOperator{\diag}{diag}
\DeclareMathOperator{\sgn}{sgn}
\DeclareMathOperator{\relu}{ReLU}
\DeclareMathOperator{\softmax}{softmax}
\DeclareMathOperator{\logsumexp}{logsumexp}
\DeclareMathOperator{\fl}{fl}
\DeclareMathOperator{\ulp}{ulp}
\DeclareMathOperator{\round}{round}

% A bare \log is FORBIDDEN in this book, and tools/check_notation.py fails the
% build on one. In Polish textbooks \log means base ten; in machine-learning
% writing it means base e. The same three characters carry opposite meanings to
% the two readerships this book serves, and they collide hardest in exactly the
% programs where it matters most -- entropy in bits and cross-entropy in nats
% are two programs apart. Write \ln for nats and \logb{2} for bits. It costs
% two characters and removes the ambiguity entirely.
\makeatletter
\let\mfa@origlog\log
\renewcommand{\log}{%
  \PackageError{mfabook}{A bare \string\log\space is forbidden in this book}%
    {Write \string\ln\space for natural logarithms or \string\logb{2}\space for
     bits. The base is never left to a convention here; see Appendix B.}%
  \mfa@origlog}
% The one legitimate form: an explicit base.
\newcommand{\logb}[1]{\mathop{\mfa@origlog}\nolimits_{#1}}
% ...and a way to write the forbidden thing when the subject IS the ban, which
% is Appendix B and nowhere else.
\newcommand{\mfalogplain}{\mathop{\mfa@origlog}\nolimits}
\makeatother

\newcommand{\R}{\mathbb{R}}
\newcommand{\N}{\mathbb{N}}
\newcommand{\Z}{\mathbb{Z}}
\newcommand{\Q}{\mathbb{Q}}
\newcommand{\C}{\mathbb{C}}
\newcommand{\Fset}{\mathbb{F}}          % the floating-point numbers, F1's subject
\newcommand{\vect}[1]{\mathbf{#1}}
\newcommand{\mat}[1]{\mathbf{#1}}
\newcommand{\T}{^{\mathsf{T}}}
\newcommand{\norm}[1]{\left\lVert #1 \right\rVert}
\newcommand{\abs}[1]{\left\lvert #1 \right\rvert}
\newcommand{\inner}[2]{\left\langle #1, #2 \right\rangle}
\newcommand{\eps}{\varepsilon}

% ============================================================
%  INDEX
%  ---------------------------------------------------------
%  Two-column, and its entries include \texttt{} names that do not hyphenate.
%  Ragged right keeps multi-word entries off the margin; it cannot help a
%  single token wider than the column, so keep verbatim index entries under
%  about 29 characters and index longer names by concept instead. Both sibling
%  books learned this the same way.
% ============================================================
\apptocmd{\theindex}{\raggedright}{}{%
  \PackageWarning{mfabook}{Could not patch theindex for ragged-right setting}}

\makeindex

% ============================================================
%  PINNED FACTS — single source of truth
%  Change these here; they propagate through both editions.
% ============================================================
\newcommand{\bookedition}{0.1}
% \pinnedon is a user-visible string and lives in lang/<lang>.tex, not here.
% It said "August 2026" on the Polish title page for exactly as long as it
% took somebody to read the Polish title page.
\newcommand{\pyver}{Python 3.13}
\newcommand{\numpyver}{2.3}
 still compiles and says what it did.
\providecommand{\booklang}{en}

\usepackage{etoolbox}   % loaded first: the language switch below needs it

% ---------- Page geometry ----------
% Same 17 x 24 cm block as the other two books in this series, so a reader who
% owns one recognises the second. Frames are short, so the measure matters
% more here than in a prose book: a frame that runs past a page turn defeats
% the cover-the-next-frame discipline the whole method rests on.
\usepackage[
  paperwidth=17cm, paperheight=24cm,
  inner=2.2cm, outer=1.8cm, top=2.2cm, bottom=2.4cm,
  head=23pt, headsep=0.7cm, footskip=1.2cm
]{geometry}

% ---------- Encoding & fonts ----------
\usepackage[T1]{fontenc}
\usepackage[utf8]{inputenc}
\IfFileExists{lmodern.sty}{\usepackage{lmodern}}{}
\IfFileExists{newtxtext.sty}{\usepackage{newtxtext}}{}
% amsmath goes here, not down with the other packages, because newtxmath
% patches amsmath's internals and must be loaded AFTER it. The wrong order does
% not fail on a machine without newtxmath -- which is most CI images -- and
% fails on the author's full TeX Live.
%
% amssymb only when newtxmath is absent: newtxmath supplies the AMS symbols
% itself (its amssymbols option is on by default) and loading both clashes.
\usepackage{amsmath}
\IfFileExists{newtxmath.sty}{\usepackage{newtxmath}}{\usepackage{amssymb}}
\IfFileExists{inconsolata.sty}{\usepackage[varqu,scaled=0.95]{inconsolata}}{}
\IfFileExists{lmodern.sty}{\usepackage{microtype}}{\usepackage[expansion=false]{microtype}}

% ---------- Language ----------
% Loading babel with a language whose .ldf is absent is a *fatal* error, not a
% warning, so both the package and each language file are probed before either
% is requested. This is inherited from the sibling books, where a minimal TeX
% installation without texlive-lang-polish aborted the run with no PDF and an
% error message that named neither babel nor the language.
%
% The Polish edition needs Polish as the main language for hyphenation; the
% English edition needs British. Both editions load the other language too,
% because the glossary appendix is bilingual in both.
\IfFileExists{babel.sty}{%
  \IfFileExists{polish.ldf}{%
    \IfFileExists{british.ldf}{%
      \ifdefstring{\booklang}{pl}%
        {\usepackage[british,main=polish]{babel}}%
        {\usepackage[polish,main=british]{babel}}%
    }{\usepackage[polish]{babel}}%
  }{%
    \IfFileExists{british.ldf}{\usepackage[british]{babel}}{}%
  }%
}{}

% babel-polish makes " an active shorthand. listings reads its body verbatim
% and an active character in a verbatim context is a category-code accident
% waiting to happen, so the shorthand is turned off globally. Polish quotation
% marks are produced with \enquote-style macros from lang/pl.tex instead.
\ifdefstring{\booklang}{pl}{\AtBeginDocument{\shorthandoff{"}}}{}

% ---------- Core packages ----------
\usepackage{graphicx}
\usepackage{xcolor}
\usepackage{booktabs}
\usepackage{tabularx}
\usepackage{longtable}
\usepackage{array}
\usepackage{enumitem}
\IfFileExists{mathtools.sty}{\usepackage{mathtools}}{}
\usepackage{listings}
\usepackage[most]{tcolorbox}
\usepackage{fancyhdr}
\usepackage{titlesec}
\usepackage{caption}
\usepackage{imakeidx}
% csquotes with autostyle follows babel's active language, so the identical
% source \enquote{x} sets 'x' in the English edition and the Polish quotation
% marks in the Polish one. This is not a nicety: babel-polish makes " an active
% shorthand, which is turned off above precisely because an active character
% near a listings body is a category-code accident waiting to happen -- so a
% quotation mark has to come from a macro rather than from a keyboard key.
\IfFileExists{csquotes.sty}{\usepackage[autostyle=true]{csquotes}}{%
  \providecommand{\enquote}[1]{``##1''}}
\usepackage[hidelinks]{hyperref}

% ---------- Palette ----------
% Deliberately the same family as the sibling books, shifted green so the three
% spines are distinguishable on a shelf.
\definecolor{mblue}{HTML}{1F4E79}
\definecolor{mgreen}{HTML}{2E6B4F}
\definecolor{mteal}{HTML}{0E7C7B}
\definecolor{mamber}{HTML}{B26A00}
\definecolor{mred}{HTML}{9B2C2C}
\definecolor{mgrey}{HTML}{5A5A5A}
\definecolor{mlight}{HTML}{F4F6F8}
\definecolor{mfaint}{HTML}{FBFCFD}
\definecolor{codebg}{HTML}{FAFAFA}
\definecolor{codeframe}{HTML}{D8DEE4}
\definecolor{kw}{HTML}{0B5394}
\definecolor{str}{HTML}{9B2C2C}
\definecolor{cmt}{HTML}{4F7A4F}
\definecolor{typ}{HTML}{0E7C7B}

\hypersetup{
  colorlinks=true,
  linkcolor=mblue,
  urlcolor=mteal,
  citecolor=mblue,
  pdfauthor={Konrad Cinkusz}
}

% \ifpl{Polish text}{English text} -- for the handful of places where a string
% is structural rather than content and belongs in the shared files:
% structure.tex's part titles, and the main files' front matter wiring.
% Anything longer than a few words belongs in lang/ or in the edition's own
% source, not here.
% \DeclareRobustCommand, not \newcommand. A fragile \ifpl inside a \part title
% is expanded one level as it is written to the .toc, which lands
% "\ifdefstring {en}{pl}{...}{...}" in the file -- and that fails on the next
% run with "Extra }, or forgotten \endgroup", an error that names neither the
% part nor this macro.
\DeclareRobustCommand{\ifpl}[2]{\ifdefstring{\booklang}{pl}{#1}{#2}}

% A robust macro cannot go into a PDF bookmark string either -- hyperref strips
% \ifdefstring and warns once per part title. Tell it which branch to take
% instead, so the bookmarks carry the right language rather than nothing.
\makeatletter
\pdfstringdefDisableCommands{%
  \ifdefstring{\booklang}{pl}{\let\ifpl\@firstoftwo}{\let\ifpl\@secondoftwo}%
}
\makeatother

% ---------- Language strings ----------
% Every user-visible word the macros below typeset comes from here. Nothing in
% programs/ or appendices/ may hard-code an English or a Polish word inside a
% macro argument that the other edition also uses.
\input{lang/\booklang}

% ============================================================
%  PROGRAM NUMBERING
%  ---------------------------------------------------------
%  Part F programs are F1..F13; the main programs restart at 1. That is
%  Stroud's scheme and it is worth reproducing, because the Foundation part is
%  explicitly optional and numbering it separately is what makes skipping it
%  feel permitted rather than delinquent.
%
%  \theHchapter is hyperref's *destination* name, and it must stay unique
%  across the whole document. Resetting the chapter counter without also
%  changing \theHchapter produces duplicate PDF destinations, a pile of
%  warnings, and bookmarks that jump to the wrong program.
% ============================================================
\newcommand{\foundationnumbering}{%
  \setcounter{chapter}{0}%
  \renewcommand{\thechapter}{F\arabic{chapter}}%
  \renewcommand{\theHchapter}{F\arabic{chapter}}%
}
\newcommand{\mainnumbering}{%
  \setcounter{chapter}{0}%
  \renewcommand{\thechapter}{\arabic{chapter}}%
  \renewcommand{\theHchapter}{M\arabic{chapter}}%
}

% \appendix resets the chapter counter and switches \thechapter to letters, but
% it knows nothing about \theHchapter, so appendix A would claim the same PDF
% destination as main program 1. Patch it once, here.
\apptocmd{\appendix}{\renewcommand{\theHchapter}{A\Alph{chapter}}}{}{%
  \PackageWarning{mfabook}{Could not patch \string\appendix\space for hyperref destinations}}

% A program is a chapter. \chaptername is set from the language file so the
% word above the title is "Program" in both editions (it happens to be the
% same word, which is a small piece of luck).
\AtBeginDocument{\renewcommand{\chaptername}{\lblProgram}}

% ---------- Headings ----------
% \raggedright on the title body is load-bearing, not cosmetic: at \Huge on a
% 13 cm text block a long title that cannot break overflows the margin badly.
\titleformat{\chapter}[display]
  {\normalfont\huge\bfseries\color{mblue}\raggedright}
  {\normalfont\normalsize\color{mgrey}\MakeUppercase{\chaptertitlename}\ \thechapter}
  {8pt}{\Huge\raggedright}
\titlespacing*{\chapter}{0pt}{10pt}{28pt}
\titleformat{\section}{\normalfont\Large\bfseries\color{mblue}}{\thesection}{0.8em}{}
\titleformat{\subsection}{\normalfont\large\bfseries\color{mgrey}}{\thesubsection}{0.7em}{}

% head=23pt in the geometry options above, not \setlength{\headheight}:
% fancyhdr computes the height it needs from the running-head font and the
% rule, wants 22.5 pt at this size, and silently overprints the top of the text
% block if it does not get it. Setting it through geometry keeps the text block
% where geometry put it instead of pushing it down.
\pagestyle{fancy}
\fancyhf{}
\fancyhead[LE]{\small\nouppercase{\leftmark}}
% The recto head carries the FRAME number, not the section. In a book of
% numbered frames "where am I" is a frame question: every Summary
% back-reference, every Quiz route and every cross-reference between programs
% is a frame number, and hunting for one by page is the single most annoying
% thing about reading a programmed text. \theframeno at shipout is the last
% frame begun on the page, which is what a reader flicking through wants.
% Guarded: the front matter and the appendices have no frames, and a head
% reading "Frame 0" over the introduction is worse than no head at all.
\fancyhead[RO]{\ifnum\value{frameno}>0\relax\small\lblFrame~\theframeno\fi}
\fancyfoot[LE,RO]{\small\thepage}
\renewcommand{\headrulewidth}{0.4pt}

% This MUST come after \pagestyle{fancy}: fancyhdr installs its own \chaptermark
% at that point and silently overwrites anything defined earlier. The symptom is
% a running-head change that appears to do nothing at all.
%
% The running head carries the *frame* number as well as the program, because
% in a book of numbered frames "where am I" is a frame question, not a page
% question, and the summary's back-references are frame numbers.
% \@chapapp, not \lblProgram. \appendix switches \@chapapp from \chaptername
% to \appendixname, and babel supplies both in the right language; hard-coding
% the word gives an appendix a running head reading "Program A. Answers" over a
% page whose own chapter head reads "APPENDIX A".
\makeatletter
\renewcommand{\chaptermark}[1]{\markboth{\@chapapp\ \thechapter.\ #1}{}}
\makeatother
\renewcommand{\sectionmark}[1]{\markright{\thesection\ #1}}

% ============================================================
%  FRAMES — the unit the whole method is built from
%  ---------------------------------------------------------
%  A program is a stream of numbered frames. Nearly every frame ends by asking
%  the reader for something, and the answer opens the *next* frame, so that the
%  reader can cover what is below and commit to an answer before seeing it.
%
%  The environment is called `fr` rather than `frame` for two reasons. First,
%  \frame is already defined by LaTeX2e (it draws a box in a picture
%  environment) and an environment named `frame` would silently redefine it.
%  Second, this is typed forty to seventy times per program, so it wants to be
%  short in the source.
%
%    \begin{fr}
%    \ans{$x = 3$}            % the answer to the previous frame, optional
%    Teaching text, ending in a question.
%    \end{fr}
% ============================================================

\newcounter{frameno}[chapter]
\renewcommand{\theframeno}{\arabic{frameno}}

% Frames are referenced constantly -- by the Summary, by the Quiz, by the
% "Can you?" checklist and by other programs. \label inside a frame must
% therefore capture the frame number, not the section number.
\makeatletter
\newcommand{\framelabelfix}{%
  \@namedef{theHframeno}{\theHchapter.\arabic{frameno}}%
}
\makeatother

\newcommand{\framerule}{%
  \par\penalty-200\vspace{9pt}%
  \noindent\textcolor{codeframe}{\rule{\linewidth}{0.5pt}}%
  \par\vspace{5pt}%
}

% The frame number as a filled badge, set in the text flow rather than in the
% margin. A margin number would have to alternate sides under `twoside` and
% would be the first thing to overflow on a 17 cm page; a badge cannot.
\newcommand{\framebadge}{%
  \begingroup
  \setlength{\fboxsep}{2.5pt}%
  \colorbox{mblue}{\textcolor{white}{\bfseries\scriptsize\theframeno}}%
  \endgroup\hspace{0.7em}%
}

\newenvironment{fr}{%
  \framerule
  \refstepcounter{frameno}%
  \nopagebreak
  \noindent\framebadge\ignorespaces
}{\par}

% The answer to the previous frame. Set apart and centred so the eye can find
% the edge to cover, and so a reader who has answered can check in one glance
% without reading on.
\newcommand{\ans}[1]{%
  \par\nobreak\vspace{-4pt}%
  \begin{tcolorbox}[
    enhanced, sharp corners, boxrule=0pt, leftrule=2.5pt,
    colback=mfaint, colframe=mgreen,
    left=8pt, right=8pt, top=4pt, bottom=4pt,
    width=0.86\linewidth, center, halign=center,
    before skip=2pt, after skip=7pt]
    #1
  \end{tcolorbox}%
  \nobreak\noindent\ignorespaces
}

% A multi-paragraph or worked answer, where a centred box would be wrong.
\newenvironment{ansblock}{%
  \par\nopagebreak
  \begin{tcolorbox}[
    enhanced, breakable, sharp corners, boxrule=0pt, leftrule=2.5pt,
    colback=mfaint, colframe=mgreen,
    left=8pt, right=8pt, top=5pt, bottom=5pt]
}{%
  \end{tcolorbox}%
  \nopagebreak\par\vspace{2pt}\noindent\ignorespaces
}

% The row of dots. Stroud's signal that the reader is to supply a word, a
% number or an expression. \blank is inline; \dotline occupies its own line,
% which is what a longer derivation step wants.
% \penalty0 offers TeX a break opportunity immediately before the dots. Without
% it the run of dots is glued to the last word of the question, and a long
% question in the Polish edition -- where the words are longer and hyphenate
% less freely -- runs into the margin.
\newcommand{\blank}[1][4em]{\penalty0\makebox[#1]{\dotfill}}
\newcommand{\dotline}{\par\vspace{2pt}\noindent\makebox[\linewidth]{\dotfill}\par\vspace{2pt}}

% "Now one for you to do." The third rung of the scaffolding gradient, and the
% one most often skipped by authors who think they have written an exercise
% when they have written another worked example.
\newcommand{\yourturn}{%
  \par\vspace{4pt}\noindent
  {\bfseries\color{mgreen}\lblYourTurn}\par\nobreak\vspace{2pt}\noindent\ignorespaces
}

% ============================================================
%  THE PROGRAM SKELETON
%  ---------------------------------------------------------
%  Every program carries the same seven parts, in the same order. They are
%  macros rather than a convention so that a missing one is a build-time
%  absence rather than something a reader notices.
% ============================================================

% ---- Learning outcomes, on entry ----------------------------------------
% Stored as well as printed, because "Can you?" at the end of the program must
% be the same list, one for one. Storing it is what makes that structural
% rather than a promise: the checklist is generated from the outcomes, so the
% two cannot drift.
\newcounter{outcomeno}[chapter]

\makeatletter
% The key under which everything belonging to the current program is stored.
% It must expand to the printed program number (F1, F2, ..., 1, 2, ...), which
% is what makes an answer findable from the program that owns it.
\newcommand{\mfa@progkey}{\thechapter}
\newcommand{\outcome}[2]{% #1 = frames that teach it, #2 = the outcome
  \stepcounter{outcomeno}%
  \csgdef{mfa@out@\mfa@progkey @\theoutcomeno}{#2}%
  \csgdef{mfa@outfr@\mfa@progkey @\theoutcomeno}{#1}%
  % The "Can you?" row is built here rather than looped over later. A \loop
  % inside a tabularx body does not survive alignment scanning: TeX needs to
  % see the & and the row break while it reads the cell, and a conditional
  % hides them. Accumulating the rows as tokens at declaration time sidesteps
  % that entirely, and has the better property anyway -- the checklist is
  % literally made of the outcomes, so the two cannot disagree.
  \csgappto{mfa@rows@\mfa@progkey}{%
    #2 & {\footnotesize\color{mgrey}#1}
      & \ratingcell & \ratingcell & \ratingcell & \ratingcell & \ratingcell
    \tabularnewline}%
  \item #2\hfill{\footnotesize\color{mgrey}[\lblFrames~#1]}%
}
\makeatother

\makeatletter
\newenvironment{outcomes}{%
  \setcounter{outcomeno}{0}%
  \csgdef{mfa@rows@\mfa@progkey}{}%
  \begin{tcolorbox}[
    enhanced, breakable, sharp corners,
    colback=mlight, colframe=mblue, boxrule=0pt, leftrule=3pt,
    left=8pt, right=8pt, top=6pt, bottom=6pt,
    fonttitle=\bfseries\small, title={\lblOutcomes}]
  \begin{itemize}[leftmargin=1.3em, itemsep=3pt, topsep=2pt]
}{%
  \end{itemize}
  \end{tcolorbox}
}
\makeatother

% ---- "Can you?" -----------------------------------------------------------
% Generated from the stored outcomes, not written out again by hand. Five
% columns rather than a tick box, because a self-assessment with a midpoint is
% one people actually use and a yes/no one is one they lie to.
\newcommand{\ratingcell}{\textcolor{codeframe}{\rule[-0.15em]{0.8em}{0.8em}}}

\makeatletter
\newcommand{\canyou}{%
  \section*{\lblCanYou}%
  \phantomsection\addcontentsline{toc}{section}{\lblCanYou}%
  \noindent\lblCanYouIntro\par\vspace{8pt}
  \noindent
  \begingroup\small
  \ifcsvoid{mfa@rows@\mfa@progkey}{%
    {\color{mred}\lblNoOutcomes}\par
  }{%
    \begin{tabularx}{\linewidth}{@{}Xcccccc@{}}
    \toprule
    \textbf{\lblOnEntryYouCould} & \textbf{\lblFrames} & 1 & 2 & 3 & 4 & 5 \\
    \midrule
    \csuse{mfa@rows@\mfa@progkey}
    \bottomrule
    \end{tabularx}%
  }%
  \endgroup
  \par\vspace{6pt}\noindent{\footnotesize\color{mgrey}\lblCanYouFooter}\par
}
\makeatother

% ============================================================
%  QUIZ, SUMMARY, EXERCISES
% ============================================================

% ---- Quiz -----------------------------------------------------------------
% Foundation programs only. It is a diagnostic on the way in and the same test
% on the way out, which is what lets one book serve a reader who needs the
% whole program and a reader who needs four frames of it. Each question names
% the frames that teach it, so a wrong answer routes the reader somewhere
% specific rather than back to the beginning.
\newlist{quizlist}{enumerate}{1}
\setlist[quizlist]{label=\arabic*., leftmargin=2.0em, itemsep=5pt, topsep=3pt}

\makeatletter
\newenvironment{quiz}{%
  \section*{\lblQuiz}%
  \phantomsection\addcontentsline{toc}{section}{\lblQuiz}%
  \def\mfa@exkey{Q\arabic{quizlisti}}%
  \noindent\lblQuizIntro\par\vspace{5pt}
  \begin{quizlist}
}{%
  \end{quizlist}
}
\makeatother

% The frames that teach the question just asked. Two forms, because a Quiz
% question that routes to a single frame reading "Frames 22" is the kind of
% small wrongness a reader notices and an author never does.
\newcommand{\teachesat}[1]{\hfill{\footnotesize\color{mgrey}(\lblFrames~#1)}}
\newcommand{\teachesatone}[1]{\hfill{\footnotesize\color{mgrey}(\lblFrame~#1)}}

% ---- Summary --------------------------------------------------------------
% Every item carries the frame number it came from, in square brackets. That
% is the eighth-edition improvement and it is the single cheapest thing in the
% format: it turns a summary into a return index, so a reader who fails one
% line of "Can you?" knows exactly which four frames to reread.
\newenvironment{summarybox}{%
  \section*{\lblSummary}%
  \phantomsection\addcontentsline{toc}{section}{\lblSummary}%
  \begin{tcolorbox}[
    enhanced, breakable, sharp corners,
    colback=mlight, colframe=mblue, boxrule=0pt, leftrule=3pt,
    left=8pt, right=8pt, top=6pt, bottom=6pt]
  \begin{itemize}[leftmargin=1.3em, itemsep=5pt, topsep=2pt]
}{%
  \end{itemize}
  \end{tcolorbox}
}

% \sumitem{12}{text} -- the frame reference is not decoration, it is the point.
\newcommand{\sumitem}[2]{\item #2~{\footnotesize\color{mgrey}[#1]}}

% ---- Answers ---------------------------------------------------------------
%  Answers are authored at the point of use and printed at the back.
%
%  They are carried there by a global macro store rather than by an auxiliary
%  file. The sibling books collect their figure and screenshot manifests with
%  \@starttoc and a custom file extension, which works because those entries
%  are one line of text; it is the wrong mechanism for an answer, because an
%  answer contains displayed maths, and material written to an .aux-style file
%  and read back is expanded at shipout, where fragile commands break in ways
%  whose error messages name neither the answer nor the program.
%
%  \include is sequential within a run, so a macro defined globally while
%  typesetting program F1 is still defined when the back matter is typeset.
%  That is all this needs.
\makeatletter
\newcommand{\mfa@storeanswer}[3]{% #1 program key, #2 item key, #3 body
  \csgdef{mfa@a@#1@#2}{#3}%
  % \listcsxadd, not \listcsgadd: the item must be EXPANDED as it is stored.
  % Storing the token \mfa@exkey means every entry in the list expands to
  % whatever the key happens to be at the time the back matter is typeset,
  % which is the last one.
  \listcsxadd{mfa@akeys@#1}{#2}%
}
% Used inside an exercise item. Prints nothing here; the body reappears at the
% back of the book under this program's heading.
\newcommand{\answerto}[1]{\mfa@storeanswer{\mfa@progkey}{\mfa@exkey}{#1}}
\providecommand{\mfa@exkey}{?}
\makeatother

\newlist{exlist}{enumerate}{1}
\setlist[exlist]{label=\arabic*., leftmargin=2.0em, itemsep=6pt, topsep=3pt}
\newlist{fplist}{enumerate}{1}
\setlist[fplist]{label=\arabic*., leftmargin=2.0em, itemsep=5pt, topsep=3pt}

% These MUST be defined inside \makeatletter. Outside it, \def\mfa@exkey{...}
% parses as \def\mfa @exkey{...} -- a definition of \mfa with a delimited
% parameter text -- which raises no error at all and quietly leaves every
% answer filed under the same key. That failure is invisible until the answers
% section at the back of the book prints one program's answers twice.
\makeatletter
\newenvironment{testexercises}{%
  \section*{\lblTestExercises}%
  \phantomsection\addcontentsline{toc}{section}{\lblTestExercises}%
  \def\mfa@exkey{T\arabic{exlisti}}%
  \noindent\lblTestIntro\par\vspace{5pt}
  \begin{exlist}
}{%
  \end{exlist}
}

\newenvironment{furtherproblems}{%
  \section*{\lblFurther}%
  \phantomsection\addcontentsline{toc}{section}{\lblFurther}%
  \def\mfa@exkey{P\arabic{fplisti}}%
  \noindent\lblFurtherIntro\par\vspace{5pt}
  \begin{fplist}
}{%
  \end{fplist}
}
\makeatother

% ---- The programs register -------------------------------------------------
% \program{Title} opens a program: it is \chapter plus the bookkeeping that
% lets the back matter reproduce this program's answers under its own name.
\makeatletter
\newcommand{\program}[1]{%
  \chapter{#1}%
  \csgdef{mfa@title@\thechapter}{#1}%
  % Expanded, for the same reason as the answer keys above.
  \listxadd{\mfaprograms}{\thechapter}%
}
\newcommand{\mfa@printoneanswer}[2]{% #1 program key, #2 item key
  \item[\textbf{#2}] \csuse{mfa@a@#1@#2}%
}
\newcommand{\mfa@answerprogram}[1]{%
  \subsection*{\lblProgram~#1\quad\textmd{\itshape\csuse{mfa@title@#1}}}%
  \ifcsdef{mfa@akeys@#1}{%
    \begin{list}{}{%
      \setlength{\leftmargin}{3.2em}\setlength{\labelwidth}{2.7em}%
      \setlength{\labelsep}{0.5em}\setlength{\itemsep}{3pt}%
      \setlength{\topsep}{3pt}\setlength{\parsep}{0pt}%
      \renewcommand{\makelabel}[1]{\hfill##1}}
      \forlistcsloop{\mfa@printoneanswer{#1}}{mfa@akeys@#1}
    \end{list}%
  }{%
    \par\noindent{\color{mred}\lblNoAnswers}\par
  }%
}
% The answers body, without a heading of its own, so the appendix that carries
% it can be a numbered appendix like any other and appear in the running heads
% and the ToC in the ordinary way.
\newcommand{\answersbody}{%
  \noindent\lblAnswersIntro\par\vspace{8pt}
  \ifdef{\mfaprograms}{\forlistloop{\mfa@answerprogram}{\mfaprograms}}{%
    \noindent{\color{mred}\lblNoAnswers}\par}%
}
% Unnumbered form, for a draft build that has no appendix wiring yet.
\newcommand{\printanswers}{%
  \chapter*{\lblAnswers}%
  \phantomsection\addcontentsline{toc}{chapter}{\lblAnswers}%
  \markboth{\lblAnswers}{\lblAnswers}%
  \answersbody
}
\makeatother

% ============================================================
%  ADMONITIONS
%  ---------------------------------------------------------
%  A deliberately small, fixed vocabulary. Each one earns its place by doing
%  something no other box does; a seventh would be a smell.
% ============================================================

\newtcolorbox{note}[1][]{
  enhanced, breakable, sharp corners,
  colback=mlight, colframe=mblue, boxrule=0pt, leftrule=3pt,
  left=8pt, right=8pt, top=6pt, bottom=6pt,
  fonttitle=\bfseries\small, title={\lblNote}, #1
}

\newtcolorbox{warning}[1][]{
  enhanced, breakable, sharp corners,
  colback=mlight, colframe=mred, boxrule=0pt, leftrule=3pt,
  left=8pt, right=8pt, top=6pt, bottom=6pt,
  fonttitle=\bfseries\small, title={\lblWarning}, #1
}

% The misconception, stated as a summary after it has been walked into. The
% trap itself belongs in the frames -- elicit the wrong answer, then say
% flatly that it is wrong -- and this box is the thing the reader marks with a
% finger and comes back to.
\newtcolorbox{trapbox}[1][]{
  enhanced, breakable, sharp corners,
  colback=white, colframe=mred, boxrule=0.8pt,
  left=8pt, right=8pt, top=6pt, bottom=6pt,
  fonttitle=\bfseries\small, title={\lblTrap}, #1
}

% Where the mathematics just done shows up in a system the reader has actually
% deployed. This is the box that separates the book from a maths textbook, and
% it is also the one most easily filled with waffle: if it cannot name a
% specific line of a specific system, it should not be written.
\newtcolorbox{aibox}[1][]{
  enhanced, breakable, sharp corners,
  colback=mlight, colframe=mteal, boxrule=0pt, leftrule=3pt,
  left=8pt, right=8pt, top=6pt, bottom=6pt,
  fonttitle=\bfseries\small, title={\lblInAI}, #1
}

% What is being asserted rather than proved, and where to read the proof.
% Stroud's method buys fluency at the cost of rigour, and the honest thing is
% to say so at each point rather than once in an introduction nobody rereads.
\newtcolorbox{rigourbox}[1][]{
  enhanced, breakable, sharp corners,
  colback=white, colframe=mgrey, boxrule=0.6pt,
  left=8pt, right=8pt, top=6pt, bottom=6pt,
  fonttitle=\bfseries\small, title={\lblRigour}, #1
}

% Notation that genuinely differs -- between the Polish and English editions,
% and between disciplines. Matrix calculus alone has two incompatible layout
% conventions, and half the confusion about backpropagation is a transpose
% that came from the other one.
\newtcolorbox{notationbox}[1][]{
  enhanced, breakable, sharp corners,
  colback=mlight, colframe=mamber, boxrule=0pt, leftrule=3pt,
  left=8pt, right=8pt, top=6pt, bottom=6pt,
  fonttitle=\bfseries\small, title={\lblNotation}, #1
}

% The maths analogue of the sibling books' "run this before you trust it".
% A number in this book is either produced by a script under code/ and pulled
% in with \val, or it carries this box. Nothing ships to a reader with one
% still attached.
\newtcolorbox{verifybox}[1][]{
  enhanced, breakable, sharp corners,
  colback=white, colframe=mamber, boxrule=0.8pt,
  left=8pt, right=8pt, top=6pt, bottom=6pt,
  fonttitle=\bfseries\small, title={\lblVerify}, #1
}

\newtcolorbox{exercisebox}[1][]{
  enhanced, breakable, sharp corners,
  colback=white, colframe=mgreen, boxrule=0pt, leftrule=3pt,
  left=8pt, right=8pt, top=6pt, bottom=6pt,
  fonttitle=\bfseries\small, title={\lblExercise}, #1
}

% ============================================================
%  CODE LISTINGS
%  ---------------------------------------------------------
%  Code appears here for one reason only: to check a number. Every listing in
%  this book is runnable and every number it prints is the number in the text.
% ============================================================
\lstdefinelanguage{pythonx}{
  language=Python,
  morekeywords={as, with, assert, match, case},
  morekeywords=[2]{
    float64,float32,float16,bfloat16,int8,uint8,
    np,numpy,array,asarray,frombuffer,nextafter,finfo,iinfo,
    exp,log,log1p,expm1,logaddexp,sum,max,abs,sqrt,mean,var
  },
  sensitive=true
}

\lstdefinestyle{mfacode}{
  basicstyle=\ttfamily\footnotesize,
  keywordstyle=\color{kw}\bfseries,
  keywordstyle=[2]\color{typ},
  stringstyle=\color{str},
  commentstyle=\color{cmt}\itshape,
  numbers=left,
  numberstyle=\ttfamily\tiny\color{mgrey},
  numbersep=8pt,
  backgroundcolor=\color{codebg},
  frame=single,
  rulecolor=\color{codeframe},
  framesep=6pt,
  breaklines=true,
  breakatwhitespace=false,
  postbreak=\mbox{\textcolor{mgrey}{$\hookrightarrow$}\space},
  showstringspaces=false,
  tabsize=4,
  columns=fullflexible,
  keepspaces=true,
  captionpos=b,
  aboveskip=1.2em,
  belowskip=1.2em,
  % Maps the characters most likely to be pasted into a listing by accident.
  %
  % Do NOT add a mapping for U+00A0 (non-breaking space). It looks like the
  % obvious next entry and it makes `listings` abort the run with "Improper
  % alphabetic constant" -- fatal, no PDF, and the message names neither the
  % character nor this line. The ASCII-in-listings convention is the guard
  % against it instead. This trap has now been hit in two of the three books
  % in this series.
  literate={ł}{{\l}}1 {Ł}{{\L}}1 {ą}{{\k a}}1 {ę}{{\k e}}1
           {ś}{{\'s}}1 {ć}{{\'c}}1 {ż}{{\.z}}1 {ź}{{\'z}}1 {ó}{{\'o}}1 {ń}{{\'n}}1
           {Ą}{{\k A}}1 {Ę}{{\k E}}1 {Ś}{{\'S}}1 {Ć}{{\'C}}1
           {Ż}{{\.Z}}1 {Ź}{{\'Z}}1 {Ó}{{\'O}}1 {Ń}{{\'N}}1
           {—}{{-{}-}}2 {–}{{-}}1 {‘}{{'}}1 {’}{{'}}1
           {“}{{"}}1 {”}{{"}}1 {°}{{\textdegree}}1
}

\lstset{style=mfacode}

\lstnewenvironment{python}[1][]
  {\lstset{language=pythonx,style=mfacode,#1}}
  {}

\lstnewenvironment{shellcmd}[1][]
  {\lstset{language=bash,style=mfacode,numbers=none,keywordstyle=\color{mteal}\bfseries,#1}}
  {}

% Terminal transcripts: no line numbers, no keyword colouring. Everything in
% one of these was copied from a real run.
\lstnewenvironment{console}[1][]
  {\lstset{style=mfacode,numbers=none,basicstyle=\ttfamily\scriptsize,
           keywordstyle=\color{black},backgroundcolor=\color{white},#1}}
  {}

\newcommand{\pyfile}[3]{%
  \lstinputlisting[language=pythonx,style=mfacode,caption={#2},label={#3}]{#1}%
}
\newcommand{\pyregion}[4]{%
  \lstinputlisting[
    language=pythonx, style=mfacode,
    linerange={--8<--\ [start:#2]--- 8<--\ [end:#2]},
    includerangemarker=false,
    caption={#3}, label={#4}]{#1}%
}

% Inline literals
\newcommand{\code}[1]{\texttt{\small #1}}
\newcommand{\pkg}[1]{\texttt{\small\color{mblue}#1}}
\newcommand{\api}[1]{\texttt{\small\color{mteal}#1}}

% ============================================================
%  COMPUTED VALUES
%  ---------------------------------------------------------
%  The house rule in the sibling books is "run the listing, or mark it". The
%  analogue here is stronger, because a wrong digit in a maths book is not a
%  broken example, it is a lie the reader will carry.
%
%  Every number that came out of a computation is written by a script under
%  code/ into figures/values/<program>.tex as
%
%      \mfaval{f01.eps64}{2.220446049250313e-16}
%
%  and pulled into the text with \val{f01.eps64}. The script is the source of
%  truth; the book cannot disagree with it, because the book does not contain
%  the digits. A value the scripts no longer produce prints a visible marker
%  and is counted by `make debt`.
%
%  \val also applies the locale's decimal separator, which is how the Polish
%  edition gets its decimal comma without a translator having to remember.
% ============================================================
\IfFileExists{siunitx.sty}{%
  \usepackage{siunitx}%
  \sisetup{
    group-digits = integer,
    group-minimum-digits = 5,
    group-separator = {\,},
    exponent-product = \cdot,
    output-decimal-marker = {\mfadecimalmarker}
  }%
  \newcommand{\mfanum}[1]{\num{##1}}%
}{%
  \newcommand{\mfanum}[1]{##1}%
}

\makeatletter
\newcommand{\mfaval}[2]{\csgdef{mfaval@#1}{#2}\listxadd{\mfavalues}{#1}}
% \num on a non-numeric body is a FATAL siunitx error -- "Invalid number" --
% and under -halt-on-error that is the whole build. A script emitting a word
% where the book expected a number is a plausible mistake.
%
% The check is NOT done here. Deciding in TeX whether a token list is a number
% means parsing it character by character, which is fragile and was got wrong
% once already. The generator knows the answer for free, so it emits \mfaval
% for a number and \mfavaltext for a word, and this end only has to enforce
% which macro reads which.
\newcommand{\val}[1]{%
  \ifcsdef{mfaval@#1}{%
    \ifcsdef{mfavaltext@#1}{%
      \PackageError{mfabook}{Value `#1' is not a number}%
        {\string\val\space passes its body to siunitx, which refuses anything
         that is not a number. Use \string\valtext{#1} instead, or fix the
         script in code/ that emits this key.}%
    }{}%
    \mfanum{\csuse{mfaval@#1}}%
  }{\textcolor{mred}{\textbf{[?\,#1]}}}%
}
% A value that is deliberately a word rather than a number -- a format name, a
% verdict, a unit. Rare, and a visibly different macro so nobody reaches for
% \val and gets an error they cannot place.
\newcommand{\valtext}[1]{%
  \ifcsdef{mfaval@#1}{\csuse{mfaval@#1}}%
    {\textcolor{mred}{\textbf{[?\,#1]}}}%
}
\newcommand{\mfavaltext}[2]{%
  \csgdef{mfaval@#1}{#2}\csgdef{mfavaltext@#1}{}\listxadd{\mfavalues}{#1}%
}

% The raw stored string, for the cases where a number must appear exactly as
% the machine printed it -- inside a console block, or where the digits
% themselves are the point and grouping them would obscure the pattern.
\newcommand{\rawval}[1]{%
  \ifcsdef{mfaval@#1}{\csuse{mfaval@#1}}%
    {\textcolor{mred}{\textbf{[?\,#1]}}}%
}
\makeatother

% figures/values/all.tex is generated by `make numbers` and does nothing but
% \input each program's value file. A build with no values yet still compiles;
% every \val{} in it prints a visible marker instead of a number, which is the
% correct behaviour for a draft and an obvious failure for a release.
\IfFileExists{figures/values/all.tex}{% Generated by `make numbers` --- do not edit.
% Generated by code/f01_numbers.py --- do not edit.
% Regenerate with `make numbers`; `make verify` fails if this file and
% the script disagree, which is what stops a number in the book drifting
% away from the computation that justifies it.

\mfaval{f01.params}{7000000000}
\mfaval{f01.weights.bytes}{14000000000}
\mfaval{f01.weights.gb}{14}
\mfaval{f01.weights.gib}{13.04}
\mfaval{f01.weights.fp32.gb}{28}
\mfaval{f01.gib.per.gb}{0.9313}
\mfaval{f01.gib.gap.pct}{6.9}
\mfaval{f01.two.ten}{1024}
\mfaval{f01.two.ten.err.pct}{2.40}
\mfaval{f01.two.eighty}{1208925819614629174706176}
\mfaval{f01.two.eighty.err.pct}{20.89}
\mfaval{f01.tokens}{2000000000000}
\mfaval{f01.train.flops}{8.40e22}
\mfaval{f01.train.flops.exp}{22}
\mfaval{f01.device.flops}{4e14}
\mfaval{f01.device.util.pct}{40}
\mfaval{f01.device.days}{6076}
\mfaval{f01.device.years}{16.6}
\mfaval{f01.devices.for.30.days}{203}
\mfaval{f01.infer.flops.per.token}{14000000000}
\mfaval{f01.base.ms}{200}
\mfaval{f01.fifty.pct.less.ms}{100}
\mfaval{f01.fifty.pct.more.rate.ms}{133.3}
\mfaval{f01.speedup.discrepancy.ms}{33.3}

}{%
  \AtBeginDocument{\PackageWarning{mfabook}{No computed values found. Run `make numbers`.}}}

% ============================================================
%  DIAGRAMS — Mermaid pipeline
%  ---------------------------------------------------------
%  \mermaidfig{key}{caption}{one-line manifest description}
%
%  Source of truth is figures/mermaid/<lang>/<key>.mmd, committed. `make
%  diagrams` renders each to figures/diagrams/<lang>/<key>.pdf, which is build
%  output and is not committed, so a diagram change reviews as a text diff.
%
%  The source is per-language because a diagram with English labels in the
%  Polish edition is exactly the kind of half-translation that makes a
%  bilingual book feel machine-made. CI checks that both languages carry the
%  same set of keys.
%
%  If the rendered PDF is absent the macro draws a placeholder AND typesets the
%  Mermaid source, so a build without mermaid-cli still shows the structure.
% ============================================================
\makeatletter
\newcommand{\listofdiagrams}{%
  \section*{\lblDiagramManifest}%
  \phantomsection\addcontentsline{toc}{section}{\lblDiagramManifest}%
  \@starttoc{dgm}%
}
\makeatother

\newcommand{\mermaidfig}[3]{%
  \addcontentsline{dgm}{subsection}{\protect\numberline{}\texttt{#1.mmd} --- #3}%
  \IfFileExists{figures/diagrams/\booklang/#1.pdf}{%
    \begin{figure}[htbp]
    \centering
    \includegraphics[width=0.95\linewidth,height=0.42\textheight,keepaspectratio]{figures/diagrams/\booklang/#1.pdf}%
    \caption{#2}\label{fig:#1}
    \end{figure}%
  }{%
    % The unrendered case deliberately does NOT float. A Mermaid source
    % typeset verbatim is often taller than a page, and a float cannot break,
    % so wrapping it in `figure` produces an overfull box per diagram and a
    % pile of tcolorbox warnings. In the flow it breaks like any other text.
    \par\medskip
    \begin{tcolorbox}[
      enhanced, breakable, width=\linewidth, sharp corners,
      colback=white, colframe=mgrey, boxrule=0.8pt,
      left=8pt, right=8pt, top=8pt, bottom=8pt]
      {\bfseries\color{mgrey}\small \lblNotRendered}\\[4pt]
      {\ttfamily\scriptsize\color{mgrey} figures/mermaid/\booklang/#1.mmd}\\[6pt]
      \IfFileExists{figures/mermaid/\booklang/#1.mmd}{%
        % ASCII only, like every other listing in this book: the fallback runs
        % the Mermaid source through `listings`, which aborts on a multi-byte
        % character it has no literate mapping for.
        \lstinputlisting[style=mfacode,numbers=none,basicstyle=\ttfamily\scriptsize,
                         backgroundcolor=\color{mlight},frame=none,
                         keywordstyle=\color{black}]{figures/mermaid/\booklang/#1.mmd}%
      }{%
        {\small\color{mred} \lblSourceMissing: \texttt{figures/mermaid/\booklang/#1.mmd}}%
      }%
    \end{tcolorbox}%
    \captionof{figure}{#2}\label{fig:#1}
    \par\medskip
  }%
}

% ============================================================
%  PROGRAM STUBS
%  ---------------------------------------------------------
%  Prints a visible marker so an unwritten program cannot be mistaken for a
%  written one and the page count never flatters the draft. `make debt` counts
%  them, and CI publishes the count on every build.
% ============================================================
\newcommand{\programstub}[1]{%
  \begin{tcolorbox}[
    enhanced, breakable, width=\linewidth, sharp corners,
    colback=mlight, colframe=mred, boxrule=1pt,
    borderline={0.6pt}{3pt}{mred, dashed},
    left=12pt, right=12pt, top=12pt, bottom=12pt]
    {\bfseries\color{mred}\large \lblNotWritten}\\[8pt]
    \small\raggedright #1
  \end{tcolorbox}%
}

% ============================================================
%  MATHEMATICAL OPERATORS
%  ---------------------------------------------------------
%  Language-invariant names live here. The ones that genuinely differ between
%  Polish and English usage -- tan/tg, Var/D^2, gcd/NWD, and the interval
%  brackets -- are defined in lang/en.tex and lang/pl.tex instead, so the
%  notation contract is executed by the build rather than remembered by a
%  translator.
% ============================================================
\DeclareMathOperator*{\argmin}{arg\,min}
\DeclareMathOperator*{\argmax}{arg\,max}
\DeclareMathOperator{\tr}{tr}
\DeclareMathOperator{\rank}{rank}
\DeclareMathOperator{\diag}{diag}
\DeclareMathOperator{\sgn}{sgn}
\DeclareMathOperator{\relu}{ReLU}
\DeclareMathOperator{\softmax}{softmax}
\DeclareMathOperator{\logsumexp}{logsumexp}
\DeclareMathOperator{\fl}{fl}
\DeclareMathOperator{\ulp}{ulp}
\DeclareMathOperator{\round}{round}

% A bare \log is FORBIDDEN in this book, and tools/check_notation.py fails the
% build on one. In Polish textbooks \log means base ten; in machine-learning
% writing it means base e. The same three characters carry opposite meanings to
% the two readerships this book serves, and they collide hardest in exactly the
% programs where it matters most -- entropy in bits and cross-entropy in nats
% are two programs apart. Write \ln for nats and \logb{2} for bits. It costs
% two characters and removes the ambiguity entirely.
\makeatletter
\let\mfa@origlog\log
\renewcommand{\log}{%
  \PackageError{mfabook}{A bare \string\log\space is forbidden in this book}%
    {Write \string\ln\space for natural logarithms or \string\logb{2}\space for
     bits. The base is never left to a convention here; see Appendix B.}%
  \mfa@origlog}
% The one legitimate form: an explicit base.
\newcommand{\logb}[1]{\mathop{\mfa@origlog}\nolimits_{#1}}
% ...and a way to write the forbidden thing when the subject IS the ban, which
% is Appendix B and nowhere else.
\newcommand{\mfalogplain}{\mathop{\mfa@origlog}\nolimits}
\makeatother

\newcommand{\R}{\mathbb{R}}
\newcommand{\N}{\mathbb{N}}
\newcommand{\Z}{\mathbb{Z}}
\newcommand{\Q}{\mathbb{Q}}
\newcommand{\C}{\mathbb{C}}
\newcommand{\Fset}{\mathbb{F}}          % the floating-point numbers, F1's subject
\newcommand{\vect}[1]{\mathbf{#1}}
\newcommand{\mat}[1]{\mathbf{#1}}
\newcommand{\T}{^{\mathsf{T}}}
\newcommand{\norm}[1]{\left\lVert #1 \right\rVert}
\newcommand{\abs}[1]{\left\lvert #1 \right\rvert}
\newcommand{\inner}[2]{\left\langle #1, #2 \right\rangle}
\newcommand{\eps}{\varepsilon}

% ============================================================
%  INDEX
%  ---------------------------------------------------------
%  Two-column, and its entries include \texttt{} names that do not hyphenate.
%  Ragged right keeps multi-word entries off the margin; it cannot help a
%  single token wider than the column, so keep verbatim index entries under
%  about 29 characters and index longer names by concept instead. Both sibling
%  books learned this the same way.
% ============================================================
\apptocmd{\theindex}{\raggedright}{}{%
  \PackageWarning{mfabook}{Could not patch theindex for ragged-right setting}}

\makeindex

% ============================================================
%  PINNED FACTS — single source of truth
%  Change these here; they propagate through both editions.
% ============================================================
\newcommand{\bookedition}{0.1}
% \pinnedon is a user-visible string and lives in lang/<lang>.tex, not here.
% It said "August 2026" on the Polish title page for exactly as long as it
% took somebody to read the Polish title page.
\newcommand{\pyver}{Python 3.13}
\newcommand{\numpyver}{2.3}
